\section*{Item 5}\label{it:5}

\textbf{Enunciado}: Calcule os seguintes parâmetros: $L_4$, $L_1$, $L_2$, $C_1$ e $D$ para o conversor Forward abaixo. Simule o circuito e verifique se os resultados são consistentes com a expectativa. Justifique eventuais discrepâncias.

\begin{itemize}
	\item O circuito opera no modo de condução contínua.
	\item Tensão de saída de $12$ V.
	\item Ripple da corrente de saída (em $L_5$) igual a $4$ A (pico a pico).
	\item Ripple da tensão de saída de $1\%$.
	\item Relação de espiras entre $L_1$ e $L_2$: $N_1 = 10\,N_2$.
	\item $L_3$ deve ser tal que garanta a desmagnetização total do núcleo durante a condução de $D_3$.
	\item A frequência de chaveamento é de $20$ kHz.
	\item $L_2 = 10~\mu\text{H}$.
	\item Utilize o transistor R6020PNJ.
	\item Para evitar problemas de convergência na simulação, utilize o resistor $R_3$ de $10~\text{M}\Omega$, mostrado na figura abaixo, bem como a configuração ``skip initial operating point solution'' (\texttt{uic}) no comando \texttt{.tran}.
\end{itemize}

\textbf{Resolução}: inicialmente, calcula-se o ciclo de trabalho a partir da característica de conversão do conversor Forward operando em modo de condução contínua~\cite{forward_notas}
\begin{equation*}
	\frac{V_o}{V_i}
	=
	\frac{N_2}{N_1}D.
\end{equation*}

Como $N_2/N_1=0{,}1$, obtém-se
\begin{equation*}
	D
	=
	\frac{V_o}{V_i(N_2/N_1)}
	=
	\frac{12}{300\cdot0{,}1}
	=
	0{,}4.
\end{equation*}

A indutância de saída $L_4$ é obtida a partir da expressão da ondulação de corrente
\begin{equation*}
	\Delta i_{L_4}
	=
	\frac{V_o}{L_4}(1-D)T_s.
\end{equation*}

Isolando $L_4$,
\begin{equation*}
	L_4
	=
	\frac{V_o(1-D)T_s}
	{\Delta i_{L_4}}.
\end{equation*}

Substituindo os valores especificados,
\begin{equation*}
	\begin{aligned}
		L_4
		&=
		\frac{12(1-0{,}4)(50\cdot10^{-6})}
		{4}
		\\
		&=
		90~\mu\mathrm{H}.
	\end{aligned}
\end{equation*}

A capacitância do filtro de saída é calculada a partir da expressão da ondulação relativa da tensão
\begin{equation*}
	\frac{\Delta V_o}{V_o}
	=
	\frac{1-D}
	{8L_4Cf_s^2}.
\end{equation*}

Como se deseja limitar o ripple da tensão de saída a $1\%$,
\begin{equation*}
	\frac{\Delta V_o}{V_o}
	=
	0{,}01.
\end{equation*}

Isolando a capacitância,
\begin{equation*}
	C
	=
	\frac{1-D}
	{8L_4f_s^2
		\left(\frac{\Delta V_o}{V_o}\right)}.
\end{equation*}

Substituindo os valores,
\begin{equation*}
	\begin{aligned}
		C
		&=
		\frac{1-0{,}4}
		{8(90\cdot10^{-6})
			(20\cdot10^3)^2
			(0{,}01)}
		\\
		&\approx
		208~\mu\mathrm{F}.
	\end{aligned}
\end{equation*}

A relação entre as indutâncias dos enrolamentos do transformador é dada por
\begin{equation*}
	\frac{L_1}{L_2}
	=
	\left(\frac{N_1}{N_2}\right)^2.
\end{equation*}

Como $N_1/N_2=10$ e $L_2=10~\mu\mathrm{H}$,
\begin{equation*}
	\begin{aligned}
		L_1
		&=
		L_2
		\left(\frac{N_1}{N_2}\right)^2
		\\
		&=
		10~\mu\mathrm{H}\times10^2
		\\
		&=1~\mathrm{mH}.
	\end{aligned}
\end{equation*}

Por fim, o enrolamento de desmagnetização deve garantir a completa remoção do fluxo armazenado durante o intervalo de condução de $D_3$. O valor máximo da indutância é dado por~\cite{forward_notas}
\begin{equation*}
	L_{3,\max}
	=
	L_1
	\left(\frac{1}{D}-1\right)^2.
\end{equation*}

Substituindo os valores obtidos,
\begin{equation*}
	\begin{aligned}
		L_{3,\max}
		&=
		1~\mathrm{mH}
		\left(\frac{1}{0{,}4}-1\right)^2
		\\
		&=2{,}25~\mathrm{mH}.
	\end{aligned}
\end{equation*}

Foi adotado um fator de segurança de $20\%$, resultando em
\begin{equation*}
	\begin{aligned}
		L_3
		&=
		0{,}8L_{3,\max}
		\\
		&=
		1{,}8~\mathrm{mH}.
	\end{aligned}
\end{equation*}

A carga limite para operação em modo de condução contínua é dada por
\begin{equation*}
	R_{\max}
	=
	\frac{2L_4}
	{(1-D)T_s},
\end{equation*}
resultando em
\begin{equation*}
	R_{\max}
	=
	\frac{2(90~\mu\mathrm{H})}
	{(1-0{,}4)(50~\mu\mathrm{s})}
	=
	6~\Omega.
\end{equation*}

Para garantir a operação em modo contínuo, adotou-se
\begin{equation*}
	R
	=
	0{,}8R_{\max}
	=
	4{,}8~\Omega.
\end{equation*}

Assim, a corrente média na carga é
\begin{equation*}
	I_o
	=
	\frac{V_o}{R}
	=
	\frac{12}{4{,}8}
	=
	2{,}5~\mathrm{A}.
\end{equation*}

A Figura~\ref{fig:ex5} apresenta, respectivamente, a tensão de saída, a corrente no indutor de saída e a tensão dreno--source do transistor, verificando que o conversor opera em CCM.
\begin{figure}[htb!]
	\centering
	%% Creator: Matplotlib, PGF backend
%%
%% To include the figure in your LaTeX document, write
%%   \input{<filename>.pgf}
%%
%% Make sure the required packages are loaded in your preamble
%%   \usepackage{pgf}
%%
%% Also ensure that all the required font packages are loaded; for instance,
%% the lmodern package is sometimes necessary when using math font.
%%   \usepackage{lmodern}
%%
%% Figures using additional raster images can only be included by \input if
%% they are in the same directory as the main LaTeX file. For loading figures
%% from other directories you can use the `import` package
%%   \usepackage{import}
%%
%% and then include the figures with
%%   \import{<path to file>}{<filename>.pgf}
%%
%% Matplotlib used the following preamble
%%   \def\mathdefault#1{#1}
%%   \everymath=\expandafter{\the\everymath\displaystyle}
%%   \IfFileExists{scrextend.sty}{
%%     \usepackage[fontsize=8.000000pt]{scrextend}
%%   }{
%%     \renewcommand{\normalsize}{\fontsize{8.000000}{9.600000}\selectfont}
%%     \normalsize
%%   }
%%   
%%           \usepackage{amsmath}
%%           \usepackage{amssymb}
%%           \usepackage{mathtools}
%%           \usepackage{dsfont}
%%           \providecommand{\mathdefault}[1]{#1}
%%       
%%   \makeatletter\@ifpackageloaded{underscore}{}{\usepackage[strings]{underscore}}\makeatother
%%
\begingroup%
\makeatletter%
\begin{pgfpicture}%
\pgfpathrectangle{\pgfpointorigin}{\pgfqpoint{3.452797in}{3.607654in}}%
\pgfusepath{use as bounding box, clip}%
\begin{pgfscope}%
\pgfsetbuttcap%
\pgfsetmiterjoin%
\definecolor{currentfill}{rgb}{1.000000,1.000000,1.000000}%
\pgfsetfillcolor{currentfill}%
\pgfsetlinewidth{0.000000pt}%
\definecolor{currentstroke}{rgb}{1.000000,1.000000,1.000000}%
\pgfsetstrokecolor{currentstroke}%
\pgfsetdash{}{0pt}%
\pgfpathmoveto{\pgfqpoint{0.000000in}{0.000000in}}%
\pgfpathlineto{\pgfqpoint{3.452797in}{0.000000in}}%
\pgfpathlineto{\pgfqpoint{3.452797in}{3.607654in}}%
\pgfpathlineto{\pgfqpoint{0.000000in}{3.607654in}}%
\pgfpathlineto{\pgfqpoint{0.000000in}{0.000000in}}%
\pgfpathclose%
\pgfusepath{fill}%
\end{pgfscope}%
\begin{pgfscope}%
\pgfsetbuttcap%
\pgfsetmiterjoin%
\definecolor{currentfill}{rgb}{1.000000,1.000000,1.000000}%
\pgfsetfillcolor{currentfill}%
\pgfsetlinewidth{0.000000pt}%
\definecolor{currentstroke}{rgb}{0.000000,0.000000,0.000000}%
\pgfsetstrokecolor{currentstroke}%
\pgfsetstrokeopacity{0.000000}%
\pgfsetdash{}{0pt}%
\pgfpathmoveto{\pgfqpoint{0.632797in}{2.601404in}}%
\pgfpathlineto{\pgfqpoint{3.352797in}{2.601404in}}%
\pgfpathlineto{\pgfqpoint{3.352797in}{3.507654in}}%
\pgfpathlineto{\pgfqpoint{0.632797in}{3.507654in}}%
\pgfpathlineto{\pgfqpoint{0.632797in}{2.601404in}}%
\pgfpathclose%
\pgfusepath{fill}%
\end{pgfscope}%
\begin{pgfscope}%
\pgfpathrectangle{\pgfqpoint{0.632797in}{2.601404in}}{\pgfqpoint{2.720000in}{0.906250in}}%
\pgfusepath{clip}%
\pgfsetbuttcap%
\pgfsetroundjoin%
\pgfsetlinewidth{0.501875pt}%
\definecolor{currentstroke}{rgb}{0.690196,0.690196,0.690196}%
\pgfsetstrokecolor{currentstroke}%
\pgfsetstrokeopacity{0.250000}%
\pgfsetdash{{1.850000pt}{0.800000pt}}{0.000000pt}%
\pgfpathmoveto{\pgfqpoint{0.756433in}{2.601404in}}%
\pgfpathlineto{\pgfqpoint{0.756433in}{3.507654in}}%
\pgfusepath{stroke}%
\end{pgfscope}%
\begin{pgfscope}%
\pgfsetbuttcap%
\pgfsetroundjoin%
\definecolor{currentfill}{rgb}{0.000000,0.000000,0.000000}%
\pgfsetfillcolor{currentfill}%
\pgfsetlinewidth{0.803000pt}%
\definecolor{currentstroke}{rgb}{0.000000,0.000000,0.000000}%
\pgfsetstrokecolor{currentstroke}%
\pgfsetdash{}{0pt}%
\pgfsys@defobject{currentmarker}{\pgfqpoint{0.000000in}{-0.048611in}}{\pgfqpoint{0.000000in}{0.000000in}}{%
\pgfpathmoveto{\pgfqpoint{0.000000in}{0.000000in}}%
\pgfpathlineto{\pgfqpoint{0.000000in}{-0.048611in}}%
\pgfusepath{stroke,fill}%
}%
\begin{pgfscope}%
\pgfsys@transformshift{0.756433in}{2.601404in}%
\pgfsys@useobject{currentmarker}{}%
\end{pgfscope}%
\end{pgfscope}%
\begin{pgfscope}%
\pgfpathrectangle{\pgfqpoint{0.632797in}{2.601404in}}{\pgfqpoint{2.720000in}{0.906250in}}%
\pgfusepath{clip}%
\pgfsetbuttcap%
\pgfsetroundjoin%
\pgfsetlinewidth{0.501875pt}%
\definecolor{currentstroke}{rgb}{0.690196,0.690196,0.690196}%
\pgfsetstrokecolor{currentstroke}%
\pgfsetstrokeopacity{0.250000}%
\pgfsetdash{{1.850000pt}{0.800000pt}}{0.000000pt}%
\pgfpathmoveto{\pgfqpoint{1.415827in}{2.601404in}}%
\pgfpathlineto{\pgfqpoint{1.415827in}{3.507654in}}%
\pgfusepath{stroke}%
\end{pgfscope}%
\begin{pgfscope}%
\pgfsetbuttcap%
\pgfsetroundjoin%
\definecolor{currentfill}{rgb}{0.000000,0.000000,0.000000}%
\pgfsetfillcolor{currentfill}%
\pgfsetlinewidth{0.803000pt}%
\definecolor{currentstroke}{rgb}{0.000000,0.000000,0.000000}%
\pgfsetstrokecolor{currentstroke}%
\pgfsetdash{}{0pt}%
\pgfsys@defobject{currentmarker}{\pgfqpoint{0.000000in}{-0.048611in}}{\pgfqpoint{0.000000in}{0.000000in}}{%
\pgfpathmoveto{\pgfqpoint{0.000000in}{0.000000in}}%
\pgfpathlineto{\pgfqpoint{0.000000in}{-0.048611in}}%
\pgfusepath{stroke,fill}%
}%
\begin{pgfscope}%
\pgfsys@transformshift{1.415827in}{2.601404in}%
\pgfsys@useobject{currentmarker}{}%
\end{pgfscope}%
\end{pgfscope}%
\begin{pgfscope}%
\pgfpathrectangle{\pgfqpoint{0.632797in}{2.601404in}}{\pgfqpoint{2.720000in}{0.906250in}}%
\pgfusepath{clip}%
\pgfsetbuttcap%
\pgfsetroundjoin%
\pgfsetlinewidth{0.501875pt}%
\definecolor{currentstroke}{rgb}{0.690196,0.690196,0.690196}%
\pgfsetstrokecolor{currentstroke}%
\pgfsetstrokeopacity{0.250000}%
\pgfsetdash{{1.850000pt}{0.800000pt}}{0.000000pt}%
\pgfpathmoveto{\pgfqpoint{2.075221in}{2.601404in}}%
\pgfpathlineto{\pgfqpoint{2.075221in}{3.507654in}}%
\pgfusepath{stroke}%
\end{pgfscope}%
\begin{pgfscope}%
\pgfsetbuttcap%
\pgfsetroundjoin%
\definecolor{currentfill}{rgb}{0.000000,0.000000,0.000000}%
\pgfsetfillcolor{currentfill}%
\pgfsetlinewidth{0.803000pt}%
\definecolor{currentstroke}{rgb}{0.000000,0.000000,0.000000}%
\pgfsetstrokecolor{currentstroke}%
\pgfsetdash{}{0pt}%
\pgfsys@defobject{currentmarker}{\pgfqpoint{0.000000in}{-0.048611in}}{\pgfqpoint{0.000000in}{0.000000in}}{%
\pgfpathmoveto{\pgfqpoint{0.000000in}{0.000000in}}%
\pgfpathlineto{\pgfqpoint{0.000000in}{-0.048611in}}%
\pgfusepath{stroke,fill}%
}%
\begin{pgfscope}%
\pgfsys@transformshift{2.075221in}{2.601404in}%
\pgfsys@useobject{currentmarker}{}%
\end{pgfscope}%
\end{pgfscope}%
\begin{pgfscope}%
\pgfpathrectangle{\pgfqpoint{0.632797in}{2.601404in}}{\pgfqpoint{2.720000in}{0.906250in}}%
\pgfusepath{clip}%
\pgfsetbuttcap%
\pgfsetroundjoin%
\pgfsetlinewidth{0.501875pt}%
\definecolor{currentstroke}{rgb}{0.690196,0.690196,0.690196}%
\pgfsetstrokecolor{currentstroke}%
\pgfsetstrokeopacity{0.250000}%
\pgfsetdash{{1.850000pt}{0.800000pt}}{0.000000pt}%
\pgfpathmoveto{\pgfqpoint{2.734615in}{2.601404in}}%
\pgfpathlineto{\pgfqpoint{2.734615in}{3.507654in}}%
\pgfusepath{stroke}%
\end{pgfscope}%
\begin{pgfscope}%
\pgfsetbuttcap%
\pgfsetroundjoin%
\definecolor{currentfill}{rgb}{0.000000,0.000000,0.000000}%
\pgfsetfillcolor{currentfill}%
\pgfsetlinewidth{0.803000pt}%
\definecolor{currentstroke}{rgb}{0.000000,0.000000,0.000000}%
\pgfsetstrokecolor{currentstroke}%
\pgfsetdash{}{0pt}%
\pgfsys@defobject{currentmarker}{\pgfqpoint{0.000000in}{-0.048611in}}{\pgfqpoint{0.000000in}{0.000000in}}{%
\pgfpathmoveto{\pgfqpoint{0.000000in}{0.000000in}}%
\pgfpathlineto{\pgfqpoint{0.000000in}{-0.048611in}}%
\pgfusepath{stroke,fill}%
}%
\begin{pgfscope}%
\pgfsys@transformshift{2.734615in}{2.601404in}%
\pgfsys@useobject{currentmarker}{}%
\end{pgfscope}%
\end{pgfscope}%
\begin{pgfscope}%
\pgfpathrectangle{\pgfqpoint{0.632797in}{2.601404in}}{\pgfqpoint{2.720000in}{0.906250in}}%
\pgfusepath{clip}%
\pgfsetbuttcap%
\pgfsetroundjoin%
\pgfsetlinewidth{0.501875pt}%
\definecolor{currentstroke}{rgb}{0.690196,0.690196,0.690196}%
\pgfsetstrokecolor{currentstroke}%
\pgfsetstrokeopacity{0.250000}%
\pgfsetdash{{1.850000pt}{0.800000pt}}{0.000000pt}%
\pgfpathmoveto{\pgfqpoint{0.632797in}{2.642598in}}%
\pgfpathlineto{\pgfqpoint{3.352797in}{2.642598in}}%
\pgfusepath{stroke}%
\end{pgfscope}%
\begin{pgfscope}%
\pgfsetbuttcap%
\pgfsetroundjoin%
\definecolor{currentfill}{rgb}{0.000000,0.000000,0.000000}%
\pgfsetfillcolor{currentfill}%
\pgfsetlinewidth{0.803000pt}%
\definecolor{currentstroke}{rgb}{0.000000,0.000000,0.000000}%
\pgfsetstrokecolor{currentstroke}%
\pgfsetdash{}{0pt}%
\pgfsys@defobject{currentmarker}{\pgfqpoint{-0.048611in}{0.000000in}}{\pgfqpoint{-0.000000in}{0.000000in}}{%
\pgfpathmoveto{\pgfqpoint{-0.000000in}{0.000000in}}%
\pgfpathlineto{\pgfqpoint{-0.048611in}{0.000000in}}%
\pgfusepath{stroke,fill}%
}%
\begin{pgfscope}%
\pgfsys@transformshift{0.632797in}{2.642598in}%
\pgfsys@useobject{currentmarker}{}%
\end{pgfscope}%
\end{pgfscope}%
\begin{pgfscope}%
\definecolor{textcolor}{rgb}{0.000000,0.000000,0.000000}%
\pgfsetstrokecolor{textcolor}%
\pgfsetfillcolor{textcolor}%
\pgftext[x=0.476546in, y=2.604017in, left, base]{\color{textcolor}{\rmfamily\fontsize{8.000000}{9.600000}\selectfont\catcode`\^=\active\def^{\ifmmode\sp\else\^{}\fi}\catcode`\%=\active\def%{\%}$\mathdefault{0}$}}%
\end{pgfscope}%
\begin{pgfscope}%
\pgfpathrectangle{\pgfqpoint{0.632797in}{2.601404in}}{\pgfqpoint{2.720000in}{0.906250in}}%
\pgfusepath{clip}%
\pgfsetbuttcap%
\pgfsetroundjoin%
\pgfsetlinewidth{0.501875pt}%
\definecolor{currentstroke}{rgb}{0.690196,0.690196,0.690196}%
\pgfsetstrokecolor{currentstroke}%
\pgfsetstrokeopacity{0.250000}%
\pgfsetdash{{1.850000pt}{0.800000pt}}{0.000000pt}%
\pgfpathmoveto{\pgfqpoint{0.632797in}{3.031500in}}%
\pgfpathlineto{\pgfqpoint{3.352797in}{3.031500in}}%
\pgfusepath{stroke}%
\end{pgfscope}%
\begin{pgfscope}%
\pgfsetbuttcap%
\pgfsetroundjoin%
\definecolor{currentfill}{rgb}{0.000000,0.000000,0.000000}%
\pgfsetfillcolor{currentfill}%
\pgfsetlinewidth{0.803000pt}%
\definecolor{currentstroke}{rgb}{0.000000,0.000000,0.000000}%
\pgfsetstrokecolor{currentstroke}%
\pgfsetdash{}{0pt}%
\pgfsys@defobject{currentmarker}{\pgfqpoint{-0.048611in}{0.000000in}}{\pgfqpoint{-0.000000in}{0.000000in}}{%
\pgfpathmoveto{\pgfqpoint{-0.000000in}{0.000000in}}%
\pgfpathlineto{\pgfqpoint{-0.048611in}{0.000000in}}%
\pgfusepath{stroke,fill}%
}%
\begin{pgfscope}%
\pgfsys@transformshift{0.632797in}{3.031500in}%
\pgfsys@useobject{currentmarker}{}%
\end{pgfscope}%
\end{pgfscope}%
\begin{pgfscope}%
\definecolor{textcolor}{rgb}{0.000000,0.000000,0.000000}%
\pgfsetstrokecolor{textcolor}%
\pgfsetfillcolor{textcolor}%
\pgftext[x=0.417518in, y=2.992919in, left, base]{\color{textcolor}{\rmfamily\fontsize{8.000000}{9.600000}\selectfont\catcode`\^=\active\def^{\ifmmode\sp\else\^{}\fi}\catcode`\%=\active\def%{\%}$\mathdefault{10}$}}%
\end{pgfscope}%
\begin{pgfscope}%
\pgfpathrectangle{\pgfqpoint{0.632797in}{2.601404in}}{\pgfqpoint{2.720000in}{0.906250in}}%
\pgfusepath{clip}%
\pgfsetbuttcap%
\pgfsetroundjoin%
\pgfsetlinewidth{0.501875pt}%
\definecolor{currentstroke}{rgb}{0.690196,0.690196,0.690196}%
\pgfsetstrokecolor{currentstroke}%
\pgfsetstrokeopacity{0.250000}%
\pgfsetdash{{1.850000pt}{0.800000pt}}{0.000000pt}%
\pgfpathmoveto{\pgfqpoint{0.632797in}{3.420401in}}%
\pgfpathlineto{\pgfqpoint{3.352797in}{3.420401in}}%
\pgfusepath{stroke}%
\end{pgfscope}%
\begin{pgfscope}%
\pgfsetbuttcap%
\pgfsetroundjoin%
\definecolor{currentfill}{rgb}{0.000000,0.000000,0.000000}%
\pgfsetfillcolor{currentfill}%
\pgfsetlinewidth{0.803000pt}%
\definecolor{currentstroke}{rgb}{0.000000,0.000000,0.000000}%
\pgfsetstrokecolor{currentstroke}%
\pgfsetdash{}{0pt}%
\pgfsys@defobject{currentmarker}{\pgfqpoint{-0.048611in}{0.000000in}}{\pgfqpoint{-0.000000in}{0.000000in}}{%
\pgfpathmoveto{\pgfqpoint{-0.000000in}{0.000000in}}%
\pgfpathlineto{\pgfqpoint{-0.048611in}{0.000000in}}%
\pgfusepath{stroke,fill}%
}%
\begin{pgfscope}%
\pgfsys@transformshift{0.632797in}{3.420401in}%
\pgfsys@useobject{currentmarker}{}%
\end{pgfscope}%
\end{pgfscope}%
\begin{pgfscope}%
\definecolor{textcolor}{rgb}{0.000000,0.000000,0.000000}%
\pgfsetstrokecolor{textcolor}%
\pgfsetfillcolor{textcolor}%
\pgftext[x=0.417518in, y=3.381821in, left, base]{\color{textcolor}{\rmfamily\fontsize{8.000000}{9.600000}\selectfont\catcode`\^=\active\def^{\ifmmode\sp\else\^{}\fi}\catcode`\%=\active\def%{\%}$\mathdefault{20}$}}%
\end{pgfscope}%
\begin{pgfscope}%
\definecolor{textcolor}{rgb}{0.000000,0.000000,0.000000}%
\pgfsetstrokecolor{textcolor}%
\pgfsetfillcolor{textcolor}%
\pgftext[x=0.361962in,y=3.054529in,,bottom,rotate=90.000000]{\color{textcolor}{\rmfamily\fontsize{8.000000}{9.600000}\selectfont\catcode`\^=\active\def^{\ifmmode\sp\else\^{}\fi}\catcode`\%=\active\def%{\%}$v_o$ (V)}}%
\end{pgfscope}%
\begin{pgfscope}%
\pgfpathrectangle{\pgfqpoint{0.632797in}{2.601404in}}{\pgfqpoint{2.720000in}{0.906250in}}%
\pgfusepath{clip}%
\pgfsetrectcap%
\pgfsetroundjoin%
\pgfsetlinewidth{1.505625pt}%
\definecolor{currentstroke}{rgb}{0.000000,0.447059,0.698039}%
\pgfsetstrokecolor{currentstroke}%
\pgfsetdash{}{0pt}%
\pgfpathmoveto{\pgfqpoint{0.756433in}{2.642598in}}%
\pgfpathlineto{\pgfqpoint{0.756434in}{2.642598in}}%
\pgfpathlineto{\pgfqpoint{0.757070in}{2.642713in}}%
\pgfpathlineto{\pgfqpoint{0.758719in}{2.644074in}}%
\pgfpathlineto{\pgfqpoint{0.760841in}{2.648067in}}%
\pgfpathlineto{\pgfqpoint{0.763028in}{2.655019in}}%
\pgfpathlineto{\pgfqpoint{0.763028in}{2.655019in}}%
\pgfpathlineto{\pgfqpoint{0.763028in}{2.654717in}}%
\pgfpathlineto{\pgfqpoint{0.763792in}{2.657634in}}%
\pgfpathlineto{\pgfqpoint{0.777508in}{2.711345in}}%
\pgfpathlineto{\pgfqpoint{0.783586in}{2.751364in}}%
\pgfpathlineto{\pgfqpoint{0.795852in}{2.836648in}}%
\pgfpathlineto{\pgfqpoint{0.806494in}{2.925710in}}%
\pgfpathlineto{\pgfqpoint{0.812112in}{2.976026in}}%
\pgfpathlineto{\pgfqpoint{0.821625in}{3.064101in}}%
\pgfpathlineto{\pgfqpoint{0.829222in}{3.132296in}}%
\pgfpathlineto{\pgfqpoint{0.836572in}{3.198230in}}%
\pgfpathlineto{\pgfqpoint{0.839416in}{3.220673in}}%
\pgfpathlineto{\pgfqpoint{0.839416in}{3.219480in}}%
\pgfpathlineto{\pgfqpoint{0.839416in}{3.212519in}}%
\pgfpathlineto{\pgfqpoint{0.840139in}{3.224811in}}%
\pgfpathlineto{\pgfqpoint{0.848669in}{3.294721in}}%
\pgfpathlineto{\pgfqpoint{0.854439in}{3.333316in}}%
\pgfpathlineto{\pgfqpoint{0.870850in}{3.421192in}}%
\pgfpathlineto{\pgfqpoint{0.875995in}{3.437243in}}%
\pgfpathlineto{\pgfqpoint{0.875995in}{3.436862in}}%
\pgfpathlineto{\pgfqpoint{0.875995in}{3.439831in}}%
\pgfpathlineto{\pgfqpoint{0.875995in}{3.436268in}}%
\pgfpathlineto{\pgfqpoint{0.876659in}{3.439156in}}%
\pgfpathlineto{\pgfqpoint{0.877004in}{3.440438in}}%
\pgfpathlineto{\pgfqpoint{0.877004in}{3.439792in}}%
\pgfpathlineto{\pgfqpoint{0.877614in}{3.445545in}}%
\pgfpathlineto{\pgfqpoint{0.877004in}{3.438576in}}%
\pgfpathlineto{\pgfqpoint{0.877768in}{3.443165in}}%
\pgfpathlineto{\pgfqpoint{0.882615in}{3.458190in}}%
\pgfpathlineto{\pgfqpoint{0.885582in}{3.462677in}}%
\pgfpathlineto{\pgfqpoint{0.887841in}{3.463652in}}%
\pgfpathlineto{\pgfqpoint{0.892665in}{3.464094in}}%
\pgfpathlineto{\pgfqpoint{0.898338in}{3.466109in}}%
\pgfpathlineto{\pgfqpoint{0.900481in}{3.464016in}}%
\pgfpathlineto{\pgfqpoint{0.903119in}{3.458881in}}%
\pgfpathlineto{\pgfqpoint{0.909495in}{3.445063in}}%
\pgfpathlineto{\pgfqpoint{0.915798in}{3.433338in}}%
\pgfpathlineto{\pgfqpoint{0.925452in}{3.412079in}}%
\pgfpathlineto{\pgfqpoint{0.933108in}{3.398897in}}%
\pgfpathlineto{\pgfqpoint{0.941759in}{3.380862in}}%
\pgfpathlineto{\pgfqpoint{0.950067in}{3.367753in}}%
\pgfpathlineto{\pgfqpoint{0.958164in}{3.351765in}}%
\pgfpathlineto{\pgfqpoint{0.962129in}{3.347890in}}%
\pgfpathlineto{\pgfqpoint{0.964767in}{3.343738in}}%
\pgfpathlineto{\pgfqpoint{0.974953in}{3.324395in}}%
\pgfpathlineto{\pgfqpoint{0.982798in}{3.314766in}}%
\pgfpathlineto{\pgfqpoint{0.990712in}{3.300245in}}%
\pgfpathlineto{\pgfqpoint{0.993874in}{3.297802in}}%
\pgfpathlineto{\pgfqpoint{0.996088in}{3.296200in}}%
\pgfpathlineto{\pgfqpoint{0.998726in}{3.292217in}}%
\pgfpathlineto{\pgfqpoint{1.007246in}{3.277183in}}%
\pgfpathlineto{\pgfqpoint{1.010175in}{3.275243in}}%
\pgfpathlineto{\pgfqpoint{1.012738in}{3.273783in}}%
\pgfpathlineto{\pgfqpoint{1.015376in}{3.270142in}}%
\pgfpathlineto{\pgfqpoint{1.024134in}{3.255709in}}%
\pgfpathlineto{\pgfqpoint{1.026951in}{3.254428in}}%
\pgfpathlineto{\pgfqpoint{1.029058in}{3.253501in}}%
\pgfpathlineto{\pgfqpoint{1.031531in}{3.250687in}}%
\pgfpathlineto{\pgfqpoint{1.036891in}{3.241053in}}%
\pgfpathlineto{\pgfqpoint{1.040263in}{3.236679in}}%
\pgfpathlineto{\pgfqpoint{1.042901in}{3.235443in}}%
\pgfpathlineto{\pgfqpoint{1.046037in}{3.234426in}}%
\pgfpathlineto{\pgfqpoint{1.048675in}{3.231374in}}%
\pgfpathlineto{\pgfqpoint{1.057884in}{3.218330in}}%
\pgfpathlineto{\pgfqpoint{1.060548in}{3.218065in}}%
\pgfpathlineto{\pgfqpoint{1.062856in}{3.217141in}}%
\pgfpathlineto{\pgfqpoint{1.065493in}{3.214212in}}%
\pgfpathlineto{\pgfqpoint{1.074289in}{3.202329in}}%
\pgfpathlineto{\pgfqpoint{1.076712in}{3.202265in}}%
\pgfpathlineto{\pgfqpoint{1.079020in}{3.201814in}}%
\pgfpathlineto{\pgfqpoint{1.081493in}{3.199674in}}%
\pgfpathlineto{\pgfqpoint{1.085079in}{3.194051in}}%
\pgfpathlineto{\pgfqpoint{1.089169in}{3.188591in}}%
\pgfpathlineto{\pgfqpoint{1.091786in}{3.187617in}}%
\pgfpathlineto{\pgfqpoint{1.096496in}{3.187140in}}%
\pgfpathlineto{\pgfqpoint{1.099134in}{3.184399in}}%
\pgfpathlineto{\pgfqpoint{1.106975in}{3.174606in}}%
\pgfpathlineto{\pgfqpoint{1.109340in}{3.174746in}}%
\pgfpathlineto{\pgfqpoint{1.111678in}{3.174954in}}%
\pgfpathlineto{\pgfqpoint{1.114151in}{3.173579in}}%
\pgfpathlineto{\pgfqpoint{1.117020in}{3.169957in}}%
\pgfpathlineto{\pgfqpoint{1.122056in}{3.163279in}}%
\pgfpathlineto{\pgfqpoint{1.124529in}{3.162562in}}%
\pgfpathlineto{\pgfqpoint{1.130104in}{3.162568in}}%
\pgfpathlineto{\pgfqpoint{1.132741in}{3.159956in}}%
\pgfpathlineto{\pgfqpoint{1.139641in}{3.151901in}}%
\pgfpathlineto{\pgfqpoint{1.141949in}{3.152106in}}%
\pgfpathlineto{\pgfqpoint{1.145109in}{3.152798in}}%
\pgfpathlineto{\pgfqpoint{1.147581in}{3.151570in}}%
\pgfpathlineto{\pgfqpoint{1.150384in}{3.148299in}}%
\pgfpathlineto{\pgfqpoint{1.154922in}{3.142652in}}%
\pgfpathlineto{\pgfqpoint{1.157291in}{3.142095in}}%
\pgfpathlineto{\pgfqpoint{1.163906in}{3.142449in}}%
\pgfpathlineto{\pgfqpoint{1.166708in}{3.139543in}}%
\pgfpathlineto{\pgfqpoint{1.171731in}{3.133661in}}%
\pgfpathlineto{\pgfqpoint{1.174039in}{3.133443in}}%
\pgfpathlineto{\pgfqpoint{1.179897in}{3.134429in}}%
\pgfpathlineto{\pgfqpoint{1.182535in}{3.132282in}}%
\pgfpathlineto{\pgfqpoint{1.189378in}{3.125523in}}%
\pgfpathlineto{\pgfqpoint{1.191635in}{3.126251in}}%
\pgfpathlineto{\pgfqpoint{1.194239in}{3.127361in}}%
\pgfpathlineto{\pgfqpoint{1.196712in}{3.126875in}}%
\pgfpathlineto{\pgfqpoint{1.199349in}{3.124716in}}%
\pgfpathlineto{\pgfqpoint{1.205613in}{3.118616in}}%
\pgfpathlineto{\pgfqpoint{1.207756in}{3.119179in}}%
\pgfpathlineto{\pgfqpoint{1.211548in}{3.120711in}}%
\pgfpathlineto{\pgfqpoint{1.214021in}{3.119874in}}%
\pgfpathlineto{\pgfqpoint{1.216823in}{3.117155in}}%
\pgfpathlineto{\pgfqpoint{1.221185in}{3.112624in}}%
\pgfpathlineto{\pgfqpoint{1.223328in}{3.112595in}}%
\pgfpathlineto{\pgfqpoint{1.230009in}{3.114252in}}%
\pgfpathlineto{\pgfqpoint{1.232647in}{3.112190in}}%
\pgfpathlineto{\pgfqpoint{1.238445in}{3.106821in}}%
\pgfpathlineto{\pgfqpoint{1.240503in}{3.107358in}}%
\pgfpathlineto{\pgfqpoint{1.245012in}{3.109311in}}%
\pgfpathlineto{\pgfqpoint{1.247649in}{3.108302in}}%
\pgfpathlineto{\pgfqpoint{1.250448in}{3.105455in}}%
\pgfpathlineto{\pgfqpoint{1.253887in}{3.102070in}}%
\pgfpathlineto{\pgfqpoint{1.256030in}{3.101963in}}%
\pgfpathlineto{\pgfqpoint{1.258614in}{3.103645in}}%
\pgfpathlineto{\pgfqpoint{1.261416in}{3.104535in}}%
\pgfpathlineto{\pgfqpoint{1.264054in}{3.103691in}}%
\pgfpathlineto{\pgfqpoint{1.266856in}{3.101031in}}%
\pgfpathlineto{\pgfqpoint{1.270426in}{3.097589in}}%
\pgfpathlineto{\pgfqpoint{1.272548in}{3.097556in}}%
\pgfpathlineto{\pgfqpoint{1.275162in}{3.099329in}}%
\pgfpathlineto{\pgfqpoint{1.277964in}{3.100254in}}%
\pgfpathlineto{\pgfqpoint{1.280602in}{3.099460in}}%
\pgfpathlineto{\pgfqpoint{1.283404in}{3.096866in}}%
\pgfpathlineto{\pgfqpoint{1.286887in}{3.093635in}}%
\pgfpathlineto{\pgfqpoint{1.289030in}{3.093676in}}%
\pgfpathlineto{\pgfqpoint{1.291664in}{3.095576in}}%
\pgfpathlineto{\pgfqpoint{1.294466in}{3.096628in}}%
\pgfpathlineto{\pgfqpoint{1.297104in}{3.095964in}}%
\pgfpathlineto{\pgfqpoint{1.299906in}{3.093521in}}%
\pgfpathlineto{\pgfqpoint{1.303281in}{3.090546in}}%
\pgfpathlineto{\pgfqpoint{1.305424in}{3.090654in}}%
\pgfpathlineto{\pgfqpoint{1.307931in}{3.092598in}}%
\pgfpathlineto{\pgfqpoint{1.310898in}{3.094020in}}%
\pgfpathlineto{\pgfqpoint{1.313536in}{3.093588in}}%
\pgfpathlineto{\pgfqpoint{1.316318in}{3.091421in}}%
\pgfpathlineto{\pgfqpoint{1.319698in}{3.088850in}}%
\pgfpathlineto{\pgfqpoint{1.321841in}{3.089212in}}%
\pgfpathlineto{\pgfqpoint{1.324308in}{3.091442in}}%
\pgfpathlineto{\pgfqpoint{1.327440in}{3.093395in}}%
\pgfpathlineto{\pgfqpoint{1.330078in}{3.093293in}}%
\pgfpathlineto{\pgfqpoint{1.332715in}{3.091608in}}%
\pgfpathlineto{\pgfqpoint{1.336237in}{3.089192in}}%
\pgfpathlineto{\pgfqpoint{1.338380in}{3.089739in}}%
\pgfpathlineto{\pgfqpoint{1.341025in}{3.092340in}}%
\pgfpathlineto{\pgfqpoint{1.344157in}{3.094335in}}%
\pgfpathlineto{\pgfqpoint{1.346794in}{3.094264in}}%
\pgfpathlineto{\pgfqpoint{1.349486in}{3.092555in}}%
\pgfpathlineto{\pgfqpoint{1.352578in}{3.090624in}}%
\pgfpathlineto{\pgfqpoint{1.354721in}{3.091189in}}%
\pgfpathlineto{\pgfqpoint{1.357276in}{3.093777in}}%
\pgfpathlineto{\pgfqpoint{1.360408in}{3.096103in}}%
\pgfpathlineto{\pgfqpoint{1.363045in}{3.096300in}}%
\pgfpathlineto{\pgfqpoint{1.365683in}{3.094902in}}%
\pgfpathlineto{\pgfqpoint{1.369209in}{3.092845in}}%
\pgfpathlineto{\pgfqpoint{1.371352in}{3.093599in}}%
\pgfpathlineto{\pgfqpoint{1.374254in}{3.096711in}}%
\pgfpathlineto{\pgfqpoint{1.377386in}{3.098776in}}%
\pgfpathlineto{\pgfqpoint{1.380024in}{3.098745in}}%
\pgfpathlineto{\pgfqpoint{1.382766in}{3.097015in}}%
\pgfpathlineto{\pgfqpoint{1.385398in}{3.095567in}}%
\pgfpathlineto{\pgfqpoint{1.387541in}{3.096140in}}%
\pgfpathlineto{\pgfqpoint{1.389914in}{3.098557in}}%
\pgfpathlineto{\pgfqpoint{1.393046in}{3.101265in}}%
\pgfpathlineto{\pgfqpoint{1.395684in}{3.101762in}}%
\pgfpathlineto{\pgfqpoint{1.398321in}{3.100643in}}%
\pgfpathlineto{\pgfqpoint{1.402550in}{3.098426in}}%
\pgfpathlineto{\pgfqpoint{1.404693in}{3.099478in}}%
\pgfpathlineto{\pgfqpoint{1.411476in}{3.104535in}}%
\pgfpathlineto{\pgfqpoint{1.414113in}{3.103783in}}%
\pgfpathlineto{\pgfqpoint{1.419838in}{3.101178in}}%
\pgfpathlineto{\pgfqpoint{1.422126in}{3.102907in}}%
\pgfpathlineto{\pgfqpoint{1.426085in}{3.106421in}}%
\pgfpathlineto{\pgfqpoint{1.428723in}{3.106734in}}%
\pgfpathlineto{\pgfqpoint{1.431360in}{3.105414in}}%
\pgfpathlineto{\pgfqpoint{1.435487in}{3.102969in}}%
\pgfpathlineto{\pgfqpoint{1.437643in}{3.103849in}}%
\pgfpathlineto{\pgfqpoint{1.444480in}{3.108372in}}%
\pgfpathlineto{\pgfqpoint{1.447117in}{3.107366in}}%
\pgfpathlineto{\pgfqpoint{1.452821in}{3.104250in}}%
\pgfpathlineto{\pgfqpoint{1.454964in}{3.105615in}}%
\pgfpathlineto{\pgfqpoint{1.459151in}{3.108884in}}%
\pgfpathlineto{\pgfqpoint{1.461788in}{3.108856in}}%
\pgfpathlineto{\pgfqpoint{1.464426in}{3.107191in}}%
\pgfpathlineto{\pgfqpoint{1.468358in}{3.104389in}}%
\pgfpathlineto{\pgfqpoint{1.470501in}{3.104951in}}%
\pgfpathlineto{\pgfqpoint{1.477650in}{3.108750in}}%
\pgfpathlineto{\pgfqpoint{1.480288in}{3.107331in}}%
\pgfpathlineto{\pgfqpoint{1.485448in}{3.103741in}}%
\pgfpathlineto{\pgfqpoint{1.487591in}{3.104640in}}%
\pgfpathlineto{\pgfqpoint{1.492779in}{3.107860in}}%
\pgfpathlineto{\pgfqpoint{1.495417in}{3.107167in}}%
\pgfpathlineto{\pgfqpoint{1.498278in}{3.104579in}}%
\pgfpathlineto{\pgfqpoint{1.500677in}{3.102616in}}%
\pgfpathlineto{\pgfqpoint{1.502964in}{3.102569in}}%
\pgfpathlineto{\pgfqpoint{1.505310in}{3.104315in}}%
\pgfpathlineto{\pgfqpoint{1.508277in}{3.106086in}}%
\pgfpathlineto{\pgfqpoint{1.510914in}{3.105916in}}%
\pgfpathlineto{\pgfqpoint{1.513552in}{3.104122in}}%
\pgfpathlineto{\pgfqpoint{1.518150in}{3.100466in}}%
\pgfpathlineto{\pgfqpoint{1.520294in}{3.101034in}}%
\pgfpathlineto{\pgfqpoint{1.526307in}{3.104097in}}%
\pgfpathlineto{\pgfqpoint{1.528945in}{3.102931in}}%
\pgfpathlineto{\pgfqpoint{1.536033in}{3.098437in}}%
\pgfpathlineto{\pgfqpoint{1.538373in}{3.100160in}}%
\pgfpathlineto{\pgfqpoint{1.541341in}{3.101789in}}%
\pgfpathlineto{\pgfqpoint{1.543978in}{3.101510in}}%
\pgfpathlineto{\pgfqpoint{1.546616in}{3.099622in}}%
\pgfpathlineto{\pgfqpoint{1.550984in}{3.096169in}}%
\pgfpathlineto{\pgfqpoint{1.553106in}{3.096698in}}%
\pgfpathlineto{\pgfqpoint{1.559641in}{3.100074in}}%
\pgfpathlineto{\pgfqpoint{1.562278in}{3.098822in}}%
\pgfpathlineto{\pgfqpoint{1.568346in}{3.094887in}}%
\pgfpathlineto{\pgfqpoint{1.570543in}{3.096200in}}%
\pgfpathlineto{\pgfqpoint{1.574380in}{3.098930in}}%
\pgfpathlineto{\pgfqpoint{1.577018in}{3.098845in}}%
\pgfpathlineto{\pgfqpoint{1.579655in}{3.097157in}}%
\pgfpathlineto{\pgfqpoint{1.583960in}{3.094065in}}%
\pgfpathlineto{\pgfqpoint{1.586164in}{3.094793in}}%
\pgfpathlineto{\pgfqpoint{1.592681in}{3.098552in}}%
\pgfpathlineto{\pgfqpoint{1.595318in}{3.097448in}}%
\pgfpathlineto{\pgfqpoint{1.601255in}{3.093936in}}%
\pgfpathlineto{\pgfqpoint{1.603398in}{3.095288in}}%
\pgfpathlineto{\pgfqpoint{1.607513in}{3.098429in}}%
\pgfpathlineto{\pgfqpoint{1.610151in}{3.098411in}}%
\pgfpathlineto{\pgfqpoint{1.612788in}{3.096793in}}%
\pgfpathlineto{\pgfqpoint{1.616722in}{3.094123in}}%
\pgfpathlineto{\pgfqpoint{1.618865in}{3.094789in}}%
\pgfpathlineto{\pgfqpoint{1.626142in}{3.099065in}}%
\pgfpathlineto{\pgfqpoint{1.628780in}{3.097810in}}%
\pgfpathlineto{\pgfqpoint{1.633694in}{3.094842in}}%
\pgfpathlineto{\pgfqpoint{1.635837in}{3.095911in}}%
\pgfpathlineto{\pgfqpoint{1.641473in}{3.099974in}}%
\pgfpathlineto{\pgfqpoint{1.644111in}{3.099467in}}%
\pgfpathlineto{\pgfqpoint{1.646993in}{3.097121in}}%
\pgfpathlineto{\pgfqpoint{1.649459in}{3.095830in}}%
\pgfpathlineto{\pgfqpoint{1.651602in}{3.096401in}}%
\pgfpathlineto{\pgfqpoint{1.654331in}{3.099087in}}%
\pgfpathlineto{\pgfqpoint{1.657299in}{3.100936in}}%
\pgfpathlineto{\pgfqpoint{1.659936in}{3.100855in}}%
\pgfpathlineto{\pgfqpoint{1.662574in}{3.099165in}}%
\pgfpathlineto{\pgfqpoint{1.665823in}{3.096993in}}%
\pgfpathlineto{\pgfqpoint{1.667966in}{3.097486in}}%
\pgfpathlineto{\pgfqpoint{1.670469in}{3.099901in}}%
\pgfpathlineto{\pgfqpoint{1.673601in}{3.102050in}}%
\pgfpathlineto{\pgfqpoint{1.676238in}{3.102078in}}%
\pgfpathlineto{\pgfqpoint{1.678876in}{3.100492in}}%
\pgfpathlineto{\pgfqpoint{1.682529in}{3.098115in}}%
\pgfpathlineto{\pgfqpoint{1.684652in}{3.098743in}}%
\pgfpathlineto{\pgfqpoint{1.687796in}{3.101816in}}%
\pgfpathlineto{\pgfqpoint{1.690763in}{3.103299in}}%
\pgfpathlineto{\pgfqpoint{1.693401in}{3.102886in}}%
\pgfpathlineto{\pgfqpoint{1.696211in}{3.100676in}}%
\pgfpathlineto{\pgfqpoint{1.698642in}{3.099166in}}%
\pgfpathlineto{\pgfqpoint{1.700785in}{3.099509in}}%
\pgfpathlineto{\pgfqpoint{1.703126in}{3.101621in}}%
\pgfpathlineto{\pgfqpoint{1.706258in}{3.103929in}}%
\pgfpathlineto{\pgfqpoint{1.708896in}{3.104083in}}%
\pgfpathlineto{\pgfqpoint{1.711534in}{3.102615in}}%
\pgfpathlineto{\pgfqpoint{1.715786in}{3.099834in}}%
\pgfpathlineto{\pgfqpoint{1.717929in}{3.100615in}}%
\pgfpathlineto{\pgfqpoint{1.724579in}{3.104769in}}%
\pgfpathlineto{\pgfqpoint{1.727217in}{3.103740in}}%
\pgfpathlineto{\pgfqpoint{1.733071in}{3.100338in}}%
\pgfpathlineto{\pgfqpoint{1.735214in}{3.101657in}}%
\pgfpathlineto{\pgfqpoint{1.739491in}{3.104872in}}%
\pgfpathlineto{\pgfqpoint{1.742129in}{3.104756in}}%
\pgfpathlineto{\pgfqpoint{1.744766in}{3.103019in}}%
\pgfpathlineto{\pgfqpoint{1.748573in}{3.100347in}}%
\pgfpathlineto{\pgfqpoint{1.750695in}{3.100907in}}%
\pgfpathlineto{\pgfqpoint{1.758004in}{3.104866in}}%
\pgfpathlineto{\pgfqpoint{1.760642in}{3.103457in}}%
\pgfpathlineto{\pgfqpoint{1.765576in}{3.100142in}}%
\pgfpathlineto{\pgfqpoint{1.767781in}{3.101101in}}%
\pgfpathlineto{\pgfqpoint{1.773188in}{3.104633in}}%
\pgfpathlineto{\pgfqpoint{1.775825in}{3.103982in}}%
\pgfpathlineto{\pgfqpoint{1.778665in}{3.101486in}}%
\pgfpathlineto{\pgfqpoint{1.781245in}{3.099741in}}%
\pgfpathlineto{\pgfqpoint{1.783388in}{3.100020in}}%
\pgfpathlineto{\pgfqpoint{1.785877in}{3.102155in}}%
\pgfpathlineto{\pgfqpoint{1.788844in}{3.103896in}}%
\pgfpathlineto{\pgfqpoint{1.791482in}{3.103709in}}%
\pgfpathlineto{\pgfqpoint{1.794119in}{3.101904in}}%
\pgfpathlineto{\pgfqpoint{1.798044in}{3.098920in}}%
\pgfpathlineto{\pgfqpoint{1.800187in}{3.099394in}}%
\pgfpathlineto{\pgfqpoint{1.807490in}{3.102997in}}%
\pgfpathlineto{\pgfqpoint{1.810128in}{3.101467in}}%
\pgfpathlineto{\pgfqpoint{1.815049in}{3.098021in}}%
\pgfpathlineto{\pgfqpoint{1.817171in}{3.098868in}}%
\pgfpathlineto{\pgfqpoint{1.822761in}{3.102351in}}%
\pgfpathlineto{\pgfqpoint{1.825398in}{3.101594in}}%
\pgfpathlineto{\pgfqpoint{1.828350in}{3.098889in}}%
\pgfpathlineto{\pgfqpoint{1.830802in}{3.097364in}}%
\pgfpathlineto{\pgfqpoint{1.833017in}{3.097756in}}%
\pgfpathlineto{\pgfqpoint{1.835600in}{3.100052in}}%
\pgfpathlineto{\pgfqpoint{1.838567in}{3.101633in}}%
\pgfpathlineto{\pgfqpoint{1.841204in}{3.101312in}}%
\pgfpathlineto{\pgfqpoint{1.843934in}{3.099291in}}%
\pgfpathlineto{\pgfqpoint{1.847256in}{3.096822in}}%
\pgfpathlineto{\pgfqpoint{1.849399in}{3.097189in}}%
\pgfpathlineto{\pgfqpoint{1.852057in}{3.099561in}}%
\pgfpathlineto{\pgfqpoint{1.855024in}{3.101212in}}%
\pgfpathlineto{\pgfqpoint{1.857661in}{3.100955in}}%
\pgfpathlineto{\pgfqpoint{1.860299in}{3.099090in}}%
\pgfpathlineto{\pgfqpoint{1.863902in}{3.096487in}}%
\pgfpathlineto{\pgfqpoint{1.866045in}{3.096979in}}%
\pgfpathlineto{\pgfqpoint{1.869479in}{3.100024in}}%
\pgfpathlineto{\pgfqpoint{1.872282in}{3.101041in}}%
\pgfpathlineto{\pgfqpoint{1.874919in}{3.100329in}}%
\pgfpathlineto{\pgfqpoint{1.877843in}{3.097738in}}%
\pgfpathlineto{\pgfqpoint{1.880315in}{3.096379in}}%
\pgfpathlineto{\pgfqpoint{1.882458in}{3.096896in}}%
\pgfpathlineto{\pgfqpoint{1.885457in}{3.099706in}}%
\pgfpathlineto{\pgfqpoint{1.888424in}{3.101176in}}%
\pgfpathlineto{\pgfqpoint{1.891061in}{3.100760in}}%
\pgfpathlineto{\pgfqpoint{1.893787in}{3.098644in}}%
\pgfpathlineto{\pgfqpoint{1.896477in}{3.096789in}}%
\pgfpathlineto{\pgfqpoint{1.898620in}{3.097084in}}%
\pgfpathlineto{\pgfqpoint{1.901042in}{3.099233in}}%
\pgfpathlineto{\pgfqpoint{1.904174in}{3.101404in}}%
\pgfpathlineto{\pgfqpoint{1.906811in}{3.101453in}}%
\pgfpathlineto{\pgfqpoint{1.909449in}{3.099890in}}%
\pgfpathlineto{\pgfqpoint{1.913460in}{3.097179in}}%
\pgfpathlineto{\pgfqpoint{1.915603in}{3.097855in}}%
\pgfpathlineto{\pgfqpoint{1.922744in}{3.102106in}}%
\pgfpathlineto{\pgfqpoint{1.925382in}{3.100894in}}%
\pgfpathlineto{\pgfqpoint{1.930553in}{3.097798in}}%
\pgfpathlineto{\pgfqpoint{1.932758in}{3.098990in}}%
\pgfpathlineto{\pgfqpoint{1.937872in}{3.102737in}}%
\pgfpathlineto{\pgfqpoint{1.940509in}{3.102369in}}%
\pgfpathlineto{\pgfqpoint{1.943274in}{3.100254in}}%
\pgfpathlineto{\pgfqpoint{1.945868in}{3.098479in}}%
\pgfpathlineto{\pgfqpoint{1.947990in}{3.098738in}}%
\pgfpathlineto{\pgfqpoint{1.950252in}{3.100680in}}%
\pgfpathlineto{\pgfqpoint{1.953366in}{3.103038in}}%
\pgfpathlineto{\pgfqpoint{1.956004in}{3.103254in}}%
\pgfpathlineto{\pgfqpoint{1.958642in}{3.101851in}}%
\pgfpathlineto{\pgfqpoint{1.963231in}{3.098881in}}%
\pgfpathlineto{\pgfqpoint{1.965374in}{3.099765in}}%
\pgfpathlineto{\pgfqpoint{1.971500in}{3.103745in}}%
\pgfpathlineto{\pgfqpoint{1.974137in}{3.102913in}}%
\pgfpathlineto{\pgfqpoint{1.980871in}{3.099421in}}%
\pgfpathlineto{\pgfqpoint{1.983151in}{3.101266in}}%
\pgfpathlineto{\pgfqpoint{1.986183in}{3.103563in}}%
\pgfpathlineto{\pgfqpoint{1.988821in}{3.103802in}}%
\pgfpathlineto{\pgfqpoint{1.991458in}{3.102421in}}%
\pgfpathlineto{\pgfqpoint{1.996362in}{3.099203in}}%
\pgfpathlineto{\pgfqpoint{1.998567in}{3.100195in}}%
\pgfpathlineto{\pgfqpoint{2.003980in}{3.103815in}}%
\pgfpathlineto{\pgfqpoint{2.006618in}{3.103207in}}%
\pgfpathlineto{\pgfqpoint{2.009462in}{3.100754in}}%
\pgfpathlineto{\pgfqpoint{2.011966in}{3.099085in}}%
\pgfpathlineto{\pgfqpoint{2.014109in}{3.099359in}}%
\pgfpathlineto{\pgfqpoint{2.016502in}{3.101422in}}%
\pgfpathlineto{\pgfqpoint{2.019469in}{3.103349in}}%
\pgfpathlineto{\pgfqpoint{2.022107in}{3.103327in}}%
\pgfpathlineto{\pgfqpoint{2.024745in}{3.101687in}}%
\pgfpathlineto{\pgfqpoint{2.029034in}{3.098636in}}%
\pgfpathlineto{\pgfqpoint{2.031156in}{3.099300in}}%
\pgfpathlineto{\pgfqpoint{2.037768in}{3.103063in}}%
\pgfpathlineto{\pgfqpoint{2.040406in}{3.101907in}}%
\pgfpathlineto{\pgfqpoint{2.046361in}{3.098215in}}%
\pgfpathlineto{\pgfqpoint{2.048484in}{3.099473in}}%
\pgfpathlineto{\pgfqpoint{2.052606in}{3.102433in}}%
\pgfpathlineto{\pgfqpoint{2.055243in}{3.102279in}}%
\pgfpathlineto{\pgfqpoint{2.057881in}{3.100512in}}%
\pgfpathlineto{\pgfqpoint{2.061822in}{3.097601in}}%
\pgfpathlineto{\pgfqpoint{2.063965in}{3.098133in}}%
\pgfpathlineto{\pgfqpoint{2.071233in}{3.101878in}}%
\pgfpathlineto{\pgfqpoint{2.073870in}{3.100417in}}%
\pgfpathlineto{\pgfqpoint{2.078767in}{3.097049in}}%
\pgfpathlineto{\pgfqpoint{2.080910in}{3.097923in}}%
\pgfpathlineto{\pgfqpoint{2.086564in}{3.101511in}}%
\pgfpathlineto{\pgfqpoint{2.089201in}{3.100769in}}%
\pgfpathlineto{\pgfqpoint{2.092154in}{3.098092in}}%
\pgfpathlineto{\pgfqpoint{2.094627in}{3.096638in}}%
\pgfpathlineto{\pgfqpoint{2.096770in}{3.097072in}}%
\pgfpathlineto{\pgfqpoint{2.099653in}{3.099672in}}%
\pgfpathlineto{\pgfqpoint{2.102620in}{3.101109in}}%
\pgfpathlineto{\pgfqpoint{2.105258in}{3.100663in}}%
\pgfpathlineto{\pgfqpoint{2.108083in}{3.098411in}}%
\pgfpathlineto{\pgfqpoint{2.110779in}{3.096543in}}%
\pgfpathlineto{\pgfqpoint{2.112984in}{3.096865in}}%
\pgfpathlineto{\pgfqpoint{2.115247in}{3.098875in}}%
\pgfpathlineto{\pgfqpoint{2.118379in}{3.101112in}}%
\pgfpathlineto{\pgfqpoint{2.121017in}{3.101217in}}%
\pgfpathlineto{\pgfqpoint{2.123655in}{3.099711in}}%
\pgfpathlineto{\pgfqpoint{2.127830in}{3.096891in}}%
\pgfpathlineto{\pgfqpoint{2.130035in}{3.097650in}}%
\pgfpathlineto{\pgfqpoint{2.136939in}{3.101797in}}%
\pgfpathlineto{\pgfqpoint{2.139576in}{3.100642in}}%
\pgfpathlineto{\pgfqpoint{2.144868in}{3.097430in}}%
\pgfpathlineto{\pgfqpoint{2.147011in}{3.098560in}}%
\pgfpathlineto{\pgfqpoint{2.152177in}{3.102318in}}%
\pgfpathlineto{\pgfqpoint{2.154814in}{3.101934in}}%
\pgfpathlineto{\pgfqpoint{2.157590in}{3.099791in}}%
\pgfpathlineto{\pgfqpoint{2.160257in}{3.097999in}}%
\pgfpathlineto{\pgfqpoint{2.162400in}{3.098324in}}%
\pgfpathlineto{\pgfqpoint{2.164867in}{3.100543in}}%
\pgfpathlineto{\pgfqpoint{2.167999in}{3.102671in}}%
\pgfpathlineto{\pgfqpoint{2.170636in}{3.102679in}}%
\pgfpathlineto{\pgfqpoint{2.173274in}{3.101071in}}%
\pgfpathlineto{\pgfqpoint{2.177199in}{3.098392in}}%
\pgfpathlineto{\pgfqpoint{2.179342in}{3.099033in}}%
\pgfpathlineto{\pgfqpoint{2.186627in}{3.103180in}}%
\pgfpathlineto{\pgfqpoint{2.189265in}{3.101850in}}%
\pgfpathlineto{\pgfqpoint{2.194163in}{3.098708in}}%
\pgfpathlineto{\pgfqpoint{2.196304in}{3.099675in}}%
\pgfpathlineto{\pgfqpoint{2.201942in}{3.103497in}}%
\pgfpathlineto{\pgfqpoint{2.204579in}{3.102864in}}%
\pgfpathlineto{\pgfqpoint{2.207416in}{3.100422in}}%
\pgfpathlineto{\pgfqpoint{2.209852in}{3.098969in}}%
\pgfpathlineto{\pgfqpoint{2.211995in}{3.099367in}}%
\pgfpathlineto{\pgfqpoint{2.214597in}{3.101757in}}%
\pgfpathlineto{\pgfqpoint{2.217564in}{3.103605in}}%
\pgfpathlineto{\pgfqpoint{2.220202in}{3.103513in}}%
\pgfpathlineto{\pgfqpoint{2.222839in}{3.101803in}}%
\pgfpathlineto{\pgfqpoint{2.226604in}{3.099180in}}%
\pgfpathlineto{\pgfqpoint{2.228806in}{3.099792in}}%
\pgfpathlineto{\pgfqpoint{2.236083in}{3.103818in}}%
\pgfpathlineto{\pgfqpoint{2.238720in}{3.102442in}}%
\pgfpathlineto{\pgfqpoint{2.243671in}{3.099218in}}%
\pgfpathlineto{\pgfqpoint{2.245814in}{3.100185in}}%
\pgfpathlineto{\pgfqpoint{2.251272in}{3.103833in}}%
\pgfpathlineto{\pgfqpoint{2.253909in}{3.103216in}}%
\pgfpathlineto{\pgfqpoint{2.256764in}{3.100744in}}%
\pgfpathlineto{\pgfqpoint{2.259309in}{3.099093in}}%
\pgfpathlineto{\pgfqpoint{2.261514in}{3.099450in}}%
\pgfpathlineto{\pgfqpoint{2.263940in}{3.101609in}}%
\pgfpathlineto{\pgfqpoint{2.266908in}{3.103420in}}%
\pgfpathlineto{\pgfqpoint{2.269545in}{3.103296in}}%
\pgfpathlineto{\pgfqpoint{2.272183in}{3.101556in}}%
\pgfpathlineto{\pgfqpoint{2.276113in}{3.098677in}}%
\pgfpathlineto{\pgfqpoint{2.278256in}{3.099213in}}%
\pgfpathlineto{\pgfqpoint{2.285548in}{3.102991in}}%
\pgfpathlineto{\pgfqpoint{2.288186in}{3.101526in}}%
\pgfpathlineto{\pgfqpoint{2.293124in}{3.098165in}}%
\pgfpathlineto{\pgfqpoint{2.295268in}{3.099073in}}%
\pgfpathlineto{\pgfqpoint{2.300703in}{3.102582in}}%
\pgfpathlineto{\pgfqpoint{2.303340in}{3.101921in}}%
\pgfpathlineto{\pgfqpoint{2.306220in}{3.099376in}}%
\pgfpathlineto{\pgfqpoint{2.308762in}{3.097679in}}%
\pgfpathlineto{\pgfqpoint{2.310905in}{3.097968in}}%
\pgfpathlineto{\pgfqpoint{2.313397in}{3.100124in}}%
\pgfpathlineto{\pgfqpoint{2.316364in}{3.101896in}}%
\pgfpathlineto{\pgfqpoint{2.319002in}{3.101743in}}%
\pgfpathlineto{\pgfqpoint{2.321639in}{3.099979in}}%
\pgfpathlineto{\pgfqpoint{2.325559in}{3.097081in}}%
\pgfpathlineto{\pgfqpoint{2.327702in}{3.097604in}}%
\pgfpathlineto{\pgfqpoint{2.334990in}{3.101394in}}%
\pgfpathlineto{\pgfqpoint{2.337628in}{3.099946in}}%
\pgfpathlineto{\pgfqpoint{2.342559in}{3.096609in}}%
\pgfpathlineto{\pgfqpoint{2.344681in}{3.097507in}}%
\pgfpathlineto{\pgfqpoint{2.350345in}{3.101134in}}%
\pgfpathlineto{\pgfqpoint{2.352983in}{3.100404in}}%
\pgfpathlineto{\pgfqpoint{2.355915in}{3.097769in}}%
\pgfpathlineto{\pgfqpoint{2.358388in}{3.096354in}}%
\pgfpathlineto{\pgfqpoint{2.360510in}{3.096810in}}%
\pgfpathlineto{\pgfqpoint{2.363344in}{3.099400in}}%
\pgfpathlineto{\pgfqpoint{2.366311in}{3.100913in}}%
\pgfpathlineto{\pgfqpoint{2.368949in}{3.100536in}}%
\pgfpathlineto{\pgfqpoint{2.371728in}{3.098402in}}%
\pgfpathlineto{\pgfqpoint{2.374570in}{3.096342in}}%
\pgfpathlineto{\pgfqpoint{2.376713in}{3.096611in}}%
\pgfpathlineto{\pgfqpoint{2.379005in}{3.098587in}}%
\pgfpathlineto{\pgfqpoint{2.382137in}{3.100760in}}%
\pgfpathlineto{\pgfqpoint{2.384775in}{3.100812in}}%
\pgfpathlineto{\pgfqpoint{2.387412in}{3.099254in}}%
\pgfpathlineto{\pgfqpoint{2.391675in}{3.096356in}}%
\pgfpathlineto{\pgfqpoint{2.393797in}{3.097082in}}%
\pgfpathlineto{\pgfqpoint{2.400624in}{3.101133in}}%
\pgfpathlineto{\pgfqpoint{2.403261in}{3.099988in}}%
\pgfpathlineto{\pgfqpoint{2.408816in}{3.096685in}}%
\pgfpathlineto{\pgfqpoint{2.411023in}{3.097986in}}%
\pgfpathlineto{\pgfqpoint{2.415453in}{3.101353in}}%
\pgfpathlineto{\pgfqpoint{2.418091in}{3.101238in}}%
\pgfpathlineto{\pgfqpoint{2.420728in}{3.099514in}}%
\pgfpathlineto{\pgfqpoint{2.424276in}{3.097022in}}%
\pgfpathlineto{\pgfqpoint{2.426419in}{3.097524in}}%
\pgfpathlineto{\pgfqpoint{2.429135in}{3.100092in}}%
\pgfpathlineto{\pgfqpoint{2.432102in}{3.101807in}}%
\pgfpathlineto{\pgfqpoint{2.434740in}{3.101604in}}%
\pgfpathlineto{\pgfqpoint{2.437509in}{3.099659in}}%
\pgfpathlineto{\pgfqpoint{2.440509in}{3.097549in}}%
\pgfpathlineto{\pgfqpoint{2.442652in}{3.097898in}}%
\pgfpathlineto{\pgfqpoint{2.445126in}{3.100155in}}%
\pgfpathlineto{\pgfqpoint{2.448258in}{3.102310in}}%
\pgfpathlineto{\pgfqpoint{2.450896in}{3.102341in}}%
\pgfpathlineto{\pgfqpoint{2.453534in}{3.100757in}}%
\pgfpathlineto{\pgfqpoint{2.457339in}{3.098232in}}%
\pgfpathlineto{\pgfqpoint{2.459461in}{3.098861in}}%
\pgfpathlineto{\pgfqpoint{2.462815in}{3.102089in}}%
\pgfpathlineto{\pgfqpoint{2.465617in}{3.103377in}}%
\pgfpathlineto{\pgfqpoint{2.468255in}{3.102910in}}%
\pgfpathlineto{\pgfqpoint{2.471061in}{3.100653in}}%
\pgfpathlineto{\pgfqpoint{2.473590in}{3.099166in}}%
\pgfpathlineto{\pgfqpoint{2.475712in}{3.099602in}}%
\pgfpathlineto{\pgfqpoint{2.478244in}{3.101978in}}%
\pgfpathlineto{\pgfqpoint{2.481211in}{3.103977in}}%
\pgfpathlineto{\pgfqpoint{2.483849in}{3.104016in}}%
\pgfpathlineto{\pgfqpoint{2.486486in}{3.102435in}}%
\pgfpathlineto{\pgfqpoint{2.490407in}{3.099786in}}%
\pgfpathlineto{\pgfqpoint{2.492530in}{3.100422in}}%
\pgfpathlineto{\pgfqpoint{2.499839in}{3.104600in}}%
\pgfpathlineto{\pgfqpoint{2.502477in}{3.103267in}}%
\pgfpathlineto{\pgfqpoint{2.507357in}{3.100105in}}%
\pgfpathlineto{\pgfqpoint{2.509479in}{3.101030in}}%
\pgfpathlineto{\pgfqpoint{2.515172in}{3.104818in}}%
\pgfpathlineto{\pgfqpoint{2.517809in}{3.104135in}}%
\pgfpathlineto{\pgfqpoint{2.520747in}{3.101526in}}%
\pgfpathlineto{\pgfqpoint{2.523220in}{3.100095in}}%
\pgfpathlineto{\pgfqpoint{2.525363in}{3.100541in}}%
\pgfpathlineto{\pgfqpoint{2.528193in}{3.103102in}}%
\pgfpathlineto{\pgfqpoint{2.531160in}{3.104572in}}%
\pgfpathlineto{\pgfqpoint{2.533798in}{3.104143in}}%
\pgfpathlineto{\pgfqpoint{2.536592in}{3.101925in}}%
\pgfpathlineto{\pgfqpoint{2.539387in}{3.099809in}}%
\pgfpathlineto{\pgfqpoint{2.541530in}{3.099985in}}%
\pgfpathlineto{\pgfqpoint{2.543922in}{3.101962in}}%
\pgfpathlineto{\pgfqpoint{2.546889in}{3.103915in}}%
\pgfpathlineto{\pgfqpoint{2.549527in}{3.103913in}}%
\pgfpathlineto{\pgfqpoint{2.552164in}{3.102292in}}%
\pgfpathlineto{\pgfqpoint{2.556710in}{3.099079in}}%
\pgfpathlineto{\pgfqpoint{2.558853in}{3.099839in}}%
\pgfpathlineto{\pgfqpoint{2.564813in}{3.103421in}}%
\pgfpathlineto{\pgfqpoint{2.567451in}{3.102512in}}%
\pgfpathlineto{\pgfqpoint{2.570718in}{3.099419in}}%
\pgfpathlineto{\pgfqpoint{2.573190in}{3.098349in}}%
\pgfpathlineto{\pgfqpoint{2.575333in}{3.099108in}}%
\pgfpathlineto{\pgfqpoint{2.581379in}{3.102701in}}%
\pgfpathlineto{\pgfqpoint{2.584017in}{3.101751in}}%
\pgfpathlineto{\pgfqpoint{2.590909in}{3.097940in}}%
\pgfpathlineto{\pgfqpoint{2.593161in}{3.099726in}}%
\pgfpathlineto{\pgfqpoint{2.596275in}{3.101881in}}%
\pgfpathlineto{\pgfqpoint{2.598913in}{3.101931in}}%
\pgfpathlineto{\pgfqpoint{2.601550in}{3.100368in}}%
\pgfpathlineto{\pgfqpoint{2.606144in}{3.097135in}}%
\pgfpathlineto{\pgfqpoint{2.608287in}{3.097908in}}%
\pgfpathlineto{\pgfqpoint{2.614421in}{3.101585in}}%
\pgfpathlineto{\pgfqpoint{2.617058in}{3.100630in}}%
\pgfpathlineto{\pgfqpoint{2.623767in}{3.096888in}}%
\pgfpathlineto{\pgfqpoint{2.626058in}{3.098664in}}%
\pgfpathlineto{\pgfqpoint{2.629083in}{3.100875in}}%
\pgfpathlineto{\pgfqpoint{2.631721in}{3.101058in}}%
\pgfpathlineto{\pgfqpoint{2.634359in}{3.099628in}}%
\pgfpathlineto{\pgfqpoint{2.639281in}{3.096332in}}%
\pgfpathlineto{\pgfqpoint{2.641362in}{3.097212in}}%
\pgfpathlineto{\pgfqpoint{2.647049in}{3.100915in}}%
\pgfpathlineto{\pgfqpoint{2.649686in}{3.100218in}}%
\pgfpathlineto{\pgfqpoint{2.652553in}{3.097680in}}%
\pgfpathlineto{\pgfqpoint{2.655026in}{3.096198in}}%
\pgfpathlineto{\pgfqpoint{2.657148in}{3.096598in}}%
\pgfpathlineto{\pgfqpoint{2.659766in}{3.098978in}}%
\pgfpathlineto{\pgfqpoint{2.662733in}{3.100721in}}%
\pgfpathlineto{\pgfqpoint{2.665371in}{3.100547in}}%
\pgfpathlineto{\pgfqpoint{2.668008in}{3.098767in}}%
\pgfpathlineto{\pgfqpoint{2.671579in}{3.096219in}}%
\pgfpathlineto{\pgfqpoint{2.673722in}{3.096711in}}%
\pgfpathlineto{\pgfqpoint{2.676565in}{3.099361in}}%
\pgfpathlineto{\pgfqpoint{2.679533in}{3.100922in}}%
\pgfpathlineto{\pgfqpoint{2.682170in}{3.100586in}}%
\pgfpathlineto{\pgfqpoint{2.684954in}{3.098491in}}%
\pgfpathlineto{\pgfqpoint{2.687811in}{3.096490in}}%
\pgfpathlineto{\pgfqpoint{2.689933in}{3.096805in}}%
\pgfpathlineto{\pgfqpoint{2.692382in}{3.098990in}}%
\pgfpathlineto{\pgfqpoint{2.695514in}{3.101092in}}%
\pgfpathlineto{\pgfqpoint{2.698152in}{3.101084in}}%
\pgfpathlineto{\pgfqpoint{2.700789in}{3.099466in}}%
\pgfpathlineto{\pgfqpoint{2.704702in}{3.096783in}}%
\pgfpathlineto{\pgfqpoint{2.706845in}{3.097418in}}%
\pgfpathlineto{\pgfqpoint{2.714143in}{3.101607in}}%
\pgfpathlineto{\pgfqpoint{2.716780in}{3.100300in}}%
\pgfpathlineto{\pgfqpoint{2.721659in}{3.097212in}}%
\pgfpathlineto{\pgfqpoint{2.723782in}{3.098179in}}%
\pgfpathlineto{\pgfqpoint{2.729556in}{3.102096in}}%
\pgfpathlineto{\pgfqpoint{2.732193in}{3.101434in}}%
\pgfpathlineto{\pgfqpoint{2.735046in}{3.098960in}}%
\pgfpathlineto{\pgfqpoint{2.737519in}{3.097607in}}%
\pgfpathlineto{\pgfqpoint{2.739723in}{3.098164in}}%
\pgfpathlineto{\pgfqpoint{2.742683in}{3.100961in}}%
\pgfpathlineto{\pgfqpoint{2.745651in}{3.102427in}}%
\pgfpathlineto{\pgfqpoint{2.748288in}{3.102003in}}%
\pgfpathlineto{\pgfqpoint{2.751102in}{3.099777in}}%
\pgfpathlineto{\pgfqpoint{2.753708in}{3.098025in}}%
\pgfpathlineto{\pgfqpoint{2.755851in}{3.098327in}}%
\pgfpathlineto{\pgfqpoint{2.758160in}{3.100355in}}%
\pgfpathlineto{\pgfqpoint{2.761292in}{3.102574in}}%
\pgfpathlineto{\pgfqpoint{2.763929in}{3.102659in}}%
\pgfpathlineto{\pgfqpoint{2.766567in}{3.101127in}}%
\pgfpathlineto{\pgfqpoint{2.770796in}{3.098262in}}%
\pgfpathlineto{\pgfqpoint{2.773001in}{3.099034in}}%
\pgfpathlineto{\pgfqpoint{2.779608in}{3.103081in}}%
\pgfpathlineto{\pgfqpoint{2.782246in}{3.102034in}}%
\pgfpathlineto{\pgfqpoint{2.788244in}{3.098653in}}%
\pgfpathlineto{\pgfqpoint{2.790387in}{3.100076in}}%
\pgfpathlineto{\pgfqpoint{2.794258in}{3.103082in}}%
\pgfpathlineto{\pgfqpoint{2.796896in}{3.103143in}}%
\pgfpathlineto{\pgfqpoint{2.799533in}{3.101586in}}%
\pgfpathlineto{\pgfqpoint{2.803856in}{3.098669in}}%
\pgfpathlineto{\pgfqpoint{2.805999in}{3.099437in}}%
\pgfpathlineto{\pgfqpoint{2.812392in}{3.103373in}}%
\pgfpathlineto{\pgfqpoint{2.815029in}{3.102394in}}%
\pgfpathlineto{\pgfqpoint{2.821434in}{3.098839in}}%
\pgfpathlineto{\pgfqpoint{2.823711in}{3.100527in}}%
\pgfpathlineto{\pgfqpoint{2.826899in}{3.102964in}}%
\pgfpathlineto{\pgfqpoint{2.829536in}{3.103177in}}%
\pgfpathlineto{\pgfqpoint{2.832174in}{3.101770in}}%
\pgfpathlineto{\pgfqpoint{2.837100in}{3.098491in}}%
\pgfpathlineto{\pgfqpoint{2.839243in}{3.099421in}}%
\pgfpathlineto{\pgfqpoint{2.844867in}{3.103053in}}%
\pgfpathlineto{\pgfqpoint{2.847504in}{3.102334in}}%
\pgfpathlineto{\pgfqpoint{2.850375in}{3.099760in}}%
\pgfpathlineto{\pgfqpoint{2.852848in}{3.098248in}}%
\pgfpathlineto{\pgfqpoint{2.854970in}{3.098618in}}%
\pgfpathlineto{\pgfqpoint{2.857560in}{3.100926in}}%
\pgfpathlineto{\pgfqpoint{2.860528in}{3.102620in}}%
\pgfpathlineto{\pgfqpoint{2.863165in}{3.102396in}}%
\pgfpathlineto{\pgfqpoint{2.865803in}{3.100559in}}%
\pgfpathlineto{\pgfqpoint{2.869400in}{3.097877in}}%
\pgfpathlineto{\pgfqpoint{2.871522in}{3.098288in}}%
\pgfpathlineto{\pgfqpoint{2.874355in}{3.100824in}}%
\pgfpathlineto{\pgfqpoint{2.877323in}{3.102291in}}%
\pgfpathlineto{\pgfqpoint{2.879960in}{3.101867in}}%
\pgfpathlineto{\pgfqpoint{2.882692in}{3.099733in}}%
\pgfpathlineto{\pgfqpoint{2.885668in}{3.097577in}}%
\pgfpathlineto{\pgfqpoint{2.887790in}{3.097877in}}%
\pgfpathlineto{\pgfqpoint{2.890268in}{3.100059in}}%
\pgfpathlineto{\pgfqpoint{2.893235in}{3.101998in}}%
\pgfpathlineto{\pgfqpoint{2.895873in}{3.101991in}}%
\pgfpathlineto{\pgfqpoint{2.898510in}{3.100371in}}%
\pgfpathlineto{\pgfqpoint{2.902616in}{3.097493in}}%
\pgfpathlineto{\pgfqpoint{2.904821in}{3.098185in}}%
\pgfpathlineto{\pgfqpoint{2.911678in}{3.102135in}}%
\pgfpathlineto{\pgfqpoint{2.914316in}{3.100930in}}%
\pgfpathlineto{\pgfqpoint{2.919834in}{3.097603in}}%
\pgfpathlineto{\pgfqpoint{2.921977in}{3.098824in}}%
\pgfpathlineto{\pgfqpoint{2.926565in}{3.102273in}}%
\pgfpathlineto{\pgfqpoint{2.929203in}{3.102099in}}%
\pgfpathlineto{\pgfqpoint{2.931832in}{3.100323in}}%
\pgfpathlineto{\pgfqpoint{2.935210in}{3.097933in}}%
\pgfpathlineto{\pgfqpoint{2.937353in}{3.098372in}}%
\pgfpathlineto{\pgfqpoint{2.939984in}{3.100830in}}%
\pgfpathlineto{\pgfqpoint{2.942951in}{3.102668in}}%
\pgfpathlineto{\pgfqpoint{2.945589in}{3.102571in}}%
\pgfpathlineto{\pgfqpoint{2.948226in}{3.100859in}}%
\pgfpathlineto{\pgfqpoint{2.951854in}{3.098339in}}%
\pgfpathlineto{\pgfqpoint{2.954049in}{3.098913in}}%
\pgfpathlineto{\pgfqpoint{2.961581in}{3.103067in}}%
\pgfpathlineto{\pgfqpoint{2.964218in}{3.101633in}}%
\pgfpathlineto{\pgfqpoint{2.968762in}{3.098613in}}%
\pgfpathlineto{\pgfqpoint{2.970884in}{3.099427in}}%
\pgfpathlineto{\pgfqpoint{2.977240in}{3.103393in}}%
\pgfpathlineto{\pgfqpoint{2.979878in}{3.102437in}}%
\pgfpathlineto{\pgfqpoint{2.986278in}{3.098935in}}%
\pgfpathlineto{\pgfqpoint{2.988560in}{3.100643in}}%
\pgfpathlineto{\pgfqpoint{2.991738in}{3.103101in}}%
\pgfpathlineto{\pgfqpoint{2.994376in}{3.103340in}}%
\pgfpathlineto{\pgfqpoint{2.997013in}{3.101958in}}%
\pgfpathlineto{\pgfqpoint{3.002029in}{3.098712in}}%
\pgfpathlineto{\pgfqpoint{3.004110in}{3.099662in}}%
\pgfpathlineto{\pgfqpoint{3.009551in}{3.103302in}}%
\pgfpathlineto{\pgfqpoint{3.012188in}{3.102697in}}%
\pgfpathlineto{\pgfqpoint{3.015030in}{3.100252in}}%
\pgfpathlineto{\pgfqpoint{3.017605in}{3.098562in}}%
\pgfpathlineto{\pgfqpoint{3.019748in}{3.098884in}}%
\pgfpathlineto{\pgfqpoint{3.022271in}{3.101105in}}%
\pgfpathlineto{\pgfqpoint{3.025238in}{3.102893in}}%
\pgfpathlineto{\pgfqpoint{3.027875in}{3.102750in}}%
\pgfpathlineto{\pgfqpoint{3.030513in}{3.100993in}}%
\pgfpathlineto{\pgfqpoint{3.034433in}{3.098205in}}%
\pgfpathlineto{\pgfqpoint{3.036576in}{3.098800in}}%
\pgfpathlineto{\pgfqpoint{3.043838in}{3.102727in}}%
\pgfpathlineto{\pgfqpoint{3.046476in}{3.101327in}}%
\pgfpathlineto{\pgfqpoint{3.051350in}{3.098058in}}%
\pgfpathlineto{\pgfqpoint{3.053473in}{3.098941in}}%
\pgfpathlineto{\pgfqpoint{3.059183in}{3.102656in}}%
\pgfpathlineto{\pgfqpoint{3.061821in}{3.101943in}}%
\pgfpathlineto{\pgfqpoint{3.064741in}{3.099329in}}%
\pgfpathlineto{\pgfqpoint{3.067214in}{3.097885in}}%
\pgfpathlineto{\pgfqpoint{3.069372in}{3.098334in}}%
\pgfpathlineto{\pgfqpoint{3.072286in}{3.100980in}}%
\pgfpathlineto{\pgfqpoint{3.075253in}{3.102414in}}%
\pgfpathlineto{\pgfqpoint{3.077890in}{3.101961in}}%
\pgfpathlineto{\pgfqpoint{3.080704in}{3.099711in}}%
\pgfpathlineto{\pgfqpoint{3.083443in}{3.097814in}}%
\pgfpathlineto{\pgfqpoint{3.085586in}{3.098118in}}%
\pgfpathlineto{\pgfqpoint{3.088024in}{3.100277in}}%
\pgfpathlineto{\pgfqpoint{3.091156in}{3.102338in}}%
\pgfpathlineto{\pgfqpoint{3.093794in}{3.102293in}}%
\pgfpathlineto{\pgfqpoint{3.096432in}{3.100633in}}%
\pgfpathlineto{\pgfqpoint{3.100383in}{3.097888in}}%
\pgfpathlineto{\pgfqpoint{3.102512in}{3.098508in}}%
\pgfpathlineto{\pgfqpoint{3.109706in}{3.102574in}}%
\pgfpathlineto{\pgfqpoint{3.112343in}{3.101265in}}%
\pgfpathlineto{\pgfqpoint{3.117337in}{3.098105in}}%
\pgfpathlineto{\pgfqpoint{3.119480in}{3.099117in}}%
\pgfpathlineto{\pgfqpoint{3.125044in}{3.102965in}}%
\pgfpathlineto{\pgfqpoint{3.127681in}{3.102386in}}%
\pgfpathlineto{\pgfqpoint{3.130487in}{3.100022in}}%
\pgfpathlineto{\pgfqpoint{3.133036in}{3.098507in}}%
\pgfpathlineto{\pgfqpoint{3.135195in}{3.098960in}}%
\pgfpathlineto{\pgfqpoint{3.137802in}{3.101400in}}%
\pgfpathlineto{\pgfqpoint{3.140769in}{3.103232in}}%
\pgfpathlineto{\pgfqpoint{3.143407in}{3.103127in}}%
\pgfpathlineto{\pgfqpoint{3.146045in}{3.101405in}}%
\pgfpathlineto{\pgfqpoint{3.149645in}{3.098873in}}%
\pgfpathlineto{\pgfqpoint{3.151788in}{3.099383in}}%
\pgfpathlineto{\pgfqpoint{3.154779in}{3.102168in}}%
\pgfpathlineto{\pgfqpoint{3.157746in}{3.103613in}}%
\pgfpathlineto{\pgfqpoint{3.160383in}{3.103166in}}%
\pgfpathlineto{\pgfqpoint{3.163187in}{3.100922in}}%
\pgfpathlineto{\pgfqpoint{3.165779in}{3.099091in}}%
\pgfpathlineto{\pgfqpoint{3.167901in}{3.099302in}}%
\pgfpathlineto{\pgfqpoint{3.170277in}{3.101309in}}%
\pgfpathlineto{\pgfqpoint{3.173409in}{3.103459in}}%
\pgfpathlineto{\pgfqpoint{3.176047in}{3.103482in}}%
\pgfpathlineto{\pgfqpoint{3.178684in}{3.101887in}}%
\pgfpathlineto{\pgfqpoint{3.182953in}{3.098907in}}%
\pgfpathlineto{\pgfqpoint{3.185096in}{3.099603in}}%
\pgfpathlineto{\pgfqpoint{3.191661in}{3.103453in}}%
\pgfpathlineto{\pgfqpoint{3.194299in}{3.102353in}}%
\pgfpathlineto{\pgfqpoint{3.200422in}{3.098739in}}%
\pgfpathlineto{\pgfqpoint{3.202650in}{3.100222in}}%
\pgfpathlineto{\pgfqpoint{3.206217in}{3.102910in}}%
\pgfpathlineto{\pgfqpoint{3.208855in}{3.102989in}}%
\pgfpathlineto{\pgfqpoint{3.211493in}{3.101451in}}%
\pgfpathlineto{\pgfqpoint{3.216109in}{3.098259in}}%
\pgfpathlineto{\pgfqpoint{3.218232in}{3.099048in}}%
\pgfpathlineto{\pgfqpoint{3.224388in}{3.102748in}}%
\pgfpathlineto{\pgfqpoint{3.227025in}{3.101782in}}%
\pgfpathlineto{\pgfqpoint{3.229161in}{3.099834in}}%
\pgfpathlineto{\pgfqpoint{3.229161in}{3.099834in}}%
\pgfusepath{stroke}%
\end{pgfscope}%
\begin{pgfscope}%
\pgfsetrectcap%
\pgfsetmiterjoin%
\pgfsetlinewidth{0.803000pt}%
\definecolor{currentstroke}{rgb}{0.000000,0.000000,0.000000}%
\pgfsetstrokecolor{currentstroke}%
\pgfsetdash{}{0pt}%
\pgfpathmoveto{\pgfqpoint{0.632797in}{2.601404in}}%
\pgfpathlineto{\pgfqpoint{0.632797in}{3.507654in}}%
\pgfusepath{stroke}%
\end{pgfscope}%
\begin{pgfscope}%
\pgfsetrectcap%
\pgfsetmiterjoin%
\pgfsetlinewidth{0.803000pt}%
\definecolor{currentstroke}{rgb}{0.000000,0.000000,0.000000}%
\pgfsetstrokecolor{currentstroke}%
\pgfsetdash{}{0pt}%
\pgfpathmoveto{\pgfqpoint{3.352797in}{2.601404in}}%
\pgfpathlineto{\pgfqpoint{3.352797in}{3.507654in}}%
\pgfusepath{stroke}%
\end{pgfscope}%
\begin{pgfscope}%
\pgfsetrectcap%
\pgfsetmiterjoin%
\pgfsetlinewidth{0.803000pt}%
\definecolor{currentstroke}{rgb}{0.000000,0.000000,0.000000}%
\pgfsetstrokecolor{currentstroke}%
\pgfsetdash{}{0pt}%
\pgfpathmoveto{\pgfqpoint{0.632797in}{2.601404in}}%
\pgfpathlineto{\pgfqpoint{3.352797in}{2.601404in}}%
\pgfusepath{stroke}%
\end{pgfscope}%
\begin{pgfscope}%
\pgfsetrectcap%
\pgfsetmiterjoin%
\pgfsetlinewidth{0.803000pt}%
\definecolor{currentstroke}{rgb}{0.000000,0.000000,0.000000}%
\pgfsetstrokecolor{currentstroke}%
\pgfsetdash{}{0pt}%
\pgfpathmoveto{\pgfqpoint{0.632797in}{3.507654in}}%
\pgfpathlineto{\pgfqpoint{3.352797in}{3.507654in}}%
\pgfusepath{stroke}%
\end{pgfscope}%
\begin{pgfscope}%
\definecolor{textcolor}{rgb}{0.000000,0.000000,0.000000}%
\pgfsetstrokecolor{textcolor}%
\pgfsetfillcolor{textcolor}%
\pgftext[x=0.687197in,y=3.444217in,left,top]{\color{textcolor}{\rmfamily\fontsize{8.000000}{9.600000}\selectfont\catcode`\^=\active\def^{\ifmmode\sp\else\^{}\fi}\catcode`\%=\active\def%{\%}(a)}}%
\end{pgfscope}%
\begin{pgfscope}%
\pgfsetbuttcap%
\pgfsetmiterjoin%
\definecolor{currentfill}{rgb}{1.000000,1.000000,1.000000}%
\pgfsetfillcolor{currentfill}%
\pgfsetlinewidth{0.000000pt}%
\definecolor{currentstroke}{rgb}{0.000000,0.000000,0.000000}%
\pgfsetstrokecolor{currentstroke}%
\pgfsetstrokeopacity{0.000000}%
\pgfsetdash{}{0pt}%
\pgfpathmoveto{\pgfqpoint{0.632797in}{1.532029in}}%
\pgfpathlineto{\pgfqpoint{3.352797in}{1.532029in}}%
\pgfpathlineto{\pgfqpoint{3.352797in}{2.438279in}}%
\pgfpathlineto{\pgfqpoint{0.632797in}{2.438279in}}%
\pgfpathlineto{\pgfqpoint{0.632797in}{1.532029in}}%
\pgfpathclose%
\pgfusepath{fill}%
\end{pgfscope}%
\begin{pgfscope}%
\pgfpathrectangle{\pgfqpoint{0.632797in}{1.532029in}}{\pgfqpoint{2.720000in}{0.906250in}}%
\pgfusepath{clip}%
\pgfsetbuttcap%
\pgfsetroundjoin%
\pgfsetlinewidth{0.501875pt}%
\definecolor{currentstroke}{rgb}{0.690196,0.690196,0.690196}%
\pgfsetstrokecolor{currentstroke}%
\pgfsetstrokeopacity{0.250000}%
\pgfsetdash{{1.850000pt}{0.800000pt}}{0.000000pt}%
\pgfpathmoveto{\pgfqpoint{0.756433in}{1.532029in}}%
\pgfpathlineto{\pgfqpoint{0.756433in}{2.438279in}}%
\pgfusepath{stroke}%
\end{pgfscope}%
\begin{pgfscope}%
\pgfsetbuttcap%
\pgfsetroundjoin%
\definecolor{currentfill}{rgb}{0.000000,0.000000,0.000000}%
\pgfsetfillcolor{currentfill}%
\pgfsetlinewidth{0.803000pt}%
\definecolor{currentstroke}{rgb}{0.000000,0.000000,0.000000}%
\pgfsetstrokecolor{currentstroke}%
\pgfsetdash{}{0pt}%
\pgfsys@defobject{currentmarker}{\pgfqpoint{0.000000in}{-0.048611in}}{\pgfqpoint{0.000000in}{0.000000in}}{%
\pgfpathmoveto{\pgfqpoint{0.000000in}{0.000000in}}%
\pgfpathlineto{\pgfqpoint{0.000000in}{-0.048611in}}%
\pgfusepath{stroke,fill}%
}%
\begin{pgfscope}%
\pgfsys@transformshift{0.756433in}{1.532029in}%
\pgfsys@useobject{currentmarker}{}%
\end{pgfscope}%
\end{pgfscope}%
\begin{pgfscope}%
\pgfpathrectangle{\pgfqpoint{0.632797in}{1.532029in}}{\pgfqpoint{2.720000in}{0.906250in}}%
\pgfusepath{clip}%
\pgfsetbuttcap%
\pgfsetroundjoin%
\pgfsetlinewidth{0.501875pt}%
\definecolor{currentstroke}{rgb}{0.690196,0.690196,0.690196}%
\pgfsetstrokecolor{currentstroke}%
\pgfsetstrokeopacity{0.250000}%
\pgfsetdash{{1.850000pt}{0.800000pt}}{0.000000pt}%
\pgfpathmoveto{\pgfqpoint{1.415827in}{1.532029in}}%
\pgfpathlineto{\pgfqpoint{1.415827in}{2.438279in}}%
\pgfusepath{stroke}%
\end{pgfscope}%
\begin{pgfscope}%
\pgfsetbuttcap%
\pgfsetroundjoin%
\definecolor{currentfill}{rgb}{0.000000,0.000000,0.000000}%
\pgfsetfillcolor{currentfill}%
\pgfsetlinewidth{0.803000pt}%
\definecolor{currentstroke}{rgb}{0.000000,0.000000,0.000000}%
\pgfsetstrokecolor{currentstroke}%
\pgfsetdash{}{0pt}%
\pgfsys@defobject{currentmarker}{\pgfqpoint{0.000000in}{-0.048611in}}{\pgfqpoint{0.000000in}{0.000000in}}{%
\pgfpathmoveto{\pgfqpoint{0.000000in}{0.000000in}}%
\pgfpathlineto{\pgfqpoint{0.000000in}{-0.048611in}}%
\pgfusepath{stroke,fill}%
}%
\begin{pgfscope}%
\pgfsys@transformshift{1.415827in}{1.532029in}%
\pgfsys@useobject{currentmarker}{}%
\end{pgfscope}%
\end{pgfscope}%
\begin{pgfscope}%
\pgfpathrectangle{\pgfqpoint{0.632797in}{1.532029in}}{\pgfqpoint{2.720000in}{0.906250in}}%
\pgfusepath{clip}%
\pgfsetbuttcap%
\pgfsetroundjoin%
\pgfsetlinewidth{0.501875pt}%
\definecolor{currentstroke}{rgb}{0.690196,0.690196,0.690196}%
\pgfsetstrokecolor{currentstroke}%
\pgfsetstrokeopacity{0.250000}%
\pgfsetdash{{1.850000pt}{0.800000pt}}{0.000000pt}%
\pgfpathmoveto{\pgfqpoint{2.075221in}{1.532029in}}%
\pgfpathlineto{\pgfqpoint{2.075221in}{2.438279in}}%
\pgfusepath{stroke}%
\end{pgfscope}%
\begin{pgfscope}%
\pgfsetbuttcap%
\pgfsetroundjoin%
\definecolor{currentfill}{rgb}{0.000000,0.000000,0.000000}%
\pgfsetfillcolor{currentfill}%
\pgfsetlinewidth{0.803000pt}%
\definecolor{currentstroke}{rgb}{0.000000,0.000000,0.000000}%
\pgfsetstrokecolor{currentstroke}%
\pgfsetdash{}{0pt}%
\pgfsys@defobject{currentmarker}{\pgfqpoint{0.000000in}{-0.048611in}}{\pgfqpoint{0.000000in}{0.000000in}}{%
\pgfpathmoveto{\pgfqpoint{0.000000in}{0.000000in}}%
\pgfpathlineto{\pgfqpoint{0.000000in}{-0.048611in}}%
\pgfusepath{stroke,fill}%
}%
\begin{pgfscope}%
\pgfsys@transformshift{2.075221in}{1.532029in}%
\pgfsys@useobject{currentmarker}{}%
\end{pgfscope}%
\end{pgfscope}%
\begin{pgfscope}%
\pgfpathrectangle{\pgfqpoint{0.632797in}{1.532029in}}{\pgfqpoint{2.720000in}{0.906250in}}%
\pgfusepath{clip}%
\pgfsetbuttcap%
\pgfsetroundjoin%
\pgfsetlinewidth{0.501875pt}%
\definecolor{currentstroke}{rgb}{0.690196,0.690196,0.690196}%
\pgfsetstrokecolor{currentstroke}%
\pgfsetstrokeopacity{0.250000}%
\pgfsetdash{{1.850000pt}{0.800000pt}}{0.000000pt}%
\pgfpathmoveto{\pgfqpoint{2.734615in}{1.532029in}}%
\pgfpathlineto{\pgfqpoint{2.734615in}{2.438279in}}%
\pgfusepath{stroke}%
\end{pgfscope}%
\begin{pgfscope}%
\pgfsetbuttcap%
\pgfsetroundjoin%
\definecolor{currentfill}{rgb}{0.000000,0.000000,0.000000}%
\pgfsetfillcolor{currentfill}%
\pgfsetlinewidth{0.803000pt}%
\definecolor{currentstroke}{rgb}{0.000000,0.000000,0.000000}%
\pgfsetstrokecolor{currentstroke}%
\pgfsetdash{}{0pt}%
\pgfsys@defobject{currentmarker}{\pgfqpoint{0.000000in}{-0.048611in}}{\pgfqpoint{0.000000in}{0.000000in}}{%
\pgfpathmoveto{\pgfqpoint{0.000000in}{0.000000in}}%
\pgfpathlineto{\pgfqpoint{0.000000in}{-0.048611in}}%
\pgfusepath{stroke,fill}%
}%
\begin{pgfscope}%
\pgfsys@transformshift{2.734615in}{1.532029in}%
\pgfsys@useobject{currentmarker}{}%
\end{pgfscope}%
\end{pgfscope}%
\begin{pgfscope}%
\pgfpathrectangle{\pgfqpoint{0.632797in}{1.532029in}}{\pgfqpoint{2.720000in}{0.906250in}}%
\pgfusepath{clip}%
\pgfsetbuttcap%
\pgfsetroundjoin%
\pgfsetlinewidth{0.501875pt}%
\definecolor{currentstroke}{rgb}{0.690196,0.690196,0.690196}%
\pgfsetstrokecolor{currentstroke}%
\pgfsetstrokeopacity{0.250000}%
\pgfsetdash{{1.850000pt}{0.800000pt}}{0.000000pt}%
\pgfpathmoveto{\pgfqpoint{0.632797in}{1.573223in}}%
\pgfpathlineto{\pgfqpoint{3.352797in}{1.573223in}}%
\pgfusepath{stroke}%
\end{pgfscope}%
\begin{pgfscope}%
\pgfsetbuttcap%
\pgfsetroundjoin%
\definecolor{currentfill}{rgb}{0.000000,0.000000,0.000000}%
\pgfsetfillcolor{currentfill}%
\pgfsetlinewidth{0.803000pt}%
\definecolor{currentstroke}{rgb}{0.000000,0.000000,0.000000}%
\pgfsetstrokecolor{currentstroke}%
\pgfsetdash{}{0pt}%
\pgfsys@defobject{currentmarker}{\pgfqpoint{-0.048611in}{0.000000in}}{\pgfqpoint{-0.000000in}{0.000000in}}{%
\pgfpathmoveto{\pgfqpoint{-0.000000in}{0.000000in}}%
\pgfpathlineto{\pgfqpoint{-0.048611in}{0.000000in}}%
\pgfusepath{stroke,fill}%
}%
\begin{pgfscope}%
\pgfsys@transformshift{0.632797in}{1.573223in}%
\pgfsys@useobject{currentmarker}{}%
\end{pgfscope}%
\end{pgfscope}%
\begin{pgfscope}%
\definecolor{textcolor}{rgb}{0.000000,0.000000,0.000000}%
\pgfsetstrokecolor{textcolor}%
\pgfsetfillcolor{textcolor}%
\pgftext[x=0.476546in, y=1.534642in, left, base]{\color{textcolor}{\rmfamily\fontsize{8.000000}{9.600000}\selectfont\catcode`\^=\active\def^{\ifmmode\sp\else\^{}\fi}\catcode`\%=\active\def%{\%}$\mathdefault{0}$}}%
\end{pgfscope}%
\begin{pgfscope}%
\pgfpathrectangle{\pgfqpoint{0.632797in}{1.532029in}}{\pgfqpoint{2.720000in}{0.906250in}}%
\pgfusepath{clip}%
\pgfsetbuttcap%
\pgfsetroundjoin%
\pgfsetlinewidth{0.501875pt}%
\definecolor{currentstroke}{rgb}{0.690196,0.690196,0.690196}%
\pgfsetstrokecolor{currentstroke}%
\pgfsetstrokeopacity{0.250000}%
\pgfsetdash{{1.850000pt}{0.800000pt}}{0.000000pt}%
\pgfpathmoveto{\pgfqpoint{0.632797in}{1.977665in}}%
\pgfpathlineto{\pgfqpoint{3.352797in}{1.977665in}}%
\pgfusepath{stroke}%
\end{pgfscope}%
\begin{pgfscope}%
\pgfsetbuttcap%
\pgfsetroundjoin%
\definecolor{currentfill}{rgb}{0.000000,0.000000,0.000000}%
\pgfsetfillcolor{currentfill}%
\pgfsetlinewidth{0.803000pt}%
\definecolor{currentstroke}{rgb}{0.000000,0.000000,0.000000}%
\pgfsetstrokecolor{currentstroke}%
\pgfsetdash{}{0pt}%
\pgfsys@defobject{currentmarker}{\pgfqpoint{-0.048611in}{0.000000in}}{\pgfqpoint{-0.000000in}{0.000000in}}{%
\pgfpathmoveto{\pgfqpoint{-0.000000in}{0.000000in}}%
\pgfpathlineto{\pgfqpoint{-0.048611in}{0.000000in}}%
\pgfusepath{stroke,fill}%
}%
\begin{pgfscope}%
\pgfsys@transformshift{0.632797in}{1.977665in}%
\pgfsys@useobject{currentmarker}{}%
\end{pgfscope}%
\end{pgfscope}%
\begin{pgfscope}%
\definecolor{textcolor}{rgb}{0.000000,0.000000,0.000000}%
\pgfsetstrokecolor{textcolor}%
\pgfsetfillcolor{textcolor}%
\pgftext[x=0.417518in, y=1.939085in, left, base]{\color{textcolor}{\rmfamily\fontsize{8.000000}{9.600000}\selectfont\catcode`\^=\active\def^{\ifmmode\sp\else\^{}\fi}\catcode`\%=\active\def%{\%}$\mathdefault{10}$}}%
\end{pgfscope}%
\begin{pgfscope}%
\pgfpathrectangle{\pgfqpoint{0.632797in}{1.532029in}}{\pgfqpoint{2.720000in}{0.906250in}}%
\pgfusepath{clip}%
\pgfsetbuttcap%
\pgfsetroundjoin%
\pgfsetlinewidth{0.501875pt}%
\definecolor{currentstroke}{rgb}{0.690196,0.690196,0.690196}%
\pgfsetstrokecolor{currentstroke}%
\pgfsetstrokeopacity{0.250000}%
\pgfsetdash{{1.850000pt}{0.800000pt}}{0.000000pt}%
\pgfpathmoveto{\pgfqpoint{0.632797in}{2.382108in}}%
\pgfpathlineto{\pgfqpoint{3.352797in}{2.382108in}}%
\pgfusepath{stroke}%
\end{pgfscope}%
\begin{pgfscope}%
\pgfsetbuttcap%
\pgfsetroundjoin%
\definecolor{currentfill}{rgb}{0.000000,0.000000,0.000000}%
\pgfsetfillcolor{currentfill}%
\pgfsetlinewidth{0.803000pt}%
\definecolor{currentstroke}{rgb}{0.000000,0.000000,0.000000}%
\pgfsetstrokecolor{currentstroke}%
\pgfsetdash{}{0pt}%
\pgfsys@defobject{currentmarker}{\pgfqpoint{-0.048611in}{0.000000in}}{\pgfqpoint{-0.000000in}{0.000000in}}{%
\pgfpathmoveto{\pgfqpoint{-0.000000in}{0.000000in}}%
\pgfpathlineto{\pgfqpoint{-0.048611in}{0.000000in}}%
\pgfusepath{stroke,fill}%
}%
\begin{pgfscope}%
\pgfsys@transformshift{0.632797in}{2.382108in}%
\pgfsys@useobject{currentmarker}{}%
\end{pgfscope}%
\end{pgfscope}%
\begin{pgfscope}%
\definecolor{textcolor}{rgb}{0.000000,0.000000,0.000000}%
\pgfsetstrokecolor{textcolor}%
\pgfsetfillcolor{textcolor}%
\pgftext[x=0.417518in, y=2.343528in, left, base]{\color{textcolor}{\rmfamily\fontsize{8.000000}{9.600000}\selectfont\catcode`\^=\active\def^{\ifmmode\sp\else\^{}\fi}\catcode`\%=\active\def%{\%}$\mathdefault{20}$}}%
\end{pgfscope}%
\begin{pgfscope}%
\definecolor{textcolor}{rgb}{0.000000,0.000000,0.000000}%
\pgfsetstrokecolor{textcolor}%
\pgfsetfillcolor{textcolor}%
\pgftext[x=0.361962in,y=1.985154in,,bottom,rotate=90.000000]{\color{textcolor}{\rmfamily\fontsize{8.000000}{9.600000}\selectfont\catcode`\^=\active\def^{\ifmmode\sp\else\^{}\fi}\catcode`\%=\active\def%{\%}$i_L$ (V)}}%
\end{pgfscope}%
\begin{pgfscope}%
\pgfpathrectangle{\pgfqpoint{0.632797in}{1.532029in}}{\pgfqpoint{2.720000in}{0.906250in}}%
\pgfusepath{clip}%
\pgfsetrectcap%
\pgfsetroundjoin%
\pgfsetlinewidth{1.505625pt}%
\definecolor{currentstroke}{rgb}{0.835294,0.368627,0.000000}%
\pgfsetstrokecolor{currentstroke}%
\pgfsetdash{}{0pt}%
\pgfpathmoveto{\pgfqpoint{0.756433in}{1.573223in}}%
\pgfpathlineto{\pgfqpoint{0.763154in}{1.842145in}}%
\pgfpathlineto{\pgfqpoint{0.763627in}{1.841658in}}%
\pgfpathlineto{\pgfqpoint{0.767748in}{1.836198in}}%
\pgfpathlineto{\pgfqpoint{0.772630in}{1.826975in}}%
\pgfpathlineto{\pgfqpoint{0.772926in}{1.826321in}}%
\pgfpathlineto{\pgfqpoint{0.772970in}{1.827655in}}%
\pgfpathlineto{\pgfqpoint{0.779656in}{2.080518in}}%
\pgfpathlineto{\pgfqpoint{0.780124in}{2.078869in}}%
\pgfpathlineto{\pgfqpoint{0.785070in}{2.058318in}}%
\pgfpathlineto{\pgfqpoint{0.789413in}{2.035715in}}%
\pgfpathlineto{\pgfqpoint{0.789565in}{2.040727in}}%
\pgfpathlineto{\pgfqpoint{0.796130in}{2.264575in}}%
\pgfpathlineto{\pgfqpoint{0.796444in}{2.262214in}}%
\pgfpathlineto{\pgfqpoint{0.802708in}{2.208738in}}%
\pgfpathlineto{\pgfqpoint{0.805898in}{2.177132in}}%
\pgfpathlineto{\pgfqpoint{0.806108in}{2.183216in}}%
\pgfpathlineto{\pgfqpoint{0.812592in}{2.373495in}}%
\pgfpathlineto{\pgfqpoint{0.812888in}{2.370107in}}%
\pgfpathlineto{\pgfqpoint{0.820471in}{2.265264in}}%
\pgfpathlineto{\pgfqpoint{0.822390in}{2.235988in}}%
\pgfpathlineto{\pgfqpoint{0.822607in}{2.241350in}}%
\pgfpathlineto{\pgfqpoint{0.829065in}{2.397086in}}%
\pgfpathlineto{\pgfqpoint{0.829319in}{2.392995in}}%
\pgfpathlineto{\pgfqpoint{0.838704in}{2.212026in}}%
\pgfpathlineto{\pgfqpoint{0.838883in}{2.208439in}}%
\pgfpathlineto{\pgfqpoint{0.839235in}{2.215465in}}%
\pgfpathlineto{\pgfqpoint{0.845556in}{2.336256in}}%
\pgfpathlineto{\pgfqpoint{0.845760in}{2.331992in}}%
\pgfpathlineto{\pgfqpoint{0.855360in}{2.102842in}}%
\pgfpathlineto{\pgfqpoint{0.856054in}{2.113726in}}%
\pgfpathlineto{\pgfqpoint{0.862040in}{2.203671in}}%
\pgfpathlineto{\pgfqpoint{0.862202in}{2.199738in}}%
\pgfpathlineto{\pgfqpoint{0.871846in}{1.936691in}}%
\pgfpathlineto{\pgfqpoint{0.872456in}{1.944479in}}%
\pgfpathlineto{\pgfqpoint{0.878509in}{2.019655in}}%
\pgfpathlineto{\pgfqpoint{0.878695in}{2.015209in}}%
\pgfpathlineto{\pgfqpoint{0.888331in}{1.734507in}}%
\pgfpathlineto{\pgfqpoint{0.888873in}{1.740713in}}%
\pgfpathlineto{\pgfqpoint{0.895007in}{1.810850in}}%
\pgfpathlineto{\pgfqpoint{0.895149in}{1.807046in}}%
\pgfpathlineto{\pgfqpoint{0.903870in}{1.573223in}}%
\pgfpathlineto{\pgfqpoint{0.904878in}{1.573947in}}%
\pgfpathlineto{\pgfqpoint{0.911491in}{1.653912in}}%
\pgfpathlineto{\pgfqpoint{0.911705in}{1.648194in}}%
\pgfpathlineto{\pgfqpoint{0.914974in}{1.573223in}}%
\pgfpathlineto{\pgfqpoint{0.921358in}{1.573976in}}%
\pgfpathlineto{\pgfqpoint{0.927976in}{1.661858in}}%
\pgfpathlineto{\pgfqpoint{0.928189in}{1.656410in}}%
\pgfpathlineto{\pgfqpoint{0.931954in}{1.573223in}}%
\pgfpathlineto{\pgfqpoint{0.937839in}{1.573936in}}%
\pgfpathlineto{\pgfqpoint{0.944461in}{1.669182in}}%
\pgfpathlineto{\pgfqpoint{0.944674in}{1.663965in}}%
\pgfpathlineto{\pgfqpoint{0.948748in}{1.573223in}}%
\pgfpathlineto{\pgfqpoint{0.954316in}{1.573926in}}%
\pgfpathlineto{\pgfqpoint{0.960946in}{1.676058in}}%
\pgfpathlineto{\pgfqpoint{0.961160in}{1.671035in}}%
\pgfpathlineto{\pgfqpoint{0.965790in}{1.573223in}}%
\pgfpathlineto{\pgfqpoint{0.970807in}{1.573965in}}%
\pgfpathlineto{\pgfqpoint{0.977420in}{1.682205in}}%
\pgfpathlineto{\pgfqpoint{0.977670in}{1.676533in}}%
\pgfpathlineto{\pgfqpoint{0.982633in}{1.573223in}}%
\pgfpathlineto{\pgfqpoint{0.987290in}{1.573976in}}%
\pgfpathlineto{\pgfqpoint{0.993916in}{1.688115in}}%
\pgfpathlineto{\pgfqpoint{0.994130in}{1.683469in}}%
\pgfpathlineto{\pgfqpoint{0.999433in}{1.573223in}}%
\pgfpathlineto{\pgfqpoint{1.003765in}{1.573916in}}%
\pgfpathlineto{\pgfqpoint{1.010401in}{1.693555in}}%
\pgfpathlineto{\pgfqpoint{1.010614in}{1.689079in}}%
\pgfpathlineto{\pgfqpoint{1.016489in}{1.573223in}}%
\pgfpathlineto{\pgfqpoint{1.020247in}{1.573934in}}%
\pgfpathlineto{\pgfqpoint{1.026886in}{1.698535in}}%
\pgfpathlineto{\pgfqpoint{1.027099in}{1.694210in}}%
\pgfpathlineto{\pgfqpoint{1.033364in}{1.573223in}}%
\pgfpathlineto{\pgfqpoint{1.036737in}{1.573988in}}%
\pgfpathlineto{\pgfqpoint{1.043370in}{1.702973in}}%
\pgfpathlineto{\pgfqpoint{1.043729in}{1.695704in}}%
\pgfpathlineto{\pgfqpoint{1.050133in}{1.573223in}}%
\pgfpathlineto{\pgfqpoint{1.051781in}{1.573223in}}%
\pgfpathlineto{\pgfqpoint{1.051946in}{1.574004in}}%
\pgfpathlineto{\pgfqpoint{1.052473in}{1.573223in}}%
\pgfpathlineto{\pgfqpoint{1.053212in}{1.573794in}}%
\pgfpathlineto{\pgfqpoint{1.059855in}{1.707080in}}%
\pgfpathlineto{\pgfqpoint{1.060218in}{1.699948in}}%
\pgfpathlineto{\pgfqpoint{1.067029in}{1.573223in}}%
\pgfpathlineto{\pgfqpoint{1.069700in}{1.573898in}}%
\pgfpathlineto{\pgfqpoint{1.076340in}{1.710868in}}%
\pgfpathlineto{\pgfqpoint{1.076712in}{1.703748in}}%
\pgfpathlineto{\pgfqpoint{1.083925in}{1.573223in}}%
\pgfpathlineto{\pgfqpoint{1.086180in}{1.573807in}}%
\pgfpathlineto{\pgfqpoint{1.092825in}{1.714300in}}%
\pgfpathlineto{\pgfqpoint{1.093199in}{1.707325in}}%
\pgfpathlineto{\pgfqpoint{1.100706in}{1.573223in}}%
\pgfpathlineto{\pgfqpoint{1.102670in}{1.573959in}}%
\pgfpathlineto{\pgfqpoint{1.109315in}{1.717472in}}%
\pgfpathlineto{\pgfqpoint{1.109700in}{1.710430in}}%
\pgfpathlineto{\pgfqpoint{1.117617in}{1.573223in}}%
\pgfpathlineto{\pgfqpoint{1.119149in}{1.573788in}}%
\pgfpathlineto{\pgfqpoint{1.125794in}{1.720152in}}%
\pgfpathlineto{\pgfqpoint{1.126147in}{1.713902in}}%
\pgfpathlineto{\pgfqpoint{1.134350in}{1.573223in}}%
\pgfpathlineto{\pgfqpoint{1.135631in}{1.573982in}}%
\pgfpathlineto{\pgfqpoint{1.142269in}{1.722797in}}%
\pgfpathlineto{\pgfqpoint{1.142636in}{1.716487in}}%
\pgfpathlineto{\pgfqpoint{1.151166in}{1.573223in}}%
\pgfpathlineto{\pgfqpoint{1.152120in}{1.573817in}}%
\pgfpathlineto{\pgfqpoint{1.158764in}{1.724949in}}%
\pgfpathlineto{\pgfqpoint{1.159125in}{1.718804in}}%
\pgfpathlineto{\pgfqpoint{1.167898in}{1.573223in}}%
\pgfpathlineto{\pgfqpoint{1.168384in}{1.573880in}}%
\pgfpathlineto{\pgfqpoint{1.168564in}{1.573223in}}%
\pgfpathlineto{\pgfqpoint{1.168720in}{1.576586in}}%
\pgfpathlineto{\pgfqpoint{1.175249in}{1.727105in}}%
\pgfpathlineto{\pgfqpoint{1.175611in}{1.721054in}}%
\pgfpathlineto{\pgfqpoint{1.184609in}{1.573223in}}%
\pgfpathlineto{\pgfqpoint{1.185074in}{1.573955in}}%
\pgfpathlineto{\pgfqpoint{1.191734in}{1.729326in}}%
\pgfpathlineto{\pgfqpoint{1.192096in}{1.723373in}}%
\pgfpathlineto{\pgfqpoint{1.201367in}{1.573223in}}%
\pgfpathlineto{\pgfqpoint{1.201849in}{1.580942in}}%
\pgfpathlineto{\pgfqpoint{1.208219in}{1.731068in}}%
\pgfpathlineto{\pgfqpoint{1.208581in}{1.725196in}}%
\pgfpathlineto{\pgfqpoint{1.218024in}{1.573223in}}%
\pgfpathlineto{\pgfqpoint{1.218084in}{1.574525in}}%
\pgfpathlineto{\pgfqpoint{1.224704in}{1.732024in}}%
\pgfpathlineto{\pgfqpoint{1.225064in}{1.726261in}}%
\pgfpathlineto{\pgfqpoint{1.234509in}{1.573223in}}%
\pgfpathlineto{\pgfqpoint{1.234618in}{1.575730in}}%
\pgfpathlineto{\pgfqpoint{1.241189in}{1.733325in}}%
\pgfpathlineto{\pgfqpoint{1.241550in}{1.727615in}}%
\pgfpathlineto{\pgfqpoint{1.250920in}{1.573223in}}%
\pgfpathlineto{\pgfqpoint{1.251186in}{1.577815in}}%
\pgfpathlineto{\pgfqpoint{1.257699in}{1.734528in}}%
\pgfpathlineto{\pgfqpoint{1.258119in}{1.727530in}}%
\pgfpathlineto{\pgfqpoint{1.267485in}{1.573223in}}%
\pgfpathlineto{\pgfqpoint{1.267696in}{1.578245in}}%
\pgfpathlineto{\pgfqpoint{1.274181in}{1.735269in}}%
\pgfpathlineto{\pgfqpoint{1.274503in}{1.729953in}}%
\pgfpathlineto{\pgfqpoint{1.283962in}{1.574387in}}%
\pgfpathlineto{\pgfqpoint{1.284166in}{1.579275in}}%
\pgfpathlineto{\pgfqpoint{1.290643in}{1.737650in}}%
\pgfpathlineto{\pgfqpoint{1.291004in}{1.732108in}}%
\pgfpathlineto{\pgfqpoint{1.300449in}{1.577816in}}%
\pgfpathlineto{\pgfqpoint{1.300696in}{1.583800in}}%
\pgfpathlineto{\pgfqpoint{1.307145in}{1.741845in}}%
\pgfpathlineto{\pgfqpoint{1.307437in}{1.737384in}}%
\pgfpathlineto{\pgfqpoint{1.316934in}{1.587393in}}%
\pgfpathlineto{\pgfqpoint{1.317082in}{1.590942in}}%
\pgfpathlineto{\pgfqpoint{1.323613in}{1.751647in}}%
\pgfpathlineto{\pgfqpoint{1.323978in}{1.746086in}}%
\pgfpathlineto{\pgfqpoint{1.333419in}{1.592393in}}%
\pgfpathlineto{\pgfqpoint{1.333657in}{1.598148in}}%
\pgfpathlineto{\pgfqpoint{1.340101in}{1.756593in}}%
\pgfpathlineto{\pgfqpoint{1.340530in}{1.749977in}}%
\pgfpathlineto{\pgfqpoint{1.349904in}{1.597037in}}%
\pgfpathlineto{\pgfqpoint{1.350144in}{1.602830in}}%
\pgfpathlineto{\pgfqpoint{1.356583in}{1.760825in}}%
\pgfpathlineto{\pgfqpoint{1.356946in}{1.755265in}}%
\pgfpathlineto{\pgfqpoint{1.366389in}{1.600561in}}%
\pgfpathlineto{\pgfqpoint{1.366648in}{1.606821in}}%
\pgfpathlineto{\pgfqpoint{1.373068in}{1.763814in}}%
\pgfpathlineto{\pgfqpoint{1.373430in}{1.758237in}}%
\pgfpathlineto{\pgfqpoint{1.382874in}{1.602672in}}%
\pgfpathlineto{\pgfqpoint{1.383130in}{1.608837in}}%
\pgfpathlineto{\pgfqpoint{1.389563in}{1.765417in}}%
\pgfpathlineto{\pgfqpoint{1.389914in}{1.759936in}}%
\pgfpathlineto{\pgfqpoint{1.399360in}{1.603431in}}%
\pgfpathlineto{\pgfqpoint{1.399627in}{1.609843in}}%
\pgfpathlineto{\pgfqpoint{1.406033in}{1.765284in}}%
\pgfpathlineto{\pgfqpoint{1.406365in}{1.760194in}}%
\pgfpathlineto{\pgfqpoint{1.415843in}{1.602219in}}%
\pgfpathlineto{\pgfqpoint{1.416086in}{1.607993in}}%
\pgfpathlineto{\pgfqpoint{1.422525in}{1.763623in}}%
\pgfpathlineto{\pgfqpoint{1.422953in}{1.756845in}}%
\pgfpathlineto{\pgfqpoint{1.432328in}{1.599832in}}%
\pgfpathlineto{\pgfqpoint{1.432612in}{1.606583in}}%
\pgfpathlineto{\pgfqpoint{1.439007in}{1.760726in}}%
\pgfpathlineto{\pgfqpoint{1.439369in}{1.755025in}}%
\pgfpathlineto{\pgfqpoint{1.448813in}{1.596360in}}%
\pgfpathlineto{\pgfqpoint{1.449063in}{1.602299in}}%
\pgfpathlineto{\pgfqpoint{1.455492in}{1.756996in}}%
\pgfpathlineto{\pgfqpoint{1.455854in}{1.751289in}}%
\pgfpathlineto{\pgfqpoint{1.465298in}{1.592407in}}%
\pgfpathlineto{\pgfqpoint{1.465569in}{1.598856in}}%
\pgfpathlineto{\pgfqpoint{1.471981in}{1.753064in}}%
\pgfpathlineto{\pgfqpoint{1.472375in}{1.746796in}}%
\pgfpathlineto{\pgfqpoint{1.481785in}{1.588875in}}%
\pgfpathlineto{\pgfqpoint{1.482033in}{1.594802in}}%
\pgfpathlineto{\pgfqpoint{1.488461in}{1.749625in}}%
\pgfpathlineto{\pgfqpoint{1.488823in}{1.743940in}}%
\pgfpathlineto{\pgfqpoint{1.498268in}{1.585491in}}%
\pgfpathlineto{\pgfqpoint{1.498542in}{1.592021in}}%
\pgfpathlineto{\pgfqpoint{1.504946in}{1.746580in}}%
\pgfpathlineto{\pgfqpoint{1.505310in}{1.740882in}}%
\pgfpathlineto{\pgfqpoint{1.514752in}{1.583036in}}%
\pgfpathlineto{\pgfqpoint{1.515027in}{1.589607in}}%
\pgfpathlineto{\pgfqpoint{1.521453in}{1.744536in}}%
\pgfpathlineto{\pgfqpoint{1.521856in}{1.737831in}}%
\pgfpathlineto{\pgfqpoint{1.531246in}{1.581601in}}%
\pgfpathlineto{\pgfqpoint{1.531496in}{1.587631in}}%
\pgfpathlineto{\pgfqpoint{1.537917in}{1.743509in}}%
\pgfpathlineto{\pgfqpoint{1.538209in}{1.739073in}}%
\pgfpathlineto{\pgfqpoint{1.547720in}{1.583898in}}%
\pgfpathlineto{\pgfqpoint{1.547934in}{1.588942in}}%
\pgfpathlineto{\pgfqpoint{1.554402in}{1.746434in}}%
\pgfpathlineto{\pgfqpoint{1.554695in}{1.741990in}}%
\pgfpathlineto{\pgfqpoint{1.564207in}{1.585980in}}%
\pgfpathlineto{\pgfqpoint{1.564441in}{1.591613in}}%
\pgfpathlineto{\pgfqpoint{1.570886in}{1.748839in}}%
\pgfpathlineto{\pgfqpoint{1.571248in}{1.743250in}}%
\pgfpathlineto{\pgfqpoint{1.580692in}{1.587703in}}%
\pgfpathlineto{\pgfqpoint{1.580905in}{1.592820in}}%
\pgfpathlineto{\pgfqpoint{1.587379in}{1.750843in}}%
\pgfpathlineto{\pgfqpoint{1.587735in}{1.745319in}}%
\pgfpathlineto{\pgfqpoint{1.597176in}{1.590116in}}%
\pgfpathlineto{\pgfqpoint{1.597465in}{1.597062in}}%
\pgfpathlineto{\pgfqpoint{1.603855in}{1.753175in}}%
\pgfpathlineto{\pgfqpoint{1.604216in}{1.747618in}}%
\pgfpathlineto{\pgfqpoint{1.613661in}{1.592177in}}%
\pgfpathlineto{\pgfqpoint{1.613932in}{1.598708in}}%
\pgfpathlineto{\pgfqpoint{1.620340in}{1.755169in}}%
\pgfpathlineto{\pgfqpoint{1.620702in}{1.749594in}}%
\pgfpathlineto{\pgfqpoint{1.630146in}{1.593989in}}%
\pgfpathlineto{\pgfqpoint{1.630416in}{1.600476in}}%
\pgfpathlineto{\pgfqpoint{1.636825in}{1.756810in}}%
\pgfpathlineto{\pgfqpoint{1.637187in}{1.751222in}}%
\pgfpathlineto{\pgfqpoint{1.646631in}{1.595326in}}%
\pgfpathlineto{\pgfqpoint{1.646871in}{1.601041in}}%
\pgfpathlineto{\pgfqpoint{1.653310in}{1.757881in}}%
\pgfpathlineto{\pgfqpoint{1.653672in}{1.752278in}}%
\pgfpathlineto{\pgfqpoint{1.663116in}{1.596017in}}%
\pgfpathlineto{\pgfqpoint{1.663399in}{1.602820in}}%
\pgfpathlineto{\pgfqpoint{1.669793in}{1.758296in}}%
\pgfpathlineto{\pgfqpoint{1.670139in}{1.752949in}}%
\pgfpathlineto{\pgfqpoint{1.679601in}{1.596017in}}%
\pgfpathlineto{\pgfqpoint{1.679917in}{1.603611in}}%
\pgfpathlineto{\pgfqpoint{1.686280in}{1.758066in}}%
\pgfpathlineto{\pgfqpoint{1.686642in}{1.752432in}}%
\pgfpathlineto{\pgfqpoint{1.696086in}{1.595431in}}%
\pgfpathlineto{\pgfqpoint{1.696353in}{1.601831in}}%
\pgfpathlineto{\pgfqpoint{1.702765in}{1.757246in}}%
\pgfpathlineto{\pgfqpoint{1.703126in}{1.751605in}}%
\pgfpathlineto{\pgfqpoint{1.712571in}{1.594317in}}%
\pgfpathlineto{\pgfqpoint{1.712839in}{1.600727in}}%
\pgfpathlineto{\pgfqpoint{1.719254in}{1.756036in}}%
\pgfpathlineto{\pgfqpoint{1.719634in}{1.750067in}}%
\pgfpathlineto{\pgfqpoint{1.729059in}{1.593090in}}%
\pgfpathlineto{\pgfqpoint{1.729334in}{1.599714in}}%
\pgfpathlineto{\pgfqpoint{1.735735in}{1.754640in}}%
\pgfpathlineto{\pgfqpoint{1.736029in}{1.750136in}}%
\pgfpathlineto{\pgfqpoint{1.745540in}{1.592003in}}%
\pgfpathlineto{\pgfqpoint{1.745821in}{1.598707in}}%
\pgfpathlineto{\pgfqpoint{1.752241in}{1.753451in}}%
\pgfpathlineto{\pgfqpoint{1.752564in}{1.748076in}}%
\pgfpathlineto{\pgfqpoint{1.762025in}{1.590609in}}%
\pgfpathlineto{\pgfqpoint{1.762292in}{1.596997in}}%
\pgfpathlineto{\pgfqpoint{1.768704in}{1.752206in}}%
\pgfpathlineto{\pgfqpoint{1.769066in}{1.746543in}}%
\pgfpathlineto{\pgfqpoint{1.778510in}{1.589155in}}%
\pgfpathlineto{\pgfqpoint{1.778782in}{1.595656in}}%
\pgfpathlineto{\pgfqpoint{1.785205in}{1.750768in}}%
\pgfpathlineto{\pgfqpoint{1.785547in}{1.745084in}}%
\pgfpathlineto{\pgfqpoint{1.794995in}{1.587788in}}%
\pgfpathlineto{\pgfqpoint{1.795285in}{1.594692in}}%
\pgfpathlineto{\pgfqpoint{1.801681in}{1.749748in}}%
\pgfpathlineto{\pgfqpoint{1.802050in}{1.743952in}}%
\pgfpathlineto{\pgfqpoint{1.811481in}{1.587826in}}%
\pgfpathlineto{\pgfqpoint{1.811768in}{1.594736in}}%
\pgfpathlineto{\pgfqpoint{1.818177in}{1.750059in}}%
\pgfpathlineto{\pgfqpoint{1.818476in}{1.745414in}}%
\pgfpathlineto{\pgfqpoint{1.827967in}{1.588311in}}%
\pgfpathlineto{\pgfqpoint{1.828232in}{1.594714in}}%
\pgfpathlineto{\pgfqpoint{1.834644in}{1.750575in}}%
\pgfpathlineto{\pgfqpoint{1.834940in}{1.746074in}}%
\pgfpathlineto{\pgfqpoint{1.844450in}{1.589332in}}%
\pgfpathlineto{\pgfqpoint{1.844682in}{1.594891in}}%
\pgfpathlineto{\pgfqpoint{1.851151in}{1.751685in}}%
\pgfpathlineto{\pgfqpoint{1.851562in}{1.744894in}}%
\pgfpathlineto{\pgfqpoint{1.860943in}{1.589680in}}%
\pgfpathlineto{\pgfqpoint{1.861160in}{1.594897in}}%
\pgfpathlineto{\pgfqpoint{1.867618in}{1.752020in}}%
\pgfpathlineto{\pgfqpoint{1.867996in}{1.746131in}}%
\pgfpathlineto{\pgfqpoint{1.877419in}{1.592005in}}%
\pgfpathlineto{\pgfqpoint{1.877688in}{1.598466in}}%
\pgfpathlineto{\pgfqpoint{1.884098in}{1.754485in}}%
\pgfpathlineto{\pgfqpoint{1.884468in}{1.748755in}}%
\pgfpathlineto{\pgfqpoint{1.893904in}{1.592598in}}%
\pgfpathlineto{\pgfqpoint{1.894177in}{1.599168in}}%
\pgfpathlineto{\pgfqpoint{1.900584in}{1.754999in}}%
\pgfpathlineto{\pgfqpoint{1.900877in}{1.750550in}}%
\pgfpathlineto{\pgfqpoint{1.910389in}{1.593337in}}%
\pgfpathlineto{\pgfqpoint{1.910668in}{1.600052in}}%
\pgfpathlineto{\pgfqpoint{1.917070in}{1.755658in}}%
\pgfpathlineto{\pgfqpoint{1.917469in}{1.749419in}}%
\pgfpathlineto{\pgfqpoint{1.926873in}{1.594176in}}%
\pgfpathlineto{\pgfqpoint{1.927119in}{1.600041in}}%
\pgfpathlineto{\pgfqpoint{1.933552in}{1.756310in}}%
\pgfpathlineto{\pgfqpoint{1.933916in}{1.750665in}}%
\pgfpathlineto{\pgfqpoint{1.943358in}{1.593869in}}%
\pgfpathlineto{\pgfqpoint{1.943620in}{1.600137in}}%
\pgfpathlineto{\pgfqpoint{1.950037in}{1.755856in}}%
\pgfpathlineto{\pgfqpoint{1.950399in}{1.750224in}}%
\pgfpathlineto{\pgfqpoint{1.959842in}{1.593224in}}%
\pgfpathlineto{\pgfqpoint{1.960114in}{1.599709in}}%
\pgfpathlineto{\pgfqpoint{1.966522in}{1.755100in}}%
\pgfpathlineto{\pgfqpoint{1.966884in}{1.749464in}}%
\pgfpathlineto{\pgfqpoint{1.976328in}{1.592342in}}%
\pgfpathlineto{\pgfqpoint{1.976599in}{1.598830in}}%
\pgfpathlineto{\pgfqpoint{1.983007in}{1.754162in}}%
\pgfpathlineto{\pgfqpoint{1.983381in}{1.748325in}}%
\pgfpathlineto{\pgfqpoint{1.992813in}{1.591356in}}%
\pgfpathlineto{\pgfqpoint{1.993081in}{1.597760in}}%
\pgfpathlineto{\pgfqpoint{1.999489in}{1.753124in}}%
\pgfpathlineto{\pgfqpoint{1.999859in}{1.747360in}}%
\pgfpathlineto{\pgfqpoint{2.009298in}{1.590373in}}%
\pgfpathlineto{\pgfqpoint{2.009543in}{1.596232in}}%
\pgfpathlineto{\pgfqpoint{2.015977in}{1.752233in}}%
\pgfpathlineto{\pgfqpoint{2.016337in}{1.746616in}}%
\pgfpathlineto{\pgfqpoint{2.025783in}{1.589564in}}%
\pgfpathlineto{\pgfqpoint{2.026093in}{1.597008in}}%
\pgfpathlineto{\pgfqpoint{2.032461in}{1.751513in}}%
\pgfpathlineto{\pgfqpoint{2.032823in}{1.745890in}}%
\pgfpathlineto{\pgfqpoint{2.042268in}{1.588991in}}%
\pgfpathlineto{\pgfqpoint{2.042557in}{1.595933in}}%
\pgfpathlineto{\pgfqpoint{2.048946in}{1.751057in}}%
\pgfpathlineto{\pgfqpoint{2.049309in}{1.745422in}}%
\pgfpathlineto{\pgfqpoint{2.058752in}{1.588724in}}%
\pgfpathlineto{\pgfqpoint{2.059027in}{1.595300in}}%
\pgfpathlineto{\pgfqpoint{2.065431in}{1.750920in}}%
\pgfpathlineto{\pgfqpoint{2.065793in}{1.745306in}}%
\pgfpathlineto{\pgfqpoint{2.075238in}{1.588788in}}%
\pgfpathlineto{\pgfqpoint{2.075520in}{1.595587in}}%
\pgfpathlineto{\pgfqpoint{2.081916in}{1.751114in}}%
\pgfpathlineto{\pgfqpoint{2.082278in}{1.745509in}}%
\pgfpathlineto{\pgfqpoint{2.091722in}{1.589140in}}%
\pgfpathlineto{\pgfqpoint{2.091996in}{1.595725in}}%
\pgfpathlineto{\pgfqpoint{2.098417in}{1.751724in}}%
\pgfpathlineto{\pgfqpoint{2.098828in}{1.745189in}}%
\pgfpathlineto{\pgfqpoint{2.108207in}{1.592389in}}%
\pgfpathlineto{\pgfqpoint{2.108400in}{1.597002in}}%
\pgfpathlineto{\pgfqpoint{2.114886in}{1.754836in}}%
\pgfpathlineto{\pgfqpoint{2.115247in}{1.749234in}}%
\pgfpathlineto{\pgfqpoint{2.124692in}{1.592897in}}%
\pgfpathlineto{\pgfqpoint{2.124931in}{1.598631in}}%
\pgfpathlineto{\pgfqpoint{2.131372in}{1.755258in}}%
\pgfpathlineto{\pgfqpoint{2.131664in}{1.750823in}}%
\pgfpathlineto{\pgfqpoint{2.141177in}{1.593507in}}%
\pgfpathlineto{\pgfqpoint{2.141447in}{1.599978in}}%
\pgfpathlineto{\pgfqpoint{2.147855in}{1.755728in}}%
\pgfpathlineto{\pgfqpoint{2.148220in}{1.750057in}}%
\pgfpathlineto{\pgfqpoint{2.157662in}{1.593439in}}%
\pgfpathlineto{\pgfqpoint{2.157974in}{1.600942in}}%
\pgfpathlineto{\pgfqpoint{2.164340in}{1.755537in}}%
\pgfpathlineto{\pgfqpoint{2.164702in}{1.749916in}}%
\pgfpathlineto{\pgfqpoint{2.174146in}{1.593080in}}%
\pgfpathlineto{\pgfqpoint{2.174438in}{1.600091in}}%
\pgfpathlineto{\pgfqpoint{2.180825in}{1.755076in}}%
\pgfpathlineto{\pgfqpoint{2.181187in}{1.749443in}}%
\pgfpathlineto{\pgfqpoint{2.190631in}{1.592483in}}%
\pgfpathlineto{\pgfqpoint{2.190902in}{1.598959in}}%
\pgfpathlineto{\pgfqpoint{2.197320in}{1.754543in}}%
\pgfpathlineto{\pgfqpoint{2.197656in}{1.749291in}}%
\pgfpathlineto{\pgfqpoint{2.207117in}{1.592835in}}%
\pgfpathlineto{\pgfqpoint{2.207416in}{1.600019in}}%
\pgfpathlineto{\pgfqpoint{2.213810in}{1.754868in}}%
\pgfpathlineto{\pgfqpoint{2.214102in}{1.750291in}}%
\pgfpathlineto{\pgfqpoint{2.223600in}{1.592333in}}%
\pgfpathlineto{\pgfqpoint{2.223850in}{1.598298in}}%
\pgfpathlineto{\pgfqpoint{2.230280in}{1.754159in}}%
\pgfpathlineto{\pgfqpoint{2.230643in}{1.748497in}}%
\pgfpathlineto{\pgfqpoint{2.240086in}{1.591348in}}%
\pgfpathlineto{\pgfqpoint{2.240398in}{1.598836in}}%
\pgfpathlineto{\pgfqpoint{2.246765in}{1.753167in}}%
\pgfpathlineto{\pgfqpoint{2.247151in}{1.747132in}}%
\pgfpathlineto{\pgfqpoint{2.256571in}{1.590398in}}%
\pgfpathlineto{\pgfqpoint{2.256845in}{1.596959in}}%
\pgfpathlineto{\pgfqpoint{2.263249in}{1.752253in}}%
\pgfpathlineto{\pgfqpoint{2.263611in}{1.746623in}}%
\pgfpathlineto{\pgfqpoint{2.273055in}{1.589568in}}%
\pgfpathlineto{\pgfqpoint{2.273351in}{1.596664in}}%
\pgfpathlineto{\pgfqpoint{2.279735in}{1.751515in}}%
\pgfpathlineto{\pgfqpoint{2.280108in}{1.745687in}}%
\pgfpathlineto{\pgfqpoint{2.289540in}{1.588991in}}%
\pgfpathlineto{\pgfqpoint{2.289851in}{1.596464in}}%
\pgfpathlineto{\pgfqpoint{2.296219in}{1.751051in}}%
\pgfpathlineto{\pgfqpoint{2.296582in}{1.745414in}}%
\pgfpathlineto{\pgfqpoint{2.306025in}{1.588708in}}%
\pgfpathlineto{\pgfqpoint{2.306298in}{1.595260in}}%
\pgfpathlineto{\pgfqpoint{2.312704in}{1.750900in}}%
\pgfpathlineto{\pgfqpoint{2.313067in}{1.745256in}}%
\pgfpathlineto{\pgfqpoint{2.322509in}{1.588756in}}%
\pgfpathlineto{\pgfqpoint{2.322771in}{1.595025in}}%
\pgfpathlineto{\pgfqpoint{2.329189in}{1.751077in}}%
\pgfpathlineto{\pgfqpoint{2.329550in}{1.745468in}}%
\pgfpathlineto{\pgfqpoint{2.338995in}{1.589099in}}%
\pgfpathlineto{\pgfqpoint{2.339269in}{1.595687in}}%
\pgfpathlineto{\pgfqpoint{2.345679in}{1.751593in}}%
\pgfpathlineto{\pgfqpoint{2.346059in}{1.745657in}}%
\pgfpathlineto{\pgfqpoint{2.355482in}{1.590096in}}%
\pgfpathlineto{\pgfqpoint{2.355787in}{1.597484in}}%
\pgfpathlineto{\pgfqpoint{2.362158in}{1.752563in}}%
\pgfpathlineto{\pgfqpoint{2.362520in}{1.746967in}}%
\pgfpathlineto{\pgfqpoint{2.371965in}{1.590775in}}%
\pgfpathlineto{\pgfqpoint{2.372183in}{1.596010in}}%
\pgfpathlineto{\pgfqpoint{2.378643in}{1.753273in}}%
\pgfpathlineto{\pgfqpoint{2.379005in}{1.747671in}}%
\pgfpathlineto{\pgfqpoint{2.388449in}{1.591466in}}%
\pgfpathlineto{\pgfqpoint{2.388720in}{1.597976in}}%
\pgfpathlineto{\pgfqpoint{2.395133in}{1.754010in}}%
\pgfpathlineto{\pgfqpoint{2.395513in}{1.748079in}}%
\pgfpathlineto{\pgfqpoint{2.404937in}{1.592297in}}%
\pgfpathlineto{\pgfqpoint{2.405229in}{1.599376in}}%
\pgfpathlineto{\pgfqpoint{2.411613in}{1.754693in}}%
\pgfpathlineto{\pgfqpoint{2.411991in}{1.748816in}}%
\pgfpathlineto{\pgfqpoint{2.421421in}{1.592711in}}%
\pgfpathlineto{\pgfqpoint{2.421690in}{1.599201in}}%
\pgfpathlineto{\pgfqpoint{2.428105in}{1.755127in}}%
\pgfpathlineto{\pgfqpoint{2.428476in}{1.749330in}}%
\pgfpathlineto{\pgfqpoint{2.437907in}{1.593937in}}%
\pgfpathlineto{\pgfqpoint{2.438122in}{1.599122in}}%
\pgfpathlineto{\pgfqpoint{2.444590in}{1.756232in}}%
\pgfpathlineto{\pgfqpoint{2.444961in}{1.750413in}}%
\pgfpathlineto{\pgfqpoint{2.454391in}{1.595595in}}%
\pgfpathlineto{\pgfqpoint{2.454607in}{1.600793in}}%
\pgfpathlineto{\pgfqpoint{2.461071in}{1.757665in}}%
\pgfpathlineto{\pgfqpoint{2.461496in}{1.750979in}}%
\pgfpathlineto{\pgfqpoint{2.470872in}{1.595081in}}%
\pgfpathlineto{\pgfqpoint{2.471160in}{1.601934in}}%
\pgfpathlineto{\pgfqpoint{2.477552in}{1.756899in}}%
\pgfpathlineto{\pgfqpoint{2.477914in}{1.751262in}}%
\pgfpathlineto{\pgfqpoint{2.487358in}{1.593975in}}%
\pgfpathlineto{\pgfqpoint{2.487625in}{1.600356in}}%
\pgfpathlineto{\pgfqpoint{2.494037in}{1.755637in}}%
\pgfpathlineto{\pgfqpoint{2.494399in}{1.749986in}}%
\pgfpathlineto{\pgfqpoint{2.503844in}{1.592566in}}%
\pgfpathlineto{\pgfqpoint{2.504110in}{1.598946in}}%
\pgfpathlineto{\pgfqpoint{2.510522in}{1.754169in}}%
\pgfpathlineto{\pgfqpoint{2.510886in}{1.748489in}}%
\pgfpathlineto{\pgfqpoint{2.520327in}{1.591054in}}%
\pgfpathlineto{\pgfqpoint{2.520624in}{1.598134in}}%
\pgfpathlineto{\pgfqpoint{2.527007in}{1.752666in}}%
\pgfpathlineto{\pgfqpoint{2.527369in}{1.747018in}}%
\pgfpathlineto{\pgfqpoint{2.536813in}{1.589623in}}%
\pgfpathlineto{\pgfqpoint{2.537087in}{1.596178in}}%
\pgfpathlineto{\pgfqpoint{2.543495in}{1.751377in}}%
\pgfpathlineto{\pgfqpoint{2.543922in}{1.744647in}}%
\pgfpathlineto{\pgfqpoint{2.553298in}{1.588572in}}%
\pgfpathlineto{\pgfqpoint{2.553612in}{1.596120in}}%
\pgfpathlineto{\pgfqpoint{2.559982in}{1.750496in}}%
\pgfpathlineto{\pgfqpoint{2.560362in}{1.744522in}}%
\pgfpathlineto{\pgfqpoint{2.569782in}{1.588385in}}%
\pgfpathlineto{\pgfqpoint{2.570078in}{1.595462in}}%
\pgfpathlineto{\pgfqpoint{2.576464in}{1.750441in}}%
\pgfpathlineto{\pgfqpoint{2.576764in}{1.745841in}}%
\pgfpathlineto{\pgfqpoint{2.586268in}{1.588740in}}%
\pgfpathlineto{\pgfqpoint{2.586472in}{1.593624in}}%
\pgfpathlineto{\pgfqpoint{2.592946in}{1.750913in}}%
\pgfpathlineto{\pgfqpoint{2.593308in}{1.745303in}}%
\pgfpathlineto{\pgfqpoint{2.602752in}{1.588733in}}%
\pgfpathlineto{\pgfqpoint{2.603052in}{1.595942in}}%
\pgfpathlineto{\pgfqpoint{2.609432in}{1.751038in}}%
\pgfpathlineto{\pgfqpoint{2.609805in}{1.745232in}}%
\pgfpathlineto{\pgfqpoint{2.619239in}{1.589048in}}%
\pgfpathlineto{\pgfqpoint{2.619543in}{1.596389in}}%
\pgfpathlineto{\pgfqpoint{2.625916in}{1.751450in}}%
\pgfpathlineto{\pgfqpoint{2.626281in}{1.745792in}}%
\pgfpathlineto{\pgfqpoint{2.635722in}{1.589587in}}%
\pgfpathlineto{\pgfqpoint{2.635995in}{1.596136in}}%
\pgfpathlineto{\pgfqpoint{2.642401in}{1.752073in}}%
\pgfpathlineto{\pgfqpoint{2.642763in}{1.746472in}}%
\pgfpathlineto{\pgfqpoint{2.652206in}{1.590303in}}%
\pgfpathlineto{\pgfqpoint{2.652438in}{1.595833in}}%
\pgfpathlineto{\pgfqpoint{2.658891in}{1.752901in}}%
\pgfpathlineto{\pgfqpoint{2.659271in}{1.746968in}}%
\pgfpathlineto{\pgfqpoint{2.668694in}{1.591354in}}%
\pgfpathlineto{\pgfqpoint{2.668987in}{1.598434in}}%
\pgfpathlineto{\pgfqpoint{2.675371in}{1.753854in}}%
\pgfpathlineto{\pgfqpoint{2.675741in}{1.748108in}}%
\pgfpathlineto{\pgfqpoint{2.685177in}{1.592045in}}%
\pgfpathlineto{\pgfqpoint{2.685413in}{1.597703in}}%
\pgfpathlineto{\pgfqpoint{2.691855in}{1.754498in}}%
\pgfpathlineto{\pgfqpoint{2.692217in}{1.748893in}}%
\pgfpathlineto{\pgfqpoint{2.701663in}{1.592587in}}%
\pgfpathlineto{\pgfqpoint{2.701942in}{1.599326in}}%
\pgfpathlineto{\pgfqpoint{2.708340in}{1.754952in}}%
\pgfpathlineto{\pgfqpoint{2.708703in}{1.749335in}}%
\pgfpathlineto{\pgfqpoint{2.718146in}{1.592899in}}%
\pgfpathlineto{\pgfqpoint{2.718418in}{1.599411in}}%
\pgfpathlineto{\pgfqpoint{2.724833in}{1.755065in}}%
\pgfpathlineto{\pgfqpoint{2.725270in}{1.747906in}}%
\pgfpathlineto{\pgfqpoint{2.734631in}{1.592707in}}%
\pgfpathlineto{\pgfqpoint{2.734898in}{1.599117in}}%
\pgfpathlineto{\pgfqpoint{2.741315in}{1.754959in}}%
\pgfpathlineto{\pgfqpoint{2.741694in}{1.749023in}}%
\pgfpathlineto{\pgfqpoint{2.751116in}{1.592802in}}%
\pgfpathlineto{\pgfqpoint{2.751415in}{1.599971in}}%
\pgfpathlineto{\pgfqpoint{2.757795in}{1.754899in}}%
\pgfpathlineto{\pgfqpoint{2.758160in}{1.749227in}}%
\pgfpathlineto{\pgfqpoint{2.767601in}{1.592456in}}%
\pgfpathlineto{\pgfqpoint{2.767868in}{1.598856in}}%
\pgfpathlineto{\pgfqpoint{2.774299in}{1.754596in}}%
\pgfpathlineto{\pgfqpoint{2.774663in}{1.748843in}}%
\pgfpathlineto{\pgfqpoint{2.784088in}{1.592443in}}%
\pgfpathlineto{\pgfqpoint{2.784335in}{1.598393in}}%
\pgfpathlineto{\pgfqpoint{2.790765in}{1.754399in}}%
\pgfpathlineto{\pgfqpoint{2.791126in}{1.748775in}}%
\pgfpathlineto{\pgfqpoint{2.800571in}{1.591780in}}%
\pgfpathlineto{\pgfqpoint{2.800797in}{1.597170in}}%
\pgfpathlineto{\pgfqpoint{2.807249in}{1.753703in}}%
\pgfpathlineto{\pgfqpoint{2.807611in}{1.748072in}}%
\pgfpathlineto{\pgfqpoint{2.817055in}{1.591074in}}%
\pgfpathlineto{\pgfqpoint{2.817327in}{1.597575in}}%
\pgfpathlineto{\pgfqpoint{2.823734in}{1.753011in}}%
\pgfpathlineto{\pgfqpoint{2.824096in}{1.747378in}}%
\pgfpathlineto{\pgfqpoint{2.833540in}{1.590421in}}%
\pgfpathlineto{\pgfqpoint{2.833813in}{1.596956in}}%
\pgfpathlineto{\pgfqpoint{2.840219in}{1.752400in}}%
\pgfpathlineto{\pgfqpoint{2.840581in}{1.746772in}}%
\pgfpathlineto{\pgfqpoint{2.850025in}{1.589901in}}%
\pgfpathlineto{\pgfqpoint{2.850266in}{1.595669in}}%
\pgfpathlineto{\pgfqpoint{2.856704in}{1.751957in}}%
\pgfpathlineto{\pgfqpoint{2.857066in}{1.746332in}}%
\pgfpathlineto{\pgfqpoint{2.866510in}{1.589564in}}%
\pgfpathlineto{\pgfqpoint{2.866775in}{1.595905in}}%
\pgfpathlineto{\pgfqpoint{2.873204in}{1.751647in}}%
\pgfpathlineto{\pgfqpoint{2.873531in}{1.746264in}}%
\pgfpathlineto{\pgfqpoint{2.882995in}{1.590967in}}%
\pgfpathlineto{\pgfqpoint{2.883238in}{1.596781in}}%
\pgfpathlineto{\pgfqpoint{2.889677in}{1.753223in}}%
\pgfpathlineto{\pgfqpoint{2.890103in}{1.746535in}}%
\pgfpathlineto{\pgfqpoint{2.899477in}{1.591039in}}%
\pgfpathlineto{\pgfqpoint{2.899717in}{1.596688in}}%
\pgfpathlineto{\pgfqpoint{2.906162in}{1.753264in}}%
\pgfpathlineto{\pgfqpoint{2.906568in}{1.746894in}}%
\pgfpathlineto{\pgfqpoint{2.915965in}{1.592309in}}%
\pgfpathlineto{\pgfqpoint{2.916250in}{1.599156in}}%
\pgfpathlineto{\pgfqpoint{2.922659in}{1.754658in}}%
\pgfpathlineto{\pgfqpoint{2.922938in}{1.750305in}}%
\pgfpathlineto{\pgfqpoint{2.932450in}{1.593293in}}%
\pgfpathlineto{\pgfqpoint{2.932668in}{1.598499in}}%
\pgfpathlineto{\pgfqpoint{2.939128in}{1.755402in}}%
\pgfpathlineto{\pgfqpoint{2.939489in}{1.749789in}}%
\pgfpathlineto{\pgfqpoint{2.948934in}{1.592957in}}%
\pgfpathlineto{\pgfqpoint{2.949240in}{1.600293in}}%
\pgfpathlineto{\pgfqpoint{2.955613in}{1.754965in}}%
\pgfpathlineto{\pgfqpoint{2.955976in}{1.749321in}}%
\pgfpathlineto{\pgfqpoint{2.965419in}{1.592396in}}%
\pgfpathlineto{\pgfqpoint{2.965667in}{1.598327in}}%
\pgfpathlineto{\pgfqpoint{2.972098in}{1.754336in}}%
\pgfpathlineto{\pgfqpoint{2.972460in}{1.748701in}}%
\pgfpathlineto{\pgfqpoint{2.981904in}{1.591700in}}%
\pgfpathlineto{\pgfqpoint{2.982172in}{1.598120in}}%
\pgfpathlineto{\pgfqpoint{2.988582in}{1.753609in}}%
\pgfpathlineto{\pgfqpoint{2.988936in}{1.748117in}}%
\pgfpathlineto{\pgfqpoint{2.998391in}{1.590959in}}%
\pgfpathlineto{\pgfqpoint{2.998618in}{1.596409in}}%
\pgfpathlineto{\pgfqpoint{3.005068in}{1.752874in}}%
\pgfpathlineto{\pgfqpoint{3.005430in}{1.747237in}}%
\pgfpathlineto{\pgfqpoint{3.014874in}{1.590274in}}%
\pgfpathlineto{\pgfqpoint{3.015145in}{1.596772in}}%
\pgfpathlineto{\pgfqpoint{3.021566in}{1.752138in}}%
\pgfpathlineto{\pgfqpoint{3.021941in}{1.745917in}}%
\pgfpathlineto{\pgfqpoint{3.031356in}{1.590883in}}%
\pgfpathlineto{\pgfqpoint{3.031555in}{1.595589in}}%
\pgfpathlineto{\pgfqpoint{3.038037in}{1.752957in}}%
\pgfpathlineto{\pgfqpoint{3.038398in}{1.747345in}}%
\pgfpathlineto{\pgfqpoint{3.047843in}{1.590518in}}%
\pgfpathlineto{\pgfqpoint{3.048109in}{1.596892in}}%
\pgfpathlineto{\pgfqpoint{3.054523in}{1.752602in}}%
\pgfpathlineto{\pgfqpoint{3.054897in}{1.746757in}}%
\pgfpathlineto{\pgfqpoint{3.064328in}{1.590231in}}%
\pgfpathlineto{\pgfqpoint{3.064593in}{1.596570in}}%
\pgfpathlineto{\pgfqpoint{3.071023in}{1.752520in}}%
\pgfpathlineto{\pgfqpoint{3.071461in}{1.745523in}}%
\pgfpathlineto{\pgfqpoint{3.080814in}{1.591900in}}%
\pgfpathlineto{\pgfqpoint{3.081039in}{1.597293in}}%
\pgfpathlineto{\pgfqpoint{3.087492in}{1.754049in}}%
\pgfpathlineto{\pgfqpoint{3.087859in}{1.748331in}}%
\pgfpathlineto{\pgfqpoint{3.097298in}{1.591715in}}%
\pgfpathlineto{\pgfqpoint{3.097603in}{1.599053in}}%
\pgfpathlineto{\pgfqpoint{3.103992in}{1.753991in}}%
\pgfpathlineto{\pgfqpoint{3.104430in}{1.747001in}}%
\pgfpathlineto{\pgfqpoint{3.113782in}{1.593181in}}%
\pgfpathlineto{\pgfqpoint{3.113974in}{1.597733in}}%
\pgfpathlineto{\pgfqpoint{3.120478in}{1.755166in}}%
\pgfpathlineto{\pgfqpoint{3.120758in}{1.750528in}}%
\pgfpathlineto{\pgfqpoint{3.130266in}{1.593419in}}%
\pgfpathlineto{\pgfqpoint{3.130487in}{1.598686in}}%
\pgfpathlineto{\pgfqpoint{3.136946in}{1.755415in}}%
\pgfpathlineto{\pgfqpoint{3.137308in}{1.749793in}}%
\pgfpathlineto{\pgfqpoint{3.146752in}{1.592777in}}%
\pgfpathlineto{\pgfqpoint{3.147022in}{1.599228in}}%
\pgfpathlineto{\pgfqpoint{3.153448in}{1.754567in}}%
\pgfpathlineto{\pgfqpoint{3.153789in}{1.748890in}}%
\pgfpathlineto{\pgfqpoint{3.163238in}{1.591695in}}%
\pgfpathlineto{\pgfqpoint{3.163510in}{1.598197in}}%
\pgfpathlineto{\pgfqpoint{3.169916in}{1.753517in}}%
\pgfpathlineto{\pgfqpoint{3.170277in}{1.747892in}}%
\pgfpathlineto{\pgfqpoint{3.179722in}{1.590791in}}%
\pgfpathlineto{\pgfqpoint{3.179995in}{1.597316in}}%
\pgfpathlineto{\pgfqpoint{3.186416in}{1.752547in}}%
\pgfpathlineto{\pgfqpoint{3.186716in}{1.747572in}}%
\pgfpathlineto{\pgfqpoint{3.196207in}{1.590159in}}%
\pgfpathlineto{\pgfqpoint{3.196504in}{1.597304in}}%
\pgfpathlineto{\pgfqpoint{3.202886in}{1.752139in}}%
\pgfpathlineto{\pgfqpoint{3.203250in}{1.746463in}}%
\pgfpathlineto{\pgfqpoint{3.212692in}{1.589602in}}%
\pgfpathlineto{\pgfqpoint{3.213003in}{1.597069in}}%
\pgfpathlineto{\pgfqpoint{3.219374in}{1.751679in}}%
\pgfpathlineto{\pgfqpoint{3.219772in}{1.745430in}}%
\pgfpathlineto{\pgfqpoint{3.229161in}{1.591706in}}%
\pgfpathlineto{\pgfqpoint{3.229161in}{1.591706in}}%
\pgfusepath{stroke}%
\end{pgfscope}%
\begin{pgfscope}%
\pgfsetrectcap%
\pgfsetmiterjoin%
\pgfsetlinewidth{0.803000pt}%
\definecolor{currentstroke}{rgb}{0.000000,0.000000,0.000000}%
\pgfsetstrokecolor{currentstroke}%
\pgfsetdash{}{0pt}%
\pgfpathmoveto{\pgfqpoint{0.632797in}{1.532029in}}%
\pgfpathlineto{\pgfqpoint{0.632797in}{2.438279in}}%
\pgfusepath{stroke}%
\end{pgfscope}%
\begin{pgfscope}%
\pgfsetrectcap%
\pgfsetmiterjoin%
\pgfsetlinewidth{0.803000pt}%
\definecolor{currentstroke}{rgb}{0.000000,0.000000,0.000000}%
\pgfsetstrokecolor{currentstroke}%
\pgfsetdash{}{0pt}%
\pgfpathmoveto{\pgfqpoint{3.352797in}{1.532029in}}%
\pgfpathlineto{\pgfqpoint{3.352797in}{2.438279in}}%
\pgfusepath{stroke}%
\end{pgfscope}%
\begin{pgfscope}%
\pgfsetrectcap%
\pgfsetmiterjoin%
\pgfsetlinewidth{0.803000pt}%
\definecolor{currentstroke}{rgb}{0.000000,0.000000,0.000000}%
\pgfsetstrokecolor{currentstroke}%
\pgfsetdash{}{0pt}%
\pgfpathmoveto{\pgfqpoint{0.632797in}{1.532029in}}%
\pgfpathlineto{\pgfqpoint{3.352797in}{1.532029in}}%
\pgfusepath{stroke}%
\end{pgfscope}%
\begin{pgfscope}%
\pgfsetrectcap%
\pgfsetmiterjoin%
\pgfsetlinewidth{0.803000pt}%
\definecolor{currentstroke}{rgb}{0.000000,0.000000,0.000000}%
\pgfsetstrokecolor{currentstroke}%
\pgfsetdash{}{0pt}%
\pgfpathmoveto{\pgfqpoint{0.632797in}{2.438279in}}%
\pgfpathlineto{\pgfqpoint{3.352797in}{2.438279in}}%
\pgfusepath{stroke}%
\end{pgfscope}%
\begin{pgfscope}%
\definecolor{textcolor}{rgb}{0.000000,0.000000,0.000000}%
\pgfsetstrokecolor{textcolor}%
\pgfsetfillcolor{textcolor}%
\pgftext[x=0.687197in,y=2.374842in,left,top]{\color{textcolor}{\rmfamily\fontsize{8.000000}{9.600000}\selectfont\catcode`\^=\active\def^{\ifmmode\sp\else\^{}\fi}\catcode`\%=\active\def%{\%}(b)}}%
\end{pgfscope}%
\begin{pgfscope}%
\pgfsetbuttcap%
\pgfsetmiterjoin%
\definecolor{currentfill}{rgb}{1.000000,1.000000,1.000000}%
\pgfsetfillcolor{currentfill}%
\pgfsetlinewidth{0.000000pt}%
\definecolor{currentstroke}{rgb}{0.000000,0.000000,0.000000}%
\pgfsetstrokecolor{currentstroke}%
\pgfsetstrokeopacity{0.000000}%
\pgfsetdash{}{0pt}%
\pgfpathmoveto{\pgfqpoint{0.632797in}{0.462654in}}%
\pgfpathlineto{\pgfqpoint{3.352797in}{0.462654in}}%
\pgfpathlineto{\pgfqpoint{3.352797in}{1.368904in}}%
\pgfpathlineto{\pgfqpoint{0.632797in}{1.368904in}}%
\pgfpathlineto{\pgfqpoint{0.632797in}{0.462654in}}%
\pgfpathclose%
\pgfusepath{fill}%
\end{pgfscope}%
\begin{pgfscope}%
\pgfpathrectangle{\pgfqpoint{0.632797in}{0.462654in}}{\pgfqpoint{2.720000in}{0.906250in}}%
\pgfusepath{clip}%
\pgfsetbuttcap%
\pgfsetroundjoin%
\pgfsetlinewidth{0.501875pt}%
\definecolor{currentstroke}{rgb}{0.690196,0.690196,0.690196}%
\pgfsetstrokecolor{currentstroke}%
\pgfsetstrokeopacity{0.250000}%
\pgfsetdash{{1.850000pt}{0.800000pt}}{0.000000pt}%
\pgfpathmoveto{\pgfqpoint{0.756433in}{0.462654in}}%
\pgfpathlineto{\pgfqpoint{0.756433in}{1.368904in}}%
\pgfusepath{stroke}%
\end{pgfscope}%
\begin{pgfscope}%
\pgfsetbuttcap%
\pgfsetroundjoin%
\definecolor{currentfill}{rgb}{0.000000,0.000000,0.000000}%
\pgfsetfillcolor{currentfill}%
\pgfsetlinewidth{0.803000pt}%
\definecolor{currentstroke}{rgb}{0.000000,0.000000,0.000000}%
\pgfsetstrokecolor{currentstroke}%
\pgfsetdash{}{0pt}%
\pgfsys@defobject{currentmarker}{\pgfqpoint{0.000000in}{-0.048611in}}{\pgfqpoint{0.000000in}{0.000000in}}{%
\pgfpathmoveto{\pgfqpoint{0.000000in}{0.000000in}}%
\pgfpathlineto{\pgfqpoint{0.000000in}{-0.048611in}}%
\pgfusepath{stroke,fill}%
}%
\begin{pgfscope}%
\pgfsys@transformshift{0.756433in}{0.462654in}%
\pgfsys@useobject{currentmarker}{}%
\end{pgfscope}%
\end{pgfscope}%
\begin{pgfscope}%
\definecolor{textcolor}{rgb}{0.000000,0.000000,0.000000}%
\pgfsetstrokecolor{textcolor}%
\pgfsetfillcolor{textcolor}%
\pgftext[x=0.756433in,y=0.365432in,,top]{\color{textcolor}{\rmfamily\fontsize{8.000000}{9.600000}\selectfont\catcode`\^=\active\def^{\ifmmode\sp\else\^{}\fi}\catcode`\%=\active\def%{\%}$\mathdefault{0}$}}%
\end{pgfscope}%
\begin{pgfscope}%
\pgfpathrectangle{\pgfqpoint{0.632797in}{0.462654in}}{\pgfqpoint{2.720000in}{0.906250in}}%
\pgfusepath{clip}%
\pgfsetbuttcap%
\pgfsetroundjoin%
\pgfsetlinewidth{0.501875pt}%
\definecolor{currentstroke}{rgb}{0.690196,0.690196,0.690196}%
\pgfsetstrokecolor{currentstroke}%
\pgfsetstrokeopacity{0.250000}%
\pgfsetdash{{1.850000pt}{0.800000pt}}{0.000000pt}%
\pgfpathmoveto{\pgfqpoint{1.415827in}{0.462654in}}%
\pgfpathlineto{\pgfqpoint{1.415827in}{1.368904in}}%
\pgfusepath{stroke}%
\end{pgfscope}%
\begin{pgfscope}%
\pgfsetbuttcap%
\pgfsetroundjoin%
\definecolor{currentfill}{rgb}{0.000000,0.000000,0.000000}%
\pgfsetfillcolor{currentfill}%
\pgfsetlinewidth{0.803000pt}%
\definecolor{currentstroke}{rgb}{0.000000,0.000000,0.000000}%
\pgfsetstrokecolor{currentstroke}%
\pgfsetdash{}{0pt}%
\pgfsys@defobject{currentmarker}{\pgfqpoint{0.000000in}{-0.048611in}}{\pgfqpoint{0.000000in}{0.000000in}}{%
\pgfpathmoveto{\pgfqpoint{0.000000in}{0.000000in}}%
\pgfpathlineto{\pgfqpoint{0.000000in}{-0.048611in}}%
\pgfusepath{stroke,fill}%
}%
\begin{pgfscope}%
\pgfsys@transformshift{1.415827in}{0.462654in}%
\pgfsys@useobject{currentmarker}{}%
\end{pgfscope}%
\end{pgfscope}%
\begin{pgfscope}%
\definecolor{textcolor}{rgb}{0.000000,0.000000,0.000000}%
\pgfsetstrokecolor{textcolor}%
\pgfsetfillcolor{textcolor}%
\pgftext[x=1.415827in,y=0.365432in,,top]{\color{textcolor}{\rmfamily\fontsize{8.000000}{9.600000}\selectfont\catcode`\^=\active\def^{\ifmmode\sp\else\^{}\fi}\catcode`\%=\active\def%{\%}$\mathdefault{2}$}}%
\end{pgfscope}%
\begin{pgfscope}%
\pgfpathrectangle{\pgfqpoint{0.632797in}{0.462654in}}{\pgfqpoint{2.720000in}{0.906250in}}%
\pgfusepath{clip}%
\pgfsetbuttcap%
\pgfsetroundjoin%
\pgfsetlinewidth{0.501875pt}%
\definecolor{currentstroke}{rgb}{0.690196,0.690196,0.690196}%
\pgfsetstrokecolor{currentstroke}%
\pgfsetstrokeopacity{0.250000}%
\pgfsetdash{{1.850000pt}{0.800000pt}}{0.000000pt}%
\pgfpathmoveto{\pgfqpoint{2.075221in}{0.462654in}}%
\pgfpathlineto{\pgfqpoint{2.075221in}{1.368904in}}%
\pgfusepath{stroke}%
\end{pgfscope}%
\begin{pgfscope}%
\pgfsetbuttcap%
\pgfsetroundjoin%
\definecolor{currentfill}{rgb}{0.000000,0.000000,0.000000}%
\pgfsetfillcolor{currentfill}%
\pgfsetlinewidth{0.803000pt}%
\definecolor{currentstroke}{rgb}{0.000000,0.000000,0.000000}%
\pgfsetstrokecolor{currentstroke}%
\pgfsetdash{}{0pt}%
\pgfsys@defobject{currentmarker}{\pgfqpoint{0.000000in}{-0.048611in}}{\pgfqpoint{0.000000in}{0.000000in}}{%
\pgfpathmoveto{\pgfqpoint{0.000000in}{0.000000in}}%
\pgfpathlineto{\pgfqpoint{0.000000in}{-0.048611in}}%
\pgfusepath{stroke,fill}%
}%
\begin{pgfscope}%
\pgfsys@transformshift{2.075221in}{0.462654in}%
\pgfsys@useobject{currentmarker}{}%
\end{pgfscope}%
\end{pgfscope}%
\begin{pgfscope}%
\definecolor{textcolor}{rgb}{0.000000,0.000000,0.000000}%
\pgfsetstrokecolor{textcolor}%
\pgfsetfillcolor{textcolor}%
\pgftext[x=2.075221in,y=0.365432in,,top]{\color{textcolor}{\rmfamily\fontsize{8.000000}{9.600000}\selectfont\catcode`\^=\active\def^{\ifmmode\sp\else\^{}\fi}\catcode`\%=\active\def%{\%}$\mathdefault{4}$}}%
\end{pgfscope}%
\begin{pgfscope}%
\pgfpathrectangle{\pgfqpoint{0.632797in}{0.462654in}}{\pgfqpoint{2.720000in}{0.906250in}}%
\pgfusepath{clip}%
\pgfsetbuttcap%
\pgfsetroundjoin%
\pgfsetlinewidth{0.501875pt}%
\definecolor{currentstroke}{rgb}{0.690196,0.690196,0.690196}%
\pgfsetstrokecolor{currentstroke}%
\pgfsetstrokeopacity{0.250000}%
\pgfsetdash{{1.850000pt}{0.800000pt}}{0.000000pt}%
\pgfpathmoveto{\pgfqpoint{2.734615in}{0.462654in}}%
\pgfpathlineto{\pgfqpoint{2.734615in}{1.368904in}}%
\pgfusepath{stroke}%
\end{pgfscope}%
\begin{pgfscope}%
\pgfsetbuttcap%
\pgfsetroundjoin%
\definecolor{currentfill}{rgb}{0.000000,0.000000,0.000000}%
\pgfsetfillcolor{currentfill}%
\pgfsetlinewidth{0.803000pt}%
\definecolor{currentstroke}{rgb}{0.000000,0.000000,0.000000}%
\pgfsetstrokecolor{currentstroke}%
\pgfsetdash{}{0pt}%
\pgfsys@defobject{currentmarker}{\pgfqpoint{0.000000in}{-0.048611in}}{\pgfqpoint{0.000000in}{0.000000in}}{%
\pgfpathmoveto{\pgfqpoint{0.000000in}{0.000000in}}%
\pgfpathlineto{\pgfqpoint{0.000000in}{-0.048611in}}%
\pgfusepath{stroke,fill}%
}%
\begin{pgfscope}%
\pgfsys@transformshift{2.734615in}{0.462654in}%
\pgfsys@useobject{currentmarker}{}%
\end{pgfscope}%
\end{pgfscope}%
\begin{pgfscope}%
\definecolor{textcolor}{rgb}{0.000000,0.000000,0.000000}%
\pgfsetstrokecolor{textcolor}%
\pgfsetfillcolor{textcolor}%
\pgftext[x=2.734615in,y=0.365432in,,top]{\color{textcolor}{\rmfamily\fontsize{8.000000}{9.600000}\selectfont\catcode`\^=\active\def^{\ifmmode\sp\else\^{}\fi}\catcode`\%=\active\def%{\%}$\mathdefault{6}$}}%
\end{pgfscope}%
\begin{pgfscope}%
\definecolor{textcolor}{rgb}{0.000000,0.000000,0.000000}%
\pgfsetstrokecolor{textcolor}%
\pgfsetfillcolor{textcolor}%
\pgftext[x=1.992797in,y=0.211111in,,top]{\color{textcolor}{\rmfamily\fontsize{8.000000}{9.600000}\selectfont\catcode`\^=\active\def^{\ifmmode\sp\else\^{}\fi}\catcode`\%=\active\def%{\%}Time (ms)}}%
\end{pgfscope}%
\begin{pgfscope}%
\pgfpathrectangle{\pgfqpoint{0.632797in}{0.462654in}}{\pgfqpoint{2.720000in}{0.906250in}}%
\pgfusepath{clip}%
\pgfsetbuttcap%
\pgfsetroundjoin%
\pgfsetlinewidth{0.501875pt}%
\definecolor{currentstroke}{rgb}{0.690196,0.690196,0.690196}%
\pgfsetstrokecolor{currentstroke}%
\pgfsetstrokeopacity{0.250000}%
\pgfsetdash{{1.850000pt}{0.800000pt}}{0.000000pt}%
\pgfpathmoveto{\pgfqpoint{0.632797in}{0.462654in}}%
\pgfpathlineto{\pgfqpoint{3.352797in}{0.462654in}}%
\pgfusepath{stroke}%
\end{pgfscope}%
\begin{pgfscope}%
\pgfsetbuttcap%
\pgfsetroundjoin%
\definecolor{currentfill}{rgb}{0.000000,0.000000,0.000000}%
\pgfsetfillcolor{currentfill}%
\pgfsetlinewidth{0.803000pt}%
\definecolor{currentstroke}{rgb}{0.000000,0.000000,0.000000}%
\pgfsetstrokecolor{currentstroke}%
\pgfsetdash{}{0pt}%
\pgfsys@defobject{currentmarker}{\pgfqpoint{-0.048611in}{0.000000in}}{\pgfqpoint{-0.000000in}{0.000000in}}{%
\pgfpathmoveto{\pgfqpoint{-0.000000in}{0.000000in}}%
\pgfpathlineto{\pgfqpoint{-0.048611in}{0.000000in}}%
\pgfusepath{stroke,fill}%
}%
\begin{pgfscope}%
\pgfsys@transformshift{0.632797in}{0.462654in}%
\pgfsys@useobject{currentmarker}{}%
\end{pgfscope}%
\end{pgfscope}%
\begin{pgfscope}%
\definecolor{textcolor}{rgb}{0.000000,0.000000,0.000000}%
\pgfsetstrokecolor{textcolor}%
\pgfsetfillcolor{textcolor}%
\pgftext[x=0.286343in, y=0.424074in, left, base]{\color{textcolor}{\rmfamily\fontsize{8.000000}{9.600000}\selectfont\catcode`\^=\active\def^{\ifmmode\sp\else\^{}\fi}\catcode`\%=\active\def%{\%}-26.2}}%
\end{pgfscope}%
\begin{pgfscope}%
\pgfpathrectangle{\pgfqpoint{0.632797in}{0.462654in}}{\pgfqpoint{2.720000in}{0.906250in}}%
\pgfusepath{clip}%
\pgfsetbuttcap%
\pgfsetroundjoin%
\pgfsetlinewidth{0.501875pt}%
\definecolor{currentstroke}{rgb}{0.690196,0.690196,0.690196}%
\pgfsetstrokecolor{currentstroke}%
\pgfsetstrokeopacity{0.250000}%
\pgfsetdash{{1.850000pt}{0.800000pt}}{0.000000pt}%
\pgfpathmoveto{\pgfqpoint{0.632797in}{0.689217in}}%
\pgfpathlineto{\pgfqpoint{3.352797in}{0.689217in}}%
\pgfusepath{stroke}%
\end{pgfscope}%
\begin{pgfscope}%
\pgfsetbuttcap%
\pgfsetroundjoin%
\definecolor{currentfill}{rgb}{0.000000,0.000000,0.000000}%
\pgfsetfillcolor{currentfill}%
\pgfsetlinewidth{0.803000pt}%
\definecolor{currentstroke}{rgb}{0.000000,0.000000,0.000000}%
\pgfsetstrokecolor{currentstroke}%
\pgfsetdash{}{0pt}%
\pgfsys@defobject{currentmarker}{\pgfqpoint{-0.048611in}{0.000000in}}{\pgfqpoint{-0.000000in}{0.000000in}}{%
\pgfpathmoveto{\pgfqpoint{-0.000000in}{0.000000in}}%
\pgfpathlineto{\pgfqpoint{-0.048611in}{0.000000in}}%
\pgfusepath{stroke,fill}%
}%
\begin{pgfscope}%
\pgfsys@transformshift{0.632797in}{0.689217in}%
\pgfsys@useobject{currentmarker}{}%
\end{pgfscope}%
\end{pgfscope}%
\begin{pgfscope}%
\definecolor{textcolor}{rgb}{0.000000,0.000000,0.000000}%
\pgfsetstrokecolor{textcolor}%
\pgfsetfillcolor{textcolor}%
\pgftext[x=0.266667in, y=0.650637in, left, base]{\color{textcolor}{\rmfamily\fontsize{8.000000}{9.600000}\selectfont\catcode`\^=\active\def^{\ifmmode\sp\else\^{}\fi}\catcode`\%=\active\def%{\%}117.9}}%
\end{pgfscope}%
\begin{pgfscope}%
\pgfpathrectangle{\pgfqpoint{0.632797in}{0.462654in}}{\pgfqpoint{2.720000in}{0.906250in}}%
\pgfusepath{clip}%
\pgfsetbuttcap%
\pgfsetroundjoin%
\pgfsetlinewidth{0.501875pt}%
\definecolor{currentstroke}{rgb}{0.690196,0.690196,0.690196}%
\pgfsetstrokecolor{currentstroke}%
\pgfsetstrokeopacity{0.250000}%
\pgfsetdash{{1.850000pt}{0.800000pt}}{0.000000pt}%
\pgfpathmoveto{\pgfqpoint{0.632797in}{0.915779in}}%
\pgfpathlineto{\pgfqpoint{3.352797in}{0.915779in}}%
\pgfusepath{stroke}%
\end{pgfscope}%
\begin{pgfscope}%
\pgfsetbuttcap%
\pgfsetroundjoin%
\definecolor{currentfill}{rgb}{0.000000,0.000000,0.000000}%
\pgfsetfillcolor{currentfill}%
\pgfsetlinewidth{0.803000pt}%
\definecolor{currentstroke}{rgb}{0.000000,0.000000,0.000000}%
\pgfsetstrokecolor{currentstroke}%
\pgfsetdash{}{0pt}%
\pgfsys@defobject{currentmarker}{\pgfqpoint{-0.048611in}{0.000000in}}{\pgfqpoint{-0.000000in}{0.000000in}}{%
\pgfpathmoveto{\pgfqpoint{-0.000000in}{0.000000in}}%
\pgfpathlineto{\pgfqpoint{-0.048611in}{0.000000in}}%
\pgfusepath{stroke,fill}%
}%
\begin{pgfscope}%
\pgfsys@transformshift{0.632797in}{0.915779in}%
\pgfsys@useobject{currentmarker}{}%
\end{pgfscope}%
\end{pgfscope}%
\begin{pgfscope}%
\definecolor{textcolor}{rgb}{0.000000,0.000000,0.000000}%
\pgfsetstrokecolor{textcolor}%
\pgfsetfillcolor{textcolor}%
\pgftext[x=0.266667in, y=0.877199in, left, base]{\color{textcolor}{\rmfamily\fontsize{8.000000}{9.600000}\selectfont\catcode`\^=\active\def^{\ifmmode\sp\else\^{}\fi}\catcode`\%=\active\def%{\%}261.9}}%
\end{pgfscope}%
\begin{pgfscope}%
\pgfpathrectangle{\pgfqpoint{0.632797in}{0.462654in}}{\pgfqpoint{2.720000in}{0.906250in}}%
\pgfusepath{clip}%
\pgfsetbuttcap%
\pgfsetroundjoin%
\pgfsetlinewidth{0.501875pt}%
\definecolor{currentstroke}{rgb}{0.690196,0.690196,0.690196}%
\pgfsetstrokecolor{currentstroke}%
\pgfsetstrokeopacity{0.250000}%
\pgfsetdash{{1.850000pt}{0.800000pt}}{0.000000pt}%
\pgfpathmoveto{\pgfqpoint{0.632797in}{1.142342in}}%
\pgfpathlineto{\pgfqpoint{3.352797in}{1.142342in}}%
\pgfusepath{stroke}%
\end{pgfscope}%
\begin{pgfscope}%
\pgfsetbuttcap%
\pgfsetroundjoin%
\definecolor{currentfill}{rgb}{0.000000,0.000000,0.000000}%
\pgfsetfillcolor{currentfill}%
\pgfsetlinewidth{0.803000pt}%
\definecolor{currentstroke}{rgb}{0.000000,0.000000,0.000000}%
\pgfsetstrokecolor{currentstroke}%
\pgfsetdash{}{0pt}%
\pgfsys@defobject{currentmarker}{\pgfqpoint{-0.048611in}{0.000000in}}{\pgfqpoint{-0.000000in}{0.000000in}}{%
\pgfpathmoveto{\pgfqpoint{-0.000000in}{0.000000in}}%
\pgfpathlineto{\pgfqpoint{-0.048611in}{0.000000in}}%
\pgfusepath{stroke,fill}%
}%
\begin{pgfscope}%
\pgfsys@transformshift{0.632797in}{1.142342in}%
\pgfsys@useobject{currentmarker}{}%
\end{pgfscope}%
\end{pgfscope}%
\begin{pgfscope}%
\definecolor{textcolor}{rgb}{0.000000,0.000000,0.000000}%
\pgfsetstrokecolor{textcolor}%
\pgfsetfillcolor{textcolor}%
\pgftext[x=0.266667in, y=1.103762in, left, base]{\color{textcolor}{\rmfamily\fontsize{8.000000}{9.600000}\selectfont\catcode`\^=\active\def^{\ifmmode\sp\else\^{}\fi}\catcode`\%=\active\def%{\%}406.0}}%
\end{pgfscope}%
\begin{pgfscope}%
\pgfpathrectangle{\pgfqpoint{0.632797in}{0.462654in}}{\pgfqpoint{2.720000in}{0.906250in}}%
\pgfusepath{clip}%
\pgfsetbuttcap%
\pgfsetroundjoin%
\pgfsetlinewidth{0.501875pt}%
\definecolor{currentstroke}{rgb}{0.690196,0.690196,0.690196}%
\pgfsetstrokecolor{currentstroke}%
\pgfsetstrokeopacity{0.250000}%
\pgfsetdash{{1.850000pt}{0.800000pt}}{0.000000pt}%
\pgfpathmoveto{\pgfqpoint{0.632797in}{1.368904in}}%
\pgfpathlineto{\pgfqpoint{3.352797in}{1.368904in}}%
\pgfusepath{stroke}%
\end{pgfscope}%
\begin{pgfscope}%
\pgfsetbuttcap%
\pgfsetroundjoin%
\definecolor{currentfill}{rgb}{0.000000,0.000000,0.000000}%
\pgfsetfillcolor{currentfill}%
\pgfsetlinewidth{0.803000pt}%
\definecolor{currentstroke}{rgb}{0.000000,0.000000,0.000000}%
\pgfsetstrokecolor{currentstroke}%
\pgfsetdash{}{0pt}%
\pgfsys@defobject{currentmarker}{\pgfqpoint{-0.048611in}{0.000000in}}{\pgfqpoint{-0.000000in}{0.000000in}}{%
\pgfpathmoveto{\pgfqpoint{-0.000000in}{0.000000in}}%
\pgfpathlineto{\pgfqpoint{-0.048611in}{0.000000in}}%
\pgfusepath{stroke,fill}%
}%
\begin{pgfscope}%
\pgfsys@transformshift{0.632797in}{1.368904in}%
\pgfsys@useobject{currentmarker}{}%
\end{pgfscope}%
\end{pgfscope}%
\begin{pgfscope}%
\definecolor{textcolor}{rgb}{0.000000,0.000000,0.000000}%
\pgfsetstrokecolor{textcolor}%
\pgfsetfillcolor{textcolor}%
\pgftext[x=0.266667in, y=1.330324in, left, base]{\color{textcolor}{\rmfamily\fontsize{8.000000}{9.600000}\selectfont\catcode`\^=\active\def^{\ifmmode\sp\else\^{}\fi}\catcode`\%=\active\def%{\%}550.1}}%
\end{pgfscope}%
\begin{pgfscope}%
\definecolor{textcolor}{rgb}{0.000000,0.000000,0.000000}%
\pgfsetstrokecolor{textcolor}%
\pgfsetfillcolor{textcolor}%
\pgftext[x=0.211111in,y=0.915779in,,bottom,rotate=90.000000]{\color{textcolor}{\rmfamily\fontsize{8.000000}{9.600000}\selectfont\catcode`\^=\active\def^{\ifmmode\sp\else\^{}\fi}\catcode`\%=\active\def%{\%}$v_{DS}$ (V)}}%
\end{pgfscope}%
\begin{pgfscope}%
\pgfpathrectangle{\pgfqpoint{0.632797in}{0.462654in}}{\pgfqpoint{2.720000in}{0.906250in}}%
\pgfusepath{clip}%
\pgfsetrectcap%
\pgfsetroundjoin%
\pgfsetlinewidth{1.505625pt}%
\definecolor{currentstroke}{rgb}{0.000000,0.619608,0.450980}%
\pgfsetstrokecolor{currentstroke}%
\pgfsetdash{}{0pt}%
\pgfpathmoveto{\pgfqpoint{0.756433in}{0.503896in}}%
\pgfpathlineto{\pgfqpoint{0.756459in}{0.506064in}}%
\pgfpathlineto{\pgfqpoint{0.756474in}{0.503848in}}%
\pgfpathlineto{\pgfqpoint{0.757235in}{0.504138in}}%
\pgfpathlineto{\pgfqpoint{0.763049in}{0.505949in}}%
\pgfpathlineto{\pgfqpoint{0.763107in}{0.520349in}}%
\pgfpathlineto{\pgfqpoint{0.763210in}{1.327709in}}%
\pgfpathlineto{\pgfqpoint{0.763957in}{1.327703in}}%
\pgfpathlineto{\pgfqpoint{0.771705in}{1.327634in}}%
\pgfpathlineto{\pgfqpoint{0.772630in}{0.976513in}}%
\pgfpathlineto{\pgfqpoint{0.772710in}{0.977411in}}%
\pgfpathlineto{\pgfqpoint{0.772926in}{0.976786in}}%
\pgfpathlineto{\pgfqpoint{0.772994in}{0.504073in}}%
\pgfpathlineto{\pgfqpoint{0.773668in}{0.504292in}}%
\pgfpathlineto{\pgfqpoint{0.779542in}{0.506225in}}%
\pgfpathlineto{\pgfqpoint{0.779599in}{0.524799in}}%
\pgfpathlineto{\pgfqpoint{0.779656in}{1.327708in}}%
\pgfpathlineto{\pgfqpoint{0.780454in}{1.327703in}}%
\pgfpathlineto{\pgfqpoint{0.788202in}{1.327621in}}%
\pgfpathlineto{\pgfqpoint{0.789421in}{0.707515in}}%
\pgfpathlineto{\pgfqpoint{0.789502in}{0.504261in}}%
\pgfpathlineto{\pgfqpoint{0.790191in}{0.504453in}}%
\pgfpathlineto{\pgfqpoint{0.796023in}{0.506303in}}%
\pgfpathlineto{\pgfqpoint{0.796079in}{0.522445in}}%
\pgfpathlineto{\pgfqpoint{0.796159in}{1.327707in}}%
\pgfpathlineto{\pgfqpoint{0.796938in}{1.327703in}}%
\pgfpathlineto{\pgfqpoint{0.804686in}{1.327616in}}%
\pgfpathlineto{\pgfqpoint{0.804851in}{1.311937in}}%
\pgfpathlineto{\pgfqpoint{0.805016in}{1.327575in}}%
\pgfpathlineto{\pgfqpoint{0.805993in}{0.504358in}}%
\pgfpathlineto{\pgfqpoint{0.806160in}{0.504395in}}%
\pgfpathlineto{\pgfqpoint{0.812509in}{0.506392in}}%
\pgfpathlineto{\pgfqpoint{0.812566in}{0.523673in}}%
\pgfpathlineto{\pgfqpoint{0.812618in}{1.327707in}}%
\pgfpathlineto{\pgfqpoint{0.813382in}{1.327704in}}%
\pgfpathlineto{\pgfqpoint{0.821295in}{1.327585in}}%
\pgfpathlineto{\pgfqpoint{0.822491in}{0.504408in}}%
\pgfpathlineto{\pgfqpoint{0.823315in}{0.504638in}}%
\pgfpathlineto{\pgfqpoint{0.829001in}{0.506637in}}%
\pgfpathlineto{\pgfqpoint{0.829065in}{0.593705in}}%
\pgfpathlineto{\pgfqpoint{0.829112in}{1.327711in}}%
\pgfpathlineto{\pgfqpoint{0.829814in}{1.327709in}}%
\pgfpathlineto{\pgfqpoint{0.836737in}{1.327669in}}%
\pgfpathlineto{\pgfqpoint{0.836902in}{0.976856in}}%
\pgfpathlineto{\pgfqpoint{0.837562in}{0.977289in}}%
\pgfpathlineto{\pgfqpoint{0.837726in}{1.327651in}}%
\pgfpathlineto{\pgfqpoint{0.838386in}{1.327630in}}%
\pgfpathlineto{\pgfqpoint{0.838977in}{0.504406in}}%
\pgfpathlineto{\pgfqpoint{0.839520in}{0.504552in}}%
\pgfpathlineto{\pgfqpoint{0.845489in}{0.506861in}}%
\pgfpathlineto{\pgfqpoint{0.845556in}{0.657556in}}%
\pgfpathlineto{\pgfqpoint{0.845582in}{1.327707in}}%
\pgfpathlineto{\pgfqpoint{0.846197in}{1.327705in}}%
\pgfpathlineto{\pgfqpoint{0.854439in}{1.327590in}}%
\pgfpathlineto{\pgfqpoint{0.855452in}{0.504320in}}%
\pgfpathlineto{\pgfqpoint{0.855573in}{0.504345in}}%
\pgfpathlineto{\pgfqpoint{0.861974in}{0.506759in}}%
\pgfpathlineto{\pgfqpoint{0.862040in}{0.648058in}}%
\pgfpathlineto{\pgfqpoint{0.862092in}{1.327707in}}%
\pgfpathlineto{\pgfqpoint{0.862772in}{1.327704in}}%
\pgfpathlineto{\pgfqpoint{0.870850in}{1.327576in}}%
\pgfpathlineto{\pgfqpoint{0.871938in}{0.504180in}}%
\pgfpathlineto{\pgfqpoint{0.873093in}{0.504497in}}%
\pgfpathlineto{\pgfqpoint{0.878459in}{0.506526in}}%
\pgfpathlineto{\pgfqpoint{0.878509in}{0.528521in}}%
\pgfpathlineto{\pgfqpoint{0.878594in}{1.327708in}}%
\pgfpathlineto{\pgfqpoint{0.879318in}{1.327702in}}%
\pgfpathlineto{\pgfqpoint{0.887066in}{1.327633in}}%
\pgfpathlineto{\pgfqpoint{0.887841in}{0.976275in}}%
\pgfpathlineto{\pgfqpoint{0.887957in}{0.976887in}}%
\pgfpathlineto{\pgfqpoint{0.888312in}{0.976234in}}%
\pgfpathlineto{\pgfqpoint{0.888320in}{0.976471in}}%
\pgfpathlineto{\pgfqpoint{0.888387in}{0.503994in}}%
\pgfpathlineto{\pgfqpoint{0.889038in}{0.504200in}}%
\pgfpathlineto{\pgfqpoint{0.894930in}{0.505930in}}%
\pgfpathlineto{\pgfqpoint{0.894984in}{0.519664in}}%
\pgfpathlineto{\pgfqpoint{0.895056in}{1.327709in}}%
\pgfpathlineto{\pgfqpoint{0.895865in}{1.327702in}}%
\pgfpathlineto{\pgfqpoint{0.903225in}{1.327647in}}%
\pgfpathlineto{\pgfqpoint{0.903328in}{1.291342in}}%
\pgfpathlineto{\pgfqpoint{0.903421in}{1.327641in}}%
\pgfpathlineto{\pgfqpoint{0.904163in}{0.880005in}}%
\pgfpathlineto{\pgfqpoint{0.904395in}{1.185412in}}%
\pgfpathlineto{\pgfqpoint{0.904811in}{0.742266in}}%
\pgfpathlineto{\pgfqpoint{0.904866in}{0.503892in}}%
\pgfpathlineto{\pgfqpoint{0.905539in}{0.504107in}}%
\pgfpathlineto{\pgfqpoint{0.911415in}{0.505832in}}%
\pgfpathlineto{\pgfqpoint{0.911469in}{0.519246in}}%
\pgfpathlineto{\pgfqpoint{0.911539in}{1.327709in}}%
\pgfpathlineto{\pgfqpoint{0.912345in}{1.327702in}}%
\pgfpathlineto{\pgfqpoint{0.919919in}{1.327646in}}%
\pgfpathlineto{\pgfqpoint{0.920084in}{0.861316in}}%
\pgfpathlineto{\pgfqpoint{0.920728in}{0.884971in}}%
\pgfpathlineto{\pgfqpoint{0.920893in}{1.327562in}}%
\pgfpathlineto{\pgfqpoint{0.921352in}{0.503886in}}%
\pgfpathlineto{\pgfqpoint{0.922155in}{0.504135in}}%
\pgfpathlineto{\pgfqpoint{0.927900in}{0.505828in}}%
\pgfpathlineto{\pgfqpoint{0.927953in}{0.519219in}}%
\pgfpathlineto{\pgfqpoint{0.928024in}{1.327709in}}%
\pgfpathlineto{\pgfqpoint{0.928828in}{1.327701in}}%
\pgfpathlineto{\pgfqpoint{0.936240in}{1.327649in}}%
\pgfpathlineto{\pgfqpoint{0.936735in}{0.970961in}}%
\pgfpathlineto{\pgfqpoint{0.937065in}{1.048533in}}%
\pgfpathlineto{\pgfqpoint{0.937230in}{1.327599in}}%
\pgfpathlineto{\pgfqpoint{0.937770in}{0.937737in}}%
\pgfpathlineto{\pgfqpoint{0.937839in}{0.503877in}}%
\pgfpathlineto{\pgfqpoint{0.938956in}{0.504221in}}%
\pgfpathlineto{\pgfqpoint{0.944385in}{0.505829in}}%
\pgfpathlineto{\pgfqpoint{0.944438in}{0.519225in}}%
\pgfpathlineto{\pgfqpoint{0.944509in}{1.327709in}}%
\pgfpathlineto{\pgfqpoint{0.945313in}{1.327701in}}%
\pgfpathlineto{\pgfqpoint{0.952870in}{1.327645in}}%
\pgfpathlineto{\pgfqpoint{0.953034in}{0.794552in}}%
\pgfpathlineto{\pgfqpoint{0.953620in}{0.878170in}}%
\pgfpathlineto{\pgfqpoint{0.953851in}{1.327491in}}%
\pgfpathlineto{\pgfqpoint{0.954273in}{0.532527in}}%
\pgfpathlineto{\pgfqpoint{0.954316in}{0.503916in}}%
\pgfpathlineto{\pgfqpoint{0.955032in}{0.504127in}}%
\pgfpathlineto{\pgfqpoint{0.960870in}{0.505856in}}%
\pgfpathlineto{\pgfqpoint{0.960923in}{0.519343in}}%
\pgfpathlineto{\pgfqpoint{0.960994in}{1.327709in}}%
\pgfpathlineto{\pgfqpoint{0.961799in}{1.327702in}}%
\pgfpathlineto{\pgfqpoint{0.969251in}{1.327650in}}%
\pgfpathlineto{\pgfqpoint{0.969746in}{0.871294in}}%
\pgfpathlineto{\pgfqpoint{0.970062in}{1.194716in}}%
\pgfpathlineto{\pgfqpoint{0.970217in}{1.327603in}}%
\pgfpathlineto{\pgfqpoint{0.970382in}{0.865400in}}%
\pgfpathlineto{\pgfqpoint{0.970689in}{1.143058in}}%
\pgfpathlineto{\pgfqpoint{0.970807in}{0.503898in}}%
\pgfpathlineto{\pgfqpoint{0.971821in}{0.504189in}}%
\pgfpathlineto{\pgfqpoint{0.977355in}{0.505835in}}%
\pgfpathlineto{\pgfqpoint{0.977420in}{0.552731in}}%
\pgfpathlineto{\pgfqpoint{0.977489in}{1.327709in}}%
\pgfpathlineto{\pgfqpoint{0.978097in}{1.327706in}}%
\pgfpathlineto{\pgfqpoint{0.985765in}{1.327643in}}%
\pgfpathlineto{\pgfqpoint{0.986842in}{0.838952in}}%
\pgfpathlineto{\pgfqpoint{0.987036in}{0.947356in}}%
\pgfpathlineto{\pgfqpoint{0.987227in}{1.140834in}}%
\pgfpathlineto{\pgfqpoint{0.987297in}{0.503882in}}%
\pgfpathlineto{\pgfqpoint{0.987475in}{0.503961in}}%
\pgfpathlineto{\pgfqpoint{0.993840in}{0.505848in}}%
\pgfpathlineto{\pgfqpoint{0.993893in}{0.519309in}}%
\pgfpathlineto{\pgfqpoint{0.993964in}{1.327709in}}%
\pgfpathlineto{\pgfqpoint{0.994770in}{1.327702in}}%
\pgfpathlineto{\pgfqpoint{1.002235in}{1.327648in}}%
\pgfpathlineto{\pgfqpoint{1.003059in}{0.973185in}}%
\pgfpathlineto{\pgfqpoint{1.003224in}{1.327586in}}%
\pgfpathlineto{\pgfqpoint{1.003721in}{0.740164in}}%
\pgfpathlineto{\pgfqpoint{1.003773in}{0.503880in}}%
\pgfpathlineto{\pgfqpoint{1.004444in}{0.504103in}}%
\pgfpathlineto{\pgfqpoint{1.010324in}{0.505855in}}%
\pgfpathlineto{\pgfqpoint{1.010378in}{0.519337in}}%
\pgfpathlineto{\pgfqpoint{1.010449in}{1.327709in}}%
\pgfpathlineto{\pgfqpoint{1.011254in}{1.327702in}}%
\pgfpathlineto{\pgfqpoint{1.018797in}{1.327646in}}%
\pgfpathlineto{\pgfqpoint{1.018962in}{0.889433in}}%
\pgfpathlineto{\pgfqpoint{1.019611in}{0.906146in}}%
\pgfpathlineto{\pgfqpoint{1.019776in}{1.327560in}}%
\pgfpathlineto{\pgfqpoint{1.020224in}{0.510061in}}%
\pgfpathlineto{\pgfqpoint{1.020255in}{0.503890in}}%
\pgfpathlineto{\pgfqpoint{1.021002in}{0.504125in}}%
\pgfpathlineto{\pgfqpoint{1.026809in}{0.505860in}}%
\pgfpathlineto{\pgfqpoint{1.026863in}{0.519358in}}%
\pgfpathlineto{\pgfqpoint{1.026934in}{1.327709in}}%
\pgfpathlineto{\pgfqpoint{1.027740in}{1.327702in}}%
\pgfpathlineto{\pgfqpoint{1.035178in}{1.327649in}}%
\pgfpathlineto{\pgfqpoint{1.035672in}{0.943047in}}%
\pgfpathlineto{\pgfqpoint{1.036000in}{1.082286in}}%
\pgfpathlineto{\pgfqpoint{1.036164in}{1.327596in}}%
\pgfpathlineto{\pgfqpoint{1.036329in}{0.942380in}}%
\pgfpathlineto{\pgfqpoint{1.036686in}{0.987558in}}%
\pgfpathlineto{\pgfqpoint{1.036745in}{0.503861in}}%
\pgfpathlineto{\pgfqpoint{1.037461in}{0.504112in}}%
\pgfpathlineto{\pgfqpoint{1.043294in}{0.505856in}}%
\pgfpathlineto{\pgfqpoint{1.043348in}{0.519344in}}%
\pgfpathlineto{\pgfqpoint{1.043419in}{1.327709in}}%
\pgfpathlineto{\pgfqpoint{1.044224in}{1.327702in}}%
\pgfpathlineto{\pgfqpoint{1.051781in}{1.327645in}}%
\pgfpathlineto{\pgfqpoint{1.051946in}{0.833203in}}%
\pgfpathlineto{\pgfqpoint{1.052473in}{1.141749in}}%
\pgfpathlineto{\pgfqpoint{1.052638in}{1.113601in}}%
\pgfpathlineto{\pgfqpoint{1.053235in}{0.503887in}}%
\pgfpathlineto{\pgfqpoint{1.053763in}{0.504072in}}%
\pgfpathlineto{\pgfqpoint{1.059779in}{0.505874in}}%
\pgfpathlineto{\pgfqpoint{1.059833in}{0.519434in}}%
\pgfpathlineto{\pgfqpoint{1.059904in}{1.327709in}}%
\pgfpathlineto{\pgfqpoint{1.060712in}{1.327702in}}%
\pgfpathlineto{\pgfqpoint{1.068183in}{1.327648in}}%
\pgfpathlineto{\pgfqpoint{1.068677in}{0.935912in}}%
\pgfpathlineto{\pgfqpoint{1.069004in}{1.092806in}}%
\pgfpathlineto{\pgfqpoint{1.069167in}{1.327594in}}%
\pgfpathlineto{\pgfqpoint{1.069332in}{0.934706in}}%
\pgfpathlineto{\pgfqpoint{1.069654in}{1.020541in}}%
\pgfpathlineto{\pgfqpoint{1.069718in}{0.503887in}}%
\pgfpathlineto{\pgfqpoint{1.070332in}{0.504086in}}%
\pgfpathlineto{\pgfqpoint{1.076264in}{0.505870in}}%
\pgfpathlineto{\pgfqpoint{1.076318in}{0.519505in}}%
\pgfpathlineto{\pgfqpoint{1.076389in}{1.327709in}}%
\pgfpathlineto{\pgfqpoint{1.077207in}{1.327702in}}%
\pgfpathlineto{\pgfqpoint{1.084584in}{1.327649in}}%
\pgfpathlineto{\pgfqpoint{1.085079in}{0.892695in}}%
\pgfpathlineto{\pgfqpoint{1.085399in}{1.159576in}}%
\pgfpathlineto{\pgfqpoint{1.085557in}{1.327598in}}%
\pgfpathlineto{\pgfqpoint{1.085721in}{0.888288in}}%
\pgfpathlineto{\pgfqpoint{1.086027in}{1.097356in}}%
\pgfpathlineto{\pgfqpoint{1.086200in}{0.503851in}}%
\pgfpathlineto{\pgfqpoint{1.087520in}{0.504283in}}%
\pgfpathlineto{\pgfqpoint{1.092749in}{0.505869in}}%
\pgfpathlineto{\pgfqpoint{1.092802in}{0.519526in}}%
\pgfpathlineto{\pgfqpoint{1.092874in}{1.327709in}}%
\pgfpathlineto{\pgfqpoint{1.093694in}{1.327701in}}%
\pgfpathlineto{\pgfqpoint{1.101036in}{1.327649in}}%
\pgfpathlineto{\pgfqpoint{1.101530in}{0.877353in}}%
\pgfpathlineto{\pgfqpoint{1.101848in}{1.184618in}}%
\pgfpathlineto{\pgfqpoint{1.102004in}{1.327600in}}%
\pgfpathlineto{\pgfqpoint{1.102168in}{0.871870in}}%
\pgfpathlineto{\pgfqpoint{1.102493in}{1.109418in}}%
\pgfpathlineto{\pgfqpoint{1.102687in}{0.503886in}}%
\pgfpathlineto{\pgfqpoint{1.103678in}{0.504189in}}%
\pgfpathlineto{\pgfqpoint{1.109237in}{0.505897in}}%
\pgfpathlineto{\pgfqpoint{1.109291in}{0.520983in}}%
\pgfpathlineto{\pgfqpoint{1.109371in}{1.327708in}}%
\pgfpathlineto{\pgfqpoint{1.110195in}{1.327703in}}%
\pgfpathlineto{\pgfqpoint{1.117781in}{1.327634in}}%
\pgfpathlineto{\pgfqpoint{1.118888in}{0.731696in}}%
\pgfpathlineto{\pgfqpoint{1.119000in}{0.854236in}}%
\pgfpathlineto{\pgfqpoint{1.119106in}{1.146282in}}%
\pgfpathlineto{\pgfqpoint{1.119185in}{0.503935in}}%
\pgfpathlineto{\pgfqpoint{1.119584in}{0.504044in}}%
\pgfpathlineto{\pgfqpoint{1.125718in}{0.505885in}}%
\pgfpathlineto{\pgfqpoint{1.125772in}{0.519335in}}%
\pgfpathlineto{\pgfqpoint{1.125843in}{1.327709in}}%
\pgfpathlineto{\pgfqpoint{1.126642in}{1.327702in}}%
\pgfpathlineto{\pgfqpoint{1.134020in}{1.327651in}}%
\pgfpathlineto{\pgfqpoint{1.134185in}{1.318008in}}%
\pgfpathlineto{\pgfqpoint{1.134350in}{1.327643in}}%
\pgfpathlineto{\pgfqpoint{1.134515in}{0.771788in}}%
\pgfpathlineto{\pgfqpoint{1.135047in}{0.980023in}}%
\pgfpathlineto{\pgfqpoint{1.135181in}{1.230995in}}%
\pgfpathlineto{\pgfqpoint{1.135603in}{0.585972in}}%
\pgfpathlineto{\pgfqpoint{1.135653in}{0.503869in}}%
\pgfpathlineto{\pgfqpoint{1.136344in}{0.504116in}}%
\pgfpathlineto{\pgfqpoint{1.142215in}{0.506206in}}%
\pgfpathlineto{\pgfqpoint{1.142269in}{0.554118in}}%
\pgfpathlineto{\pgfqpoint{1.142314in}{1.327710in}}%
\pgfpathlineto{\pgfqpoint{1.142966in}{1.327702in}}%
\pgfpathlineto{\pgfqpoint{1.150384in}{1.327657in}}%
\pgfpathlineto{\pgfqpoint{1.150549in}{1.111872in}}%
\pgfpathlineto{\pgfqpoint{1.151118in}{1.284398in}}%
\pgfpathlineto{\pgfqpoint{1.151166in}{1.327615in}}%
\pgfpathlineto{\pgfqpoint{1.151368in}{1.185184in}}%
\pgfpathlineto{\pgfqpoint{1.151466in}{1.214016in}}%
\pgfpathlineto{\pgfqpoint{1.152142in}{0.503915in}}%
\pgfpathlineto{\pgfqpoint{1.152614in}{0.504053in}}%
\pgfpathlineto{\pgfqpoint{1.158688in}{0.505882in}}%
\pgfpathlineto{\pgfqpoint{1.158742in}{0.519455in}}%
\pgfpathlineto{\pgfqpoint{1.158813in}{1.327709in}}%
\pgfpathlineto{\pgfqpoint{1.159620in}{1.327702in}}%
\pgfpathlineto{\pgfqpoint{1.167038in}{1.327649in}}%
\pgfpathlineto{\pgfqpoint{1.167898in}{1.010343in}}%
\pgfpathlineto{\pgfqpoint{1.168022in}{1.297590in}}%
\pgfpathlineto{\pgfqpoint{1.168384in}{0.732165in}}%
\pgfpathlineto{\pgfqpoint{1.168568in}{0.831743in}}%
\pgfpathlineto{\pgfqpoint{1.168626in}{0.503885in}}%
\pgfpathlineto{\pgfqpoint{1.169259in}{0.504111in}}%
\pgfpathlineto{\pgfqpoint{1.175173in}{0.505897in}}%
\pgfpathlineto{\pgfqpoint{1.175227in}{0.519518in}}%
\pgfpathlineto{\pgfqpoint{1.175298in}{1.327709in}}%
\pgfpathlineto{\pgfqpoint{1.176106in}{1.327702in}}%
\pgfpathlineto{\pgfqpoint{1.183524in}{1.327650in}}%
\pgfpathlineto{\pgfqpoint{1.183689in}{1.304951in}}%
\pgfpathlineto{\pgfqpoint{1.183853in}{1.327642in}}%
\pgfpathlineto{\pgfqpoint{1.184018in}{0.780926in}}%
\pgfpathlineto{\pgfqpoint{1.184609in}{1.278633in}}%
\pgfpathlineto{\pgfqpoint{1.185111in}{0.503900in}}%
\pgfpathlineto{\pgfqpoint{1.185937in}{0.504163in}}%
\pgfpathlineto{\pgfqpoint{1.191658in}{0.505894in}}%
\pgfpathlineto{\pgfqpoint{1.191711in}{0.519504in}}%
\pgfpathlineto{\pgfqpoint{1.191783in}{1.327709in}}%
\pgfpathlineto{\pgfqpoint{1.192590in}{1.327702in}}%
\pgfpathlineto{\pgfqpoint{1.200009in}{1.327650in}}%
\pgfpathlineto{\pgfqpoint{1.200503in}{0.834820in}}%
\pgfpathlineto{\pgfqpoint{1.200796in}{1.274748in}}%
\pgfpathlineto{\pgfqpoint{1.200888in}{1.327605in}}%
\pgfpathlineto{\pgfqpoint{1.200970in}{1.105888in}}%
\pgfpathlineto{\pgfqpoint{1.201170in}{1.122415in}}%
\pgfpathlineto{\pgfqpoint{1.201593in}{0.503918in}}%
\pgfpathlineto{\pgfqpoint{1.202152in}{0.504087in}}%
\pgfpathlineto{\pgfqpoint{1.208143in}{0.505897in}}%
\pgfpathlineto{\pgfqpoint{1.208196in}{0.519519in}}%
\pgfpathlineto{\pgfqpoint{1.208268in}{1.327709in}}%
\pgfpathlineto{\pgfqpoint{1.209076in}{1.327702in}}%
\pgfpathlineto{\pgfqpoint{1.216494in}{1.327650in}}%
\pgfpathlineto{\pgfqpoint{1.216659in}{1.277059in}}%
\pgfpathlineto{\pgfqpoint{1.216823in}{1.327641in}}%
\pgfpathlineto{\pgfqpoint{1.216988in}{0.835915in}}%
\pgfpathlineto{\pgfqpoint{1.217663in}{0.966811in}}%
\pgfpathlineto{\pgfqpoint{1.217816in}{1.163921in}}%
\pgfpathlineto{\pgfqpoint{1.218080in}{0.503899in}}%
\pgfpathlineto{\pgfqpoint{1.218149in}{0.503924in}}%
\pgfpathlineto{\pgfqpoint{1.224627in}{0.505873in}}%
\pgfpathlineto{\pgfqpoint{1.224681in}{0.519415in}}%
\pgfpathlineto{\pgfqpoint{1.224752in}{1.327709in}}%
\pgfpathlineto{\pgfqpoint{1.225558in}{1.327701in}}%
\pgfpathlineto{\pgfqpoint{1.232977in}{1.327648in}}%
\pgfpathlineto{\pgfqpoint{1.233471in}{0.992330in}}%
\pgfpathlineto{\pgfqpoint{1.233801in}{1.026027in}}%
\pgfpathlineto{\pgfqpoint{1.233966in}{1.327587in}}%
\pgfpathlineto{\pgfqpoint{1.234376in}{0.949426in}}%
\pgfpathlineto{\pgfqpoint{1.234501in}{0.967677in}}%
\pgfpathlineto{\pgfqpoint{1.234561in}{0.503914in}}%
\pgfpathlineto{\pgfqpoint{1.235251in}{0.504112in}}%
\pgfpathlineto{\pgfqpoint{1.241112in}{0.505886in}}%
\pgfpathlineto{\pgfqpoint{1.241166in}{0.519472in}}%
\pgfpathlineto{\pgfqpoint{1.241237in}{1.327709in}}%
\pgfpathlineto{\pgfqpoint{1.242044in}{1.327702in}}%
\pgfpathlineto{\pgfqpoint{1.249463in}{1.327649in}}%
\pgfpathlineto{\pgfqpoint{1.249957in}{0.965283in}}%
\pgfpathlineto{\pgfqpoint{1.250287in}{0.976649in}}%
\pgfpathlineto{\pgfqpoint{1.250448in}{1.327594in}}%
\pgfpathlineto{\pgfqpoint{1.250991in}{0.868041in}}%
\pgfpathlineto{\pgfqpoint{1.251045in}{0.503881in}}%
\pgfpathlineto{\pgfqpoint{1.251744in}{0.504111in}}%
\pgfpathlineto{\pgfqpoint{1.257607in}{0.506057in}}%
\pgfpathlineto{\pgfqpoint{1.257672in}{0.585010in}}%
\pgfpathlineto{\pgfqpoint{1.257724in}{1.327711in}}%
\pgfpathlineto{\pgfqpoint{1.258449in}{1.327709in}}%
\pgfpathlineto{\pgfqpoint{1.265372in}{1.327670in}}%
\pgfpathlineto{\pgfqpoint{1.265537in}{0.974492in}}%
\pgfpathlineto{\pgfqpoint{1.266197in}{0.975880in}}%
\pgfpathlineto{\pgfqpoint{1.266361in}{1.327653in}}%
\pgfpathlineto{\pgfqpoint{1.267021in}{1.327633in}}%
\pgfpathlineto{\pgfqpoint{1.267542in}{0.503893in}}%
\pgfpathlineto{\pgfqpoint{1.268118in}{0.504096in}}%
\pgfpathlineto{\pgfqpoint{1.274092in}{0.506093in}}%
\pgfpathlineto{\pgfqpoint{1.274153in}{0.588242in}}%
\pgfpathlineto{\pgfqpoint{1.274295in}{1.327710in}}%
\pgfpathlineto{\pgfqpoint{1.274832in}{1.327708in}}%
\pgfpathlineto{\pgfqpoint{1.282415in}{1.327654in}}%
\pgfpathlineto{\pgfqpoint{1.282580in}{0.977886in}}%
\pgfpathlineto{\pgfqpoint{1.283240in}{0.978413in}}%
\pgfpathlineto{\pgfqpoint{1.283404in}{1.327616in}}%
\pgfpathlineto{\pgfqpoint{1.283952in}{0.855801in}}%
\pgfpathlineto{\pgfqpoint{1.284022in}{0.503875in}}%
\pgfpathlineto{\pgfqpoint{1.284744in}{0.504107in}}%
\pgfpathlineto{\pgfqpoint{1.290567in}{0.505872in}}%
\pgfpathlineto{\pgfqpoint{1.290620in}{0.519413in}}%
\pgfpathlineto{\pgfqpoint{1.290692in}{1.327709in}}%
\pgfpathlineto{\pgfqpoint{1.291499in}{1.327701in}}%
\pgfpathlineto{\pgfqpoint{1.298917in}{1.327647in}}%
\pgfpathlineto{\pgfqpoint{1.299082in}{0.976864in}}%
\pgfpathlineto{\pgfqpoint{1.299741in}{0.984188in}}%
\pgfpathlineto{\pgfqpoint{1.299906in}{1.327579in}}%
\pgfpathlineto{\pgfqpoint{1.300434in}{0.956980in}}%
\pgfpathlineto{\pgfqpoint{1.300441in}{0.957395in}}%
\pgfpathlineto{\pgfqpoint{1.300509in}{0.503900in}}%
\pgfpathlineto{\pgfqpoint{1.301138in}{0.504095in}}%
\pgfpathlineto{\pgfqpoint{1.307063in}{0.506157in}}%
\pgfpathlineto{\pgfqpoint{1.307110in}{0.523280in}}%
\pgfpathlineto{\pgfqpoint{1.307178in}{1.327707in}}%
\pgfpathlineto{\pgfqpoint{1.307931in}{1.327705in}}%
\pgfpathlineto{\pgfqpoint{1.316009in}{1.327461in}}%
\pgfpathlineto{\pgfqpoint{1.316174in}{1.316479in}}%
\pgfpathlineto{\pgfqpoint{1.316986in}{0.503891in}}%
\pgfpathlineto{\pgfqpoint{1.317159in}{0.503969in}}%
\pgfpathlineto{\pgfqpoint{1.323537in}{0.505898in}}%
\pgfpathlineto{\pgfqpoint{1.323590in}{0.519543in}}%
\pgfpathlineto{\pgfqpoint{1.323662in}{1.327709in}}%
\pgfpathlineto{\pgfqpoint{1.324473in}{1.327702in}}%
\pgfpathlineto{\pgfqpoint{1.331891in}{1.327648in}}%
\pgfpathlineto{\pgfqpoint{1.332385in}{0.976001in}}%
\pgfpathlineto{\pgfqpoint{1.332715in}{0.978484in}}%
\pgfpathlineto{\pgfqpoint{1.332879in}{1.327589in}}%
\pgfpathlineto{\pgfqpoint{1.333403in}{0.975702in}}%
\pgfpathlineto{\pgfqpoint{1.333411in}{0.976101in}}%
\pgfpathlineto{\pgfqpoint{1.333477in}{0.503890in}}%
\pgfpathlineto{\pgfqpoint{1.334094in}{0.504102in}}%
\pgfpathlineto{\pgfqpoint{1.340036in}{0.506530in}}%
\pgfpathlineto{\pgfqpoint{1.340101in}{0.645474in}}%
\pgfpathlineto{\pgfqpoint{1.340155in}{1.327707in}}%
\pgfpathlineto{\pgfqpoint{1.340860in}{1.327704in}}%
\pgfpathlineto{\pgfqpoint{1.348937in}{1.327593in}}%
\pgfpathlineto{\pgfqpoint{1.349724in}{0.976221in}}%
\pgfpathlineto{\pgfqpoint{1.349888in}{0.976503in}}%
\pgfpathlineto{\pgfqpoint{1.349896in}{0.976417in}}%
\pgfpathlineto{\pgfqpoint{1.349964in}{0.503913in}}%
\pgfpathlineto{\pgfqpoint{1.350600in}{0.504112in}}%
\pgfpathlineto{\pgfqpoint{1.356506in}{0.505903in}}%
\pgfpathlineto{\pgfqpoint{1.356560in}{0.519550in}}%
\pgfpathlineto{\pgfqpoint{1.356632in}{1.327709in}}%
\pgfpathlineto{\pgfqpoint{1.357441in}{1.327702in}}%
\pgfpathlineto{\pgfqpoint{1.364859in}{1.327648in}}%
\pgfpathlineto{\pgfqpoint{1.365353in}{0.976065in}}%
\pgfpathlineto{\pgfqpoint{1.365683in}{0.977185in}}%
\pgfpathlineto{\pgfqpoint{1.365847in}{1.327588in}}%
\pgfpathlineto{\pgfqpoint{1.366373in}{0.975956in}}%
\pgfpathlineto{\pgfqpoint{1.366381in}{0.976316in}}%
\pgfpathlineto{\pgfqpoint{1.366447in}{0.503897in}}%
\pgfpathlineto{\pgfqpoint{1.367086in}{0.504114in}}%
\pgfpathlineto{\pgfqpoint{1.372991in}{0.505904in}}%
\pgfpathlineto{\pgfqpoint{1.373045in}{0.519549in}}%
\pgfpathlineto{\pgfqpoint{1.373117in}{1.327709in}}%
\pgfpathlineto{\pgfqpoint{1.373925in}{1.327702in}}%
\pgfpathlineto{\pgfqpoint{1.381343in}{1.327648in}}%
\pgfpathlineto{\pgfqpoint{1.381837in}{0.976127in}}%
\pgfpathlineto{\pgfqpoint{1.382167in}{0.977773in}}%
\pgfpathlineto{\pgfqpoint{1.382331in}{1.327587in}}%
\pgfpathlineto{\pgfqpoint{1.382766in}{0.976212in}}%
\pgfpathlineto{\pgfqpoint{1.382866in}{0.976319in}}%
\pgfpathlineto{\pgfqpoint{1.382931in}{0.503897in}}%
\pgfpathlineto{\pgfqpoint{1.383585in}{0.504119in}}%
\pgfpathlineto{\pgfqpoint{1.389481in}{0.505971in}}%
\pgfpathlineto{\pgfqpoint{1.389537in}{0.522767in}}%
\pgfpathlineto{\pgfqpoint{1.389618in}{1.327707in}}%
\pgfpathlineto{\pgfqpoint{1.390408in}{1.327704in}}%
\pgfpathlineto{\pgfqpoint{1.398486in}{1.327566in}}%
\pgfpathlineto{\pgfqpoint{1.399417in}{0.503899in}}%
\pgfpathlineto{\pgfqpoint{1.399627in}{0.503982in}}%
\pgfpathlineto{\pgfqpoint{1.405971in}{0.506122in}}%
\pgfpathlineto{\pgfqpoint{1.406033in}{0.596433in}}%
\pgfpathlineto{\pgfqpoint{1.406083in}{1.327709in}}%
\pgfpathlineto{\pgfqpoint{1.406695in}{1.327707in}}%
\pgfpathlineto{\pgfqpoint{1.414278in}{1.327650in}}%
\pgfpathlineto{\pgfqpoint{1.414443in}{0.976036in}}%
\pgfpathlineto{\pgfqpoint{1.415102in}{1.019884in}}%
\pgfpathlineto{\pgfqpoint{1.415267in}{1.327598in}}%
\pgfpathlineto{\pgfqpoint{1.415747in}{0.975429in}}%
\pgfpathlineto{\pgfqpoint{1.415836in}{0.975820in}}%
\pgfpathlineto{\pgfqpoint{1.415903in}{0.503910in}}%
\pgfpathlineto{\pgfqpoint{1.416542in}{0.504111in}}%
\pgfpathlineto{\pgfqpoint{1.422460in}{0.506530in}}%
\pgfpathlineto{\pgfqpoint{1.422525in}{0.645583in}}%
\pgfpathlineto{\pgfqpoint{1.422580in}{1.327707in}}%
\pgfpathlineto{\pgfqpoint{1.423283in}{1.327704in}}%
\pgfpathlineto{\pgfqpoint{1.431360in}{1.327589in}}%
\pgfpathlineto{\pgfqpoint{1.432169in}{0.976414in}}%
\pgfpathlineto{\pgfqpoint{1.432317in}{0.976651in}}%
\pgfpathlineto{\pgfqpoint{1.432321in}{0.976619in}}%
\pgfpathlineto{\pgfqpoint{1.432384in}{0.503910in}}%
\pgfpathlineto{\pgfqpoint{1.433055in}{0.504124in}}%
\pgfpathlineto{\pgfqpoint{1.438930in}{0.505903in}}%
\pgfpathlineto{\pgfqpoint{1.438984in}{0.519546in}}%
\pgfpathlineto{\pgfqpoint{1.439056in}{1.327709in}}%
\pgfpathlineto{\pgfqpoint{1.439864in}{1.327702in}}%
\pgfpathlineto{\pgfqpoint{1.447282in}{1.327648in}}%
\pgfpathlineto{\pgfqpoint{1.447777in}{0.975958in}}%
\pgfpathlineto{\pgfqpoint{1.448106in}{0.976783in}}%
\pgfpathlineto{\pgfqpoint{1.448270in}{1.327589in}}%
\pgfpathlineto{\pgfqpoint{1.448797in}{0.976028in}}%
\pgfpathlineto{\pgfqpoint{1.448805in}{0.976274in}}%
\pgfpathlineto{\pgfqpoint{1.448871in}{0.503893in}}%
\pgfpathlineto{\pgfqpoint{1.449524in}{0.504115in}}%
\pgfpathlineto{\pgfqpoint{1.455415in}{0.505899in}}%
\pgfpathlineto{\pgfqpoint{1.455469in}{0.519527in}}%
\pgfpathlineto{\pgfqpoint{1.455541in}{1.327709in}}%
\pgfpathlineto{\pgfqpoint{1.456348in}{1.327702in}}%
\pgfpathlineto{\pgfqpoint{1.463767in}{1.327648in}}%
\pgfpathlineto{\pgfqpoint{1.464261in}{0.976028in}}%
\pgfpathlineto{\pgfqpoint{1.464591in}{0.980210in}}%
\pgfpathlineto{\pgfqpoint{1.464755in}{1.327588in}}%
\pgfpathlineto{\pgfqpoint{1.465282in}{0.975704in}}%
\pgfpathlineto{\pgfqpoint{1.465290in}{0.976138in}}%
\pgfpathlineto{\pgfqpoint{1.465356in}{0.503890in}}%
\pgfpathlineto{\pgfqpoint{1.466050in}{0.504124in}}%
\pgfpathlineto{\pgfqpoint{1.471903in}{0.505931in}}%
\pgfpathlineto{\pgfqpoint{1.471957in}{0.521000in}}%
\pgfpathlineto{\pgfqpoint{1.472037in}{1.327708in}}%
\pgfpathlineto{\pgfqpoint{1.472870in}{1.327702in}}%
\pgfpathlineto{\pgfqpoint{1.480617in}{1.327628in}}%
\pgfpathlineto{\pgfqpoint{1.481543in}{0.948591in}}%
\pgfpathlineto{\pgfqpoint{1.481655in}{0.975149in}}%
\pgfpathlineto{\pgfqpoint{1.481776in}{0.974982in}}%
\pgfpathlineto{\pgfqpoint{1.481841in}{0.503891in}}%
\pgfpathlineto{\pgfqpoint{1.482481in}{0.504099in}}%
\pgfpathlineto{\pgfqpoint{1.488385in}{0.505886in}}%
\pgfpathlineto{\pgfqpoint{1.488439in}{0.519471in}}%
\pgfpathlineto{\pgfqpoint{1.488510in}{1.327709in}}%
\pgfpathlineto{\pgfqpoint{1.489317in}{1.327701in}}%
\pgfpathlineto{\pgfqpoint{1.496735in}{1.327647in}}%
\pgfpathlineto{\pgfqpoint{1.497230in}{0.985936in}}%
\pgfpathlineto{\pgfqpoint{1.497560in}{1.032663in}}%
\pgfpathlineto{\pgfqpoint{1.497725in}{1.327584in}}%
\pgfpathlineto{\pgfqpoint{1.498177in}{0.975950in}}%
\pgfpathlineto{\pgfqpoint{1.498260in}{0.976055in}}%
\pgfpathlineto{\pgfqpoint{1.498325in}{0.503885in}}%
\pgfpathlineto{\pgfqpoint{1.499029in}{0.504122in}}%
\pgfpathlineto{\pgfqpoint{1.504870in}{0.505892in}}%
\pgfpathlineto{\pgfqpoint{1.504924in}{0.519508in}}%
\pgfpathlineto{\pgfqpoint{1.504995in}{1.327709in}}%
\pgfpathlineto{\pgfqpoint{1.505804in}{1.327702in}}%
\pgfpathlineto{\pgfqpoint{1.513222in}{1.327648in}}%
\pgfpathlineto{\pgfqpoint{1.513717in}{0.975922in}}%
\pgfpathlineto{\pgfqpoint{1.514046in}{0.987330in}}%
\pgfpathlineto{\pgfqpoint{1.514211in}{1.327590in}}%
\pgfpathlineto{\pgfqpoint{1.514736in}{0.975651in}}%
\pgfpathlineto{\pgfqpoint{1.514745in}{0.975810in}}%
\pgfpathlineto{\pgfqpoint{1.514810in}{0.503882in}}%
\pgfpathlineto{\pgfqpoint{1.515513in}{0.504120in}}%
\pgfpathlineto{\pgfqpoint{1.521364in}{0.506040in}}%
\pgfpathlineto{\pgfqpoint{1.521427in}{0.583548in}}%
\pgfpathlineto{\pgfqpoint{1.521476in}{1.327711in}}%
\pgfpathlineto{\pgfqpoint{1.522186in}{1.327709in}}%
\pgfpathlineto{\pgfqpoint{1.529110in}{1.327670in}}%
\pgfpathlineto{\pgfqpoint{1.529274in}{0.975584in}}%
\pgfpathlineto{\pgfqpoint{1.529934in}{0.977371in}}%
\pgfpathlineto{\pgfqpoint{1.530099in}{1.327653in}}%
\pgfpathlineto{\pgfqpoint{1.530758in}{1.327633in}}%
\pgfpathlineto{\pgfqpoint{1.531295in}{0.503923in}}%
\pgfpathlineto{\pgfqpoint{1.531932in}{0.504112in}}%
\pgfpathlineto{\pgfqpoint{1.537851in}{0.506220in}}%
\pgfpathlineto{\pgfqpoint{1.537917in}{0.631609in}}%
\pgfpathlineto{\pgfqpoint{1.537975in}{1.327708in}}%
\pgfpathlineto{\pgfqpoint{1.538703in}{1.327704in}}%
\pgfpathlineto{\pgfqpoint{1.546781in}{1.327608in}}%
\pgfpathlineto{\pgfqpoint{1.547782in}{0.503897in}}%
\pgfpathlineto{\pgfqpoint{1.547934in}{0.503964in}}%
\pgfpathlineto{\pgfqpoint{1.554336in}{0.506210in}}%
\pgfpathlineto{\pgfqpoint{1.554402in}{0.631263in}}%
\pgfpathlineto{\pgfqpoint{1.554460in}{1.327708in}}%
\pgfpathlineto{\pgfqpoint{1.555190in}{1.327704in}}%
\pgfpathlineto{\pgfqpoint{1.563267in}{1.327603in}}%
\pgfpathlineto{\pgfqpoint{1.563858in}{0.934909in}}%
\pgfpathlineto{\pgfqpoint{1.564192in}{0.965251in}}%
\pgfpathlineto{\pgfqpoint{1.564199in}{0.966453in}}%
\pgfpathlineto{\pgfqpoint{1.564267in}{0.503898in}}%
\pgfpathlineto{\pgfqpoint{1.564885in}{0.504093in}}%
\pgfpathlineto{\pgfqpoint{1.570809in}{0.505887in}}%
\pgfpathlineto{\pgfqpoint{1.570863in}{0.519482in}}%
\pgfpathlineto{\pgfqpoint{1.570934in}{1.327709in}}%
\pgfpathlineto{\pgfqpoint{1.571743in}{1.327701in}}%
\pgfpathlineto{\pgfqpoint{1.579161in}{1.327647in}}%
\pgfpathlineto{\pgfqpoint{1.579655in}{0.977717in}}%
\pgfpathlineto{\pgfqpoint{1.579985in}{1.040119in}}%
\pgfpathlineto{\pgfqpoint{1.580150in}{1.327586in}}%
\pgfpathlineto{\pgfqpoint{1.580363in}{0.976083in}}%
\pgfpathlineto{\pgfqpoint{1.580684in}{0.976308in}}%
\pgfpathlineto{\pgfqpoint{1.580745in}{0.503868in}}%
\pgfpathlineto{\pgfqpoint{1.581322in}{0.504087in}}%
\pgfpathlineto{\pgfqpoint{1.587299in}{0.505950in}}%
\pgfpathlineto{\pgfqpoint{1.587354in}{0.522440in}}%
\pgfpathlineto{\pgfqpoint{1.587434in}{1.327707in}}%
\pgfpathlineto{\pgfqpoint{1.588230in}{1.327704in}}%
\pgfpathlineto{\pgfqpoint{1.596307in}{1.327580in}}%
\pgfpathlineto{\pgfqpoint{1.597236in}{0.503889in}}%
\pgfpathlineto{\pgfqpoint{1.597380in}{0.503955in}}%
\pgfpathlineto{\pgfqpoint{1.603779in}{0.505887in}}%
\pgfpathlineto{\pgfqpoint{1.603833in}{0.519473in}}%
\pgfpathlineto{\pgfqpoint{1.603904in}{1.327709in}}%
\pgfpathlineto{\pgfqpoint{1.604711in}{1.327701in}}%
\pgfpathlineto{\pgfqpoint{1.612129in}{1.327647in}}%
\pgfpathlineto{\pgfqpoint{1.612624in}{1.001301in}}%
\pgfpathlineto{\pgfqpoint{1.612953in}{1.016848in}}%
\pgfpathlineto{\pgfqpoint{1.613118in}{1.327582in}}%
\pgfpathlineto{\pgfqpoint{1.613646in}{0.975901in}}%
\pgfpathlineto{\pgfqpoint{1.613654in}{0.976024in}}%
\pgfpathlineto{\pgfqpoint{1.613719in}{0.503888in}}%
\pgfpathlineto{\pgfqpoint{1.614414in}{0.504123in}}%
\pgfpathlineto{\pgfqpoint{1.620264in}{0.505896in}}%
\pgfpathlineto{\pgfqpoint{1.620317in}{0.519517in}}%
\pgfpathlineto{\pgfqpoint{1.620389in}{1.327709in}}%
\pgfpathlineto{\pgfqpoint{1.621197in}{1.327702in}}%
\pgfpathlineto{\pgfqpoint{1.628615in}{1.327648in}}%
\pgfpathlineto{\pgfqpoint{1.629109in}{0.976142in}}%
\pgfpathlineto{\pgfqpoint{1.629439in}{0.990070in}}%
\pgfpathlineto{\pgfqpoint{1.629604in}{1.327588in}}%
\pgfpathlineto{\pgfqpoint{1.630130in}{0.975930in}}%
\pgfpathlineto{\pgfqpoint{1.630139in}{0.976280in}}%
\pgfpathlineto{\pgfqpoint{1.630204in}{0.503892in}}%
\pgfpathlineto{\pgfqpoint{1.630892in}{0.504124in}}%
\pgfpathlineto{\pgfqpoint{1.636749in}{0.505899in}}%
\pgfpathlineto{\pgfqpoint{1.636802in}{0.519528in}}%
\pgfpathlineto{\pgfqpoint{1.636874in}{1.327709in}}%
\pgfpathlineto{\pgfqpoint{1.637682in}{1.327702in}}%
\pgfpathlineto{\pgfqpoint{1.645100in}{1.327648in}}%
\pgfpathlineto{\pgfqpoint{1.645594in}{0.976038in}}%
\pgfpathlineto{\pgfqpoint{1.645924in}{0.977854in}}%
\pgfpathlineto{\pgfqpoint{1.646088in}{1.327588in}}%
\pgfpathlineto{\pgfqpoint{1.646615in}{0.975764in}}%
\pgfpathlineto{\pgfqpoint{1.646623in}{0.976183in}}%
\pgfpathlineto{\pgfqpoint{1.646691in}{0.503903in}}%
\pgfpathlineto{\pgfqpoint{1.647316in}{0.504106in}}%
\pgfpathlineto{\pgfqpoint{1.653233in}{0.505899in}}%
\pgfpathlineto{\pgfqpoint{1.653287in}{0.519527in}}%
\pgfpathlineto{\pgfqpoint{1.653359in}{1.327709in}}%
\pgfpathlineto{\pgfqpoint{1.654167in}{1.327702in}}%
\pgfpathlineto{\pgfqpoint{1.661585in}{1.327648in}}%
\pgfpathlineto{\pgfqpoint{1.662079in}{0.976122in}}%
\pgfpathlineto{\pgfqpoint{1.662409in}{0.985111in}}%
\pgfpathlineto{\pgfqpoint{1.662574in}{1.327588in}}%
\pgfpathlineto{\pgfqpoint{1.663100in}{0.975873in}}%
\pgfpathlineto{\pgfqpoint{1.663108in}{0.976275in}}%
\pgfpathlineto{\pgfqpoint{1.663168in}{0.503925in}}%
\pgfpathlineto{\pgfqpoint{1.663845in}{0.504120in}}%
\pgfpathlineto{\pgfqpoint{1.669717in}{0.505885in}}%
\pgfpathlineto{\pgfqpoint{1.669770in}{0.518925in}}%
\pgfpathlineto{\pgfqpoint{1.669840in}{1.327709in}}%
\pgfpathlineto{\pgfqpoint{1.670633in}{1.327701in}}%
\pgfpathlineto{\pgfqpoint{1.678052in}{1.327651in}}%
\pgfpathlineto{\pgfqpoint{1.678216in}{0.976143in}}%
\pgfpathlineto{\pgfqpoint{1.678876in}{1.039239in}}%
\pgfpathlineto{\pgfqpoint{1.679041in}{1.327605in}}%
\pgfpathlineto{\pgfqpoint{1.679554in}{0.968852in}}%
\pgfpathlineto{\pgfqpoint{1.679593in}{0.975324in}}%
\pgfpathlineto{\pgfqpoint{1.679661in}{0.503902in}}%
\pgfpathlineto{\pgfqpoint{1.680222in}{0.504080in}}%
\pgfpathlineto{\pgfqpoint{1.686203in}{0.505891in}}%
\pgfpathlineto{\pgfqpoint{1.686257in}{0.519494in}}%
\pgfpathlineto{\pgfqpoint{1.686328in}{1.327709in}}%
\pgfpathlineto{\pgfqpoint{1.687136in}{1.327701in}}%
\pgfpathlineto{\pgfqpoint{1.694554in}{1.327647in}}%
\pgfpathlineto{\pgfqpoint{1.695049in}{0.996723in}}%
\pgfpathlineto{\pgfqpoint{1.695379in}{1.021512in}}%
\pgfpathlineto{\pgfqpoint{1.695544in}{1.327581in}}%
\pgfpathlineto{\pgfqpoint{1.696003in}{0.976028in}}%
\pgfpathlineto{\pgfqpoint{1.696078in}{0.976095in}}%
\pgfpathlineto{\pgfqpoint{1.696143in}{0.503891in}}%
\pgfpathlineto{\pgfqpoint{1.696828in}{0.504123in}}%
\pgfpathlineto{\pgfqpoint{1.702688in}{0.505898in}}%
\pgfpathlineto{\pgfqpoint{1.702742in}{0.519523in}}%
\pgfpathlineto{\pgfqpoint{1.702813in}{1.327709in}}%
\pgfpathlineto{\pgfqpoint{1.703621in}{1.327702in}}%
\pgfpathlineto{\pgfqpoint{1.711039in}{1.327648in}}%
\pgfpathlineto{\pgfqpoint{1.711534in}{0.976141in}}%
\pgfpathlineto{\pgfqpoint{1.711863in}{0.989247in}}%
\pgfpathlineto{\pgfqpoint{1.712028in}{1.327588in}}%
\pgfpathlineto{\pgfqpoint{1.712555in}{0.975909in}}%
\pgfpathlineto{\pgfqpoint{1.712563in}{0.976282in}}%
\pgfpathlineto{\pgfqpoint{1.712628in}{0.503892in}}%
\pgfpathlineto{\pgfqpoint{1.713313in}{0.504123in}}%
\pgfpathlineto{\pgfqpoint{1.719176in}{0.505931in}}%
\pgfpathlineto{\pgfqpoint{1.719230in}{0.521159in}}%
\pgfpathlineto{\pgfqpoint{1.719310in}{1.327708in}}%
\pgfpathlineto{\pgfqpoint{1.720128in}{1.327703in}}%
\pgfpathlineto{\pgfqpoint{1.727876in}{1.327627in}}%
\pgfpathlineto{\pgfqpoint{1.728823in}{0.967883in}}%
\pgfpathlineto{\pgfqpoint{1.728931in}{0.976181in}}%
\pgfpathlineto{\pgfqpoint{1.729040in}{0.966078in}}%
\pgfpathlineto{\pgfqpoint{1.729050in}{0.974418in}}%
\pgfpathlineto{\pgfqpoint{1.729112in}{0.503897in}}%
\pgfpathlineto{\pgfqpoint{1.729774in}{0.504109in}}%
\pgfpathlineto{\pgfqpoint{1.735670in}{0.506207in}}%
\pgfpathlineto{\pgfqpoint{1.735735in}{0.631097in}}%
\pgfpathlineto{\pgfqpoint{1.735793in}{1.327708in}}%
\pgfpathlineto{\pgfqpoint{1.736524in}{1.327703in}}%
\pgfpathlineto{\pgfqpoint{1.744272in}{1.327624in}}%
\pgfpathlineto{\pgfqpoint{1.744437in}{1.307497in}}%
\pgfpathlineto{\pgfqpoint{1.744601in}{1.327599in}}%
\pgfpathlineto{\pgfqpoint{1.744766in}{0.940251in}}%
\pgfpathlineto{\pgfqpoint{1.745401in}{0.975066in}}%
\pgfpathlineto{\pgfqpoint{1.745524in}{0.975237in}}%
\pgfpathlineto{\pgfqpoint{1.745532in}{0.975149in}}%
\pgfpathlineto{\pgfqpoint{1.745598in}{0.503882in}}%
\pgfpathlineto{\pgfqpoint{1.746285in}{0.504114in}}%
\pgfpathlineto{\pgfqpoint{1.752152in}{0.506078in}}%
\pgfpathlineto{\pgfqpoint{1.752213in}{0.587710in}}%
\pgfpathlineto{\pgfqpoint{1.752356in}{1.327709in}}%
\pgfpathlineto{\pgfqpoint{1.752894in}{1.327707in}}%
\pgfpathlineto{\pgfqpoint{1.760477in}{1.327651in}}%
\pgfpathlineto{\pgfqpoint{1.760642in}{0.974646in}}%
\pgfpathlineto{\pgfqpoint{1.761301in}{0.975617in}}%
\pgfpathlineto{\pgfqpoint{1.761466in}{1.327605in}}%
\pgfpathlineto{\pgfqpoint{1.762021in}{0.909724in}}%
\pgfpathlineto{\pgfqpoint{1.762082in}{0.503876in}}%
\pgfpathlineto{\pgfqpoint{1.762773in}{0.504115in}}%
\pgfpathlineto{\pgfqpoint{1.768627in}{0.505888in}}%
\pgfpathlineto{\pgfqpoint{1.768681in}{0.519489in}}%
\pgfpathlineto{\pgfqpoint{1.768753in}{1.327709in}}%
\pgfpathlineto{\pgfqpoint{1.769561in}{1.327701in}}%
\pgfpathlineto{\pgfqpoint{1.776979in}{1.327647in}}%
\pgfpathlineto{\pgfqpoint{1.777474in}{0.988077in}}%
\pgfpathlineto{\pgfqpoint{1.777803in}{1.030432in}}%
\pgfpathlineto{\pgfqpoint{1.777968in}{1.327584in}}%
\pgfpathlineto{\pgfqpoint{1.778422in}{0.976002in}}%
\pgfpathlineto{\pgfqpoint{1.778502in}{0.976105in}}%
\pgfpathlineto{\pgfqpoint{1.778568in}{0.503887in}}%
\pgfpathlineto{\pgfqpoint{1.779267in}{0.504123in}}%
\pgfpathlineto{\pgfqpoint{1.785112in}{0.505894in}}%
\pgfpathlineto{\pgfqpoint{1.785162in}{0.518069in}}%
\pgfpathlineto{\pgfqpoint{1.785259in}{1.327710in}}%
\pgfpathlineto{\pgfqpoint{1.786042in}{1.327700in}}%
\pgfpathlineto{\pgfqpoint{1.793130in}{1.327661in}}%
\pgfpathlineto{\pgfqpoint{1.793625in}{0.976924in}}%
\pgfpathlineto{\pgfqpoint{1.793954in}{1.032577in}}%
\pgfpathlineto{\pgfqpoint{1.794119in}{1.327636in}}%
\pgfpathlineto{\pgfqpoint{1.794284in}{0.985990in}}%
\pgfpathlineto{\pgfqpoint{1.794756in}{1.327582in}}%
\pgfpathlineto{\pgfqpoint{1.795049in}{0.503857in}}%
\pgfpathlineto{\pgfqpoint{1.796066in}{0.504204in}}%
\pgfpathlineto{\pgfqpoint{1.801601in}{0.505933in}}%
\pgfpathlineto{\pgfqpoint{1.801657in}{0.522110in}}%
\pgfpathlineto{\pgfqpoint{1.801736in}{1.327708in}}%
\pgfpathlineto{\pgfqpoint{1.802545in}{1.327703in}}%
\pgfpathlineto{\pgfqpoint{1.810293in}{1.327617in}}%
\pgfpathlineto{\pgfqpoint{1.811536in}{0.503889in}}%
\pgfpathlineto{\pgfqpoint{1.812082in}{0.504070in}}%
\pgfpathlineto{\pgfqpoint{1.818092in}{0.506099in}}%
\pgfpathlineto{\pgfqpoint{1.818141in}{0.522542in}}%
\pgfpathlineto{\pgfqpoint{1.818211in}{1.327707in}}%
\pgfpathlineto{\pgfqpoint{1.818969in}{1.327704in}}%
\pgfpathlineto{\pgfqpoint{1.827047in}{1.327567in}}%
\pgfpathlineto{\pgfqpoint{1.828022in}{0.503890in}}%
\pgfpathlineto{\pgfqpoint{1.828139in}{0.503949in}}%
\pgfpathlineto{\pgfqpoint{1.834579in}{0.506207in}}%
\pgfpathlineto{\pgfqpoint{1.834644in}{0.630906in}}%
\pgfpathlineto{\pgfqpoint{1.834702in}{1.327708in}}%
\pgfpathlineto{\pgfqpoint{1.835435in}{1.327703in}}%
\pgfpathlineto{\pgfqpoint{1.843183in}{1.327624in}}%
\pgfpathlineto{\pgfqpoint{1.843347in}{1.306650in}}%
\pgfpathlineto{\pgfqpoint{1.843512in}{1.327599in}}%
\pgfpathlineto{\pgfqpoint{1.843677in}{0.923973in}}%
\pgfpathlineto{\pgfqpoint{1.844317in}{0.975886in}}%
\pgfpathlineto{\pgfqpoint{1.844433in}{0.971526in}}%
\pgfpathlineto{\pgfqpoint{1.844442in}{0.975422in}}%
\pgfpathlineto{\pgfqpoint{1.844510in}{0.503904in}}%
\pgfpathlineto{\pgfqpoint{1.845113in}{0.504092in}}%
\pgfpathlineto{\pgfqpoint{1.851061in}{0.506046in}}%
\pgfpathlineto{\pgfqpoint{1.851125in}{0.583772in}}%
\pgfpathlineto{\pgfqpoint{1.851174in}{1.327711in}}%
\pgfpathlineto{\pgfqpoint{1.851892in}{1.327709in}}%
\pgfpathlineto{\pgfqpoint{1.858815in}{1.327670in}}%
\pgfpathlineto{\pgfqpoint{1.858980in}{0.975058in}}%
\pgfpathlineto{\pgfqpoint{1.859640in}{0.975699in}}%
\pgfpathlineto{\pgfqpoint{1.859804in}{1.327653in}}%
\pgfpathlineto{\pgfqpoint{1.860464in}{1.327632in}}%
\pgfpathlineto{\pgfqpoint{1.860998in}{0.503901in}}%
\pgfpathlineto{\pgfqpoint{1.861594in}{0.504110in}}%
\pgfpathlineto{\pgfqpoint{1.867539in}{0.505943in}}%
\pgfpathlineto{\pgfqpoint{1.867594in}{0.521226in}}%
\pgfpathlineto{\pgfqpoint{1.867674in}{1.327708in}}%
\pgfpathlineto{\pgfqpoint{1.868490in}{1.327703in}}%
\pgfpathlineto{\pgfqpoint{1.876238in}{1.327630in}}%
\pgfpathlineto{\pgfqpoint{1.877171in}{0.855736in}}%
\pgfpathlineto{\pgfqpoint{1.877346in}{0.964531in}}%
\pgfpathlineto{\pgfqpoint{1.877403in}{0.976088in}}%
\pgfpathlineto{\pgfqpoint{1.877411in}{0.975612in}}%
\pgfpathlineto{\pgfqpoint{1.877476in}{0.503883in}}%
\pgfpathlineto{\pgfqpoint{1.878172in}{0.504120in}}%
\pgfpathlineto{\pgfqpoint{1.884022in}{0.505896in}}%
\pgfpathlineto{\pgfqpoint{1.884075in}{0.519661in}}%
\pgfpathlineto{\pgfqpoint{1.884148in}{1.327709in}}%
\pgfpathlineto{\pgfqpoint{1.884962in}{1.327702in}}%
\pgfpathlineto{\pgfqpoint{1.892380in}{1.327647in}}%
\pgfpathlineto{\pgfqpoint{1.892875in}{0.976151in}}%
\pgfpathlineto{\pgfqpoint{1.893205in}{0.992914in}}%
\pgfpathlineto{\pgfqpoint{1.893369in}{1.327579in}}%
\pgfpathlineto{\pgfqpoint{1.893888in}{0.975520in}}%
\pgfpathlineto{\pgfqpoint{1.893896in}{0.975845in}}%
\pgfpathlineto{\pgfqpoint{1.893962in}{0.503887in}}%
\pgfpathlineto{\pgfqpoint{1.894664in}{0.504124in}}%
\pgfpathlineto{\pgfqpoint{1.900518in}{0.506213in}}%
\pgfpathlineto{\pgfqpoint{1.900584in}{0.631328in}}%
\pgfpathlineto{\pgfqpoint{1.900641in}{1.327708in}}%
\pgfpathlineto{\pgfqpoint{1.901371in}{1.327704in}}%
\pgfpathlineto{\pgfqpoint{1.909449in}{1.327601in}}%
\pgfpathlineto{\pgfqpoint{1.909614in}{0.954066in}}%
\pgfpathlineto{\pgfqpoint{1.910375in}{0.974747in}}%
\pgfpathlineto{\pgfqpoint{1.910381in}{0.974650in}}%
\pgfpathlineto{\pgfqpoint{1.910446in}{0.503880in}}%
\pgfpathlineto{\pgfqpoint{1.911152in}{0.504117in}}%
\pgfpathlineto{\pgfqpoint{1.916994in}{0.505914in}}%
\pgfpathlineto{\pgfqpoint{1.917047in}{0.520490in}}%
\pgfpathlineto{\pgfqpoint{1.917123in}{1.327708in}}%
\pgfpathlineto{\pgfqpoint{1.917963in}{1.327702in}}%
\pgfpathlineto{\pgfqpoint{1.925382in}{1.327642in}}%
\pgfpathlineto{\pgfqpoint{1.925876in}{0.975322in}}%
\pgfpathlineto{\pgfqpoint{1.926206in}{0.975709in}}%
\pgfpathlineto{\pgfqpoint{1.926370in}{1.238478in}}%
\pgfpathlineto{\pgfqpoint{1.926872in}{0.847804in}}%
\pgfpathlineto{\pgfqpoint{1.926926in}{0.503893in}}%
\pgfpathlineto{\pgfqpoint{1.927586in}{0.504106in}}%
\pgfpathlineto{\pgfqpoint{1.933476in}{0.505891in}}%
\pgfpathlineto{\pgfqpoint{1.933530in}{0.519503in}}%
\pgfpathlineto{\pgfqpoint{1.933601in}{1.327709in}}%
\pgfpathlineto{\pgfqpoint{1.934410in}{1.327701in}}%
\pgfpathlineto{\pgfqpoint{1.941828in}{1.327647in}}%
\pgfpathlineto{\pgfqpoint{1.942323in}{0.991186in}}%
\pgfpathlineto{\pgfqpoint{1.942652in}{1.027208in}}%
\pgfpathlineto{\pgfqpoint{1.942817in}{1.327582in}}%
\pgfpathlineto{\pgfqpoint{1.943274in}{0.976023in}}%
\pgfpathlineto{\pgfqpoint{1.943351in}{0.976127in}}%
\pgfpathlineto{\pgfqpoint{1.943416in}{0.503890in}}%
\pgfpathlineto{\pgfqpoint{1.944055in}{0.504109in}}%
\pgfpathlineto{\pgfqpoint{1.949961in}{0.505897in}}%
\pgfpathlineto{\pgfqpoint{1.950014in}{0.519520in}}%
\pgfpathlineto{\pgfqpoint{1.950086in}{1.327709in}}%
\pgfpathlineto{\pgfqpoint{1.950894in}{1.327702in}}%
\pgfpathlineto{\pgfqpoint{1.958312in}{1.327648in}}%
\pgfpathlineto{\pgfqpoint{1.958806in}{0.976096in}}%
\pgfpathlineto{\pgfqpoint{1.959136in}{0.986427in}}%
\pgfpathlineto{\pgfqpoint{1.959301in}{1.327588in}}%
\pgfpathlineto{\pgfqpoint{1.959827in}{0.975819in}}%
\pgfpathlineto{\pgfqpoint{1.959835in}{0.976229in}}%
\pgfpathlineto{\pgfqpoint{1.959901in}{0.503891in}}%
\pgfpathlineto{\pgfqpoint{1.960593in}{0.504124in}}%
\pgfpathlineto{\pgfqpoint{1.966446in}{0.505898in}}%
\pgfpathlineto{\pgfqpoint{1.966499in}{0.519521in}}%
\pgfpathlineto{\pgfqpoint{1.966571in}{1.327709in}}%
\pgfpathlineto{\pgfqpoint{1.967378in}{1.327702in}}%
\pgfpathlineto{\pgfqpoint{1.974797in}{1.327648in}}%
\pgfpathlineto{\pgfqpoint{1.975291in}{0.976024in}}%
\pgfpathlineto{\pgfqpoint{1.975621in}{0.980265in}}%
\pgfpathlineto{\pgfqpoint{1.975785in}{1.327589in}}%
\pgfpathlineto{\pgfqpoint{1.976312in}{0.975715in}}%
\pgfpathlineto{\pgfqpoint{1.976320in}{0.976131in}}%
\pgfpathlineto{\pgfqpoint{1.976386in}{0.503890in}}%
\pgfpathlineto{\pgfqpoint{1.977080in}{0.504124in}}%
\pgfpathlineto{\pgfqpoint{1.982931in}{0.505903in}}%
\pgfpathlineto{\pgfqpoint{1.982984in}{0.519655in}}%
\pgfpathlineto{\pgfqpoint{1.983056in}{1.327709in}}%
\pgfpathlineto{\pgfqpoint{1.983875in}{1.327702in}}%
\pgfpathlineto{\pgfqpoint{1.991294in}{1.327647in}}%
\pgfpathlineto{\pgfqpoint{1.991788in}{0.976110in}}%
\pgfpathlineto{\pgfqpoint{1.992118in}{0.990737in}}%
\pgfpathlineto{\pgfqpoint{1.992283in}{1.327585in}}%
\pgfpathlineto{\pgfqpoint{1.992577in}{0.975839in}}%
\pgfpathlineto{\pgfqpoint{1.992805in}{0.976124in}}%
\pgfpathlineto{\pgfqpoint{1.992871in}{0.503888in}}%
\pgfpathlineto{\pgfqpoint{1.993560in}{0.504122in}}%
\pgfpathlineto{\pgfqpoint{1.999417in}{0.505911in}}%
\pgfpathlineto{\pgfqpoint{1.999467in}{0.518874in}}%
\pgfpathlineto{\pgfqpoint{1.999631in}{1.327710in}}%
\pgfpathlineto{\pgfqpoint{2.000353in}{1.327701in}}%
\pgfpathlineto{\pgfqpoint{2.007442in}{1.327659in}}%
\pgfpathlineto{\pgfqpoint{2.007607in}{1.290940in}}%
\pgfpathlineto{\pgfqpoint{2.007772in}{1.327652in}}%
\pgfpathlineto{\pgfqpoint{2.007937in}{0.974814in}}%
\pgfpathlineto{\pgfqpoint{2.008596in}{0.975705in}}%
\pgfpathlineto{\pgfqpoint{2.008760in}{1.327610in}}%
\pgfpathlineto{\pgfqpoint{2.009295in}{0.867902in}}%
\pgfpathlineto{\pgfqpoint{2.009355in}{0.503882in}}%
\pgfpathlineto{\pgfqpoint{2.009987in}{0.504091in}}%
\pgfpathlineto{\pgfqpoint{2.015900in}{0.505881in}}%
\pgfpathlineto{\pgfqpoint{2.015954in}{0.519447in}}%
\pgfpathlineto{\pgfqpoint{2.016025in}{1.327709in}}%
\pgfpathlineto{\pgfqpoint{2.016832in}{1.327701in}}%
\pgfpathlineto{\pgfqpoint{2.024250in}{1.327646in}}%
\pgfpathlineto{\pgfqpoint{2.024415in}{0.976875in}}%
\pgfpathlineto{\pgfqpoint{2.025074in}{0.987470in}}%
\pgfpathlineto{\pgfqpoint{2.025239in}{1.327574in}}%
\pgfpathlineto{\pgfqpoint{2.025695in}{0.975341in}}%
\pgfpathlineto{\pgfqpoint{2.025775in}{0.975453in}}%
\pgfpathlineto{\pgfqpoint{2.025842in}{0.503901in}}%
\pgfpathlineto{\pgfqpoint{2.026561in}{0.504125in}}%
\pgfpathlineto{\pgfqpoint{2.032385in}{0.505890in}}%
\pgfpathlineto{\pgfqpoint{2.032439in}{0.519488in}}%
\pgfpathlineto{\pgfqpoint{2.032510in}{1.327709in}}%
\pgfpathlineto{\pgfqpoint{2.033317in}{1.327702in}}%
\pgfpathlineto{\pgfqpoint{2.040736in}{1.327647in}}%
\pgfpathlineto{\pgfqpoint{2.041230in}{0.976647in}}%
\pgfpathlineto{\pgfqpoint{2.041560in}{1.026393in}}%
\pgfpathlineto{\pgfqpoint{2.041725in}{1.327586in}}%
\pgfpathlineto{\pgfqpoint{2.042252in}{0.976183in}}%
\pgfpathlineto{\pgfqpoint{2.042260in}{0.976424in}}%
\pgfpathlineto{\pgfqpoint{2.042323in}{0.503924in}}%
\pgfpathlineto{\pgfqpoint{2.042899in}{0.504088in}}%
\pgfpathlineto{\pgfqpoint{2.048870in}{0.505896in}}%
\pgfpathlineto{\pgfqpoint{2.048924in}{0.519518in}}%
\pgfpathlineto{\pgfqpoint{2.048995in}{1.327709in}}%
\pgfpathlineto{\pgfqpoint{2.049803in}{1.327702in}}%
\pgfpathlineto{\pgfqpoint{2.057222in}{1.327648in}}%
\pgfpathlineto{\pgfqpoint{2.057716in}{0.975894in}}%
\pgfpathlineto{\pgfqpoint{2.058046in}{0.977127in}}%
\pgfpathlineto{\pgfqpoint{2.058210in}{1.327590in}}%
\pgfpathlineto{\pgfqpoint{2.058736in}{0.975747in}}%
\pgfpathlineto{\pgfqpoint{2.058745in}{0.975976in}}%
\pgfpathlineto{\pgfqpoint{2.058810in}{0.503886in}}%
\pgfpathlineto{\pgfqpoint{2.059515in}{0.504124in}}%
\pgfpathlineto{\pgfqpoint{2.065355in}{0.505894in}}%
\pgfpathlineto{\pgfqpoint{2.065408in}{0.519504in}}%
\pgfpathlineto{\pgfqpoint{2.065480in}{1.327709in}}%
\pgfpathlineto{\pgfqpoint{2.066287in}{1.327702in}}%
\pgfpathlineto{\pgfqpoint{2.073706in}{1.327648in}}%
\pgfpathlineto{\pgfqpoint{2.074200in}{0.976065in}}%
\pgfpathlineto{\pgfqpoint{2.074530in}{0.990286in}}%
\pgfpathlineto{\pgfqpoint{2.074695in}{1.327589in}}%
\pgfpathlineto{\pgfqpoint{2.075221in}{0.975893in}}%
\pgfpathlineto{\pgfqpoint{2.075229in}{0.976134in}}%
\pgfpathlineto{\pgfqpoint{2.075289in}{0.503907in}}%
\pgfpathlineto{\pgfqpoint{2.075964in}{0.504115in}}%
\pgfpathlineto{\pgfqpoint{2.081840in}{0.505895in}}%
\pgfpathlineto{\pgfqpoint{2.081893in}{0.519509in}}%
\pgfpathlineto{\pgfqpoint{2.081965in}{1.327709in}}%
\pgfpathlineto{\pgfqpoint{2.082772in}{1.327702in}}%
\pgfpathlineto{\pgfqpoint{2.090190in}{1.327648in}}%
\pgfpathlineto{\pgfqpoint{2.090685in}{0.975975in}}%
\pgfpathlineto{\pgfqpoint{2.091015in}{0.980831in}}%
\pgfpathlineto{\pgfqpoint{2.091179in}{1.327589in}}%
\pgfpathlineto{\pgfqpoint{2.091706in}{0.975716in}}%
\pgfpathlineto{\pgfqpoint{2.091714in}{0.976022in}}%
\pgfpathlineto{\pgfqpoint{2.091780in}{0.503887in}}%
\pgfpathlineto{\pgfqpoint{2.092484in}{0.504124in}}%
\pgfpathlineto{\pgfqpoint{2.098330in}{0.505978in}}%
\pgfpathlineto{\pgfqpoint{2.098388in}{0.525065in}}%
\pgfpathlineto{\pgfqpoint{2.098445in}{1.327708in}}%
\pgfpathlineto{\pgfqpoint{2.099158in}{1.327704in}}%
\pgfpathlineto{\pgfqpoint{2.107236in}{1.327610in}}%
\pgfpathlineto{\pgfqpoint{2.107849in}{0.849407in}}%
\pgfpathlineto{\pgfqpoint{2.108137in}{0.976414in}}%
\pgfpathlineto{\pgfqpoint{2.108191in}{0.974354in}}%
\pgfpathlineto{\pgfqpoint{2.108199in}{0.975588in}}%
\pgfpathlineto{\pgfqpoint{2.108264in}{0.503884in}}%
\pgfpathlineto{\pgfqpoint{2.108966in}{0.504122in}}%
\pgfpathlineto{\pgfqpoint{2.114809in}{0.505893in}}%
\pgfpathlineto{\pgfqpoint{2.114863in}{0.519499in}}%
\pgfpathlineto{\pgfqpoint{2.114935in}{1.327709in}}%
\pgfpathlineto{\pgfqpoint{2.115742in}{1.327702in}}%
\pgfpathlineto{\pgfqpoint{2.123160in}{1.327647in}}%
\pgfpathlineto{\pgfqpoint{2.123655in}{0.976638in}}%
\pgfpathlineto{\pgfqpoint{2.123984in}{1.024163in}}%
\pgfpathlineto{\pgfqpoint{2.124149in}{1.327586in}}%
\pgfpathlineto{\pgfqpoint{2.124392in}{0.976186in}}%
\pgfpathlineto{\pgfqpoint{2.124683in}{0.976455in}}%
\pgfpathlineto{\pgfqpoint{2.124746in}{0.503929in}}%
\pgfpathlineto{\pgfqpoint{2.125357in}{0.504101in}}%
\pgfpathlineto{\pgfqpoint{2.131306in}{0.506217in}}%
\pgfpathlineto{\pgfqpoint{2.131372in}{0.631524in}}%
\pgfpathlineto{\pgfqpoint{2.131429in}{1.327708in}}%
\pgfpathlineto{\pgfqpoint{2.132158in}{1.327704in}}%
\pgfpathlineto{\pgfqpoint{2.140236in}{1.327602in}}%
\pgfpathlineto{\pgfqpoint{2.141229in}{0.503879in}}%
\pgfpathlineto{\pgfqpoint{2.141357in}{0.503947in}}%
\pgfpathlineto{\pgfqpoint{2.147779in}{0.505887in}}%
\pgfpathlineto{\pgfqpoint{2.147833in}{0.519500in}}%
\pgfpathlineto{\pgfqpoint{2.147904in}{1.327709in}}%
\pgfpathlineto{\pgfqpoint{2.148715in}{1.327701in}}%
\pgfpathlineto{\pgfqpoint{2.156133in}{1.327646in}}%
\pgfpathlineto{\pgfqpoint{2.156298in}{0.977938in}}%
\pgfpathlineto{\pgfqpoint{2.156957in}{0.978804in}}%
\pgfpathlineto{\pgfqpoint{2.157122in}{1.327574in}}%
\pgfpathlineto{\pgfqpoint{2.157590in}{0.975407in}}%
\pgfpathlineto{\pgfqpoint{2.157654in}{0.975627in}}%
\pgfpathlineto{\pgfqpoint{2.157721in}{0.503904in}}%
\pgfpathlineto{\pgfqpoint{2.158279in}{0.504082in}}%
\pgfpathlineto{\pgfqpoint{2.164264in}{0.505893in}}%
\pgfpathlineto{\pgfqpoint{2.164317in}{0.519501in}}%
\pgfpathlineto{\pgfqpoint{2.164389in}{1.327709in}}%
\pgfpathlineto{\pgfqpoint{2.165196in}{1.327702in}}%
\pgfpathlineto{\pgfqpoint{2.172615in}{1.327647in}}%
\pgfpathlineto{\pgfqpoint{2.173109in}{0.976647in}}%
\pgfpathlineto{\pgfqpoint{2.173439in}{1.024422in}}%
\pgfpathlineto{\pgfqpoint{2.173604in}{1.327586in}}%
\pgfpathlineto{\pgfqpoint{2.173848in}{0.976174in}}%
\pgfpathlineto{\pgfqpoint{2.174138in}{0.976441in}}%
\pgfpathlineto{\pgfqpoint{2.174205in}{0.503899in}}%
\pgfpathlineto{\pgfqpoint{2.174891in}{0.504124in}}%
\pgfpathlineto{\pgfqpoint{2.180749in}{0.505899in}}%
\pgfpathlineto{\pgfqpoint{2.180802in}{0.519526in}}%
\pgfpathlineto{\pgfqpoint{2.180874in}{1.327709in}}%
\pgfpathlineto{\pgfqpoint{2.181682in}{1.327702in}}%
\pgfpathlineto{\pgfqpoint{2.189100in}{1.327648in}}%
\pgfpathlineto{\pgfqpoint{2.189594in}{0.975944in}}%
\pgfpathlineto{\pgfqpoint{2.189924in}{0.977018in}}%
\pgfpathlineto{\pgfqpoint{2.190088in}{1.327589in}}%
\pgfpathlineto{\pgfqpoint{2.190615in}{0.975860in}}%
\pgfpathlineto{\pgfqpoint{2.190623in}{0.976122in}}%
\pgfpathlineto{\pgfqpoint{2.190689in}{0.503890in}}%
\pgfpathlineto{\pgfqpoint{2.191382in}{0.504123in}}%
\pgfpathlineto{\pgfqpoint{2.197238in}{0.505951in}}%
\pgfpathlineto{\pgfqpoint{2.197295in}{0.523303in}}%
\pgfpathlineto{\pgfqpoint{2.197374in}{1.327707in}}%
\pgfpathlineto{\pgfqpoint{2.198150in}{1.327704in}}%
\pgfpathlineto{\pgfqpoint{2.206228in}{1.327527in}}%
\pgfpathlineto{\pgfqpoint{2.207176in}{0.503899in}}%
\pgfpathlineto{\pgfqpoint{2.207416in}{0.503983in}}%
\pgfpathlineto{\pgfqpoint{2.213724in}{0.505967in}}%
\pgfpathlineto{\pgfqpoint{2.213782in}{0.524902in}}%
\pgfpathlineto{\pgfqpoint{2.213839in}{1.327708in}}%
\pgfpathlineto{\pgfqpoint{2.214597in}{1.327703in}}%
\pgfpathlineto{\pgfqpoint{2.222345in}{1.327624in}}%
\pgfpathlineto{\pgfqpoint{2.223608in}{0.553145in}}%
\pgfpathlineto{\pgfqpoint{2.223653in}{0.503874in}}%
\pgfpathlineto{\pgfqpoint{2.224296in}{0.504097in}}%
\pgfpathlineto{\pgfqpoint{2.230203in}{0.505886in}}%
\pgfpathlineto{\pgfqpoint{2.230257in}{0.519485in}}%
\pgfpathlineto{\pgfqpoint{2.230328in}{1.327709in}}%
\pgfpathlineto{\pgfqpoint{2.231137in}{1.327701in}}%
\pgfpathlineto{\pgfqpoint{2.238556in}{1.327647in}}%
\pgfpathlineto{\pgfqpoint{2.239380in}{0.995376in}}%
\pgfpathlineto{\pgfqpoint{2.239545in}{1.327578in}}%
\pgfpathlineto{\pgfqpoint{2.240000in}{0.975568in}}%
\pgfpathlineto{\pgfqpoint{2.240078in}{0.975762in}}%
\pgfpathlineto{\pgfqpoint{2.240145in}{0.503904in}}%
\pgfpathlineto{\pgfqpoint{2.240704in}{0.504082in}}%
\pgfpathlineto{\pgfqpoint{2.246689in}{0.505906in}}%
\pgfpathlineto{\pgfqpoint{2.246742in}{0.519842in}}%
\pgfpathlineto{\pgfqpoint{2.246815in}{1.327709in}}%
\pgfpathlineto{\pgfqpoint{2.247645in}{1.327702in}}%
\pgfpathlineto{\pgfqpoint{2.255063in}{1.327646in}}%
\pgfpathlineto{\pgfqpoint{2.255558in}{0.976089in}}%
\pgfpathlineto{\pgfqpoint{2.255888in}{0.990574in}}%
\pgfpathlineto{\pgfqpoint{2.256052in}{1.327575in}}%
\pgfpathlineto{\pgfqpoint{2.256555in}{0.975363in}}%
\pgfpathlineto{\pgfqpoint{2.256563in}{0.975512in}}%
\pgfpathlineto{\pgfqpoint{2.256628in}{0.503883in}}%
\pgfpathlineto{\pgfqpoint{2.257331in}{0.504121in}}%
\pgfpathlineto{\pgfqpoint{2.263173in}{0.505891in}}%
\pgfpathlineto{\pgfqpoint{2.263227in}{0.519491in}}%
\pgfpathlineto{\pgfqpoint{2.263298in}{1.327709in}}%
\pgfpathlineto{\pgfqpoint{2.264105in}{1.327702in}}%
\pgfpathlineto{\pgfqpoint{2.271524in}{1.327647in}}%
\pgfpathlineto{\pgfqpoint{2.272018in}{0.976618in}}%
\pgfpathlineto{\pgfqpoint{2.272348in}{1.025012in}}%
\pgfpathlineto{\pgfqpoint{2.272513in}{1.327586in}}%
\pgfpathlineto{\pgfqpoint{2.273040in}{0.976210in}}%
\pgfpathlineto{\pgfqpoint{2.273048in}{0.976430in}}%
\pgfpathlineto{\pgfqpoint{2.273112in}{0.503904in}}%
\pgfpathlineto{\pgfqpoint{2.273805in}{0.504123in}}%
\pgfpathlineto{\pgfqpoint{2.279658in}{0.505902in}}%
\pgfpathlineto{\pgfqpoint{2.279712in}{0.519653in}}%
\pgfpathlineto{\pgfqpoint{2.279784in}{1.327709in}}%
\pgfpathlineto{\pgfqpoint{2.280603in}{1.327702in}}%
\pgfpathlineto{\pgfqpoint{2.288021in}{1.327647in}}%
\pgfpathlineto{\pgfqpoint{2.288515in}{0.975938in}}%
\pgfpathlineto{\pgfqpoint{2.288845in}{0.978027in}}%
\pgfpathlineto{\pgfqpoint{2.289009in}{1.327586in}}%
\pgfpathlineto{\pgfqpoint{2.289442in}{0.975735in}}%
\pgfpathlineto{\pgfqpoint{2.289533in}{0.975983in}}%
\pgfpathlineto{\pgfqpoint{2.289600in}{0.503904in}}%
\pgfpathlineto{\pgfqpoint{2.290157in}{0.504082in}}%
\pgfpathlineto{\pgfqpoint{2.296143in}{0.505894in}}%
\pgfpathlineto{\pgfqpoint{2.296196in}{0.519510in}}%
\pgfpathlineto{\pgfqpoint{2.296268in}{1.327709in}}%
\pgfpathlineto{\pgfqpoint{2.297076in}{1.327702in}}%
\pgfpathlineto{\pgfqpoint{2.304494in}{1.327648in}}%
\pgfpathlineto{\pgfqpoint{2.304989in}{0.976090in}}%
\pgfpathlineto{\pgfqpoint{2.305319in}{0.992679in}}%
\pgfpathlineto{\pgfqpoint{2.305483in}{1.327589in}}%
\pgfpathlineto{\pgfqpoint{2.306009in}{0.975931in}}%
\pgfpathlineto{\pgfqpoint{2.306017in}{0.976160in}}%
\pgfpathlineto{\pgfqpoint{2.306083in}{0.503888in}}%
\pgfpathlineto{\pgfqpoint{2.306784in}{0.504124in}}%
\pgfpathlineto{\pgfqpoint{2.312627in}{0.505895in}}%
\pgfpathlineto{\pgfqpoint{2.312681in}{0.519521in}}%
\pgfpathlineto{\pgfqpoint{2.312753in}{1.327709in}}%
\pgfpathlineto{\pgfqpoint{2.313562in}{1.327702in}}%
\pgfpathlineto{\pgfqpoint{2.320980in}{1.327648in}}%
\pgfpathlineto{\pgfqpoint{2.321475in}{0.976007in}}%
\pgfpathlineto{\pgfqpoint{2.321804in}{0.984773in}}%
\pgfpathlineto{\pgfqpoint{2.321969in}{1.327589in}}%
\pgfpathlineto{\pgfqpoint{2.322494in}{0.975745in}}%
\pgfpathlineto{\pgfqpoint{2.322501in}{0.976067in}}%
\pgfpathlineto{\pgfqpoint{2.322569in}{0.503897in}}%
\pgfpathlineto{\pgfqpoint{2.323251in}{0.504118in}}%
\pgfpathlineto{\pgfqpoint{2.329112in}{0.505894in}}%
\pgfpathlineto{\pgfqpoint{2.329166in}{0.519507in}}%
\pgfpathlineto{\pgfqpoint{2.329238in}{1.327709in}}%
\pgfpathlineto{\pgfqpoint{2.330045in}{1.327702in}}%
\pgfpathlineto{\pgfqpoint{2.337463in}{1.327648in}}%
\pgfpathlineto{\pgfqpoint{2.337958in}{0.976020in}}%
\pgfpathlineto{\pgfqpoint{2.338287in}{0.985429in}}%
\pgfpathlineto{\pgfqpoint{2.338452in}{1.327589in}}%
\pgfpathlineto{\pgfqpoint{2.338979in}{0.975796in}}%
\pgfpathlineto{\pgfqpoint{2.338987in}{0.976079in}}%
\pgfpathlineto{\pgfqpoint{2.339053in}{0.503887in}}%
\pgfpathlineto{\pgfqpoint{2.339756in}{0.504124in}}%
\pgfpathlineto{\pgfqpoint{2.345600in}{0.505927in}}%
\pgfpathlineto{\pgfqpoint{2.345654in}{0.521133in}}%
\pgfpathlineto{\pgfqpoint{2.345735in}{1.327708in}}%
\pgfpathlineto{\pgfqpoint{2.346553in}{1.327703in}}%
\pgfpathlineto{\pgfqpoint{2.354301in}{1.327627in}}%
\pgfpathlineto{\pgfqpoint{2.355238in}{0.946090in}}%
\pgfpathlineto{\pgfqpoint{2.355351in}{0.975766in}}%
\pgfpathlineto{\pgfqpoint{2.355474in}{0.975390in}}%
\pgfpathlineto{\pgfqpoint{2.355537in}{0.503892in}}%
\pgfpathlineto{\pgfqpoint{2.356245in}{0.504121in}}%
\pgfpathlineto{\pgfqpoint{2.362082in}{0.505890in}}%
\pgfpathlineto{\pgfqpoint{2.362136in}{0.519490in}}%
\pgfpathlineto{\pgfqpoint{2.362207in}{1.327709in}}%
\pgfpathlineto{\pgfqpoint{2.363014in}{1.327702in}}%
\pgfpathlineto{\pgfqpoint{2.370433in}{1.327647in}}%
\pgfpathlineto{\pgfqpoint{2.370927in}{0.976825in}}%
\pgfpathlineto{\pgfqpoint{2.371257in}{1.031038in}}%
\pgfpathlineto{\pgfqpoint{2.371422in}{1.327586in}}%
\pgfpathlineto{\pgfqpoint{2.371949in}{0.976202in}}%
\pgfpathlineto{\pgfqpoint{2.371956in}{0.976381in}}%
\pgfpathlineto{\pgfqpoint{2.372020in}{0.503906in}}%
\pgfpathlineto{\pgfqpoint{2.372592in}{0.504088in}}%
\pgfpathlineto{\pgfqpoint{2.378567in}{0.505897in}}%
\pgfpathlineto{\pgfqpoint{2.378620in}{0.519520in}}%
\pgfpathlineto{\pgfqpoint{2.378692in}{1.327709in}}%
\pgfpathlineto{\pgfqpoint{2.379500in}{1.327702in}}%
\pgfpathlineto{\pgfqpoint{2.386918in}{1.327648in}}%
\pgfpathlineto{\pgfqpoint{2.387412in}{0.975935in}}%
\pgfpathlineto{\pgfqpoint{2.387742in}{0.977082in}}%
\pgfpathlineto{\pgfqpoint{2.387906in}{1.327589in}}%
\pgfpathlineto{\pgfqpoint{2.388433in}{0.975848in}}%
\pgfpathlineto{\pgfqpoint{2.388442in}{0.976080in}}%
\pgfpathlineto{\pgfqpoint{2.388507in}{0.503889in}}%
\pgfpathlineto{\pgfqpoint{2.389202in}{0.504123in}}%
\pgfpathlineto{\pgfqpoint{2.395055in}{0.505929in}}%
\pgfpathlineto{\pgfqpoint{2.395109in}{0.521142in}}%
\pgfpathlineto{\pgfqpoint{2.395189in}{1.327708in}}%
\pgfpathlineto{\pgfqpoint{2.396008in}{1.327703in}}%
\pgfpathlineto{\pgfqpoint{2.403756in}{1.327627in}}%
\pgfpathlineto{\pgfqpoint{2.404669in}{0.974273in}}%
\pgfpathlineto{\pgfqpoint{2.404794in}{0.975946in}}%
\pgfpathlineto{\pgfqpoint{2.404918in}{0.959747in}}%
\pgfpathlineto{\pgfqpoint{2.404929in}{0.967091in}}%
\pgfpathlineto{\pgfqpoint{2.404991in}{0.503894in}}%
\pgfpathlineto{\pgfqpoint{2.405684in}{0.504117in}}%
\pgfpathlineto{\pgfqpoint{2.411537in}{0.505899in}}%
\pgfpathlineto{\pgfqpoint{2.411591in}{0.519699in}}%
\pgfpathlineto{\pgfqpoint{2.411663in}{1.327709in}}%
\pgfpathlineto{\pgfqpoint{2.412486in}{1.327701in}}%
\pgfpathlineto{\pgfqpoint{2.419904in}{1.327646in}}%
\pgfpathlineto{\pgfqpoint{2.420399in}{0.980157in}}%
\pgfpathlineto{\pgfqpoint{2.420728in}{1.038729in}}%
\pgfpathlineto{\pgfqpoint{2.420893in}{1.327577in}}%
\pgfpathlineto{\pgfqpoint{2.421138in}{0.975347in}}%
\pgfpathlineto{\pgfqpoint{2.421413in}{0.975787in}}%
\pgfpathlineto{\pgfqpoint{2.421471in}{0.503924in}}%
\pgfpathlineto{\pgfqpoint{2.422133in}{0.504111in}}%
\pgfpathlineto{\pgfqpoint{2.428027in}{0.505960in}}%
\pgfpathlineto{\pgfqpoint{2.428082in}{0.522658in}}%
\pgfpathlineto{\pgfqpoint{2.428159in}{1.327707in}}%
\pgfpathlineto{\pgfqpoint{2.428970in}{1.327704in}}%
\pgfpathlineto{\pgfqpoint{2.437048in}{1.327566in}}%
\pgfpathlineto{\pgfqpoint{2.437958in}{0.503917in}}%
\pgfpathlineto{\pgfqpoint{2.438122in}{0.503960in}}%
\pgfpathlineto{\pgfqpoint{2.444511in}{0.505952in}}%
\pgfpathlineto{\pgfqpoint{2.444567in}{0.522623in}}%
\pgfpathlineto{\pgfqpoint{2.444644in}{1.327707in}}%
\pgfpathlineto{\pgfqpoint{2.445456in}{1.327703in}}%
\pgfpathlineto{\pgfqpoint{2.453204in}{1.327612in}}%
\pgfpathlineto{\pgfqpoint{2.453369in}{1.314752in}}%
\pgfpathlineto{\pgfqpoint{2.453534in}{1.327549in}}%
\pgfpathlineto{\pgfqpoint{2.454454in}{0.503920in}}%
\pgfpathlineto{\pgfqpoint{2.454607in}{0.503959in}}%
\pgfpathlineto{\pgfqpoint{2.461005in}{0.506507in}}%
\pgfpathlineto{\pgfqpoint{2.461071in}{0.645740in}}%
\pgfpathlineto{\pgfqpoint{2.461125in}{1.327707in}}%
\pgfpathlineto{\pgfqpoint{2.461826in}{1.327704in}}%
\pgfpathlineto{\pgfqpoint{2.469903in}{1.327580in}}%
\pgfpathlineto{\pgfqpoint{2.470726in}{0.975546in}}%
\pgfpathlineto{\pgfqpoint{2.470864in}{0.975578in}}%
\pgfpathlineto{\pgfqpoint{2.470926in}{0.503881in}}%
\pgfpathlineto{\pgfqpoint{2.471612in}{0.504116in}}%
\pgfpathlineto{\pgfqpoint{2.477476in}{0.505893in}}%
\pgfpathlineto{\pgfqpoint{2.477530in}{0.519502in}}%
\pgfpathlineto{\pgfqpoint{2.477601in}{1.327709in}}%
\pgfpathlineto{\pgfqpoint{2.478409in}{1.327702in}}%
\pgfpathlineto{\pgfqpoint{2.485827in}{1.327647in}}%
\pgfpathlineto{\pgfqpoint{2.486321in}{0.976952in}}%
\pgfpathlineto{\pgfqpoint{2.486651in}{1.032991in}}%
\pgfpathlineto{\pgfqpoint{2.486816in}{1.327585in}}%
\pgfpathlineto{\pgfqpoint{2.487268in}{0.976219in}}%
\pgfpathlineto{\pgfqpoint{2.487351in}{0.976332in}}%
\pgfpathlineto{\pgfqpoint{2.487416in}{0.503891in}}%
\pgfpathlineto{\pgfqpoint{2.488099in}{0.504123in}}%
\pgfpathlineto{\pgfqpoint{2.493961in}{0.505898in}}%
\pgfpathlineto{\pgfqpoint{2.494014in}{0.519524in}}%
\pgfpathlineto{\pgfqpoint{2.494086in}{1.327709in}}%
\pgfpathlineto{\pgfqpoint{2.494894in}{1.327702in}}%
\pgfpathlineto{\pgfqpoint{2.502312in}{1.327648in}}%
\pgfpathlineto{\pgfqpoint{2.502807in}{0.975995in}}%
\pgfpathlineto{\pgfqpoint{2.503136in}{0.977713in}}%
\pgfpathlineto{\pgfqpoint{2.503300in}{1.327589in}}%
\pgfpathlineto{\pgfqpoint{2.503827in}{0.975714in}}%
\pgfpathlineto{\pgfqpoint{2.503835in}{0.976101in}}%
\pgfpathlineto{\pgfqpoint{2.503895in}{0.503926in}}%
\pgfpathlineto{\pgfqpoint{2.504554in}{0.504112in}}%
\pgfpathlineto{\pgfqpoint{2.510446in}{0.505897in}}%
\pgfpathlineto{\pgfqpoint{2.510499in}{0.519529in}}%
\pgfpathlineto{\pgfqpoint{2.510571in}{1.327709in}}%
\pgfpathlineto{\pgfqpoint{2.511380in}{1.327702in}}%
\pgfpathlineto{\pgfqpoint{2.518798in}{1.327648in}}%
\pgfpathlineto{\pgfqpoint{2.519293in}{0.976088in}}%
\pgfpathlineto{\pgfqpoint{2.519623in}{0.988785in}}%
\pgfpathlineto{\pgfqpoint{2.519787in}{1.327588in}}%
\pgfpathlineto{\pgfqpoint{2.520312in}{0.975796in}}%
\pgfpathlineto{\pgfqpoint{2.520319in}{0.976203in}}%
\pgfpathlineto{\pgfqpoint{2.520385in}{0.503904in}}%
\pgfpathlineto{\pgfqpoint{2.521077in}{0.504123in}}%
\pgfpathlineto{\pgfqpoint{2.526930in}{0.505896in}}%
\pgfpathlineto{\pgfqpoint{2.526984in}{0.519513in}}%
\pgfpathlineto{\pgfqpoint{2.527056in}{1.327709in}}%
\pgfpathlineto{\pgfqpoint{2.527863in}{1.327702in}}%
\pgfpathlineto{\pgfqpoint{2.535281in}{1.327648in}}%
\pgfpathlineto{\pgfqpoint{2.535776in}{0.975988in}}%
\pgfpathlineto{\pgfqpoint{2.536106in}{0.981112in}}%
\pgfpathlineto{\pgfqpoint{2.536270in}{1.327589in}}%
\pgfpathlineto{\pgfqpoint{2.536797in}{0.975707in}}%
\pgfpathlineto{\pgfqpoint{2.536805in}{0.976046in}}%
\pgfpathlineto{\pgfqpoint{2.536871in}{0.503888in}}%
\pgfpathlineto{\pgfqpoint{2.537574in}{0.504124in}}%
\pgfpathlineto{\pgfqpoint{2.543430in}{0.506523in}}%
\pgfpathlineto{\pgfqpoint{2.543495in}{0.645674in}}%
\pgfpathlineto{\pgfqpoint{2.543549in}{1.327707in}}%
\pgfpathlineto{\pgfqpoint{2.544251in}{1.327704in}}%
\pgfpathlineto{\pgfqpoint{2.552329in}{1.327592in}}%
\pgfpathlineto{\pgfqpoint{2.553139in}{0.975915in}}%
\pgfpathlineto{\pgfqpoint{2.553282in}{0.976331in}}%
\pgfpathlineto{\pgfqpoint{2.553289in}{0.976155in}}%
\pgfpathlineto{\pgfqpoint{2.553351in}{0.503868in}}%
\pgfpathlineto{\pgfqpoint{2.554072in}{0.504129in}}%
\pgfpathlineto{\pgfqpoint{2.559903in}{0.505927in}}%
\pgfpathlineto{\pgfqpoint{2.559957in}{0.521127in}}%
\pgfpathlineto{\pgfqpoint{2.560038in}{1.327708in}}%
\pgfpathlineto{\pgfqpoint{2.560857in}{1.327703in}}%
\pgfpathlineto{\pgfqpoint{2.568605in}{1.327627in}}%
\pgfpathlineto{\pgfqpoint{2.569547in}{0.928511in}}%
\pgfpathlineto{\pgfqpoint{2.569712in}{0.975499in}}%
\pgfpathlineto{\pgfqpoint{2.569782in}{0.786229in}}%
\pgfpathlineto{\pgfqpoint{2.569838in}{0.503864in}}%
\pgfpathlineto{\pgfqpoint{2.570553in}{0.504122in}}%
\pgfpathlineto{\pgfqpoint{2.576381in}{0.505856in}}%
\pgfpathlineto{\pgfqpoint{2.576441in}{0.520725in}}%
\pgfpathlineto{\pgfqpoint{2.576520in}{1.327708in}}%
\pgfpathlineto{\pgfqpoint{2.577258in}{1.327704in}}%
\pgfpathlineto{\pgfqpoint{2.585336in}{1.327597in}}%
\pgfpathlineto{\pgfqpoint{2.585501in}{0.930465in}}%
\pgfpathlineto{\pgfqpoint{2.586253in}{0.975349in}}%
\pgfpathlineto{\pgfqpoint{2.586256in}{0.975409in}}%
\pgfpathlineto{\pgfqpoint{2.586260in}{0.975382in}}%
\pgfpathlineto{\pgfqpoint{2.586328in}{0.503902in}}%
\pgfpathlineto{\pgfqpoint{2.586884in}{0.504078in}}%
\pgfpathlineto{\pgfqpoint{2.592870in}{0.505889in}}%
\pgfpathlineto{\pgfqpoint{2.592923in}{0.519486in}}%
\pgfpathlineto{\pgfqpoint{2.592995in}{1.327709in}}%
\pgfpathlineto{\pgfqpoint{2.593802in}{1.327702in}}%
\pgfpathlineto{\pgfqpoint{2.601220in}{1.327647in}}%
\pgfpathlineto{\pgfqpoint{2.601715in}{0.976684in}}%
\pgfpathlineto{\pgfqpoint{2.602045in}{1.027786in}}%
\pgfpathlineto{\pgfqpoint{2.602209in}{1.327586in}}%
\pgfpathlineto{\pgfqpoint{2.602736in}{0.976176in}}%
\pgfpathlineto{\pgfqpoint{2.602745in}{0.976425in}}%
\pgfpathlineto{\pgfqpoint{2.602811in}{0.503904in}}%
\pgfpathlineto{\pgfqpoint{2.603507in}{0.504124in}}%
\pgfpathlineto{\pgfqpoint{2.609355in}{0.505902in}}%
\pgfpathlineto{\pgfqpoint{2.609409in}{0.519650in}}%
\pgfpathlineto{\pgfqpoint{2.609481in}{1.327709in}}%
\pgfpathlineto{\pgfqpoint{2.610299in}{1.327702in}}%
\pgfpathlineto{\pgfqpoint{2.617718in}{1.327647in}}%
\pgfpathlineto{\pgfqpoint{2.618212in}{0.975950in}}%
\pgfpathlineto{\pgfqpoint{2.618542in}{0.978862in}}%
\pgfpathlineto{\pgfqpoint{2.618706in}{1.327587in}}%
\pgfpathlineto{\pgfqpoint{2.619221in}{0.975918in}}%
\pgfpathlineto{\pgfqpoint{2.619230in}{0.975994in}}%
\pgfpathlineto{\pgfqpoint{2.619291in}{0.503868in}}%
\pgfpathlineto{\pgfqpoint{2.619996in}{0.504124in}}%
\pgfpathlineto{\pgfqpoint{2.625840in}{0.505896in}}%
\pgfpathlineto{\pgfqpoint{2.625893in}{0.519535in}}%
\pgfpathlineto{\pgfqpoint{2.625965in}{1.327709in}}%
\pgfpathlineto{\pgfqpoint{2.626776in}{1.327702in}}%
\pgfpathlineto{\pgfqpoint{2.634194in}{1.327648in}}%
\pgfpathlineto{\pgfqpoint{2.634688in}{0.976121in}}%
\pgfpathlineto{\pgfqpoint{2.635018in}{0.994297in}}%
\pgfpathlineto{\pgfqpoint{2.635183in}{1.327588in}}%
\pgfpathlineto{\pgfqpoint{2.635706in}{0.975851in}}%
\pgfpathlineto{\pgfqpoint{2.635714in}{0.976216in}}%
\pgfpathlineto{\pgfqpoint{2.635780in}{0.503889in}}%
\pgfpathlineto{\pgfqpoint{2.636478in}{0.504124in}}%
\pgfpathlineto{\pgfqpoint{2.642324in}{0.505896in}}%
\pgfpathlineto{\pgfqpoint{2.642378in}{0.519513in}}%
\pgfpathlineto{\pgfqpoint{2.642450in}{1.327709in}}%
\pgfpathlineto{\pgfqpoint{2.643257in}{1.327702in}}%
\pgfpathlineto{\pgfqpoint{2.650675in}{1.327648in}}%
\pgfpathlineto{\pgfqpoint{2.651170in}{0.975979in}}%
\pgfpathlineto{\pgfqpoint{2.651500in}{0.979215in}}%
\pgfpathlineto{\pgfqpoint{2.651664in}{1.327589in}}%
\pgfpathlineto{\pgfqpoint{2.652191in}{0.975706in}}%
\pgfpathlineto{\pgfqpoint{2.652198in}{0.976046in}}%
\pgfpathlineto{\pgfqpoint{2.652265in}{0.503897in}}%
\pgfpathlineto{\pgfqpoint{2.652883in}{0.504100in}}%
\pgfpathlineto{\pgfqpoint{2.658812in}{0.505928in}}%
\pgfpathlineto{\pgfqpoint{2.658867in}{0.521132in}}%
\pgfpathlineto{\pgfqpoint{2.658947in}{1.327708in}}%
\pgfpathlineto{\pgfqpoint{2.659766in}{1.327703in}}%
\pgfpathlineto{\pgfqpoint{2.667514in}{1.327627in}}%
\pgfpathlineto{\pgfqpoint{2.668423in}{0.969992in}}%
\pgfpathlineto{\pgfqpoint{2.668554in}{0.975890in}}%
\pgfpathlineto{\pgfqpoint{2.668676in}{0.959448in}}%
\pgfpathlineto{\pgfqpoint{2.668686in}{0.966419in}}%
\pgfpathlineto{\pgfqpoint{2.668749in}{0.503891in}}%
\pgfpathlineto{\pgfqpoint{2.669436in}{0.504115in}}%
\pgfpathlineto{\pgfqpoint{2.675295in}{0.505894in}}%
\pgfpathlineto{\pgfqpoint{2.675348in}{0.519582in}}%
\pgfpathlineto{\pgfqpoint{2.675420in}{1.327709in}}%
\pgfpathlineto{\pgfqpoint{2.676236in}{1.327701in}}%
\pgfpathlineto{\pgfqpoint{2.683654in}{1.327647in}}%
\pgfpathlineto{\pgfqpoint{2.684148in}{0.986958in}}%
\pgfpathlineto{\pgfqpoint{2.684478in}{1.031598in}}%
\pgfpathlineto{\pgfqpoint{2.684643in}{1.327582in}}%
\pgfpathlineto{\pgfqpoint{2.685161in}{0.975796in}}%
\pgfpathlineto{\pgfqpoint{2.685169in}{0.976022in}}%
\pgfpathlineto{\pgfqpoint{2.685239in}{0.503931in}}%
\pgfpathlineto{\pgfqpoint{2.685833in}{0.504095in}}%
\pgfpathlineto{\pgfqpoint{2.691779in}{0.505896in}}%
\pgfpathlineto{\pgfqpoint{2.691833in}{0.519514in}}%
\pgfpathlineto{\pgfqpoint{2.691904in}{1.327709in}}%
\pgfpathlineto{\pgfqpoint{2.692712in}{1.327702in}}%
\pgfpathlineto{\pgfqpoint{2.700130in}{1.327648in}}%
\pgfpathlineto{\pgfqpoint{2.700625in}{0.976132in}}%
\pgfpathlineto{\pgfqpoint{2.700954in}{0.990937in}}%
\pgfpathlineto{\pgfqpoint{2.701119in}{1.327588in}}%
\pgfpathlineto{\pgfqpoint{2.701646in}{0.975931in}}%
\pgfpathlineto{\pgfqpoint{2.701654in}{0.976243in}}%
\pgfpathlineto{\pgfqpoint{2.701714in}{0.503882in}}%
\pgfpathlineto{\pgfqpoint{2.702394in}{0.504119in}}%
\pgfpathlineto{\pgfqpoint{2.708264in}{0.505898in}}%
\pgfpathlineto{\pgfqpoint{2.708317in}{0.519522in}}%
\pgfpathlineto{\pgfqpoint{2.708389in}{1.327709in}}%
\pgfpathlineto{\pgfqpoint{2.709197in}{1.327702in}}%
\pgfpathlineto{\pgfqpoint{2.716615in}{1.327648in}}%
\pgfpathlineto{\pgfqpoint{2.717110in}{0.976022in}}%
\pgfpathlineto{\pgfqpoint{2.717440in}{0.979255in}}%
\pgfpathlineto{\pgfqpoint{2.717604in}{1.327589in}}%
\pgfpathlineto{\pgfqpoint{2.718130in}{0.975719in}}%
\pgfpathlineto{\pgfqpoint{2.718139in}{0.976132in}}%
\pgfpathlineto{\pgfqpoint{2.718204in}{0.503890in}}%
\pgfpathlineto{\pgfqpoint{2.718898in}{0.504124in}}%
\pgfpathlineto{\pgfqpoint{2.724760in}{0.506178in}}%
\pgfpathlineto{\pgfqpoint{2.724811in}{0.541255in}}%
\pgfpathlineto{\pgfqpoint{2.724907in}{1.327708in}}%
\pgfpathlineto{\pgfqpoint{2.725599in}{1.327706in}}%
\pgfpathlineto{\pgfqpoint{2.733512in}{1.327625in}}%
\pgfpathlineto{\pgfqpoint{2.733677in}{0.975177in}}%
\pgfpathlineto{\pgfqpoint{2.734349in}{0.976106in}}%
\pgfpathlineto{\pgfqpoint{2.734619in}{0.976226in}}%
\pgfpathlineto{\pgfqpoint{2.734623in}{0.976205in}}%
\pgfpathlineto{\pgfqpoint{2.734689in}{0.503889in}}%
\pgfpathlineto{\pgfqpoint{2.735376in}{0.504122in}}%
\pgfpathlineto{\pgfqpoint{2.741236in}{0.505931in}}%
\pgfpathlineto{\pgfqpoint{2.741291in}{0.521168in}}%
\pgfpathlineto{\pgfqpoint{2.741371in}{1.327708in}}%
\pgfpathlineto{\pgfqpoint{2.742189in}{1.327703in}}%
\pgfpathlineto{\pgfqpoint{2.749937in}{1.327627in}}%
\pgfpathlineto{\pgfqpoint{2.750881in}{0.971230in}}%
\pgfpathlineto{\pgfqpoint{2.750991in}{0.976190in}}%
\pgfpathlineto{\pgfqpoint{2.751100in}{0.962537in}}%
\pgfpathlineto{\pgfqpoint{2.751105in}{0.967485in}}%
\pgfpathlineto{\pgfqpoint{2.751172in}{0.503875in}}%
\pgfpathlineto{\pgfqpoint{2.751895in}{0.504126in}}%
\pgfpathlineto{\pgfqpoint{2.757718in}{0.505893in}}%
\pgfpathlineto{\pgfqpoint{2.757772in}{0.519522in}}%
\pgfpathlineto{\pgfqpoint{2.757844in}{1.327709in}}%
\pgfpathlineto{\pgfqpoint{2.758654in}{1.327701in}}%
\pgfpathlineto{\pgfqpoint{2.766072in}{1.327647in}}%
\pgfpathlineto{\pgfqpoint{2.766567in}{0.978029in}}%
\pgfpathlineto{\pgfqpoint{2.766897in}{1.040341in}}%
\pgfpathlineto{\pgfqpoint{2.767061in}{1.327584in}}%
\pgfpathlineto{\pgfqpoint{2.767283in}{0.975965in}}%
\pgfpathlineto{\pgfqpoint{2.767593in}{0.976239in}}%
\pgfpathlineto{\pgfqpoint{2.767659in}{0.503889in}}%
\pgfpathlineto{\pgfqpoint{2.768324in}{0.504116in}}%
\pgfpathlineto{\pgfqpoint{2.774214in}{0.506116in}}%
\pgfpathlineto{\pgfqpoint{2.774262in}{0.523193in}}%
\pgfpathlineto{\pgfqpoint{2.774333in}{1.327707in}}%
\pgfpathlineto{\pgfqpoint{2.774993in}{1.327705in}}%
\pgfpathlineto{\pgfqpoint{2.783235in}{1.327580in}}%
\pgfpathlineto{\pgfqpoint{2.784144in}{0.503891in}}%
\pgfpathlineto{\pgfqpoint{2.784335in}{0.503969in}}%
\pgfpathlineto{\pgfqpoint{2.790688in}{0.505887in}}%
\pgfpathlineto{\pgfqpoint{2.790742in}{0.519475in}}%
\pgfpathlineto{\pgfqpoint{2.790813in}{1.327709in}}%
\pgfpathlineto{\pgfqpoint{2.791620in}{1.327701in}}%
\pgfpathlineto{\pgfqpoint{2.799039in}{1.327647in}}%
\pgfpathlineto{\pgfqpoint{2.799533in}{1.007772in}}%
\pgfpathlineto{\pgfqpoint{2.799863in}{1.010324in}}%
\pgfpathlineto{\pgfqpoint{2.800028in}{1.327581in}}%
\pgfpathlineto{\pgfqpoint{2.800567in}{0.922619in}}%
\pgfpathlineto{\pgfqpoint{2.800629in}{0.503894in}}%
\pgfpathlineto{\pgfqpoint{2.801239in}{0.504098in}}%
\pgfpathlineto{\pgfqpoint{2.807173in}{0.505895in}}%
\pgfpathlineto{\pgfqpoint{2.807227in}{0.519508in}}%
\pgfpathlineto{\pgfqpoint{2.807298in}{1.327709in}}%
\pgfpathlineto{\pgfqpoint{2.808106in}{1.327702in}}%
\pgfpathlineto{\pgfqpoint{2.815524in}{1.327648in}}%
\pgfpathlineto{\pgfqpoint{2.816018in}{0.976193in}}%
\pgfpathlineto{\pgfqpoint{2.816348in}{0.998397in}}%
\pgfpathlineto{\pgfqpoint{2.816513in}{1.327588in}}%
\pgfpathlineto{\pgfqpoint{2.817040in}{0.975895in}}%
\pgfpathlineto{\pgfqpoint{2.817048in}{0.976309in}}%
\pgfpathlineto{\pgfqpoint{2.817113in}{0.503890in}}%
\pgfpathlineto{\pgfqpoint{2.817807in}{0.504124in}}%
\pgfpathlineto{\pgfqpoint{2.823658in}{0.505897in}}%
\pgfpathlineto{\pgfqpoint{2.823711in}{0.519519in}}%
\pgfpathlineto{\pgfqpoint{2.823783in}{1.327709in}}%
\pgfpathlineto{\pgfqpoint{2.824591in}{1.327702in}}%
\pgfpathlineto{\pgfqpoint{2.832009in}{1.327648in}}%
\pgfpathlineto{\pgfqpoint{2.832503in}{0.975951in}}%
\pgfpathlineto{\pgfqpoint{2.832833in}{0.977509in}}%
\pgfpathlineto{\pgfqpoint{2.832997in}{1.327589in}}%
\pgfpathlineto{\pgfqpoint{2.833524in}{0.975689in}}%
\pgfpathlineto{\pgfqpoint{2.833532in}{0.976037in}}%
\pgfpathlineto{\pgfqpoint{2.833598in}{0.503888in}}%
\pgfpathlineto{\pgfqpoint{2.834298in}{0.504124in}}%
\pgfpathlineto{\pgfqpoint{2.840143in}{0.505895in}}%
\pgfpathlineto{\pgfqpoint{2.840196in}{0.519509in}}%
\pgfpathlineto{\pgfqpoint{2.840268in}{1.327709in}}%
\pgfpathlineto{\pgfqpoint{2.841075in}{1.327702in}}%
\pgfpathlineto{\pgfqpoint{2.848494in}{1.327648in}}%
\pgfpathlineto{\pgfqpoint{2.848988in}{0.976070in}}%
\pgfpathlineto{\pgfqpoint{2.849318in}{0.988930in}}%
\pgfpathlineto{\pgfqpoint{2.849483in}{1.327589in}}%
\pgfpathlineto{\pgfqpoint{2.850009in}{0.975857in}}%
\pgfpathlineto{\pgfqpoint{2.850017in}{0.976157in}}%
\pgfpathlineto{\pgfqpoint{2.850077in}{0.503924in}}%
\pgfpathlineto{\pgfqpoint{2.850705in}{0.504102in}}%
\pgfpathlineto{\pgfqpoint{2.856627in}{0.505895in}}%
\pgfpathlineto{\pgfqpoint{2.856681in}{0.519513in}}%
\pgfpathlineto{\pgfqpoint{2.856753in}{1.327709in}}%
\pgfpathlineto{\pgfqpoint{2.857560in}{1.327702in}}%
\pgfpathlineto{\pgfqpoint{2.864978in}{1.327648in}}%
\pgfpathlineto{\pgfqpoint{2.865473in}{0.975981in}}%
\pgfpathlineto{\pgfqpoint{2.865803in}{0.980641in}}%
\pgfpathlineto{\pgfqpoint{2.865967in}{1.327589in}}%
\pgfpathlineto{\pgfqpoint{2.866494in}{0.975682in}}%
\pgfpathlineto{\pgfqpoint{2.866502in}{0.976040in}}%
\pgfpathlineto{\pgfqpoint{2.866568in}{0.503887in}}%
\pgfpathlineto{\pgfqpoint{2.867257in}{0.504120in}}%
\pgfpathlineto{\pgfqpoint{2.873127in}{0.506494in}}%
\pgfpathlineto{\pgfqpoint{2.873178in}{0.540604in}}%
\pgfpathlineto{\pgfqpoint{2.873225in}{1.327710in}}%
\pgfpathlineto{\pgfqpoint{2.873861in}{1.327702in}}%
\pgfpathlineto{\pgfqpoint{2.881279in}{1.327656in}}%
\pgfpathlineto{\pgfqpoint{2.881773in}{0.976081in}}%
\pgfpathlineto{\pgfqpoint{2.882038in}{1.132172in}}%
\pgfpathlineto{\pgfqpoint{2.882092in}{1.327624in}}%
\pgfpathlineto{\pgfqpoint{2.882692in}{0.975948in}}%
\pgfpathlineto{\pgfqpoint{2.883055in}{0.503899in}}%
\pgfpathlineto{\pgfqpoint{2.883854in}{0.504150in}}%
\pgfpathlineto{\pgfqpoint{2.889611in}{0.506518in}}%
\pgfpathlineto{\pgfqpoint{2.889677in}{0.645734in}}%
\pgfpathlineto{\pgfqpoint{2.889731in}{1.327707in}}%
\pgfpathlineto{\pgfqpoint{2.890433in}{1.327704in}}%
\pgfpathlineto{\pgfqpoint{2.898510in}{1.327589in}}%
\pgfpathlineto{\pgfqpoint{2.899339in}{0.976831in}}%
\pgfpathlineto{\pgfqpoint{2.899464in}{0.976737in}}%
\pgfpathlineto{\pgfqpoint{2.899472in}{0.977146in}}%
\pgfpathlineto{\pgfqpoint{2.899532in}{0.503920in}}%
\pgfpathlineto{\pgfqpoint{2.900144in}{0.504100in}}%
\pgfpathlineto{\pgfqpoint{2.906085in}{0.505930in}}%
\pgfpathlineto{\pgfqpoint{2.906138in}{0.520771in}}%
\pgfpathlineto{\pgfqpoint{2.906219in}{1.327709in}}%
\pgfpathlineto{\pgfqpoint{2.906898in}{1.327706in}}%
\pgfpathlineto{\pgfqpoint{2.914811in}{1.327630in}}%
\pgfpathlineto{\pgfqpoint{2.915743in}{0.878607in}}%
\pgfpathlineto{\pgfqpoint{2.915949in}{0.975563in}}%
\pgfpathlineto{\pgfqpoint{2.915971in}{0.574427in}}%
\pgfpathlineto{\pgfqpoint{2.916017in}{0.503882in}}%
\pgfpathlineto{\pgfqpoint{2.916702in}{0.504114in}}%
\pgfpathlineto{\pgfqpoint{2.922573in}{0.505973in}}%
\pgfpathlineto{\pgfqpoint{2.922630in}{0.524975in}}%
\pgfpathlineto{\pgfqpoint{2.922687in}{1.327708in}}%
\pgfpathlineto{\pgfqpoint{2.923433in}{1.327704in}}%
\pgfpathlineto{\pgfqpoint{2.931510in}{1.327604in}}%
\pgfpathlineto{\pgfqpoint{2.931675in}{0.934245in}}%
\pgfpathlineto{\pgfqpoint{2.932345in}{0.968557in}}%
\pgfpathlineto{\pgfqpoint{2.932433in}{0.975032in}}%
\pgfpathlineto{\pgfqpoint{2.932442in}{0.974856in}}%
\pgfpathlineto{\pgfqpoint{2.932509in}{0.503900in}}%
\pgfpathlineto{\pgfqpoint{2.933087in}{0.504082in}}%
\pgfpathlineto{\pgfqpoint{2.939052in}{0.505887in}}%
\pgfpathlineto{\pgfqpoint{2.939105in}{0.519477in}}%
\pgfpathlineto{\pgfqpoint{2.939177in}{1.327709in}}%
\pgfpathlineto{\pgfqpoint{2.939984in}{1.327701in}}%
\pgfpathlineto{\pgfqpoint{2.947402in}{1.327647in}}%
\pgfpathlineto{\pgfqpoint{2.948226in}{1.007360in}}%
\pgfpathlineto{\pgfqpoint{2.948391in}{1.327580in}}%
\pgfpathlineto{\pgfqpoint{2.948811in}{0.975825in}}%
\pgfpathlineto{\pgfqpoint{2.948926in}{0.975914in}}%
\pgfpathlineto{\pgfqpoint{2.948994in}{0.503906in}}%
\pgfpathlineto{\pgfqpoint{2.949711in}{0.504130in}}%
\pgfpathlineto{\pgfqpoint{2.955537in}{0.505896in}}%
\pgfpathlineto{\pgfqpoint{2.955590in}{0.519518in}}%
\pgfpathlineto{\pgfqpoint{2.955662in}{1.327709in}}%
\pgfpathlineto{\pgfqpoint{2.956470in}{1.327702in}}%
\pgfpathlineto{\pgfqpoint{2.963888in}{1.327648in}}%
\pgfpathlineto{\pgfqpoint{2.964383in}{0.976254in}}%
\pgfpathlineto{\pgfqpoint{2.964713in}{1.001897in}}%
\pgfpathlineto{\pgfqpoint{2.964878in}{1.327587in}}%
\pgfpathlineto{\pgfqpoint{2.965403in}{0.976051in}}%
\pgfpathlineto{\pgfqpoint{2.965411in}{0.976372in}}%
\pgfpathlineto{\pgfqpoint{2.965479in}{0.503909in}}%
\pgfpathlineto{\pgfqpoint{2.966124in}{0.504112in}}%
\pgfpathlineto{\pgfqpoint{2.972021in}{0.505898in}}%
\pgfpathlineto{\pgfqpoint{2.972075in}{0.519524in}}%
\pgfpathlineto{\pgfqpoint{2.972147in}{1.327709in}}%
\pgfpathlineto{\pgfqpoint{2.972954in}{1.327702in}}%
\pgfpathlineto{\pgfqpoint{2.980373in}{1.327648in}}%
\pgfpathlineto{\pgfqpoint{2.980867in}{0.975952in}}%
\pgfpathlineto{\pgfqpoint{2.981197in}{0.977206in}}%
\pgfpathlineto{\pgfqpoint{2.981361in}{1.327589in}}%
\pgfpathlineto{\pgfqpoint{2.981888in}{0.975763in}}%
\pgfpathlineto{\pgfqpoint{2.981896in}{0.976089in}}%
\pgfpathlineto{\pgfqpoint{2.981961in}{0.503888in}}%
\pgfpathlineto{\pgfqpoint{2.982652in}{0.504122in}}%
\pgfpathlineto{\pgfqpoint{2.988506in}{0.505893in}}%
\pgfpathlineto{\pgfqpoint{2.988560in}{0.519367in}}%
\pgfpathlineto{\pgfqpoint{2.988631in}{1.327709in}}%
\pgfpathlineto{\pgfqpoint{2.989430in}{1.327702in}}%
\pgfpathlineto{\pgfqpoint{2.996849in}{1.327648in}}%
\pgfpathlineto{\pgfqpoint{2.997343in}{0.976287in}}%
\pgfpathlineto{\pgfqpoint{2.997673in}{1.005810in}}%
\pgfpathlineto{\pgfqpoint{2.997838in}{1.327593in}}%
\pgfpathlineto{\pgfqpoint{2.998373in}{0.975829in}}%
\pgfpathlineto{\pgfqpoint{2.998382in}{0.976087in}}%
\pgfpathlineto{\pgfqpoint{2.998445in}{0.503901in}}%
\pgfpathlineto{\pgfqpoint{2.999041in}{0.504094in}}%
\pgfpathlineto{\pgfqpoint{3.004991in}{0.505896in}}%
\pgfpathlineto{\pgfqpoint{3.005045in}{0.519514in}}%
\pgfpathlineto{\pgfqpoint{3.005116in}{1.327709in}}%
\pgfpathlineto{\pgfqpoint{3.005924in}{1.327702in}}%
\pgfpathlineto{\pgfqpoint{3.013342in}{1.327648in}}%
\pgfpathlineto{\pgfqpoint{3.013837in}{0.976061in}}%
\pgfpathlineto{\pgfqpoint{3.014167in}{0.987453in}}%
\pgfpathlineto{\pgfqpoint{3.014331in}{1.327589in}}%
\pgfpathlineto{\pgfqpoint{3.014858in}{0.975819in}}%
\pgfpathlineto{\pgfqpoint{3.014866in}{0.976150in}}%
\pgfpathlineto{\pgfqpoint{3.014931in}{0.503889in}}%
\pgfpathlineto{\pgfqpoint{3.015627in}{0.504123in}}%
\pgfpathlineto{\pgfqpoint{3.021489in}{0.506258in}}%
\pgfpathlineto{\pgfqpoint{3.021544in}{0.576528in}}%
\pgfpathlineto{\pgfqpoint{3.021585in}{1.327709in}}%
\pgfpathlineto{\pgfqpoint{3.022271in}{1.327702in}}%
\pgfpathlineto{\pgfqpoint{3.029689in}{1.327654in}}%
\pgfpathlineto{\pgfqpoint{3.030183in}{0.976227in}}%
\pgfpathlineto{\pgfqpoint{3.030513in}{1.000404in}}%
\pgfpathlineto{\pgfqpoint{3.030658in}{1.327617in}}%
\pgfpathlineto{\pgfqpoint{3.031206in}{0.892867in}}%
\pgfpathlineto{\pgfqpoint{3.031274in}{0.890966in}}%
\pgfpathlineto{\pgfqpoint{3.031349in}{0.918657in}}%
\pgfpathlineto{\pgfqpoint{3.031360in}{0.674504in}}%
\pgfpathlineto{\pgfqpoint{3.031412in}{0.503870in}}%
\pgfpathlineto{\pgfqpoint{3.032125in}{0.504122in}}%
\pgfpathlineto{\pgfqpoint{3.037961in}{0.505890in}}%
\pgfpathlineto{\pgfqpoint{3.038014in}{0.519490in}}%
\pgfpathlineto{\pgfqpoint{3.038086in}{1.327709in}}%
\pgfpathlineto{\pgfqpoint{3.038893in}{1.327702in}}%
\pgfpathlineto{\pgfqpoint{3.046311in}{1.327647in}}%
\pgfpathlineto{\pgfqpoint{3.046805in}{0.976927in}}%
\pgfpathlineto{\pgfqpoint{3.047135in}{1.033360in}}%
\pgfpathlineto{\pgfqpoint{3.047300in}{1.327585in}}%
\pgfpathlineto{\pgfqpoint{3.047827in}{0.976210in}}%
\pgfpathlineto{\pgfqpoint{3.047835in}{0.976332in}}%
\pgfpathlineto{\pgfqpoint{3.047901in}{0.503888in}}%
\pgfpathlineto{\pgfqpoint{3.048548in}{0.504110in}}%
\pgfpathlineto{\pgfqpoint{3.054446in}{0.505903in}}%
\pgfpathlineto{\pgfqpoint{3.054500in}{0.519678in}}%
\pgfpathlineto{\pgfqpoint{3.054572in}{1.327709in}}%
\pgfpathlineto{\pgfqpoint{3.055392in}{1.327702in}}%
\pgfpathlineto{\pgfqpoint{3.062810in}{1.327647in}}%
\pgfpathlineto{\pgfqpoint{3.063305in}{0.975969in}}%
\pgfpathlineto{\pgfqpoint{3.063634in}{0.978625in}}%
\pgfpathlineto{\pgfqpoint{3.063798in}{1.327585in}}%
\pgfpathlineto{\pgfqpoint{3.064105in}{0.975693in}}%
\pgfpathlineto{\pgfqpoint{3.064320in}{0.975996in}}%
\pgfpathlineto{\pgfqpoint{3.064386in}{0.503885in}}%
\pgfpathlineto{\pgfqpoint{3.065071in}{0.504119in}}%
\pgfpathlineto{\pgfqpoint{3.070936in}{0.505976in}}%
\pgfpathlineto{\pgfqpoint{3.070994in}{0.525013in}}%
\pgfpathlineto{\pgfqpoint{3.071051in}{1.327708in}}%
\pgfpathlineto{\pgfqpoint{3.071791in}{1.327704in}}%
\pgfpathlineto{\pgfqpoint{3.079869in}{1.327607in}}%
\pgfpathlineto{\pgfqpoint{3.080480in}{0.896847in}}%
\pgfpathlineto{\pgfqpoint{3.080798in}{0.975249in}}%
\pgfpathlineto{\pgfqpoint{3.080805in}{0.975139in}}%
\pgfpathlineto{\pgfqpoint{3.080873in}{0.503903in}}%
\pgfpathlineto{\pgfqpoint{3.081465in}{0.504087in}}%
\pgfpathlineto{\pgfqpoint{3.087416in}{0.505892in}}%
\pgfpathlineto{\pgfqpoint{3.087469in}{0.519541in}}%
\pgfpathlineto{\pgfqpoint{3.087541in}{1.327709in}}%
\pgfpathlineto{\pgfqpoint{3.088354in}{1.327701in}}%
\pgfpathlineto{\pgfqpoint{3.095772in}{1.327647in}}%
\pgfpathlineto{\pgfqpoint{3.096267in}{0.999024in}}%
\pgfpathlineto{\pgfqpoint{3.096596in}{1.019163in}}%
\pgfpathlineto{\pgfqpoint{3.096761in}{1.327581in}}%
\pgfpathlineto{\pgfqpoint{3.097282in}{0.975875in}}%
\pgfpathlineto{\pgfqpoint{3.097290in}{0.975965in}}%
\pgfpathlineto{\pgfqpoint{3.097358in}{0.503906in}}%
\pgfpathlineto{\pgfqpoint{3.098075in}{0.504130in}}%
\pgfpathlineto{\pgfqpoint{3.103906in}{0.505976in}}%
\pgfpathlineto{\pgfqpoint{3.103964in}{0.525014in}}%
\pgfpathlineto{\pgfqpoint{3.104021in}{1.327708in}}%
\pgfpathlineto{\pgfqpoint{3.104760in}{1.327704in}}%
\pgfpathlineto{\pgfqpoint{3.112838in}{1.327607in}}%
\pgfpathlineto{\pgfqpoint{3.113457in}{0.902023in}}%
\pgfpathlineto{\pgfqpoint{3.113769in}{0.975120in}}%
\pgfpathlineto{\pgfqpoint{3.113774in}{0.975062in}}%
\pgfpathlineto{\pgfqpoint{3.113837in}{0.503865in}}%
\pgfpathlineto{\pgfqpoint{3.114535in}{0.504116in}}%
\pgfpathlineto{\pgfqpoint{3.120398in}{0.506246in}}%
\pgfpathlineto{\pgfqpoint{3.120456in}{0.592175in}}%
\pgfpathlineto{\pgfqpoint{3.120498in}{1.327709in}}%
\pgfpathlineto{\pgfqpoint{3.121252in}{1.327701in}}%
\pgfpathlineto{\pgfqpoint{3.128670in}{1.327652in}}%
\pgfpathlineto{\pgfqpoint{3.128835in}{0.975484in}}%
\pgfpathlineto{\pgfqpoint{3.129495in}{0.976438in}}%
\pgfpathlineto{\pgfqpoint{3.129659in}{1.327609in}}%
\pgfpathlineto{\pgfqpoint{3.130172in}{0.898257in}}%
\pgfpathlineto{\pgfqpoint{3.130262in}{0.906688in}}%
\pgfpathlineto{\pgfqpoint{3.130263in}{0.889626in}}%
\pgfpathlineto{\pgfqpoint{3.130328in}{0.503923in}}%
\pgfpathlineto{\pgfqpoint{3.130893in}{0.504079in}}%
\pgfpathlineto{\pgfqpoint{3.136870in}{0.505888in}}%
\pgfpathlineto{\pgfqpoint{3.136923in}{0.519480in}}%
\pgfpathlineto{\pgfqpoint{3.136995in}{1.327709in}}%
\pgfpathlineto{\pgfqpoint{3.137802in}{1.327701in}}%
\pgfpathlineto{\pgfqpoint{3.145220in}{1.327647in}}%
\pgfpathlineto{\pgfqpoint{3.145715in}{1.002645in}}%
\pgfpathlineto{\pgfqpoint{3.146045in}{1.015486in}}%
\pgfpathlineto{\pgfqpoint{3.146209in}{1.327581in}}%
\pgfpathlineto{\pgfqpoint{3.146736in}{0.975890in}}%
\pgfpathlineto{\pgfqpoint{3.146745in}{0.975996in}}%
\pgfpathlineto{\pgfqpoint{3.146810in}{0.503889in}}%
\pgfpathlineto{\pgfqpoint{3.147502in}{0.504122in}}%
\pgfpathlineto{\pgfqpoint{3.153355in}{0.505896in}}%
\pgfpathlineto{\pgfqpoint{3.153405in}{0.518074in}}%
\pgfpathlineto{\pgfqpoint{3.153501in}{1.327710in}}%
\pgfpathlineto{\pgfqpoint{3.154284in}{1.327700in}}%
\pgfpathlineto{\pgfqpoint{3.161373in}{1.327661in}}%
\pgfpathlineto{\pgfqpoint{3.161867in}{0.979780in}}%
\pgfpathlineto{\pgfqpoint{3.162197in}{1.039134in}}%
\pgfpathlineto{\pgfqpoint{3.162362in}{1.327636in}}%
\pgfpathlineto{\pgfqpoint{3.162526in}{0.979760in}}%
\pgfpathlineto{\pgfqpoint{3.162997in}{1.327580in}}%
\pgfpathlineto{\pgfqpoint{3.163295in}{0.503882in}}%
\pgfpathlineto{\pgfqpoint{3.164295in}{0.504203in}}%
\pgfpathlineto{\pgfqpoint{3.169840in}{0.505887in}}%
\pgfpathlineto{\pgfqpoint{3.169893in}{0.519474in}}%
\pgfpathlineto{\pgfqpoint{3.169965in}{1.327709in}}%
\pgfpathlineto{\pgfqpoint{3.170772in}{1.327701in}}%
\pgfpathlineto{\pgfqpoint{3.178190in}{1.327647in}}%
\pgfpathlineto{\pgfqpoint{3.178684in}{1.002426in}}%
\pgfpathlineto{\pgfqpoint{3.179014in}{1.015709in}}%
\pgfpathlineto{\pgfqpoint{3.179179in}{1.327581in}}%
\pgfpathlineto{\pgfqpoint{3.179706in}{0.975835in}}%
\pgfpathlineto{\pgfqpoint{3.179714in}{0.975966in}}%
\pgfpathlineto{\pgfqpoint{3.179780in}{0.503887in}}%
\pgfpathlineto{\pgfqpoint{3.180480in}{0.504124in}}%
\pgfpathlineto{\pgfqpoint{3.186328in}{0.505941in}}%
\pgfpathlineto{\pgfqpoint{3.186373in}{0.517633in}}%
\pgfpathlineto{\pgfqpoint{3.186436in}{1.327710in}}%
\pgfpathlineto{\pgfqpoint{3.187375in}{1.327707in}}%
\pgfpathlineto{\pgfqpoint{3.194299in}{1.327660in}}%
\pgfpathlineto{\pgfqpoint{3.194464in}{1.300147in}}%
\pgfpathlineto{\pgfqpoint{3.194629in}{1.327654in}}%
\pgfpathlineto{\pgfqpoint{3.194793in}{0.974714in}}%
\pgfpathlineto{\pgfqpoint{3.195453in}{0.975648in}}%
\pgfpathlineto{\pgfqpoint{3.195618in}{1.327616in}}%
\pgfpathlineto{\pgfqpoint{3.196193in}{0.913998in}}%
\pgfpathlineto{\pgfqpoint{3.196208in}{0.672600in}}%
\pgfpathlineto{\pgfqpoint{3.196265in}{0.503918in}}%
\pgfpathlineto{\pgfqpoint{3.196960in}{0.504112in}}%
\pgfpathlineto{\pgfqpoint{3.202809in}{0.505885in}}%
\pgfpathlineto{\pgfqpoint{3.202863in}{0.519489in}}%
\pgfpathlineto{\pgfqpoint{3.202934in}{1.327709in}}%
\pgfpathlineto{\pgfqpoint{3.203745in}{1.327701in}}%
\pgfpathlineto{\pgfqpoint{3.211163in}{1.327647in}}%
\pgfpathlineto{\pgfqpoint{3.211987in}{0.984548in}}%
\pgfpathlineto{\pgfqpoint{3.212152in}{1.327577in}}%
\pgfpathlineto{\pgfqpoint{3.212615in}{0.975281in}}%
\pgfpathlineto{\pgfqpoint{3.212684in}{0.975593in}}%
\pgfpathlineto{\pgfqpoint{3.212752in}{0.503902in}}%
\pgfpathlineto{\pgfqpoint{3.213472in}{0.504127in}}%
\pgfpathlineto{\pgfqpoint{3.219297in}{0.505919in}}%
\pgfpathlineto{\pgfqpoint{3.219350in}{0.520510in}}%
\pgfpathlineto{\pgfqpoint{3.219426in}{1.327709in}}%
\pgfpathlineto{\pgfqpoint{3.220266in}{1.327702in}}%
\pgfpathlineto{\pgfqpoint{3.227685in}{1.327642in}}%
\pgfpathlineto{\pgfqpoint{3.227849in}{1.298297in}}%
\pgfpathlineto{\pgfqpoint{3.228014in}{1.327631in}}%
\pgfpathlineto{\pgfqpoint{3.228939in}{0.852666in}}%
\pgfpathlineto{\pgfqpoint{3.229161in}{0.975614in}}%
\pgfpathlineto{\pgfqpoint{3.229161in}{0.975614in}}%
\pgfusepath{stroke}%
\end{pgfscope}%
\begin{pgfscope}%
\pgfsetrectcap%
\pgfsetmiterjoin%
\pgfsetlinewidth{0.803000pt}%
\definecolor{currentstroke}{rgb}{0.000000,0.000000,0.000000}%
\pgfsetstrokecolor{currentstroke}%
\pgfsetdash{}{0pt}%
\pgfpathmoveto{\pgfqpoint{0.632797in}{0.462654in}}%
\pgfpathlineto{\pgfqpoint{0.632797in}{1.368904in}}%
\pgfusepath{stroke}%
\end{pgfscope}%
\begin{pgfscope}%
\pgfsetrectcap%
\pgfsetmiterjoin%
\pgfsetlinewidth{0.803000pt}%
\definecolor{currentstroke}{rgb}{0.000000,0.000000,0.000000}%
\pgfsetstrokecolor{currentstroke}%
\pgfsetdash{}{0pt}%
\pgfpathmoveto{\pgfqpoint{3.352797in}{0.462654in}}%
\pgfpathlineto{\pgfqpoint{3.352797in}{1.368904in}}%
\pgfusepath{stroke}%
\end{pgfscope}%
\begin{pgfscope}%
\pgfsetrectcap%
\pgfsetmiterjoin%
\pgfsetlinewidth{0.803000pt}%
\definecolor{currentstroke}{rgb}{0.000000,0.000000,0.000000}%
\pgfsetstrokecolor{currentstroke}%
\pgfsetdash{}{0pt}%
\pgfpathmoveto{\pgfqpoint{0.632797in}{0.462654in}}%
\pgfpathlineto{\pgfqpoint{3.352797in}{0.462654in}}%
\pgfusepath{stroke}%
\end{pgfscope}%
\begin{pgfscope}%
\pgfsetrectcap%
\pgfsetmiterjoin%
\pgfsetlinewidth{0.803000pt}%
\definecolor{currentstroke}{rgb}{0.000000,0.000000,0.000000}%
\pgfsetstrokecolor{currentstroke}%
\pgfsetdash{}{0pt}%
\pgfpathmoveto{\pgfqpoint{0.632797in}{1.368904in}}%
\pgfpathlineto{\pgfqpoint{3.352797in}{1.368904in}}%
\pgfusepath{stroke}%
\end{pgfscope}%
\begin{pgfscope}%
\definecolor{textcolor}{rgb}{0.000000,0.000000,0.000000}%
\pgfsetstrokecolor{textcolor}%
\pgfsetfillcolor{textcolor}%
\pgftext[x=0.687197in,y=1.305467in,left,top]{\color{textcolor}{\rmfamily\fontsize{8.000000}{9.600000}\selectfont\catcode`\^=\active\def^{\ifmmode\sp\else\^{}\fi}\catcode`\%=\active\def%{\%}(c)}}%
\end{pgfscope}%
\begin{pgfscope}%
\pgfsetbuttcap%
\pgfsetmiterjoin%
\definecolor{currentfill}{rgb}{1.000000,1.000000,1.000000}%
\pgfsetfillcolor{currentfill}%
\pgfsetlinewidth{0.000000pt}%
\definecolor{currentstroke}{rgb}{0.000000,0.000000,0.000000}%
\pgfsetstrokecolor{currentstroke}%
\pgfsetstrokeopacity{0.000000}%
\pgfsetdash{}{0pt}%
\pgfpathmoveto{\pgfqpoint{2.256930in}{2.718071in}}%
\pgfpathlineto{\pgfqpoint{3.236130in}{2.718071in}}%
\pgfpathlineto{\pgfqpoint{3.236130in}{3.044321in}}%
\pgfpathlineto{\pgfqpoint{2.256930in}{3.044321in}}%
\pgfpathlineto{\pgfqpoint{2.256930in}{2.718071in}}%
\pgfpathclose%
\pgfusepath{fill}%
\end{pgfscope}%
\begin{pgfscope}%
\pgfpathrectangle{\pgfqpoint{2.256930in}{2.718071in}}{\pgfqpoint{0.979200in}{0.326250in}}%
\pgfusepath{clip}%
\pgfsetbuttcap%
\pgfsetroundjoin%
\pgfsetlinewidth{0.401500pt}%
\definecolor{currentstroke}{rgb}{0.690196,0.690196,0.690196}%
\pgfsetstrokecolor{currentstroke}%
\pgfsetstrokeopacity{0.250000}%
\pgfsetdash{{1.480000pt}{0.640000pt}}{0.000000pt}%
\pgfpathmoveto{\pgfqpoint{2.256930in}{2.718071in}}%
\pgfpathlineto{\pgfqpoint{2.256930in}{3.044321in}}%
\pgfusepath{stroke}%
\end{pgfscope}%
\begin{pgfscope}%
\pgfsetbuttcap%
\pgfsetroundjoin%
\definecolor{currentfill}{rgb}{0.000000,0.000000,0.000000}%
\pgfsetfillcolor{currentfill}%
\pgfsetlinewidth{0.803000pt}%
\definecolor{currentstroke}{rgb}{0.000000,0.000000,0.000000}%
\pgfsetstrokecolor{currentstroke}%
\pgfsetdash{}{0pt}%
\pgfsys@defobject{currentmarker}{\pgfqpoint{0.000000in}{0.000000in}}{\pgfqpoint{0.000000in}{0.027778in}}{%
\pgfpathmoveto{\pgfqpoint{0.000000in}{0.000000in}}%
\pgfpathlineto{\pgfqpoint{0.000000in}{0.027778in}}%
\pgfusepath{stroke,fill}%
}%
\begin{pgfscope}%
\pgfsys@transformshift{2.256930in}{2.718071in}%
\pgfsys@useobject{currentmarker}{}%
\end{pgfscope}%
\end{pgfscope}%
\begin{pgfscope}%
\definecolor{textcolor}{rgb}{0.000000,0.000000,0.000000}%
\pgfsetstrokecolor{textcolor}%
\pgfsetfillcolor{textcolor}%
\pgftext[x=2.256930in,y=2.704182in,,top]{\color{textcolor}{\rmfamily\fontsize{5.000000}{6.000000}\selectfont\catcode`\^=\active\def^{\ifmmode\sp\else\^{}\fi}\catcode`\%=\active\def%{\%}7.000}}%
\end{pgfscope}%
\begin{pgfscope}%
\pgfpathrectangle{\pgfqpoint{2.256930in}{2.718071in}}{\pgfqpoint{0.979200in}{0.326250in}}%
\pgfusepath{clip}%
\pgfsetbuttcap%
\pgfsetroundjoin%
\pgfsetlinewidth{0.401500pt}%
\definecolor{currentstroke}{rgb}{0.690196,0.690196,0.690196}%
\pgfsetstrokecolor{currentstroke}%
\pgfsetstrokeopacity{0.250000}%
\pgfsetdash{{1.480000pt}{0.640000pt}}{0.000000pt}%
\pgfpathmoveto{\pgfqpoint{2.501730in}{2.718071in}}%
\pgfpathlineto{\pgfqpoint{2.501730in}{3.044321in}}%
\pgfusepath{stroke}%
\end{pgfscope}%
\begin{pgfscope}%
\pgfsetbuttcap%
\pgfsetroundjoin%
\definecolor{currentfill}{rgb}{0.000000,0.000000,0.000000}%
\pgfsetfillcolor{currentfill}%
\pgfsetlinewidth{0.803000pt}%
\definecolor{currentstroke}{rgb}{0.000000,0.000000,0.000000}%
\pgfsetstrokecolor{currentstroke}%
\pgfsetdash{}{0pt}%
\pgfsys@defobject{currentmarker}{\pgfqpoint{0.000000in}{0.000000in}}{\pgfqpoint{0.000000in}{0.027778in}}{%
\pgfpathmoveto{\pgfqpoint{0.000000in}{0.000000in}}%
\pgfpathlineto{\pgfqpoint{0.000000in}{0.027778in}}%
\pgfusepath{stroke,fill}%
}%
\begin{pgfscope}%
\pgfsys@transformshift{2.501730in}{2.718071in}%
\pgfsys@useobject{currentmarker}{}%
\end{pgfscope}%
\end{pgfscope}%
\begin{pgfscope}%
\definecolor{textcolor}{rgb}{0.000000,0.000000,0.000000}%
\pgfsetstrokecolor{textcolor}%
\pgfsetfillcolor{textcolor}%
\pgftext[x=2.501730in,y=2.704182in,,top]{\color{textcolor}{\rmfamily\fontsize{5.000000}{6.000000}\selectfont\catcode`\^=\active\def^{\ifmmode\sp\else\^{}\fi}\catcode`\%=\active\def%{\%}7.125}}%
\end{pgfscope}%
\begin{pgfscope}%
\pgfpathrectangle{\pgfqpoint{2.256930in}{2.718071in}}{\pgfqpoint{0.979200in}{0.326250in}}%
\pgfusepath{clip}%
\pgfsetbuttcap%
\pgfsetroundjoin%
\pgfsetlinewidth{0.401500pt}%
\definecolor{currentstroke}{rgb}{0.690196,0.690196,0.690196}%
\pgfsetstrokecolor{currentstroke}%
\pgfsetstrokeopacity{0.250000}%
\pgfsetdash{{1.480000pt}{0.640000pt}}{0.000000pt}%
\pgfpathmoveto{\pgfqpoint{2.746530in}{2.718071in}}%
\pgfpathlineto{\pgfqpoint{2.746530in}{3.044321in}}%
\pgfusepath{stroke}%
\end{pgfscope}%
\begin{pgfscope}%
\pgfsetbuttcap%
\pgfsetroundjoin%
\definecolor{currentfill}{rgb}{0.000000,0.000000,0.000000}%
\pgfsetfillcolor{currentfill}%
\pgfsetlinewidth{0.803000pt}%
\definecolor{currentstroke}{rgb}{0.000000,0.000000,0.000000}%
\pgfsetstrokecolor{currentstroke}%
\pgfsetdash{}{0pt}%
\pgfsys@defobject{currentmarker}{\pgfqpoint{0.000000in}{0.000000in}}{\pgfqpoint{0.000000in}{0.027778in}}{%
\pgfpathmoveto{\pgfqpoint{0.000000in}{0.000000in}}%
\pgfpathlineto{\pgfqpoint{0.000000in}{0.027778in}}%
\pgfusepath{stroke,fill}%
}%
\begin{pgfscope}%
\pgfsys@transformshift{2.746530in}{2.718071in}%
\pgfsys@useobject{currentmarker}{}%
\end{pgfscope}%
\end{pgfscope}%
\begin{pgfscope}%
\definecolor{textcolor}{rgb}{0.000000,0.000000,0.000000}%
\pgfsetstrokecolor{textcolor}%
\pgfsetfillcolor{textcolor}%
\pgftext[x=2.746530in,y=2.704182in,,top]{\color{textcolor}{\rmfamily\fontsize{5.000000}{6.000000}\selectfont\catcode`\^=\active\def^{\ifmmode\sp\else\^{}\fi}\catcode`\%=\active\def%{\%}7.250}}%
\end{pgfscope}%
\begin{pgfscope}%
\pgfpathrectangle{\pgfqpoint{2.256930in}{2.718071in}}{\pgfqpoint{0.979200in}{0.326250in}}%
\pgfusepath{clip}%
\pgfsetbuttcap%
\pgfsetroundjoin%
\pgfsetlinewidth{0.401500pt}%
\definecolor{currentstroke}{rgb}{0.690196,0.690196,0.690196}%
\pgfsetstrokecolor{currentstroke}%
\pgfsetstrokeopacity{0.250000}%
\pgfsetdash{{1.480000pt}{0.640000pt}}{0.000000pt}%
\pgfpathmoveto{\pgfqpoint{2.991330in}{2.718071in}}%
\pgfpathlineto{\pgfqpoint{2.991330in}{3.044321in}}%
\pgfusepath{stroke}%
\end{pgfscope}%
\begin{pgfscope}%
\pgfsetbuttcap%
\pgfsetroundjoin%
\definecolor{currentfill}{rgb}{0.000000,0.000000,0.000000}%
\pgfsetfillcolor{currentfill}%
\pgfsetlinewidth{0.803000pt}%
\definecolor{currentstroke}{rgb}{0.000000,0.000000,0.000000}%
\pgfsetstrokecolor{currentstroke}%
\pgfsetdash{}{0pt}%
\pgfsys@defobject{currentmarker}{\pgfqpoint{0.000000in}{0.000000in}}{\pgfqpoint{0.000000in}{0.027778in}}{%
\pgfpathmoveto{\pgfqpoint{0.000000in}{0.000000in}}%
\pgfpathlineto{\pgfqpoint{0.000000in}{0.027778in}}%
\pgfusepath{stroke,fill}%
}%
\begin{pgfscope}%
\pgfsys@transformshift{2.991330in}{2.718071in}%
\pgfsys@useobject{currentmarker}{}%
\end{pgfscope}%
\end{pgfscope}%
\begin{pgfscope}%
\definecolor{textcolor}{rgb}{0.000000,0.000000,0.000000}%
\pgfsetstrokecolor{textcolor}%
\pgfsetfillcolor{textcolor}%
\pgftext[x=2.991330in,y=2.704182in,,top]{\color{textcolor}{\rmfamily\fontsize{5.000000}{6.000000}\selectfont\catcode`\^=\active\def^{\ifmmode\sp\else\^{}\fi}\catcode`\%=\active\def%{\%}7.375}}%
\end{pgfscope}%
\begin{pgfscope}%
\pgfpathrectangle{\pgfqpoint{2.256930in}{2.718071in}}{\pgfqpoint{0.979200in}{0.326250in}}%
\pgfusepath{clip}%
\pgfsetbuttcap%
\pgfsetroundjoin%
\pgfsetlinewidth{0.401500pt}%
\definecolor{currentstroke}{rgb}{0.690196,0.690196,0.690196}%
\pgfsetstrokecolor{currentstroke}%
\pgfsetstrokeopacity{0.250000}%
\pgfsetdash{{1.480000pt}{0.640000pt}}{0.000000pt}%
\pgfpathmoveto{\pgfqpoint{3.236130in}{2.718071in}}%
\pgfpathlineto{\pgfqpoint{3.236130in}{3.044321in}}%
\pgfusepath{stroke}%
\end{pgfscope}%
\begin{pgfscope}%
\pgfsetbuttcap%
\pgfsetroundjoin%
\definecolor{currentfill}{rgb}{0.000000,0.000000,0.000000}%
\pgfsetfillcolor{currentfill}%
\pgfsetlinewidth{0.803000pt}%
\definecolor{currentstroke}{rgb}{0.000000,0.000000,0.000000}%
\pgfsetstrokecolor{currentstroke}%
\pgfsetdash{}{0pt}%
\pgfsys@defobject{currentmarker}{\pgfqpoint{0.000000in}{0.000000in}}{\pgfqpoint{0.000000in}{0.027778in}}{%
\pgfpathmoveto{\pgfqpoint{0.000000in}{0.000000in}}%
\pgfpathlineto{\pgfqpoint{0.000000in}{0.027778in}}%
\pgfusepath{stroke,fill}%
}%
\begin{pgfscope}%
\pgfsys@transformshift{3.236130in}{2.718071in}%
\pgfsys@useobject{currentmarker}{}%
\end{pgfscope}%
\end{pgfscope}%
\begin{pgfscope}%
\definecolor{textcolor}{rgb}{0.000000,0.000000,0.000000}%
\pgfsetstrokecolor{textcolor}%
\pgfsetfillcolor{textcolor}%
\pgftext[x=3.236130in,y=2.704182in,,top]{\color{textcolor}{\rmfamily\fontsize{5.000000}{6.000000}\selectfont\catcode`\^=\active\def^{\ifmmode\sp\else\^{}\fi}\catcode`\%=\active\def%{\%}7.500}}%
\end{pgfscope}%
\begin{pgfscope}%
\pgfpathrectangle{\pgfqpoint{2.256930in}{2.718071in}}{\pgfqpoint{0.979200in}{0.326250in}}%
\pgfusepath{clip}%
\pgfsetbuttcap%
\pgfsetroundjoin%
\pgfsetlinewidth{0.401500pt}%
\definecolor{currentstroke}{rgb}{0.690196,0.690196,0.690196}%
\pgfsetstrokecolor{currentstroke}%
\pgfsetstrokeopacity{0.250000}%
\pgfsetdash{{1.480000pt}{0.640000pt}}{0.000000pt}%
\pgfpathmoveto{\pgfqpoint{2.256930in}{2.718071in}}%
\pgfpathlineto{\pgfqpoint{3.236130in}{2.718071in}}%
\pgfusepath{stroke}%
\end{pgfscope}%
\begin{pgfscope}%
\pgfsetbuttcap%
\pgfsetroundjoin%
\definecolor{currentfill}{rgb}{0.000000,0.000000,0.000000}%
\pgfsetfillcolor{currentfill}%
\pgfsetlinewidth{0.803000pt}%
\definecolor{currentstroke}{rgb}{0.000000,0.000000,0.000000}%
\pgfsetstrokecolor{currentstroke}%
\pgfsetdash{}{0pt}%
\pgfsys@defobject{currentmarker}{\pgfqpoint{0.000000in}{0.000000in}}{\pgfqpoint{0.027778in}{0.000000in}}{%
\pgfpathmoveto{\pgfqpoint{0.000000in}{0.000000in}}%
\pgfpathlineto{\pgfqpoint{0.027778in}{0.000000in}}%
\pgfusepath{stroke,fill}%
}%
\begin{pgfscope}%
\pgfsys@transformshift{2.256930in}{2.718071in}%
\pgfsys@useobject{currentmarker}{}%
\end{pgfscope}%
\end{pgfscope}%
\begin{pgfscope}%
\definecolor{textcolor}{rgb}{0.000000,0.000000,0.000000}%
\pgfsetstrokecolor{textcolor}%
\pgfsetfillcolor{textcolor}%
\pgftext[x=2.026024in, y=2.693959in, left, base]{\color{textcolor}{\rmfamily\fontsize{5.000000}{6.000000}\selectfont\catcode`\^=\active\def^{\ifmmode\sp\else\^{}\fi}\catcode`\%=\active\def%{\%}11.69}}%
\end{pgfscope}%
\begin{pgfscope}%
\pgfpathrectangle{\pgfqpoint{2.256930in}{2.718071in}}{\pgfqpoint{0.979200in}{0.326250in}}%
\pgfusepath{clip}%
\pgfsetbuttcap%
\pgfsetroundjoin%
\pgfsetlinewidth{0.401500pt}%
\definecolor{currentstroke}{rgb}{0.690196,0.690196,0.690196}%
\pgfsetstrokecolor{currentstroke}%
\pgfsetstrokeopacity{0.250000}%
\pgfsetdash{{1.480000pt}{0.640000pt}}{0.000000pt}%
\pgfpathmoveto{\pgfqpoint{2.256930in}{2.881196in}}%
\pgfpathlineto{\pgfqpoint{3.236130in}{2.881196in}}%
\pgfusepath{stroke}%
\end{pgfscope}%
\begin{pgfscope}%
\pgfsetbuttcap%
\pgfsetroundjoin%
\definecolor{currentfill}{rgb}{0.000000,0.000000,0.000000}%
\pgfsetfillcolor{currentfill}%
\pgfsetlinewidth{0.803000pt}%
\definecolor{currentstroke}{rgb}{0.000000,0.000000,0.000000}%
\pgfsetstrokecolor{currentstroke}%
\pgfsetdash{}{0pt}%
\pgfsys@defobject{currentmarker}{\pgfqpoint{0.000000in}{0.000000in}}{\pgfqpoint{0.027778in}{0.000000in}}{%
\pgfpathmoveto{\pgfqpoint{0.000000in}{0.000000in}}%
\pgfpathlineto{\pgfqpoint{0.027778in}{0.000000in}}%
\pgfusepath{stroke,fill}%
}%
\begin{pgfscope}%
\pgfsys@transformshift{2.256930in}{2.881196in}%
\pgfsys@useobject{currentmarker}{}%
\end{pgfscope}%
\end{pgfscope}%
\begin{pgfscope}%
\definecolor{textcolor}{rgb}{0.000000,0.000000,0.000000}%
\pgfsetstrokecolor{textcolor}%
\pgfsetfillcolor{textcolor}%
\pgftext[x=2.026024in, y=2.857084in, left, base]{\color{textcolor}{\rmfamily\fontsize{5.000000}{6.000000}\selectfont\catcode`\^=\active\def^{\ifmmode\sp\else\^{}\fi}\catcode`\%=\active\def%{\%}11.78}}%
\end{pgfscope}%
\begin{pgfscope}%
\pgfpathrectangle{\pgfqpoint{2.256930in}{2.718071in}}{\pgfqpoint{0.979200in}{0.326250in}}%
\pgfusepath{clip}%
\pgfsetbuttcap%
\pgfsetroundjoin%
\pgfsetlinewidth{0.401500pt}%
\definecolor{currentstroke}{rgb}{0.690196,0.690196,0.690196}%
\pgfsetstrokecolor{currentstroke}%
\pgfsetstrokeopacity{0.250000}%
\pgfsetdash{{1.480000pt}{0.640000pt}}{0.000000pt}%
\pgfpathmoveto{\pgfqpoint{2.256930in}{3.044321in}}%
\pgfpathlineto{\pgfqpoint{3.236130in}{3.044321in}}%
\pgfusepath{stroke}%
\end{pgfscope}%
\begin{pgfscope}%
\pgfsetbuttcap%
\pgfsetroundjoin%
\definecolor{currentfill}{rgb}{0.000000,0.000000,0.000000}%
\pgfsetfillcolor{currentfill}%
\pgfsetlinewidth{0.803000pt}%
\definecolor{currentstroke}{rgb}{0.000000,0.000000,0.000000}%
\pgfsetstrokecolor{currentstroke}%
\pgfsetdash{}{0pt}%
\pgfsys@defobject{currentmarker}{\pgfqpoint{0.000000in}{0.000000in}}{\pgfqpoint{0.027778in}{0.000000in}}{%
\pgfpathmoveto{\pgfqpoint{0.000000in}{0.000000in}}%
\pgfpathlineto{\pgfqpoint{0.027778in}{0.000000in}}%
\pgfusepath{stroke,fill}%
}%
\begin{pgfscope}%
\pgfsys@transformshift{2.256930in}{3.044321in}%
\pgfsys@useobject{currentmarker}{}%
\end{pgfscope}%
\end{pgfscope}%
\begin{pgfscope}%
\definecolor{textcolor}{rgb}{0.000000,0.000000,0.000000}%
\pgfsetstrokecolor{textcolor}%
\pgfsetfillcolor{textcolor}%
\pgftext[x=2.026024in, y=3.020209in, left, base]{\color{textcolor}{\rmfamily\fontsize{5.000000}{6.000000}\selectfont\catcode`\^=\active\def^{\ifmmode\sp\else\^{}\fi}\catcode`\%=\active\def%{\%}11.87}}%
\end{pgfscope}%
\begin{pgfscope}%
\pgfpathrectangle{\pgfqpoint{2.256930in}{2.718071in}}{\pgfqpoint{0.979200in}{0.326250in}}%
\pgfusepath{clip}%
\pgfsetrectcap%
\pgfsetroundjoin%
\pgfsetlinewidth{1.505625pt}%
\definecolor{currentstroke}{rgb}{0.000000,0.447059,0.698039}%
\pgfsetstrokecolor{currentstroke}%
\pgfsetdash{}{0pt}%
\pgfpathmoveto{\pgfqpoint{2.256642in}{2.841882in}}%
\pgfpathlineto{\pgfqpoint{2.262416in}{2.797110in}}%
\pgfpathlineto{\pgfqpoint{2.267312in}{2.770780in}}%
\pgfpathlineto{\pgfqpoint{2.271228in}{2.757432in}}%
\pgfpathlineto{\pgfqpoint{2.274166in}{2.751885in}}%
\pgfpathlineto{\pgfqpoint{2.276124in}{2.750327in}}%
\pgfpathlineto{\pgfqpoint{2.278083in}{2.750457in}}%
\pgfpathlineto{\pgfqpoint{2.280041in}{2.752291in}}%
\pgfpathlineto{\pgfqpoint{2.282979in}{2.758223in}}%
\pgfpathlineto{\pgfqpoint{2.286498in}{2.770335in}}%
\pgfpathlineto{\pgfqpoint{2.289925in}{2.787362in}}%
\pgfpathlineto{\pgfqpoint{2.294821in}{2.820618in}}%
\pgfpathlineto{\pgfqpoint{2.305271in}{2.899824in}}%
\pgfpathlineto{\pgfqpoint{2.311147in}{2.930493in}}%
\pgfpathlineto{\pgfqpoint{2.316043in}{2.948008in}}%
\pgfpathlineto{\pgfqpoint{2.319959in}{2.956767in}}%
\pgfpathlineto{\pgfqpoint{2.322897in}{2.960277in}}%
\pgfpathlineto{\pgfqpoint{2.325835in}{2.961170in}}%
\pgfpathlineto{\pgfqpoint{2.327793in}{2.960312in}}%
\pgfpathlineto{\pgfqpoint{2.330731in}{2.956851in}}%
\pgfpathlineto{\pgfqpoint{2.333668in}{2.950784in}}%
\pgfpathlineto{\pgfqpoint{2.337585in}{2.938650in}}%
\pgfpathlineto{\pgfqpoint{2.342481in}{2.917003in}}%
\pgfpathlineto{\pgfqpoint{2.347377in}{2.888182in}}%
\pgfpathlineto{\pgfqpoint{2.353740in}{2.840222in}}%
\pgfpathlineto{\pgfqpoint{2.360778in}{2.786819in}}%
\pgfpathlineto{\pgfqpoint{2.365674in}{2.762364in}}%
\pgfpathlineto{\pgfqpoint{2.369590in}{2.750496in}}%
\pgfpathlineto{\pgfqpoint{2.372528in}{2.746083in}}%
\pgfpathlineto{\pgfqpoint{2.374486in}{2.745259in}}%
\pgfpathlineto{\pgfqpoint{2.376445in}{2.746144in}}%
\pgfpathlineto{\pgfqpoint{2.378403in}{2.748719in}}%
\pgfpathlineto{\pgfqpoint{2.381341in}{2.755769in}}%
\pgfpathlineto{\pgfqpoint{2.385258in}{2.771072in}}%
\pgfpathlineto{\pgfqpoint{2.389174in}{2.793115in}}%
\pgfpathlineto{\pgfqpoint{2.394482in}{2.833706in}}%
\pgfpathlineto{\pgfqpoint{2.400718in}{2.883422in}}%
\pgfpathlineto{\pgfqpoint{2.406593in}{2.919400in}}%
\pgfpathlineto{\pgfqpoint{2.411489in}{2.941330in}}%
\pgfpathlineto{\pgfqpoint{2.415406in}{2.953614in}}%
\pgfpathlineto{\pgfqpoint{2.419322in}{2.961230in}}%
\pgfpathlineto{\pgfqpoint{2.422260in}{2.963884in}}%
\pgfpathlineto{\pgfqpoint{2.424218in}{2.964199in}}%
\pgfpathlineto{\pgfqpoint{2.426177in}{2.963353in}}%
\pgfpathlineto{\pgfqpoint{2.429114in}{2.959907in}}%
\pgfpathlineto{\pgfqpoint{2.432052in}{2.953856in}}%
\pgfpathlineto{\pgfqpoint{2.435969in}{2.941743in}}%
\pgfpathlineto{\pgfqpoint{2.440865in}{2.920119in}}%
\pgfpathlineto{\pgfqpoint{2.445761in}{2.891321in}}%
\pgfpathlineto{\pgfqpoint{2.452148in}{2.843024in}}%
\pgfpathlineto{\pgfqpoint{2.458459in}{2.795041in}}%
\pgfpathlineto{\pgfqpoint{2.463355in}{2.769956in}}%
\pgfpathlineto{\pgfqpoint{2.467272in}{2.757573in}}%
\pgfpathlineto{\pgfqpoint{2.470210in}{2.752769in}}%
\pgfpathlineto{\pgfqpoint{2.472168in}{2.751699in}}%
\pgfpathlineto{\pgfqpoint{2.474126in}{2.752319in}}%
\pgfpathlineto{\pgfqpoint{2.476613in}{2.755562in}}%
\pgfpathlineto{\pgfqpoint{2.479428in}{2.762526in}}%
\pgfpathlineto{\pgfqpoint{2.482855in}{2.775713in}}%
\pgfpathlineto{\pgfqpoint{2.486772in}{2.797106in}}%
\pgfpathlineto{\pgfqpoint{2.492118in}{2.837148in}}%
\pgfpathlineto{\pgfqpoint{2.500129in}{2.899509in}}%
\pgfpathlineto{\pgfqpoint{2.506004in}{2.932850in}}%
\pgfpathlineto{\pgfqpoint{2.510900in}{2.952586in}}%
\pgfpathlineto{\pgfqpoint{2.514817in}{2.963116in}}%
\pgfpathlineto{\pgfqpoint{2.517754in}{2.967951in}}%
\pgfpathlineto{\pgfqpoint{2.520692in}{2.970164in}}%
\pgfpathlineto{\pgfqpoint{2.522650in}{2.970186in}}%
\pgfpathlineto{\pgfqpoint{2.524609in}{2.969046in}}%
\pgfpathlineto{\pgfqpoint{2.527546in}{2.965161in}}%
\pgfpathlineto{\pgfqpoint{2.530484in}{2.958669in}}%
\pgfpathlineto{\pgfqpoint{2.534401in}{2.945970in}}%
\pgfpathlineto{\pgfqpoint{2.539297in}{2.923614in}}%
\pgfpathlineto{\pgfqpoint{2.544193in}{2.894085in}}%
\pgfpathlineto{\pgfqpoint{2.550890in}{2.842470in}}%
\pgfpathlineto{\pgfqpoint{2.555252in}{2.809540in}}%
\pgfpathlineto{\pgfqpoint{2.560148in}{2.782732in}}%
\pgfpathlineto{\pgfqpoint{2.564065in}{2.768996in}}%
\pgfpathlineto{\pgfqpoint{2.567003in}{2.763155in}}%
\pgfpathlineto{\pgfqpoint{2.568961in}{2.761405in}}%
\pgfpathlineto{\pgfqpoint{2.570920in}{2.761345in}}%
\pgfpathlineto{\pgfqpoint{2.572878in}{2.762983in}}%
\pgfpathlineto{\pgfqpoint{2.575816in}{2.768617in}}%
\pgfpathlineto{\pgfqpoint{2.578753in}{2.778059in}}%
\pgfpathlineto{\pgfqpoint{2.582670in}{2.796563in}}%
\pgfpathlineto{\pgfqpoint{2.587566in}{2.829140in}}%
\pgfpathlineto{\pgfqpoint{2.601030in}{2.931199in}}%
\pgfpathlineto{\pgfqpoint{2.605926in}{2.955699in}}%
\pgfpathlineto{\pgfqpoint{2.610822in}{2.972881in}}%
\pgfpathlineto{\pgfqpoint{2.614738in}{2.981368in}}%
\pgfpathlineto{\pgfqpoint{2.617676in}{2.984671in}}%
\pgfpathlineto{\pgfqpoint{2.620614in}{2.985353in}}%
\pgfpathlineto{\pgfqpoint{2.622572in}{2.984354in}}%
\pgfpathlineto{\pgfqpoint{2.625510in}{2.980679in}}%
\pgfpathlineto{\pgfqpoint{2.628447in}{2.974396in}}%
\pgfpathlineto{\pgfqpoint{2.632364in}{2.961970in}}%
\pgfpathlineto{\pgfqpoint{2.637260in}{2.939951in}}%
\pgfpathlineto{\pgfqpoint{2.642156in}{2.910752in}}%
\pgfpathlineto{\pgfqpoint{2.648673in}{2.860850in}}%
\pgfpathlineto{\pgfqpoint{2.653401in}{2.825353in}}%
\pgfpathlineto{\pgfqpoint{2.658297in}{2.799136in}}%
\pgfpathlineto{\pgfqpoint{2.662214in}{2.785857in}}%
\pgfpathlineto{\pgfqpoint{2.665152in}{2.780370in}}%
\pgfpathlineto{\pgfqpoint{2.667110in}{2.778840in}}%
\pgfpathlineto{\pgfqpoint{2.669068in}{2.779016in}}%
\pgfpathlineto{\pgfqpoint{2.671027in}{2.780880in}}%
\pgfpathlineto{\pgfqpoint{2.673964in}{2.786854in}}%
\pgfpathlineto{\pgfqpoint{2.677483in}{2.799025in}}%
\pgfpathlineto{\pgfqpoint{2.680911in}{2.816107in}}%
\pgfpathlineto{\pgfqpoint{2.685684in}{2.848453in}}%
\pgfpathlineto{\pgfqpoint{2.698356in}{2.945816in}}%
\pgfpathlineto{\pgfqpoint{2.704231in}{2.975549in}}%
\pgfpathlineto{\pgfqpoint{2.709127in}{2.992271in}}%
\pgfpathlineto{\pgfqpoint{2.713044in}{3.000386in}}%
\pgfpathlineto{\pgfqpoint{2.715982in}{3.003407in}}%
\pgfpathlineto{\pgfqpoint{2.717940in}{3.003965in}}%
\pgfpathlineto{\pgfqpoint{2.719899in}{3.003358in}}%
\pgfpathlineto{\pgfqpoint{2.722836in}{3.000268in}}%
\pgfpathlineto{\pgfqpoint{2.725774in}{2.994566in}}%
\pgfpathlineto{\pgfqpoint{2.729691in}{2.982911in}}%
\pgfpathlineto{\pgfqpoint{2.734587in}{2.961846in}}%
\pgfpathlineto{\pgfqpoint{2.739483in}{2.933590in}}%
\pgfpathlineto{\pgfqpoint{2.745789in}{2.886651in}}%
\pgfpathlineto{\pgfqpoint{2.752056in}{2.839259in}}%
\pgfpathlineto{\pgfqpoint{2.756952in}{2.814245in}}%
\pgfpathlineto{\pgfqpoint{2.760869in}{2.801913in}}%
\pgfpathlineto{\pgfqpoint{2.763807in}{2.797138in}}%
\pgfpathlineto{\pgfqpoint{2.765765in}{2.796078in}}%
\pgfpathlineto{\pgfqpoint{2.767724in}{2.796714in}}%
\pgfpathlineto{\pgfqpoint{2.769682in}{2.799053in}}%
\pgfpathlineto{\pgfqpoint{2.772620in}{2.805733in}}%
\pgfpathlineto{\pgfqpoint{2.775557in}{2.816199in}}%
\pgfpathlineto{\pgfqpoint{2.779474in}{2.836059in}}%
\pgfpathlineto{\pgfqpoint{2.784370in}{2.870330in}}%
\pgfpathlineto{\pgfqpoint{2.795280in}{2.954299in}}%
\pgfpathlineto{\pgfqpoint{2.801155in}{2.985164in}}%
\pgfpathlineto{\pgfqpoint{2.806051in}{3.002822in}}%
\pgfpathlineto{\pgfqpoint{2.809968in}{3.011682in}}%
\pgfpathlineto{\pgfqpoint{2.812905in}{3.015260in}}%
\pgfpathlineto{\pgfqpoint{2.815843in}{3.016213in}}%
\pgfpathlineto{\pgfqpoint{2.817801in}{3.015392in}}%
\pgfpathlineto{\pgfqpoint{2.820739in}{3.011980in}}%
\pgfpathlineto{\pgfqpoint{2.823677in}{3.005956in}}%
\pgfpathlineto{\pgfqpoint{2.827593in}{2.993870in}}%
\pgfpathlineto{\pgfqpoint{2.832489in}{2.972264in}}%
\pgfpathlineto{\pgfqpoint{2.837385in}{2.943466in}}%
\pgfpathlineto{\pgfqpoint{2.843818in}{2.894755in}}%
\pgfpathlineto{\pgfqpoint{2.849850in}{2.848302in}}%
\pgfpathlineto{\pgfqpoint{2.854746in}{2.822388in}}%
\pgfpathlineto{\pgfqpoint{2.858663in}{2.809351in}}%
\pgfpathlineto{\pgfqpoint{2.861601in}{2.804039in}}%
\pgfpathlineto{\pgfqpoint{2.863559in}{2.802630in}}%
\pgfpathlineto{\pgfqpoint{2.865885in}{2.803154in}}%
\pgfpathlineto{\pgfqpoint{2.868333in}{2.806302in}}%
\pgfpathlineto{\pgfqpoint{2.871270in}{2.813552in}}%
\pgfpathlineto{\pgfqpoint{2.875187in}{2.829125in}}%
\pgfpathlineto{\pgfqpoint{2.879104in}{2.851404in}}%
\pgfpathlineto{\pgfqpoint{2.884506in}{2.893131in}}%
\pgfpathlineto{\pgfqpoint{2.890279in}{2.938319in}}%
\pgfpathlineto{\pgfqpoint{2.896154in}{2.973822in}}%
\pgfpathlineto{\pgfqpoint{2.901050in}{2.995339in}}%
\pgfpathlineto{\pgfqpoint{2.904967in}{3.007280in}}%
\pgfpathlineto{\pgfqpoint{2.908884in}{3.014542in}}%
\pgfpathlineto{\pgfqpoint{2.911821in}{3.016923in}}%
\pgfpathlineto{\pgfqpoint{2.913780in}{3.017053in}}%
\pgfpathlineto{\pgfqpoint{2.915738in}{3.016019in}}%
\pgfpathlineto{\pgfqpoint{2.918676in}{3.012287in}}%
\pgfpathlineto{\pgfqpoint{2.921613in}{3.005944in}}%
\pgfpathlineto{\pgfqpoint{2.925530in}{2.993432in}}%
\pgfpathlineto{\pgfqpoint{2.930426in}{2.971296in}}%
\pgfpathlineto{\pgfqpoint{2.935322in}{2.941970in}}%
\pgfpathlineto{\pgfqpoint{2.941626in}{2.893670in}}%
\pgfpathlineto{\pgfqpoint{2.947946in}{2.844570in}}%
\pgfpathlineto{\pgfqpoint{2.952842in}{2.818591in}}%
\pgfpathlineto{\pgfqpoint{2.956759in}{2.805494in}}%
\pgfpathlineto{\pgfqpoint{2.959697in}{2.800147in}}%
\pgfpathlineto{\pgfqpoint{2.961655in}{2.798707in}}%
\pgfpathlineto{\pgfqpoint{2.963614in}{2.798973in}}%
\pgfpathlineto{\pgfqpoint{2.965572in}{2.800923in}}%
\pgfpathlineto{\pgfqpoint{2.968510in}{2.807026in}}%
\pgfpathlineto{\pgfqpoint{2.971447in}{2.816933in}}%
\pgfpathlineto{\pgfqpoint{2.975364in}{2.836028in}}%
\pgfpathlineto{\pgfqpoint{2.980138in}{2.868384in}}%
\pgfpathlineto{\pgfqpoint{2.990863in}{2.949015in}}%
\pgfpathlineto{\pgfqpoint{2.996738in}{2.979063in}}%
\pgfpathlineto{\pgfqpoint{3.001634in}{2.996046in}}%
\pgfpathlineto{\pgfqpoint{3.005551in}{3.004368in}}%
\pgfpathlineto{\pgfqpoint{3.008488in}{3.007544in}}%
\pgfpathlineto{\pgfqpoint{3.010447in}{3.008204in}}%
\pgfpathlineto{\pgfqpoint{3.012405in}{3.007699in}}%
\pgfpathlineto{\pgfqpoint{3.015343in}{3.004762in}}%
\pgfpathlineto{\pgfqpoint{3.018280in}{2.999212in}}%
\pgfpathlineto{\pgfqpoint{3.022197in}{2.987758in}}%
\pgfpathlineto{\pgfqpoint{3.027093in}{2.966942in}}%
\pgfpathlineto{\pgfqpoint{3.031989in}{2.938934in}}%
\pgfpathlineto{\pgfqpoint{3.037858in}{2.895938in}}%
\pgfpathlineto{\pgfqpoint{3.046816in}{2.826236in}}%
\pgfpathlineto{\pgfqpoint{3.051712in}{2.802057in}}%
\pgfpathlineto{\pgfqpoint{3.055628in}{2.790399in}}%
\pgfpathlineto{\pgfqpoint{3.058566in}{2.786125in}}%
\pgfpathlineto{\pgfqpoint{3.060524in}{2.785393in}}%
\pgfpathlineto{\pgfqpoint{3.062483in}{2.786379in}}%
\pgfpathlineto{\pgfqpoint{3.064441in}{2.789049in}}%
\pgfpathlineto{\pgfqpoint{3.067379in}{2.796213in}}%
\pgfpathlineto{\pgfqpoint{3.070997in}{2.810281in}}%
\pgfpathlineto{\pgfqpoint{3.074914in}{2.831964in}}%
\pgfpathlineto{\pgfqpoint{3.080202in}{2.871884in}}%
\pgfpathlineto{\pgfqpoint{3.086139in}{2.917654in}}%
\pgfpathlineto{\pgfqpoint{3.092014in}{2.952324in}}%
\pgfpathlineto{\pgfqpoint{3.096910in}{2.973157in}}%
\pgfpathlineto{\pgfqpoint{3.100827in}{2.984557in}}%
\pgfpathlineto{\pgfqpoint{3.104744in}{2.991284in}}%
\pgfpathlineto{\pgfqpoint{3.107681in}{2.993269in}}%
\pgfpathlineto{\pgfqpoint{3.109640in}{2.993136in}}%
\pgfpathlineto{\pgfqpoint{3.111598in}{2.991841in}}%
\pgfpathlineto{\pgfqpoint{3.114536in}{2.987721in}}%
\pgfpathlineto{\pgfqpoint{3.117473in}{2.980992in}}%
\pgfpathlineto{\pgfqpoint{3.121390in}{2.967973in}}%
\pgfpathlineto{\pgfqpoint{3.126286in}{2.945212in}}%
\pgfpathlineto{\pgfqpoint{3.132161in}{2.908425in}}%
\pgfpathlineto{\pgfqpoint{3.138910in}{2.853629in}}%
\pgfpathlineto{\pgfqpoint{3.143918in}{2.815698in}}%
\pgfpathlineto{\pgfqpoint{3.148814in}{2.789468in}}%
\pgfpathlineto{\pgfqpoint{3.152730in}{2.776176in}}%
\pgfpathlineto{\pgfqpoint{3.155668in}{2.770688in}}%
\pgfpathlineto{\pgfqpoint{3.157626in}{2.769147in}}%
\pgfpathlineto{\pgfqpoint{3.159585in}{2.769315in}}%
\pgfpathlineto{\pgfqpoint{3.161543in}{2.771183in}}%
\pgfpathlineto{\pgfqpoint{3.164603in}{2.777488in}}%
\pgfpathlineto{\pgfqpoint{3.167296in}{2.786442in}}%
\pgfpathlineto{\pgfqpoint{3.171213in}{2.805175in}}%
\pgfpathlineto{\pgfqpoint{3.176109in}{2.838019in}}%
\pgfpathlineto{\pgfqpoint{3.187215in}{2.919936in}}%
\pgfpathlineto{\pgfqpoint{3.193090in}{2.948735in}}%
\pgfpathlineto{\pgfqpoint{3.197986in}{2.964692in}}%
\pgfpathlineto{\pgfqpoint{3.201903in}{2.972202in}}%
\pgfpathlineto{\pgfqpoint{3.204841in}{2.974775in}}%
\pgfpathlineto{\pgfqpoint{3.206799in}{2.975036in}}%
\pgfpathlineto{\pgfqpoint{3.208757in}{2.974134in}}%
\pgfpathlineto{\pgfqpoint{3.211695in}{2.970605in}}%
\pgfpathlineto{\pgfqpoint{3.214633in}{2.964470in}}%
\pgfpathlineto{\pgfqpoint{3.218549in}{2.952242in}}%
\pgfpathlineto{\pgfqpoint{3.223445in}{2.930472in}}%
\pgfpathlineto{\pgfqpoint{3.228341in}{2.901525in}}%
\pgfpathlineto{\pgfqpoint{3.234816in}{2.852441in}}%
\pgfpathlineto{\pgfqpoint{3.236130in}{2.841172in}}%
\pgfpathlineto{\pgfqpoint{3.236130in}{2.841172in}}%
\pgfusepath{stroke}%
\end{pgfscope}%
\begin{pgfscope}%
\pgfsetrectcap%
\pgfsetmiterjoin%
\pgfsetlinewidth{0.803000pt}%
\definecolor{currentstroke}{rgb}{0.000000,0.000000,0.000000}%
\pgfsetstrokecolor{currentstroke}%
\pgfsetdash{}{0pt}%
\pgfpathmoveto{\pgfqpoint{2.256930in}{2.718071in}}%
\pgfpathlineto{\pgfqpoint{2.256930in}{3.044321in}}%
\pgfusepath{stroke}%
\end{pgfscope}%
\begin{pgfscope}%
\pgfsetrectcap%
\pgfsetmiterjoin%
\pgfsetlinewidth{0.803000pt}%
\definecolor{currentstroke}{rgb}{0.000000,0.000000,0.000000}%
\pgfsetstrokecolor{currentstroke}%
\pgfsetdash{}{0pt}%
\pgfpathmoveto{\pgfqpoint{3.236130in}{2.718071in}}%
\pgfpathlineto{\pgfqpoint{3.236130in}{3.044321in}}%
\pgfusepath{stroke}%
\end{pgfscope}%
\begin{pgfscope}%
\pgfsetrectcap%
\pgfsetmiterjoin%
\pgfsetlinewidth{0.803000pt}%
\definecolor{currentstroke}{rgb}{0.000000,0.000000,0.000000}%
\pgfsetstrokecolor{currentstroke}%
\pgfsetdash{}{0pt}%
\pgfpathmoveto{\pgfqpoint{2.256930in}{2.718071in}}%
\pgfpathlineto{\pgfqpoint{3.236130in}{2.718071in}}%
\pgfusepath{stroke}%
\end{pgfscope}%
\begin{pgfscope}%
\pgfsetrectcap%
\pgfsetmiterjoin%
\pgfsetlinewidth{0.803000pt}%
\definecolor{currentstroke}{rgb}{0.000000,0.000000,0.000000}%
\pgfsetstrokecolor{currentstroke}%
\pgfsetdash{}{0pt}%
\pgfpathmoveto{\pgfqpoint{2.256930in}{3.044321in}}%
\pgfpathlineto{\pgfqpoint{3.236130in}{3.044321in}}%
\pgfusepath{stroke}%
\end{pgfscope}%
\begin{pgfscope}%
\pgfsetbuttcap%
\pgfsetmiterjoin%
\definecolor{currentfill}{rgb}{1.000000,1.000000,1.000000}%
\pgfsetfillcolor{currentfill}%
\pgfsetlinewidth{0.000000pt}%
\definecolor{currentstroke}{rgb}{0.000000,0.000000,0.000000}%
\pgfsetstrokecolor{currentstroke}%
\pgfsetstrokeopacity{0.000000}%
\pgfsetdash{}{0pt}%
\pgfpathmoveto{\pgfqpoint{2.256930in}{1.995363in}}%
\pgfpathlineto{\pgfqpoint{3.236130in}{1.995363in}}%
\pgfpathlineto{\pgfqpoint{3.236130in}{2.321613in}}%
\pgfpathlineto{\pgfqpoint{2.256930in}{2.321613in}}%
\pgfpathlineto{\pgfqpoint{2.256930in}{1.995363in}}%
\pgfpathclose%
\pgfusepath{fill}%
\end{pgfscope}%
\begin{pgfscope}%
\pgfpathrectangle{\pgfqpoint{2.256930in}{1.995363in}}{\pgfqpoint{0.979200in}{0.326250in}}%
\pgfusepath{clip}%
\pgfsetrectcap%
\pgfsetroundjoin%
\pgfsetlinewidth{0.702625pt}%
\definecolor{currentstroke}{rgb}{0.600000,0.600000,0.600000}%
\pgfsetstrokecolor{currentstroke}%
\pgfsetdash{}{0pt}%
\pgfpathmoveto{\pgfqpoint{2.256930in}{2.025022in}}%
\pgfpathlineto{\pgfqpoint{3.236130in}{2.025022in}}%
\pgfusepath{stroke}%
\end{pgfscope}%
\begin{pgfscope}%
\pgfpathrectangle{\pgfqpoint{2.256930in}{1.995363in}}{\pgfqpoint{0.979200in}{0.326250in}}%
\pgfusepath{clip}%
\pgfsetbuttcap%
\pgfsetroundjoin%
\pgfsetlinewidth{0.401500pt}%
\definecolor{currentstroke}{rgb}{0.690196,0.690196,0.690196}%
\pgfsetstrokecolor{currentstroke}%
\pgfsetstrokeopacity{0.250000}%
\pgfsetdash{{1.480000pt}{0.640000pt}}{0.000000pt}%
\pgfpathmoveto{\pgfqpoint{2.256930in}{1.995363in}}%
\pgfpathlineto{\pgfqpoint{2.256930in}{2.321613in}}%
\pgfusepath{stroke}%
\end{pgfscope}%
\begin{pgfscope}%
\pgfsetbuttcap%
\pgfsetroundjoin%
\definecolor{currentfill}{rgb}{0.000000,0.000000,0.000000}%
\pgfsetfillcolor{currentfill}%
\pgfsetlinewidth{0.803000pt}%
\definecolor{currentstroke}{rgb}{0.000000,0.000000,0.000000}%
\pgfsetstrokecolor{currentstroke}%
\pgfsetdash{}{0pt}%
\pgfsys@defobject{currentmarker}{\pgfqpoint{0.000000in}{0.000000in}}{\pgfqpoint{0.000000in}{0.027778in}}{%
\pgfpathmoveto{\pgfqpoint{0.000000in}{0.000000in}}%
\pgfpathlineto{\pgfqpoint{0.000000in}{0.027778in}}%
\pgfusepath{stroke,fill}%
}%
\begin{pgfscope}%
\pgfsys@transformshift{2.256930in}{1.995363in}%
\pgfsys@useobject{currentmarker}{}%
\end{pgfscope}%
\end{pgfscope}%
\begin{pgfscope}%
\definecolor{textcolor}{rgb}{0.000000,0.000000,0.000000}%
\pgfsetstrokecolor{textcolor}%
\pgfsetfillcolor{textcolor}%
\pgftext[x=2.256930in,y=1.981474in,,top]{\color{textcolor}{\rmfamily\fontsize{5.000000}{6.000000}\selectfont\catcode`\^=\active\def^{\ifmmode\sp\else\^{}\fi}\catcode`\%=\active\def%{\%}7.000}}%
\end{pgfscope}%
\begin{pgfscope}%
\pgfpathrectangle{\pgfqpoint{2.256930in}{1.995363in}}{\pgfqpoint{0.979200in}{0.326250in}}%
\pgfusepath{clip}%
\pgfsetbuttcap%
\pgfsetroundjoin%
\pgfsetlinewidth{0.401500pt}%
\definecolor{currentstroke}{rgb}{0.690196,0.690196,0.690196}%
\pgfsetstrokecolor{currentstroke}%
\pgfsetstrokeopacity{0.250000}%
\pgfsetdash{{1.480000pt}{0.640000pt}}{0.000000pt}%
\pgfpathmoveto{\pgfqpoint{2.501730in}{1.995363in}}%
\pgfpathlineto{\pgfqpoint{2.501730in}{2.321613in}}%
\pgfusepath{stroke}%
\end{pgfscope}%
\begin{pgfscope}%
\pgfsetbuttcap%
\pgfsetroundjoin%
\definecolor{currentfill}{rgb}{0.000000,0.000000,0.000000}%
\pgfsetfillcolor{currentfill}%
\pgfsetlinewidth{0.803000pt}%
\definecolor{currentstroke}{rgb}{0.000000,0.000000,0.000000}%
\pgfsetstrokecolor{currentstroke}%
\pgfsetdash{}{0pt}%
\pgfsys@defobject{currentmarker}{\pgfqpoint{0.000000in}{0.000000in}}{\pgfqpoint{0.000000in}{0.027778in}}{%
\pgfpathmoveto{\pgfqpoint{0.000000in}{0.000000in}}%
\pgfpathlineto{\pgfqpoint{0.000000in}{0.027778in}}%
\pgfusepath{stroke,fill}%
}%
\begin{pgfscope}%
\pgfsys@transformshift{2.501730in}{1.995363in}%
\pgfsys@useobject{currentmarker}{}%
\end{pgfscope}%
\end{pgfscope}%
\begin{pgfscope}%
\definecolor{textcolor}{rgb}{0.000000,0.000000,0.000000}%
\pgfsetstrokecolor{textcolor}%
\pgfsetfillcolor{textcolor}%
\pgftext[x=2.501730in,y=1.981474in,,top]{\color{textcolor}{\rmfamily\fontsize{5.000000}{6.000000}\selectfont\catcode`\^=\active\def^{\ifmmode\sp\else\^{}\fi}\catcode`\%=\active\def%{\%}7.125}}%
\end{pgfscope}%
\begin{pgfscope}%
\pgfpathrectangle{\pgfqpoint{2.256930in}{1.995363in}}{\pgfqpoint{0.979200in}{0.326250in}}%
\pgfusepath{clip}%
\pgfsetbuttcap%
\pgfsetroundjoin%
\pgfsetlinewidth{0.401500pt}%
\definecolor{currentstroke}{rgb}{0.690196,0.690196,0.690196}%
\pgfsetstrokecolor{currentstroke}%
\pgfsetstrokeopacity{0.250000}%
\pgfsetdash{{1.480000pt}{0.640000pt}}{0.000000pt}%
\pgfpathmoveto{\pgfqpoint{2.746530in}{1.995363in}}%
\pgfpathlineto{\pgfqpoint{2.746530in}{2.321613in}}%
\pgfusepath{stroke}%
\end{pgfscope}%
\begin{pgfscope}%
\pgfsetbuttcap%
\pgfsetroundjoin%
\definecolor{currentfill}{rgb}{0.000000,0.000000,0.000000}%
\pgfsetfillcolor{currentfill}%
\pgfsetlinewidth{0.803000pt}%
\definecolor{currentstroke}{rgb}{0.000000,0.000000,0.000000}%
\pgfsetstrokecolor{currentstroke}%
\pgfsetdash{}{0pt}%
\pgfsys@defobject{currentmarker}{\pgfqpoint{0.000000in}{0.000000in}}{\pgfqpoint{0.000000in}{0.027778in}}{%
\pgfpathmoveto{\pgfqpoint{0.000000in}{0.000000in}}%
\pgfpathlineto{\pgfqpoint{0.000000in}{0.027778in}}%
\pgfusepath{stroke,fill}%
}%
\begin{pgfscope}%
\pgfsys@transformshift{2.746530in}{1.995363in}%
\pgfsys@useobject{currentmarker}{}%
\end{pgfscope}%
\end{pgfscope}%
\begin{pgfscope}%
\definecolor{textcolor}{rgb}{0.000000,0.000000,0.000000}%
\pgfsetstrokecolor{textcolor}%
\pgfsetfillcolor{textcolor}%
\pgftext[x=2.746530in,y=1.981474in,,top]{\color{textcolor}{\rmfamily\fontsize{5.000000}{6.000000}\selectfont\catcode`\^=\active\def^{\ifmmode\sp\else\^{}\fi}\catcode`\%=\active\def%{\%}7.250}}%
\end{pgfscope}%
\begin{pgfscope}%
\pgfpathrectangle{\pgfqpoint{2.256930in}{1.995363in}}{\pgfqpoint{0.979200in}{0.326250in}}%
\pgfusepath{clip}%
\pgfsetbuttcap%
\pgfsetroundjoin%
\pgfsetlinewidth{0.401500pt}%
\definecolor{currentstroke}{rgb}{0.690196,0.690196,0.690196}%
\pgfsetstrokecolor{currentstroke}%
\pgfsetstrokeopacity{0.250000}%
\pgfsetdash{{1.480000pt}{0.640000pt}}{0.000000pt}%
\pgfpathmoveto{\pgfqpoint{2.991330in}{1.995363in}}%
\pgfpathlineto{\pgfqpoint{2.991330in}{2.321613in}}%
\pgfusepath{stroke}%
\end{pgfscope}%
\begin{pgfscope}%
\pgfsetbuttcap%
\pgfsetroundjoin%
\definecolor{currentfill}{rgb}{0.000000,0.000000,0.000000}%
\pgfsetfillcolor{currentfill}%
\pgfsetlinewidth{0.803000pt}%
\definecolor{currentstroke}{rgb}{0.000000,0.000000,0.000000}%
\pgfsetstrokecolor{currentstroke}%
\pgfsetdash{}{0pt}%
\pgfsys@defobject{currentmarker}{\pgfqpoint{0.000000in}{0.000000in}}{\pgfqpoint{0.000000in}{0.027778in}}{%
\pgfpathmoveto{\pgfqpoint{0.000000in}{0.000000in}}%
\pgfpathlineto{\pgfqpoint{0.000000in}{0.027778in}}%
\pgfusepath{stroke,fill}%
}%
\begin{pgfscope}%
\pgfsys@transformshift{2.991330in}{1.995363in}%
\pgfsys@useobject{currentmarker}{}%
\end{pgfscope}%
\end{pgfscope}%
\begin{pgfscope}%
\definecolor{textcolor}{rgb}{0.000000,0.000000,0.000000}%
\pgfsetstrokecolor{textcolor}%
\pgfsetfillcolor{textcolor}%
\pgftext[x=2.991330in,y=1.981474in,,top]{\color{textcolor}{\rmfamily\fontsize{5.000000}{6.000000}\selectfont\catcode`\^=\active\def^{\ifmmode\sp\else\^{}\fi}\catcode`\%=\active\def%{\%}7.375}}%
\end{pgfscope}%
\begin{pgfscope}%
\pgfpathrectangle{\pgfqpoint{2.256930in}{1.995363in}}{\pgfqpoint{0.979200in}{0.326250in}}%
\pgfusepath{clip}%
\pgfsetbuttcap%
\pgfsetroundjoin%
\pgfsetlinewidth{0.401500pt}%
\definecolor{currentstroke}{rgb}{0.690196,0.690196,0.690196}%
\pgfsetstrokecolor{currentstroke}%
\pgfsetstrokeopacity{0.250000}%
\pgfsetdash{{1.480000pt}{0.640000pt}}{0.000000pt}%
\pgfpathmoveto{\pgfqpoint{3.236130in}{1.995363in}}%
\pgfpathlineto{\pgfqpoint{3.236130in}{2.321613in}}%
\pgfusepath{stroke}%
\end{pgfscope}%
\begin{pgfscope}%
\pgfsetbuttcap%
\pgfsetroundjoin%
\definecolor{currentfill}{rgb}{0.000000,0.000000,0.000000}%
\pgfsetfillcolor{currentfill}%
\pgfsetlinewidth{0.803000pt}%
\definecolor{currentstroke}{rgb}{0.000000,0.000000,0.000000}%
\pgfsetstrokecolor{currentstroke}%
\pgfsetdash{}{0pt}%
\pgfsys@defobject{currentmarker}{\pgfqpoint{0.000000in}{0.000000in}}{\pgfqpoint{0.000000in}{0.027778in}}{%
\pgfpathmoveto{\pgfqpoint{0.000000in}{0.000000in}}%
\pgfpathlineto{\pgfqpoint{0.000000in}{0.027778in}}%
\pgfusepath{stroke,fill}%
}%
\begin{pgfscope}%
\pgfsys@transformshift{3.236130in}{1.995363in}%
\pgfsys@useobject{currentmarker}{}%
\end{pgfscope}%
\end{pgfscope}%
\begin{pgfscope}%
\definecolor{textcolor}{rgb}{0.000000,0.000000,0.000000}%
\pgfsetstrokecolor{textcolor}%
\pgfsetfillcolor{textcolor}%
\pgftext[x=3.236130in,y=1.981474in,,top]{\color{textcolor}{\rmfamily\fontsize{5.000000}{6.000000}\selectfont\catcode`\^=\active\def^{\ifmmode\sp\else\^{}\fi}\catcode`\%=\active\def%{\%}7.500}}%
\end{pgfscope}%
\begin{pgfscope}%
\pgfpathrectangle{\pgfqpoint{2.256930in}{1.995363in}}{\pgfqpoint{0.979200in}{0.326250in}}%
\pgfusepath{clip}%
\pgfsetbuttcap%
\pgfsetroundjoin%
\pgfsetlinewidth{0.401500pt}%
\definecolor{currentstroke}{rgb}{0.690196,0.690196,0.690196}%
\pgfsetstrokecolor{currentstroke}%
\pgfsetstrokeopacity{0.250000}%
\pgfsetdash{{1.480000pt}{0.640000pt}}{0.000000pt}%
\pgfpathmoveto{\pgfqpoint{2.256930in}{1.995363in}}%
\pgfpathlineto{\pgfqpoint{3.236130in}{1.995363in}}%
\pgfusepath{stroke}%
\end{pgfscope}%
\begin{pgfscope}%
\pgfsetbuttcap%
\pgfsetroundjoin%
\definecolor{currentfill}{rgb}{0.000000,0.000000,0.000000}%
\pgfsetfillcolor{currentfill}%
\pgfsetlinewidth{0.803000pt}%
\definecolor{currentstroke}{rgb}{0.000000,0.000000,0.000000}%
\pgfsetstrokecolor{currentstroke}%
\pgfsetdash{}{0pt}%
\pgfsys@defobject{currentmarker}{\pgfqpoint{0.000000in}{0.000000in}}{\pgfqpoint{0.027778in}{0.000000in}}{%
\pgfpathmoveto{\pgfqpoint{0.000000in}{0.000000in}}%
\pgfpathlineto{\pgfqpoint{0.027778in}{0.000000in}}%
\pgfusepath{stroke,fill}%
}%
\begin{pgfscope}%
\pgfsys@transformshift{2.256930in}{1.995363in}%
\pgfsys@useobject{currentmarker}{}%
\end{pgfscope}%
\end{pgfscope}%
\begin{pgfscope}%
\definecolor{textcolor}{rgb}{0.000000,0.000000,0.000000}%
\pgfsetstrokecolor{textcolor}%
\pgfsetfillcolor{textcolor}%
\pgftext[x=2.040492in, y=1.971250in, left, base]{\color{textcolor}{\rmfamily\fontsize{5.000000}{6.000000}\selectfont\catcode`\^=\active\def^{\ifmmode\sp\else\^{}\fi}\catcode`\%=\active\def%{\%}-0.50}}%
\end{pgfscope}%
\begin{pgfscope}%
\pgfpathrectangle{\pgfqpoint{2.256930in}{1.995363in}}{\pgfqpoint{0.979200in}{0.326250in}}%
\pgfusepath{clip}%
\pgfsetbuttcap%
\pgfsetroundjoin%
\pgfsetlinewidth{0.401500pt}%
\definecolor{currentstroke}{rgb}{0.690196,0.690196,0.690196}%
\pgfsetstrokecolor{currentstroke}%
\pgfsetstrokeopacity{0.250000}%
\pgfsetdash{{1.480000pt}{0.640000pt}}{0.000000pt}%
\pgfpathmoveto{\pgfqpoint{2.256930in}{2.158488in}}%
\pgfpathlineto{\pgfqpoint{3.236130in}{2.158488in}}%
\pgfusepath{stroke}%
\end{pgfscope}%
\begin{pgfscope}%
\pgfsetbuttcap%
\pgfsetroundjoin%
\definecolor{currentfill}{rgb}{0.000000,0.000000,0.000000}%
\pgfsetfillcolor{currentfill}%
\pgfsetlinewidth{0.803000pt}%
\definecolor{currentstroke}{rgb}{0.000000,0.000000,0.000000}%
\pgfsetstrokecolor{currentstroke}%
\pgfsetdash{}{0pt}%
\pgfsys@defobject{currentmarker}{\pgfqpoint{0.000000in}{0.000000in}}{\pgfqpoint{0.027778in}{0.000000in}}{%
\pgfpathmoveto{\pgfqpoint{0.000000in}{0.000000in}}%
\pgfpathlineto{\pgfqpoint{0.027778in}{0.000000in}}%
\pgfusepath{stroke,fill}%
}%
\begin{pgfscope}%
\pgfsys@transformshift{2.256930in}{2.158488in}%
\pgfsys@useobject{currentmarker}{}%
\end{pgfscope}%
\end{pgfscope}%
\begin{pgfscope}%
\definecolor{textcolor}{rgb}{0.000000,0.000000,0.000000}%
\pgfsetstrokecolor{textcolor}%
\pgfsetfillcolor{textcolor}%
\pgftext[x=2.073286in, y=2.134375in, left, base]{\color{textcolor}{\rmfamily\fontsize{5.000000}{6.000000}\selectfont\catcode`\^=\active\def^{\ifmmode\sp\else\^{}\fi}\catcode`\%=\active\def%{\%}2.25}}%
\end{pgfscope}%
\begin{pgfscope}%
\pgfpathrectangle{\pgfqpoint{2.256930in}{1.995363in}}{\pgfqpoint{0.979200in}{0.326250in}}%
\pgfusepath{clip}%
\pgfsetbuttcap%
\pgfsetroundjoin%
\pgfsetlinewidth{0.401500pt}%
\definecolor{currentstroke}{rgb}{0.690196,0.690196,0.690196}%
\pgfsetstrokecolor{currentstroke}%
\pgfsetstrokeopacity{0.250000}%
\pgfsetdash{{1.480000pt}{0.640000pt}}{0.000000pt}%
\pgfpathmoveto{\pgfqpoint{2.256930in}{2.321613in}}%
\pgfpathlineto{\pgfqpoint{3.236130in}{2.321613in}}%
\pgfusepath{stroke}%
\end{pgfscope}%
\begin{pgfscope}%
\pgfsetbuttcap%
\pgfsetroundjoin%
\definecolor{currentfill}{rgb}{0.000000,0.000000,0.000000}%
\pgfsetfillcolor{currentfill}%
\pgfsetlinewidth{0.803000pt}%
\definecolor{currentstroke}{rgb}{0.000000,0.000000,0.000000}%
\pgfsetstrokecolor{currentstroke}%
\pgfsetdash{}{0pt}%
\pgfsys@defobject{currentmarker}{\pgfqpoint{0.000000in}{0.000000in}}{\pgfqpoint{0.027778in}{0.000000in}}{%
\pgfpathmoveto{\pgfqpoint{0.000000in}{0.000000in}}%
\pgfpathlineto{\pgfqpoint{0.027778in}{0.000000in}}%
\pgfusepath{stroke,fill}%
}%
\begin{pgfscope}%
\pgfsys@transformshift{2.256930in}{2.321613in}%
\pgfsys@useobject{currentmarker}{}%
\end{pgfscope}%
\end{pgfscope}%
\begin{pgfscope}%
\definecolor{textcolor}{rgb}{0.000000,0.000000,0.000000}%
\pgfsetstrokecolor{textcolor}%
\pgfsetfillcolor{textcolor}%
\pgftext[x=2.073286in, y=2.297500in, left, base]{\color{textcolor}{\rmfamily\fontsize{5.000000}{6.000000}\selectfont\catcode`\^=\active\def^{\ifmmode\sp\else\^{}\fi}\catcode`\%=\active\def%{\%}5.00}}%
\end{pgfscope}%
\begin{pgfscope}%
\pgfpathrectangle{\pgfqpoint{2.256930in}{1.995363in}}{\pgfqpoint{0.979200in}{0.326250in}}%
\pgfusepath{clip}%
\pgfsetrectcap%
\pgfsetroundjoin%
\pgfsetlinewidth{1.505625pt}%
\definecolor{currentstroke}{rgb}{0.835294,0.368627,0.000000}%
\pgfsetstrokecolor{currentstroke}%
\pgfsetdash{}{0pt}%
\pgfpathmoveto{\pgfqpoint{2.256642in}{2.051421in}}%
\pgfpathlineto{\pgfqpoint{2.257025in}{2.049967in}}%
\pgfpathlineto{\pgfqpoint{2.257318in}{2.051562in}}%
\pgfpathlineto{\pgfqpoint{2.296791in}{2.287992in}}%
\pgfpathlineto{\pgfqpoint{2.297125in}{2.287021in}}%
\pgfpathlineto{\pgfqpoint{2.354948in}{2.052416in}}%
\pgfpathlineto{\pgfqpoint{2.355083in}{2.053083in}}%
\pgfpathlineto{\pgfqpoint{2.394618in}{2.290234in}}%
\pgfpathlineto{\pgfqpoint{2.395221in}{2.288309in}}%
\pgfpathlineto{\pgfqpoint{2.452866in}{2.052145in}}%
\pgfpathlineto{\pgfqpoint{2.453220in}{2.054113in}}%
\pgfpathlineto{\pgfqpoint{2.492631in}{2.290148in}}%
\pgfpathlineto{\pgfqpoint{2.492965in}{2.289178in}}%
\pgfpathlineto{\pgfqpoint{2.550780in}{2.054294in}}%
\pgfpathlineto{\pgfqpoint{2.550938in}{2.055041in}}%
\pgfpathlineto{\pgfqpoint{2.590555in}{2.291872in}}%
\pgfpathlineto{\pgfqpoint{2.591035in}{2.289907in}}%
\pgfpathlineto{\pgfqpoint{2.648695in}{2.054644in}}%
\pgfpathlineto{\pgfqpoint{2.648928in}{2.055850in}}%
\pgfpathlineto{\pgfqpoint{2.688377in}{2.292237in}}%
\pgfpathlineto{\pgfqpoint{2.688856in}{2.290821in}}%
\pgfpathlineto{\pgfqpoint{2.746625in}{2.053702in}}%
\pgfpathlineto{\pgfqpoint{2.746968in}{2.055594in}}%
\pgfpathlineto{\pgfqpoint{2.786395in}{2.290993in}}%
\pgfpathlineto{\pgfqpoint{2.786614in}{2.290093in}}%
\pgfpathlineto{\pgfqpoint{2.844549in}{2.052114in}}%
\pgfpathlineto{\pgfqpoint{2.844932in}{2.054217in}}%
\pgfpathlineto{\pgfqpoint{2.884217in}{2.289454in}}%
\pgfpathlineto{\pgfqpoint{2.884605in}{2.288408in}}%
\pgfpathlineto{\pgfqpoint{2.942465in}{2.050788in}}%
\pgfpathlineto{\pgfqpoint{2.942808in}{2.052679in}}%
\pgfpathlineto{\pgfqpoint{2.982228in}{2.288031in}}%
\pgfpathlineto{\pgfqpoint{2.982445in}{2.287143in}}%
\pgfpathlineto{\pgfqpoint{3.040386in}{2.049861in}}%
\pgfpathlineto{\pgfqpoint{3.040694in}{2.051589in}}%
\pgfpathlineto{\pgfqpoint{3.080057in}{2.287432in}}%
\pgfpathlineto{\pgfqpoint{3.080446in}{2.286385in}}%
\pgfpathlineto{\pgfqpoint{3.138305in}{2.049044in}}%
\pgfpathlineto{\pgfqpoint{3.138660in}{2.051010in}}%
\pgfpathlineto{\pgfqpoint{3.177994in}{2.286757in}}%
\pgfpathlineto{\pgfqpoint{3.178305in}{2.286009in}}%
\pgfpathlineto{\pgfqpoint{3.236130in}{2.052131in}}%
\pgfpathlineto{\pgfqpoint{3.236130in}{2.052131in}}%
\pgfusepath{stroke}%
\end{pgfscope}%
\begin{pgfscope}%
\pgfsetrectcap%
\pgfsetmiterjoin%
\pgfsetlinewidth{0.803000pt}%
\definecolor{currentstroke}{rgb}{0.000000,0.000000,0.000000}%
\pgfsetstrokecolor{currentstroke}%
\pgfsetdash{}{0pt}%
\pgfpathmoveto{\pgfqpoint{2.256930in}{1.995363in}}%
\pgfpathlineto{\pgfqpoint{2.256930in}{2.321613in}}%
\pgfusepath{stroke}%
\end{pgfscope}%
\begin{pgfscope}%
\pgfsetrectcap%
\pgfsetmiterjoin%
\pgfsetlinewidth{0.803000pt}%
\definecolor{currentstroke}{rgb}{0.000000,0.000000,0.000000}%
\pgfsetstrokecolor{currentstroke}%
\pgfsetdash{}{0pt}%
\pgfpathmoveto{\pgfqpoint{3.236130in}{1.995363in}}%
\pgfpathlineto{\pgfqpoint{3.236130in}{2.321613in}}%
\pgfusepath{stroke}%
\end{pgfscope}%
\begin{pgfscope}%
\pgfsetrectcap%
\pgfsetmiterjoin%
\pgfsetlinewidth{0.803000pt}%
\definecolor{currentstroke}{rgb}{0.000000,0.000000,0.000000}%
\pgfsetstrokecolor{currentstroke}%
\pgfsetdash{}{0pt}%
\pgfpathmoveto{\pgfqpoint{2.256930in}{1.995363in}}%
\pgfpathlineto{\pgfqpoint{3.236130in}{1.995363in}}%
\pgfusepath{stroke}%
\end{pgfscope}%
\begin{pgfscope}%
\pgfsetrectcap%
\pgfsetmiterjoin%
\pgfsetlinewidth{0.803000pt}%
\definecolor{currentstroke}{rgb}{0.000000,0.000000,0.000000}%
\pgfsetstrokecolor{currentstroke}%
\pgfsetdash{}{0pt}%
\pgfpathmoveto{\pgfqpoint{2.256930in}{2.321613in}}%
\pgfpathlineto{\pgfqpoint{3.236130in}{2.321613in}}%
\pgfusepath{stroke}%
\end{pgfscope}%
\begin{pgfscope}%
\pgfsetbuttcap%
\pgfsetmiterjoin%
\definecolor{currentfill}{rgb}{1.000000,1.000000,1.000000}%
\pgfsetfillcolor{currentfill}%
\pgfsetlinewidth{0.000000pt}%
\definecolor{currentstroke}{rgb}{0.000000,0.000000,0.000000}%
\pgfsetstrokecolor{currentstroke}%
\pgfsetstrokeopacity{0.000000}%
\pgfsetdash{}{0pt}%
\pgfpathmoveto{\pgfqpoint{2.256930in}{0.925988in}}%
\pgfpathlineto{\pgfqpoint{3.236130in}{0.925988in}}%
\pgfpathlineto{\pgfqpoint{3.236130in}{1.252238in}}%
\pgfpathlineto{\pgfqpoint{2.256930in}{1.252238in}}%
\pgfpathlineto{\pgfqpoint{2.256930in}{0.925988in}}%
\pgfpathclose%
\pgfusepath{fill}%
\end{pgfscope}%
\begin{pgfscope}%
\pgfpathrectangle{\pgfqpoint{2.256930in}{0.925988in}}{\pgfqpoint{0.979200in}{0.326250in}}%
\pgfusepath{clip}%
\pgfsetbuttcap%
\pgfsetroundjoin%
\pgfsetlinewidth{0.401500pt}%
\definecolor{currentstroke}{rgb}{0.690196,0.690196,0.690196}%
\pgfsetstrokecolor{currentstroke}%
\pgfsetstrokeopacity{0.250000}%
\pgfsetdash{{1.480000pt}{0.640000pt}}{0.000000pt}%
\pgfpathmoveto{\pgfqpoint{2.256930in}{0.925988in}}%
\pgfpathlineto{\pgfqpoint{2.256930in}{1.252238in}}%
\pgfusepath{stroke}%
\end{pgfscope}%
\begin{pgfscope}%
\pgfsetbuttcap%
\pgfsetroundjoin%
\definecolor{currentfill}{rgb}{0.000000,0.000000,0.000000}%
\pgfsetfillcolor{currentfill}%
\pgfsetlinewidth{0.803000pt}%
\definecolor{currentstroke}{rgb}{0.000000,0.000000,0.000000}%
\pgfsetstrokecolor{currentstroke}%
\pgfsetdash{}{0pt}%
\pgfsys@defobject{currentmarker}{\pgfqpoint{0.000000in}{0.000000in}}{\pgfqpoint{0.000000in}{0.027778in}}{%
\pgfpathmoveto{\pgfqpoint{0.000000in}{0.000000in}}%
\pgfpathlineto{\pgfqpoint{0.000000in}{0.027778in}}%
\pgfusepath{stroke,fill}%
}%
\begin{pgfscope}%
\pgfsys@transformshift{2.256930in}{0.925988in}%
\pgfsys@useobject{currentmarker}{}%
\end{pgfscope}%
\end{pgfscope}%
\begin{pgfscope}%
\definecolor{textcolor}{rgb}{0.000000,0.000000,0.000000}%
\pgfsetstrokecolor{textcolor}%
\pgfsetfillcolor{textcolor}%
\pgftext[x=2.256930in,y=0.912099in,,top]{\color{textcolor}{\rmfamily\fontsize{5.000000}{6.000000}\selectfont\catcode`\^=\active\def^{\ifmmode\sp\else\^{}\fi}\catcode`\%=\active\def%{\%}7.000}}%
\end{pgfscope}%
\begin{pgfscope}%
\pgfpathrectangle{\pgfqpoint{2.256930in}{0.925988in}}{\pgfqpoint{0.979200in}{0.326250in}}%
\pgfusepath{clip}%
\pgfsetbuttcap%
\pgfsetroundjoin%
\pgfsetlinewidth{0.401500pt}%
\definecolor{currentstroke}{rgb}{0.690196,0.690196,0.690196}%
\pgfsetstrokecolor{currentstroke}%
\pgfsetstrokeopacity{0.250000}%
\pgfsetdash{{1.480000pt}{0.640000pt}}{0.000000pt}%
\pgfpathmoveto{\pgfqpoint{2.501730in}{0.925988in}}%
\pgfpathlineto{\pgfqpoint{2.501730in}{1.252238in}}%
\pgfusepath{stroke}%
\end{pgfscope}%
\begin{pgfscope}%
\pgfsetbuttcap%
\pgfsetroundjoin%
\definecolor{currentfill}{rgb}{0.000000,0.000000,0.000000}%
\pgfsetfillcolor{currentfill}%
\pgfsetlinewidth{0.803000pt}%
\definecolor{currentstroke}{rgb}{0.000000,0.000000,0.000000}%
\pgfsetstrokecolor{currentstroke}%
\pgfsetdash{}{0pt}%
\pgfsys@defobject{currentmarker}{\pgfqpoint{0.000000in}{0.000000in}}{\pgfqpoint{0.000000in}{0.027778in}}{%
\pgfpathmoveto{\pgfqpoint{0.000000in}{0.000000in}}%
\pgfpathlineto{\pgfqpoint{0.000000in}{0.027778in}}%
\pgfusepath{stroke,fill}%
}%
\begin{pgfscope}%
\pgfsys@transformshift{2.501730in}{0.925988in}%
\pgfsys@useobject{currentmarker}{}%
\end{pgfscope}%
\end{pgfscope}%
\begin{pgfscope}%
\definecolor{textcolor}{rgb}{0.000000,0.000000,0.000000}%
\pgfsetstrokecolor{textcolor}%
\pgfsetfillcolor{textcolor}%
\pgftext[x=2.501730in,y=0.912099in,,top]{\color{textcolor}{\rmfamily\fontsize{5.000000}{6.000000}\selectfont\catcode`\^=\active\def^{\ifmmode\sp\else\^{}\fi}\catcode`\%=\active\def%{\%}7.125}}%
\end{pgfscope}%
\begin{pgfscope}%
\pgfpathrectangle{\pgfqpoint{2.256930in}{0.925988in}}{\pgfqpoint{0.979200in}{0.326250in}}%
\pgfusepath{clip}%
\pgfsetbuttcap%
\pgfsetroundjoin%
\pgfsetlinewidth{0.401500pt}%
\definecolor{currentstroke}{rgb}{0.690196,0.690196,0.690196}%
\pgfsetstrokecolor{currentstroke}%
\pgfsetstrokeopacity{0.250000}%
\pgfsetdash{{1.480000pt}{0.640000pt}}{0.000000pt}%
\pgfpathmoveto{\pgfqpoint{2.746530in}{0.925988in}}%
\pgfpathlineto{\pgfqpoint{2.746530in}{1.252238in}}%
\pgfusepath{stroke}%
\end{pgfscope}%
\begin{pgfscope}%
\pgfsetbuttcap%
\pgfsetroundjoin%
\definecolor{currentfill}{rgb}{0.000000,0.000000,0.000000}%
\pgfsetfillcolor{currentfill}%
\pgfsetlinewidth{0.803000pt}%
\definecolor{currentstroke}{rgb}{0.000000,0.000000,0.000000}%
\pgfsetstrokecolor{currentstroke}%
\pgfsetdash{}{0pt}%
\pgfsys@defobject{currentmarker}{\pgfqpoint{0.000000in}{0.000000in}}{\pgfqpoint{0.000000in}{0.027778in}}{%
\pgfpathmoveto{\pgfqpoint{0.000000in}{0.000000in}}%
\pgfpathlineto{\pgfqpoint{0.000000in}{0.027778in}}%
\pgfusepath{stroke,fill}%
}%
\begin{pgfscope}%
\pgfsys@transformshift{2.746530in}{0.925988in}%
\pgfsys@useobject{currentmarker}{}%
\end{pgfscope}%
\end{pgfscope}%
\begin{pgfscope}%
\definecolor{textcolor}{rgb}{0.000000,0.000000,0.000000}%
\pgfsetstrokecolor{textcolor}%
\pgfsetfillcolor{textcolor}%
\pgftext[x=2.746530in,y=0.912099in,,top]{\color{textcolor}{\rmfamily\fontsize{5.000000}{6.000000}\selectfont\catcode`\^=\active\def^{\ifmmode\sp\else\^{}\fi}\catcode`\%=\active\def%{\%}7.250}}%
\end{pgfscope}%
\begin{pgfscope}%
\pgfpathrectangle{\pgfqpoint{2.256930in}{0.925988in}}{\pgfqpoint{0.979200in}{0.326250in}}%
\pgfusepath{clip}%
\pgfsetbuttcap%
\pgfsetroundjoin%
\pgfsetlinewidth{0.401500pt}%
\definecolor{currentstroke}{rgb}{0.690196,0.690196,0.690196}%
\pgfsetstrokecolor{currentstroke}%
\pgfsetstrokeopacity{0.250000}%
\pgfsetdash{{1.480000pt}{0.640000pt}}{0.000000pt}%
\pgfpathmoveto{\pgfqpoint{2.991330in}{0.925988in}}%
\pgfpathlineto{\pgfqpoint{2.991330in}{1.252238in}}%
\pgfusepath{stroke}%
\end{pgfscope}%
\begin{pgfscope}%
\pgfsetbuttcap%
\pgfsetroundjoin%
\definecolor{currentfill}{rgb}{0.000000,0.000000,0.000000}%
\pgfsetfillcolor{currentfill}%
\pgfsetlinewidth{0.803000pt}%
\definecolor{currentstroke}{rgb}{0.000000,0.000000,0.000000}%
\pgfsetstrokecolor{currentstroke}%
\pgfsetdash{}{0pt}%
\pgfsys@defobject{currentmarker}{\pgfqpoint{0.000000in}{0.000000in}}{\pgfqpoint{0.000000in}{0.027778in}}{%
\pgfpathmoveto{\pgfqpoint{0.000000in}{0.000000in}}%
\pgfpathlineto{\pgfqpoint{0.000000in}{0.027778in}}%
\pgfusepath{stroke,fill}%
}%
\begin{pgfscope}%
\pgfsys@transformshift{2.991330in}{0.925988in}%
\pgfsys@useobject{currentmarker}{}%
\end{pgfscope}%
\end{pgfscope}%
\begin{pgfscope}%
\definecolor{textcolor}{rgb}{0.000000,0.000000,0.000000}%
\pgfsetstrokecolor{textcolor}%
\pgfsetfillcolor{textcolor}%
\pgftext[x=2.991330in,y=0.912099in,,top]{\color{textcolor}{\rmfamily\fontsize{5.000000}{6.000000}\selectfont\catcode`\^=\active\def^{\ifmmode\sp\else\^{}\fi}\catcode`\%=\active\def%{\%}7.375}}%
\end{pgfscope}%
\begin{pgfscope}%
\pgfpathrectangle{\pgfqpoint{2.256930in}{0.925988in}}{\pgfqpoint{0.979200in}{0.326250in}}%
\pgfusepath{clip}%
\pgfsetbuttcap%
\pgfsetroundjoin%
\pgfsetlinewidth{0.401500pt}%
\definecolor{currentstroke}{rgb}{0.690196,0.690196,0.690196}%
\pgfsetstrokecolor{currentstroke}%
\pgfsetstrokeopacity{0.250000}%
\pgfsetdash{{1.480000pt}{0.640000pt}}{0.000000pt}%
\pgfpathmoveto{\pgfqpoint{3.236130in}{0.925988in}}%
\pgfpathlineto{\pgfqpoint{3.236130in}{1.252238in}}%
\pgfusepath{stroke}%
\end{pgfscope}%
\begin{pgfscope}%
\pgfsetbuttcap%
\pgfsetroundjoin%
\definecolor{currentfill}{rgb}{0.000000,0.000000,0.000000}%
\pgfsetfillcolor{currentfill}%
\pgfsetlinewidth{0.803000pt}%
\definecolor{currentstroke}{rgb}{0.000000,0.000000,0.000000}%
\pgfsetstrokecolor{currentstroke}%
\pgfsetdash{}{0pt}%
\pgfsys@defobject{currentmarker}{\pgfqpoint{0.000000in}{0.000000in}}{\pgfqpoint{0.000000in}{0.027778in}}{%
\pgfpathmoveto{\pgfqpoint{0.000000in}{0.000000in}}%
\pgfpathlineto{\pgfqpoint{0.000000in}{0.027778in}}%
\pgfusepath{stroke,fill}%
}%
\begin{pgfscope}%
\pgfsys@transformshift{3.236130in}{0.925988in}%
\pgfsys@useobject{currentmarker}{}%
\end{pgfscope}%
\end{pgfscope}%
\begin{pgfscope}%
\definecolor{textcolor}{rgb}{0.000000,0.000000,0.000000}%
\pgfsetstrokecolor{textcolor}%
\pgfsetfillcolor{textcolor}%
\pgftext[x=3.236130in,y=0.912099in,,top]{\color{textcolor}{\rmfamily\fontsize{5.000000}{6.000000}\selectfont\catcode`\^=\active\def^{\ifmmode\sp\else\^{}\fi}\catcode`\%=\active\def%{\%}7.500}}%
\end{pgfscope}%
\begin{pgfscope}%
\pgfpathrectangle{\pgfqpoint{2.256930in}{0.925988in}}{\pgfqpoint{0.979200in}{0.326250in}}%
\pgfusepath{clip}%
\pgfsetbuttcap%
\pgfsetroundjoin%
\pgfsetlinewidth{0.401500pt}%
\definecolor{currentstroke}{rgb}{0.690196,0.690196,0.690196}%
\pgfsetstrokecolor{currentstroke}%
\pgfsetstrokeopacity{0.250000}%
\pgfsetdash{{1.480000pt}{0.640000pt}}{0.000000pt}%
\pgfpathmoveto{\pgfqpoint{2.256930in}{0.925988in}}%
\pgfpathlineto{\pgfqpoint{3.236130in}{0.925988in}}%
\pgfusepath{stroke}%
\end{pgfscope}%
\begin{pgfscope}%
\pgfsetbuttcap%
\pgfsetroundjoin%
\definecolor{currentfill}{rgb}{0.000000,0.000000,0.000000}%
\pgfsetfillcolor{currentfill}%
\pgfsetlinewidth{0.803000pt}%
\definecolor{currentstroke}{rgb}{0.000000,0.000000,0.000000}%
\pgfsetstrokecolor{currentstroke}%
\pgfsetdash{}{0pt}%
\pgfsys@defobject{currentmarker}{\pgfqpoint{0.000000in}{0.000000in}}{\pgfqpoint{0.027778in}{0.000000in}}{%
\pgfpathmoveto{\pgfqpoint{0.000000in}{0.000000in}}%
\pgfpathlineto{\pgfqpoint{0.027778in}{0.000000in}}%
\pgfusepath{stroke,fill}%
}%
\begin{pgfscope}%
\pgfsys@transformshift{2.256930in}{0.925988in}%
\pgfsys@useobject{currentmarker}{}%
\end{pgfscope}%
\end{pgfscope}%
\begin{pgfscope}%
\definecolor{textcolor}{rgb}{0.000000,0.000000,0.000000}%
\pgfsetstrokecolor{textcolor}%
\pgfsetfillcolor{textcolor}%
\pgftext[x=1.993231in, y=0.901875in, left, base]{\color{textcolor}{\rmfamily\fontsize{5.000000}{6.000000}\selectfont\catcode`\^=\active\def^{\ifmmode\sp\else\^{}\fi}\catcode`\%=\active\def%{\%}-52.41}}%
\end{pgfscope}%
\begin{pgfscope}%
\pgfpathrectangle{\pgfqpoint{2.256930in}{0.925988in}}{\pgfqpoint{0.979200in}{0.326250in}}%
\pgfusepath{clip}%
\pgfsetbuttcap%
\pgfsetroundjoin%
\pgfsetlinewidth{0.401500pt}%
\definecolor{currentstroke}{rgb}{0.690196,0.690196,0.690196}%
\pgfsetstrokecolor{currentstroke}%
\pgfsetstrokeopacity{0.250000}%
\pgfsetdash{{1.480000pt}{0.640000pt}}{0.000000pt}%
\pgfpathmoveto{\pgfqpoint{2.256930in}{1.089113in}}%
\pgfpathlineto{\pgfqpoint{3.236130in}{1.089113in}}%
\pgfusepath{stroke}%
\end{pgfscope}%
\begin{pgfscope}%
\pgfsetbuttcap%
\pgfsetroundjoin%
\definecolor{currentfill}{rgb}{0.000000,0.000000,0.000000}%
\pgfsetfillcolor{currentfill}%
\pgfsetlinewidth{0.803000pt}%
\definecolor{currentstroke}{rgb}{0.000000,0.000000,0.000000}%
\pgfsetstrokecolor{currentstroke}%
\pgfsetdash{}{0pt}%
\pgfsys@defobject{currentmarker}{\pgfqpoint{0.000000in}{0.000000in}}{\pgfqpoint{0.027778in}{0.000000in}}{%
\pgfpathmoveto{\pgfqpoint{0.000000in}{0.000000in}}%
\pgfpathlineto{\pgfqpoint{0.027778in}{0.000000in}}%
\pgfusepath{stroke,fill}%
}%
\begin{pgfscope}%
\pgfsys@transformshift{2.256930in}{1.089113in}%
\pgfsys@useobject{currentmarker}{}%
\end{pgfscope}%
\end{pgfscope}%
\begin{pgfscope}%
\definecolor{textcolor}{rgb}{0.000000,0.000000,0.000000}%
\pgfsetstrokecolor{textcolor}%
\pgfsetfillcolor{textcolor}%
\pgftext[x=1.978763in, y=1.065000in, left, base]{\color{textcolor}{\rmfamily\fontsize{5.000000}{6.000000}\selectfont\catcode`\^=\active\def^{\ifmmode\sp\else\^{}\fi}\catcode`\%=\active\def%{\%}261.94}}%
\end{pgfscope}%
\begin{pgfscope}%
\pgfpathrectangle{\pgfqpoint{2.256930in}{0.925988in}}{\pgfqpoint{0.979200in}{0.326250in}}%
\pgfusepath{clip}%
\pgfsetbuttcap%
\pgfsetroundjoin%
\pgfsetlinewidth{0.401500pt}%
\definecolor{currentstroke}{rgb}{0.690196,0.690196,0.690196}%
\pgfsetstrokecolor{currentstroke}%
\pgfsetstrokeopacity{0.250000}%
\pgfsetdash{{1.480000pt}{0.640000pt}}{0.000000pt}%
\pgfpathmoveto{\pgfqpoint{2.256930in}{1.252238in}}%
\pgfpathlineto{\pgfqpoint{3.236130in}{1.252238in}}%
\pgfusepath{stroke}%
\end{pgfscope}%
\begin{pgfscope}%
\pgfsetbuttcap%
\pgfsetroundjoin%
\definecolor{currentfill}{rgb}{0.000000,0.000000,0.000000}%
\pgfsetfillcolor{currentfill}%
\pgfsetlinewidth{0.803000pt}%
\definecolor{currentstroke}{rgb}{0.000000,0.000000,0.000000}%
\pgfsetstrokecolor{currentstroke}%
\pgfsetdash{}{0pt}%
\pgfsys@defobject{currentmarker}{\pgfqpoint{0.000000in}{0.000000in}}{\pgfqpoint{0.027778in}{0.000000in}}{%
\pgfpathmoveto{\pgfqpoint{0.000000in}{0.000000in}}%
\pgfpathlineto{\pgfqpoint{0.027778in}{0.000000in}}%
\pgfusepath{stroke,fill}%
}%
\begin{pgfscope}%
\pgfsys@transformshift{2.256930in}{1.252238in}%
\pgfsys@useobject{currentmarker}{}%
\end{pgfscope}%
\end{pgfscope}%
\begin{pgfscope}%
\definecolor{textcolor}{rgb}{0.000000,0.000000,0.000000}%
\pgfsetstrokecolor{textcolor}%
\pgfsetfillcolor{textcolor}%
\pgftext[x=1.978763in, y=1.228125in, left, base]{\color{textcolor}{\rmfamily\fontsize{5.000000}{6.000000}\selectfont\catcode`\^=\active\def^{\ifmmode\sp\else\^{}\fi}\catcode`\%=\active\def%{\%}576.30}}%
\end{pgfscope}%
\begin{pgfscope}%
\pgfpathrectangle{\pgfqpoint{2.256930in}{0.925988in}}{\pgfqpoint{0.979200in}{0.326250in}}%
\pgfusepath{clip}%
\pgfsetrectcap%
\pgfsetroundjoin%
\pgfsetlinewidth{1.505625pt}%
\definecolor{currentstroke}{rgb}{0.000000,0.619608,0.450980}%
\pgfsetstrokecolor{currentstroke}%
\pgfsetdash{}{0pt}%
\pgfpathmoveto{\pgfqpoint{2.256642in}{1.108883in}}%
\pgfpathlineto{\pgfqpoint{2.256978in}{1.108982in}}%
\pgfpathlineto{\pgfqpoint{2.257366in}{0.953182in}}%
\pgfpathlineto{\pgfqpoint{2.257822in}{0.953202in}}%
\pgfpathlineto{\pgfqpoint{2.296278in}{0.953872in}}%
\pgfpathlineto{\pgfqpoint{2.296621in}{0.960155in}}%
\pgfpathlineto{\pgfqpoint{2.297927in}{1.225049in}}%
\pgfpathlineto{\pgfqpoint{2.349335in}{1.225016in}}%
\pgfpathlineto{\pgfqpoint{2.350315in}{1.103628in}}%
\pgfpathlineto{\pgfqpoint{2.351276in}{1.193758in}}%
\pgfpathlineto{\pgfqpoint{2.352966in}{1.082862in}}%
\pgfpathlineto{\pgfqpoint{2.354850in}{1.108774in}}%
\pgfpathlineto{\pgfqpoint{2.354900in}{1.108699in}}%
\pgfpathlineto{\pgfqpoint{2.355304in}{0.953188in}}%
\pgfpathlineto{\pgfqpoint{2.355646in}{0.953198in}}%
\pgfpathlineto{\pgfqpoint{2.394263in}{0.954250in}}%
\pgfpathlineto{\pgfqpoint{2.394482in}{0.958349in}}%
\pgfpathlineto{\pgfqpoint{2.395911in}{1.225048in}}%
\pgfpathlineto{\pgfqpoint{2.443802in}{1.225029in}}%
\pgfpathlineto{\pgfqpoint{2.444782in}{1.123276in}}%
\pgfpathlineto{\pgfqpoint{2.445761in}{1.225026in}}%
\pgfpathlineto{\pgfqpoint{2.446740in}{1.116582in}}%
\pgfpathlineto{\pgfqpoint{2.447719in}{1.225020in}}%
\pgfpathlineto{\pgfqpoint{2.448698in}{1.123227in}}%
\pgfpathlineto{\pgfqpoint{2.449678in}{1.225008in}}%
\pgfpathlineto{\pgfqpoint{2.451445in}{1.108978in}}%
\pgfpathlineto{\pgfqpoint{2.452819in}{1.108972in}}%
\pgfpathlineto{\pgfqpoint{2.453220in}{0.953189in}}%
\pgfpathlineto{\pgfqpoint{2.453606in}{0.953201in}}%
\pgfpathlineto{\pgfqpoint{2.492118in}{0.953872in}}%
\pgfpathlineto{\pgfqpoint{2.492461in}{0.960155in}}%
\pgfpathlineto{\pgfqpoint{2.493766in}{1.225049in}}%
\pgfpathlineto{\pgfqpoint{2.545172in}{1.225016in}}%
\pgfpathlineto{\pgfqpoint{2.546151in}{1.105153in}}%
\pgfpathlineto{\pgfqpoint{2.547119in}{1.193498in}}%
\pgfpathlineto{\pgfqpoint{2.548853in}{1.084570in}}%
\pgfpathlineto{\pgfqpoint{2.550690in}{1.108761in}}%
\pgfpathlineto{\pgfqpoint{2.550732in}{1.108674in}}%
\pgfpathlineto{\pgfqpoint{2.551106in}{0.953175in}}%
\pgfpathlineto{\pgfqpoint{2.551608in}{0.953200in}}%
\pgfpathlineto{\pgfqpoint{2.590108in}{0.954270in}}%
\pgfpathlineto{\pgfqpoint{2.590310in}{0.957969in}}%
\pgfpathlineto{\pgfqpoint{2.590427in}{0.982318in}}%
\pgfpathlineto{\pgfqpoint{2.591225in}{1.225050in}}%
\pgfpathlineto{\pgfqpoint{2.639218in}{1.225031in}}%
\pgfpathlineto{\pgfqpoint{2.640198in}{1.108813in}}%
\pgfpathlineto{\pgfqpoint{2.641177in}{1.225028in}}%
\pgfpathlineto{\pgfqpoint{2.642156in}{1.108924in}}%
\pgfpathlineto{\pgfqpoint{2.643135in}{1.225024in}}%
\pgfpathlineto{\pgfqpoint{2.644114in}{1.109128in}}%
\pgfpathlineto{\pgfqpoint{2.645088in}{1.225017in}}%
\pgfpathlineto{\pgfqpoint{2.645950in}{1.136325in}}%
\pgfpathlineto{\pgfqpoint{2.646704in}{1.224981in}}%
\pgfpathlineto{\pgfqpoint{2.647041in}{1.198440in}}%
\pgfpathlineto{\pgfqpoint{2.648138in}{1.083328in}}%
\pgfpathlineto{\pgfqpoint{2.648673in}{1.086110in}}%
\pgfpathlineto{\pgfqpoint{2.649064in}{0.953194in}}%
\pgfpathlineto{\pgfqpoint{2.649560in}{0.953200in}}%
\pgfpathlineto{\pgfqpoint{2.688021in}{0.954229in}}%
\pgfpathlineto{\pgfqpoint{2.688241in}{0.958328in}}%
\pgfpathlineto{\pgfqpoint{2.688377in}{0.992504in}}%
\pgfpathlineto{\pgfqpoint{2.689218in}{1.225050in}}%
\pgfpathlineto{\pgfqpoint{2.737524in}{1.225029in}}%
\pgfpathlineto{\pgfqpoint{2.738503in}{1.122062in}}%
\pgfpathlineto{\pgfqpoint{2.739483in}{1.225026in}}%
\pgfpathlineto{\pgfqpoint{2.740462in}{1.117777in}}%
\pgfpathlineto{\pgfqpoint{2.741441in}{1.225020in}}%
\pgfpathlineto{\pgfqpoint{2.742420in}{1.122014in}}%
\pgfpathlineto{\pgfqpoint{2.743399in}{1.225008in}}%
\pgfpathlineto{\pgfqpoint{2.745015in}{1.108980in}}%
\pgfpathlineto{\pgfqpoint{2.745145in}{1.109010in}}%
\pgfpathlineto{\pgfqpoint{2.746579in}{1.108982in}}%
\pgfpathlineto{\pgfqpoint{2.746968in}{0.953183in}}%
\pgfpathlineto{\pgfqpoint{2.747434in}{0.953202in}}%
\pgfpathlineto{\pgfqpoint{2.785955in}{0.954397in}}%
\pgfpathlineto{\pgfqpoint{2.786140in}{0.957864in}}%
\pgfpathlineto{\pgfqpoint{2.786262in}{0.982123in}}%
\pgfpathlineto{\pgfqpoint{2.787118in}{1.225050in}}%
\pgfpathlineto{\pgfqpoint{2.833469in}{1.225034in}}%
\pgfpathlineto{\pgfqpoint{2.834448in}{1.129881in}}%
\pgfpathlineto{\pgfqpoint{2.835427in}{1.225032in}}%
\pgfpathlineto{\pgfqpoint{2.836406in}{1.110231in}}%
\pgfpathlineto{\pgfqpoint{2.837385in}{1.225029in}}%
\pgfpathlineto{\pgfqpoint{2.838365in}{1.129818in}}%
\pgfpathlineto{\pgfqpoint{2.839344in}{1.225026in}}%
\pgfpathlineto{\pgfqpoint{2.840323in}{1.110224in}}%
\pgfpathlineto{\pgfqpoint{2.841302in}{1.225020in}}%
\pgfpathlineto{\pgfqpoint{2.842281in}{1.129796in}}%
\pgfpathlineto{\pgfqpoint{2.843120in}{1.225007in}}%
\pgfpathlineto{\pgfqpoint{2.844549in}{1.047887in}}%
\pgfpathlineto{\pgfqpoint{2.844890in}{0.953181in}}%
\pgfpathlineto{\pgfqpoint{2.845350in}{0.953199in}}%
\pgfpathlineto{\pgfqpoint{2.883861in}{0.954229in}}%
\pgfpathlineto{\pgfqpoint{2.884081in}{0.958327in}}%
\pgfpathlineto{\pgfqpoint{2.884217in}{0.992476in}}%
\pgfpathlineto{\pgfqpoint{2.885058in}{1.225050in}}%
\pgfpathlineto{\pgfqpoint{2.933364in}{1.225029in}}%
\pgfpathlineto{\pgfqpoint{2.934343in}{1.122135in}}%
\pgfpathlineto{\pgfqpoint{2.935322in}{1.225026in}}%
\pgfpathlineto{\pgfqpoint{2.936301in}{1.117704in}}%
\pgfpathlineto{\pgfqpoint{2.937280in}{1.225020in}}%
\pgfpathlineto{\pgfqpoint{2.938260in}{1.122088in}}%
\pgfpathlineto{\pgfqpoint{2.939239in}{1.225008in}}%
\pgfpathlineto{\pgfqpoint{2.940843in}{1.108963in}}%
\pgfpathlineto{\pgfqpoint{2.940973in}{1.109012in}}%
\pgfpathlineto{\pgfqpoint{2.942419in}{1.108972in}}%
\pgfpathlineto{\pgfqpoint{2.942808in}{0.953183in}}%
\pgfpathlineto{\pgfqpoint{2.943270in}{0.953202in}}%
\pgfpathlineto{\pgfqpoint{2.981783in}{0.954231in}}%
\pgfpathlineto{\pgfqpoint{2.981971in}{0.957719in}}%
\pgfpathlineto{\pgfqpoint{2.982090in}{0.979399in}}%
\pgfpathlineto{\pgfqpoint{2.982919in}{1.225050in}}%
\pgfpathlineto{\pgfqpoint{3.029051in}{1.225034in}}%
\pgfpathlineto{\pgfqpoint{3.030031in}{1.215954in}}%
\pgfpathlineto{\pgfqpoint{3.031010in}{1.225032in}}%
\pgfpathlineto{\pgfqpoint{3.031989in}{1.108559in}}%
\pgfpathlineto{\pgfqpoint{3.032968in}{1.225029in}}%
\pgfpathlineto{\pgfqpoint{3.033947in}{1.108730in}}%
\pgfpathlineto{\pgfqpoint{3.034927in}{1.225025in}}%
\pgfpathlineto{\pgfqpoint{3.035906in}{1.108867in}}%
\pgfpathlineto{\pgfqpoint{3.036885in}{1.225019in}}%
\pgfpathlineto{\pgfqpoint{3.037858in}{1.109111in}}%
\pgfpathlineto{\pgfqpoint{3.038728in}{1.224999in}}%
\pgfpathlineto{\pgfqpoint{3.039109in}{1.209288in}}%
\pgfpathlineto{\pgfqpoint{3.040409in}{0.983528in}}%
\pgfpathlineto{\pgfqpoint{3.040768in}{0.953193in}}%
\pgfpathlineto{\pgfqpoint{3.041019in}{0.953196in}}%
\pgfpathlineto{\pgfqpoint{3.079702in}{0.954238in}}%
\pgfpathlineto{\pgfqpoint{3.079921in}{0.958331in}}%
\pgfpathlineto{\pgfqpoint{3.080057in}{0.992558in}}%
\pgfpathlineto{\pgfqpoint{3.080901in}{1.225050in}}%
\pgfpathlineto{\pgfqpoint{3.129224in}{1.225029in}}%
\pgfpathlineto{\pgfqpoint{3.130203in}{1.111841in}}%
\pgfpathlineto{\pgfqpoint{3.131182in}{1.225026in}}%
\pgfpathlineto{\pgfqpoint{3.132161in}{1.128166in}}%
\pgfpathlineto{\pgfqpoint{3.133140in}{1.225020in}}%
\pgfpathlineto{\pgfqpoint{3.134120in}{1.111804in}}%
\pgfpathlineto{\pgfqpoint{3.135099in}{1.225006in}}%
\pgfpathlineto{\pgfqpoint{3.136829in}{1.108776in}}%
\pgfpathlineto{\pgfqpoint{3.137031in}{1.108981in}}%
\pgfpathlineto{\pgfqpoint{3.137846in}{1.108746in}}%
\pgfpathlineto{\pgfqpoint{3.138259in}{1.108849in}}%
\pgfpathlineto{\pgfqpoint{3.138660in}{0.953188in}}%
\pgfpathlineto{\pgfqpoint{3.139034in}{0.953200in}}%
\pgfpathlineto{\pgfqpoint{3.177635in}{0.954410in}}%
\pgfpathlineto{\pgfqpoint{3.177855in}{0.958668in}}%
\pgfpathlineto{\pgfqpoint{3.179381in}{1.225048in}}%
\pgfpathlineto{\pgfqpoint{3.227362in}{1.225028in}}%
\pgfpathlineto{\pgfqpoint{3.228341in}{1.215344in}}%
\pgfpathlineto{\pgfqpoint{3.229321in}{1.225024in}}%
\pgfpathlineto{\pgfqpoint{3.230300in}{1.102155in}}%
\pgfpathlineto{\pgfqpoint{3.231279in}{1.225017in}}%
\pgfpathlineto{\pgfqpoint{3.232258in}{1.108765in}}%
\pgfpathlineto{\pgfqpoint{3.233204in}{1.215815in}}%
\pgfpathlineto{\pgfqpoint{3.234816in}{1.068282in}}%
\pgfpathlineto{\pgfqpoint{3.235769in}{1.097345in}}%
\pgfpathlineto{\pgfqpoint{3.236130in}{1.108856in}}%
\pgfpathlineto{\pgfqpoint{3.236130in}{1.108856in}}%
\pgfusepath{stroke}%
\end{pgfscope}%
\begin{pgfscope}%
\pgfsetrectcap%
\pgfsetmiterjoin%
\pgfsetlinewidth{0.803000pt}%
\definecolor{currentstroke}{rgb}{0.000000,0.000000,0.000000}%
\pgfsetstrokecolor{currentstroke}%
\pgfsetdash{}{0pt}%
\pgfpathmoveto{\pgfqpoint{2.256930in}{0.925988in}}%
\pgfpathlineto{\pgfqpoint{2.256930in}{1.252238in}}%
\pgfusepath{stroke}%
\end{pgfscope}%
\begin{pgfscope}%
\pgfsetrectcap%
\pgfsetmiterjoin%
\pgfsetlinewidth{0.803000pt}%
\definecolor{currentstroke}{rgb}{0.000000,0.000000,0.000000}%
\pgfsetstrokecolor{currentstroke}%
\pgfsetdash{}{0pt}%
\pgfpathmoveto{\pgfqpoint{3.236130in}{0.925988in}}%
\pgfpathlineto{\pgfqpoint{3.236130in}{1.252238in}}%
\pgfusepath{stroke}%
\end{pgfscope}%
\begin{pgfscope}%
\pgfsetrectcap%
\pgfsetmiterjoin%
\pgfsetlinewidth{0.803000pt}%
\definecolor{currentstroke}{rgb}{0.000000,0.000000,0.000000}%
\pgfsetstrokecolor{currentstroke}%
\pgfsetdash{}{0pt}%
\pgfpathmoveto{\pgfqpoint{2.256930in}{0.925988in}}%
\pgfpathlineto{\pgfqpoint{3.236130in}{0.925988in}}%
\pgfusepath{stroke}%
\end{pgfscope}%
\begin{pgfscope}%
\pgfsetrectcap%
\pgfsetmiterjoin%
\pgfsetlinewidth{0.803000pt}%
\definecolor{currentstroke}{rgb}{0.000000,0.000000,0.000000}%
\pgfsetstrokecolor{currentstroke}%
\pgfsetdash{}{0pt}%
\pgfpathmoveto{\pgfqpoint{2.256930in}{1.252238in}}%
\pgfpathlineto{\pgfqpoint{3.236130in}{1.252238in}}%
\pgfusepath{stroke}%
\end{pgfscope}%
\end{pgfpicture}%
\makeatother%
\endgroup%

	\caption{(a) tensão de saída, (b) corrente no indutor de saída e (c) tensão de saídatensão dreno--source do transistor do circuito Forward.}
	\label{fig:ex5}
\end{figure}

Pode-se observar que a tesão de saída apresentou $\approx -220$ mV de diferença para o valor de $12$ V especificado no enunciado. Além disso, ao final da simulação, ainda persiste uma oscilação associada ao transitório de ligamento do conversor. Apesar disso, é possível verificar que o \textit{ripple} de saída é $\approx 1$\%. Ademais, com o auxílio de ferramentas de análise do \textit{software} ngspice, verificou-se que a corrente no indutor de saída possui média $ \approx 2{,}4$ A e não toca o zero, portanto o conversor opera no CCM. Por fim, observou-se que $\Delta i_L \approx 3.96$ A, atendendo as especificação do enunciado.